% !TEX program = xelatex
% ============================================================
% SPSS算法实战：数模美赛快速入门教材
% 主文档 - 使用 xelatex + ctex 编译
% ============================================================
\documentclass[12pt,a4paper,openany,fontset=none]{ctexbook}

% ---------- 中文字体配置（统一使用宋体 SimSun） ----------
\usepackage{fontspec}
\setCJKmainfont{SimSun}[AutoFakeBold=2.5]
\setCJKsansfont{SimSun}[AutoFakeBold=2.5]
\setCJKmonofont{SimSun}
\setmainfont{Times New Roman}
\setsansfont{Times New Roman}

% ---------- 基础宏包 ----------
\usepackage[margin=2.5cm]{geometry}
\usepackage{amsmath,amssymb,amsfonts}
\usepackage{graphicx}
\usepackage{booktabs}
\usepackage{longtable}
\usepackage{array}
\usepackage{multirow}
\usepackage{makecell}
\usepackage[table]{xcolor}
\usepackage{enumitem}
\usepackage{hyperref}
\usepackage{caption}
\usepackage{subcaption}
\usepackage{float}
\usepackage{titlesec}
\usepackage{fancyhdr}
\usepackage{tocloft}
\usepackage{etoolbox}
\usepackage{tikz}
\usepackage{tcolorbox}
\tcbuselibrary{skins,}

% ---------- 超链接 ----------
\hypersetup{
    colorlinks=true,
    linkcolor=blue!60!black,
    citecolor=blue!60!black,
    urlcolor=blue!60!black,
    bookmarksnumbered=true,
}

% ---------- 颜色定义 ----------
\definecolor{algoBlue}{RGB}{30,100,180}
\definecolor{algoGreen}{RGB}{30,130,80}
\definecolor{algoOrange}{RGB}{200,120,30}
\definecolor{algoRed}{RGB}{180,50,50}
\definecolor{algoPurple}{RGB}{120,60,150}
\definecolor{algoTeal}{RGB}{20,120,120}
\definecolor{lightGray}{RGB}{245,245,245}
\definecolor{boxBlue}{RGB}{230,240,252}
\definecolor{boxGreen}{RGB}{230,248,235}
\definecolor{boxOrange}{RGB}{253,243,230}
\definecolor{boxRed}{RGB}{252,235,235}
\definecolor{boxPurple}{RGB}{243,235,250}
\definecolor{boxTeal}{RGB}{230,247,247}

% ---------- 章节标题样式 ----------
\titleformat{\chapter}[display]
    {\normalfont\bfseries\huge\color{algoBlue}}
    {\chaptertitlename\ \thechapter}{20pt}{\Huge}
\titlespacing*{\chapter}{0pt}{-20pt}{30pt}

\titleformat{\section}
    {\normalfont\bfseries\Large\color{algoBlue!80!black}}
    {\thesection}{1em}{}

\titleformat{\subsection}
    {\normalfont\bfseries\large\color{algoBlue!70!black}}
    {\thesubsection}{1em}{}

\titleformat{\subsubsection}
    {\normalfont\bfseries\normalsize\color{algoBlue!60!black}}
    {\thesubsubsection}{1em}{}

% ---------- 自定义模块环境 ----------
\newtcolorbox{algointro}{colback=boxBlue,colframe=algoBlue,boxrule=0.5pt,arc=3pt,left=8pt,right=8pt,top=6pt,bottom=6pt,title={\textbf{\textcolor{white}{一句话讲清}}},colbacktitle=algoBlue,fonttitle=\bfseries}
\newtcolorbox{algoscene}{colback=boxGreen,colframe=algoGreen,boxrule=0.5pt,arc=3pt,left=8pt,right=8pt,top=6pt,bottom=6pt,title={\textbf{\textcolor{white}{美赛适用场景}}},colbacktitle=algoGreen,fonttitle=\bfseries}
\newtcolorbox{algoprinciple}{colback=boxPurple,colframe=algoPurple,boxrule=0.5pt,arc=3pt,left=8pt,right=8pt,top=6pt,bottom=6pt,title={\textbf{\textcolor{white}{通俗原理}}},colbacktitle=algoPurple,fonttitle=\bfseries}
\newtcolorbox{algosposs}{colback=boxOrange,colframe=algoOrange,boxrule=0.5pt,arc=3pt,left=8pt,right=8pt,top=6pt,bottom=6pt,title={\textbf{\textcolor{white}{SPSS操作步骤}}},colbacktitle=algoOrange,fonttitle=\bfseries}
\newtcolorbox{algoresult}{colback=boxTeal,colframe=algoTeal,boxrule=0.5pt,arc=3pt,left=8pt,right=8pt,top=6pt,bottom=6pt,title={\textbf{\textcolor{white}{结果怎么看}}},colbacktitle=algoTeal,fonttitle=\bfseries}
\newtcolorbox{algonotes}{colback=boxRed,colframe=algoRed,boxrule=0.5pt,arc=3pt,left=8pt,right=8pt,top=6pt,bottom=6pt,title={\textbf{\textcolor{white}{注意事项 \& 决策指引}}},colbacktitle=algoRed,fonttitle=\bfseries}
\newtcolorbox{algobrief}{colback=lightGray,colframe=gray!60,boxrule=0.5pt,arc=3pt,left=8pt,right=8pt,top=6pt,bottom=6pt,title={\textbf{\textcolor{white}{进阶算法速览}}},colbacktitle=gray!70,fonttitle=\bfseries}
\newtcolorbox{chapteroverview}{colback=boxBlue!50,colframe=algoBlue,boxrule=1pt,arc=4pt,left=10pt,right=10pt,top=8pt,bottom=8pt,title={\textbf{\textcolor{white}{本章导览}}},colbacktitle=algoBlue!80,fonttitle=\bfseries\large}

% ---------- 算法标题命令 ----------
% 基础算法标题：\basicalgo{算法名称}{序号}
\newcommand{\basicalgo}[2]{
    \subsection[#1]{#1 \hfill \small \textcolor{gray}{基础级}}
    \label{algo:#2}
}

% 进阶算法标题：\advancedalgo{算法名称}{序号}
\newcommand{\advancedalgo}[2]{
    \subsection[#1]{#1 \hfill \small \textcolor{gray!70}{进阶级}}
    \label{algo:#2}
}

% 星标命令
\newcommand{\starmark}{\textcolor{algoOrange}{$\bigstar$}\,}

% ---------- 列表样式 ----------
\setlist[itemize]{leftmargin=*, itemsep=2pt, topsep=4pt}
\setlist[enumerate]{leftmargin=*, itemsep=2pt, topsep=4pt}

% ---------- 页眉页脚 ----------
\setlength{\headheight}{14pt}
\pagestyle{fancy}
\fancyhf{}
\fancyhead[L]{\small \leftmark}
\fancyhead[R]{\small \thepage}
\fancyfoot[C]{}
\renewcommand{\headrulewidth}{0.4pt}
\renewcommand{\footrulewidth}{0pt}

% ---------- 目录设置 ----------
\setcounter{tocdepth}{2}
\setlength{\cftbeforechapskip}{8pt}

% ---------- 表格样式辅助 ----------
\newcommand{\tableheader}[1]{\rowcolor{algoBlue!15}\textbf{#1}}

% ============================================================
% 文档开始
% ============================================================
\begin{document}

% ---------- 封面 ----------
\begin{titlepage}
\centering
\vspace*{3cm}
{\Huge\bfseries\textcolor{algoBlue}{SPSS算法实战}\par}
\vspace{0.5cm}
{\LARGE\bfseries\textcolor{algoBlue!80}{数模美赛快速入门教材}\par}
\vspace{1cm}
{\Large 覆盖18大分类 · 252个算法 · 从入门到实战\par}
\vspace{2cm}
\begin{tikzpicture}[remember picture, overlay]
    \draw[algoBlue!30, line width=2pt] (current page.west) ++(0,-2) -- ++(3,0);
    \draw[algoBlue!30, line width=2pt] (current page.east) ++(0,-2) -- ++(-3,0);
\end{tikzpicture}
\vspace{3cm}
{\large\itshape 写给大一数模新手的SPSS算法百科全书\par}
\vspace{0.5cm}
{\large 2026年9月\par}
\vfill
{\small 基于Dabbit AI SPSS算法工具箱252个真实案例编写\par}
\end{titlepage}

% ---------- 前言 ----------
\frontmatter
\chapter*{前言}
\addcontentsline{toc}{chapter}{前言}
% 前言
\section*{这本书是给谁看的？}

如果你是一名刚上大一、对数模美赛（MCM/ICM）充满好奇但又不知从何下手的同学，这本书就是为你写的。

你可能听说过"回归分析""层次分析法""时间序列预测"这些名词，但不知道它们到底是干什么的、什么时候该用、在SPSS里怎么点出来。你可能看过一些获奖论文，里面满是公式和图表，却看不懂作者为什么选这个模型而不是那个。

这本书的目标很简单：\textbf{让你在最短时间内搞懂252个常用统计与建模算法，知道每个算法是干什么的、美赛什么题能用、SPSS怎么操作、结果怎么看。}

\section*{这本书怎么用？}

\begin{itemize}
    \item \textbf{按章节顺序读}：全书18章，从最基础的描述性统计开始，逐步深入到回归、时间序列、机器学习等高级方法。每章开头有"本章导览"，告诉你这个分类是干什么的、美赛什么题型常用、建议怎么学。
    \item \textbf{基础级 vs 进阶级}：每章内的算法分为两级。\textbf{基础级}算法有完整的六大模块讲解，适合新手仔细阅读；\textbf{进阶级}算法以"速览"形式呈现核心要点，适合有一定基础后查阅。
    \item \textbf{星标算法}：标有{\starmark}的算法是美赛中最常用的核心算法，建议优先掌握。
    \item \textbf{当工具书查}：不需要从头读到尾。比赛时遇到问题，直接翻到对应章节查找合适的算法即可。
\end{itemize}

\section*{每个算法讲什么？}

每个基础级算法都包含以下六个模块，用不同颜色的框体区分，一目了然：

\begin{center}
\begin{tabular}{cl}
\toprule
\textbf{模块} & \textbf{内容} \\
\midrule
\textcolor{algoBlue}{$\blacksquare$} 一句话讲清 & 用大白话告诉你这个算法是干什么的 \\
\textcolor{algoGreen}{$\blacksquare$} 美赛适用场景 & 什么题型、什么问题适合用它 \\
\textcolor{algoPurple}{$\blacksquare$} 通俗原理 & 用生活类比讲清原理，公式配文字说明 \\
\textcolor{algoOrange}{$\blacksquare$} SPSS操作步骤 & 具体点哪个菜单、选什么变量、参数怎么填 \\
\textcolor{algoTeal}{$\blacksquare$} 结果怎么看 & 输出表格中每个关键指标什么意思、怎么判断好坏 \\
\textcolor{algoRed}{$\blacksquare$} 注意事项 \& 决策指引 & 什么时候用、什么时候不用、常见坑、和相似算法的区别 \\
\bottomrule
\end{tabular}
\end{center}

\section*{阅读约定}

\begin{itemize}
    \item 菜单路径用 \textbf{分析 → 回归 → 线性} 的形式表示，即依次点击"分析"菜单、"回归"子菜单、"线性"选项。
    \item 变量名和参数值用\textbf{粗体}标注。
    \item 公式中的专业术语首次出现时会用通俗语言解释。
    \item 书中的SPSS操作基于SPSS 26及以上版本，不同版本界面可能略有差异。
\end{itemize}

\vspace{1em}
\begin{center}
\textit{祝你在数模美赛的旅程中，从"看不懂"到"用得溜"，加油！}
\end{center}


% ---------- 目录 ----------
\tableofcontents
\cleardoublepage

% ---------- 正文 ----------
\mainmatter

%% 各章节由子任务生成，按顺序 include
% 第1章 描述性统计
\chapter{描述性统计}

\begin{chapteroverview}
\textbf{这个分类是干什么的？}

描述性统计是所有数据分析的起点。它不做预测、不做推断，只回答一个问题：\textbf{你的数据长什么样？}通过计算均值、中位数、标准差、频数等指标，帮你快速了解数据的集中趋势、离散程度和分布形态。

\textbf{美赛什么题型常用？}

几乎所有题目都需要！无论是A题的连续型问题、B题的离散型问题，还是C题的大数据题，拿到数据后的第一步永远是描述性统计。它是后续所有建模的基础——你连数据长什么样都不知道，怎么选模型？

\textbf{学习路径建议}

先学\textbf{描述性统计}和\textbf{频数分析}（数值变量和分类变量的基础画像），再学\textbf{正态性检验}（判断数据分布，决定后续用参数还是非参数方法），最后了解\textbf{数据探索}和\textbf{列联分析}等进阶工具。
\end{chapteroverview}

\section{基础级算法}

%% ========== 描述性统计 ==========
\basicalgo{{\starmark}描述性统计}{desc_stats}

\begin{algointro}
描述性统计就是给你的数据做一次"全面体检"——算出均值、中位数、标准差、最大值、最小值等指标，让你一眼看清数据的"长相"。
\end{algointro}

\begin{algoscene}
\begin{itemize}
    \item 拿到任何数据集后的\textbf{第一步}，必须做！
    \item 美赛C题（大数据题）中，对题目给出的CSV数据进行初步探索
    \item 论文中"数据描述"部分的核心内容
    \item 判断数据是否有异常值、缺失值，决定后续清洗策略
\end{itemize}
\end{algoscene}

\begin{algoprinciple}
想象你是班主任，拿到全班的考试成绩单。你会怎么看？

\begin{itemize}
    \item \textbf{集中趋势}：平均分多少？（均值）中间水平是多少？（中位数）出现最多的分数是？（众数）——这告诉你班级整体水平
    \item \textbf{离散程度}：最高分和最低分差多少？（极差）大家分数差距大不大？（标准差、方差）——这告诉你成绩是否稳定
    \item \textbf{分布形态}：分数是对称的还是偏的？（偏度）是集中在中间还是两头多？（峰度）——这告诉你分布形状
    \item \textbf{分位数}：前25\%的人考了多少分？（四分位数）——这告诉你数据的分段情况
\end{itemize}

\textbf{关键公式：}
\begin{itemize}
    \item 均值：$\bar{x} = \frac{1}{n}\sum_{i=1}^{n}x_i$，就是所有数加起来除以个数
    \item 标准差：$s = \sqrt{\frac{1}{n-1}\sum_{i=1}^{n}(x_i-\bar{x})^2}$，衡量数据偏离均值的平均程度，越大说明数据越分散
    \item 偏度（Skewness）：大于0表示右偏（长尾在右），小于0表示左偏，接近0表示对称
    \item 峰度（Kurtosis）：大于3表示尖峰（数据集中在中间），小于3表示平峰
\end{itemize}
\end{algoprinciple}

\begin{algosposs}
\begin{enumerate}
    \item 打开SPSS，导入数据文件（\textbf{文件 → 打开 → 数据}）
    \item 点击菜单 \textbf{分析 → 描述统计 → 描述}
    \item 在弹出的对话框中，将需要分析的变量选入\textbf{变量}框
    \item 点击右侧\textbf{选项}按钮，勾选需要的统计量：
    \begin{itemize}
        \item 集中趋势：均值、中位数、众数
        \item 离散程度：标准差、方差、极差、最小值、最大值
        \item 分布形态：偏度、峰度
        \item 分位数：点击"百分位数"可添加自定义分位点
    \end{itemize}
    \item 如有分组变量（如按性别分别统计），先使用 \textbf{数据 → 拆分文件} 按组拆分，再执行上述操作
    \item 点击\textbf{确定}，输出结果
\end{enumerate}

\textbf{默认参数参考}：集中趋势选均值+中位数，离散程度选标准差+四分位距，分布形态选偏度+峰度，离群判定倍数1.5。
\end{algosposs}

\begin{algoresult}
SPSS会输出一个"描述统计量"表格，每一行是一个变量，每一列是一个指标：

\begin{itemize}
    \item \textbf{N（样本量）}：有效数据个数。如果远小于总行数，说明缺失值多
    \item \textbf{均值（Mean）}：平均值。注意：如果数据偏态严重，均值会被极端值拉偏，此时看中位数更可靠
    \item \textbf{标准差（Std. Deviation）}：越大越分散。一般来说，标准差/均值 $>$ 1 说明离散程度较大
    \item \textbf{最小值/最大值}：检查是否有不合理的值（如年龄出现200岁）
    \item \textbf{偏度（Skewness）}：绝对值 $>$ 1 表示严重偏态，此时均值不可靠，应使用中位数
    \item \textbf{峰度（Kurtosis）}：SPSS中默认输出的是"超额峰度"（已减3），接近0表示正态，正值表示尖峰
\end{itemize}

\textbf{判断好坏}：描述性统计没有"好坏"之分，关键是通过这些指标\textbf{发现问题}——有没有异常值？数据偏不偏？缺失多不多？这些发现决定了你下一步该怎么做。
\end{algoresult}

\begin{algonotes}
\textbf{什么时候用？}
\begin{itemize}
    \item 拿到任何数据后的\textbf{第一步}，没有例外
    \item 论文中需要报告数据基本情况时
\end{itemize}

\textbf{什么时候不用？}
\begin{itemize}
    \item 不要用描述性统计的结果直接下结论（如"A组均值大于B组，所以A组更好"）——组间差异需要做假设检验（见第3章差异比较）
    \item 数据严重偏态时，不要只看均值，必须结合中位数
\end{itemize}

\textbf{常见坑}
\begin{itemize}
    \item \textbf{坑1}：忽略缺失值。SPSS默认剔除缺失值后计算，N值会变小，务必注意
    \item \textbf{坑2}：把分类变量（如性别编码为1、2）当数值变量算均值，得到无意义的"1.5"
    \item \textbf{坑3}：只看均值不看标准差。两个班平均分都是80，但一个班标准差5，另一个标准差20，差别巨大
\end{itemize}

\textbf{和相似算法的区别}
\begin{itemize}
    \item vs \textbf{频数分析}：描述性统计用于\textbf{数值变量}，频数分析用于\textbf{分类变量}
    \item vs \textbf{数据探索}：描述性统计只算指标，数据探索还会自动检测异常值、缺失模式等更深入的质量问题
\end{itemize}
\end{algonotes}

%% ========== 正态性检验 ==========
\basicalgo{{\starmark}正态性检验}{normality_test}

\begin{algointro}
正态性检验就是判断你的数据"长得像不像正态分布（钟形曲线）"。这很重要，因为很多高级统计方法（如t检验、方差分析、线性回归）都要求数据近似正态，不符合的话就得换方法。
\end{algointro}

\begin{algoscene}
\begin{itemize}
    \item 在做t检验、方差分析、线性回归等\textbf{参数检验之前}，必须先做正态性检验
    \item 美赛中需要验证模型假设时（如回归残差是否正态）
    \item 判断该用参数方法还是非参数方法（非正态时用非参数检验）
\end{itemize}
\end{algoscene}

\begin{algoprinciple}
想象你在猜一个人的身高。大多数人身高集中在中间（170cm左右），特别高和特别矮的人很少——画出来就是一个钟形曲线，这就是\textbf{正态分布}。

正态性检验的思路是：\textbf{先假设数据服从正态分布，然后看实际数据和正态分布的差距有多大}。如果差距小到可以用"随机误差"解释，就认为是正态的；如果差距太大，就拒绝正态假设。

\textbf{两种常用检验：}
\begin{itemize}
    \item \textbf{Shapiro-Wilk检验}：适合小样本（$n \leq 5000$），检验功效更高
    \item \textbf{Kolmogorov-Smirnov检验（K-S检验）}：适合大样本，对尾部差异更敏感
\end{itemize}

\textbf{核心概念——p值}：p值 $> 0.05$ 表示"不能拒绝正态假设"，即认为数据正态；p值 $\leq 0.05$ 表示数据显著偏离正态。

\textbf{辅助判断}：除了检验，还要看\textbf{偏度}和\textbf{峰度}，以及\textbf{Q-Q图}（数据点是否大致在一条直线上）。三者结合判断最可靠。
\end{algoprinciple}

\begin{algosposs}
\begin{enumerate}
    \item 点击菜单 \textbf{分析 → 描述统计 → 探索}
    \item 将待检验变量选入\textbf{因变量列表}
    \item 如有分组变量（如按性别分别检验），选入\textbf{因子列表}
    \item 点击右侧\textbf{绘制}按钮，勾选\textbf{含检验的正态图}（同时输出Q-Q图和检验结果）
    \item 点击\textbf{继续}，再点击\textbf{确定}
\end{enumerate}

\textbf{替代路径}：也可通过 \textbf{分析 → 非参数检验 → 单样本}，选择"正态性检验"。

\textbf{默认参数}：显著性水平 $\alpha=0.05$，方法自动选择（小样本Shapiro-Wilk，大样本K-S）。
\end{algosposs}

\begin{algoresult}
输出"正态性检验"表格，包含Shapiro-Wilk和K-S两行：

\begin{itemize}
    \item \textbf{统计量（Statistic）}：越接近1表示越接近正态
    \item \textbf{显著性（Sig. / p值）}：\textbf{最关键的指标}
    \begin{itemize}
        \item $p > 0.05$：数据服从正态分布，可以使用参数检验
        \item $p \leq 0.05$：数据不服从正态分布，应考虑非参数检验或数据变换
    \end{itemize}
\end{itemize}

同时输出\textbf{正态Q-Q图}：如果数据点大致落在对角线上，说明正态性良好；如果点明显偏离直线（尤其是两端），说明偏离正态。

\textbf{注意}：样本量很大时（如 $n>1000$），即使数据轻微偏离正态，p值也会小于0.05。此时应结合Q-Q图和偏度峰度综合判断，不要只看p值。
\end{algoresult}

\begin{algonotes}
\textbf{什么时候用？}
\begin{itemize}
    \item 参数检验（t检验、ANOVA、线性回归等）之前的必做步骤
    \item 论文中需要报告"数据满足正态性假设"时
\end{itemize}

\textbf{什么时候不用？}
\begin{itemize}
    \item 非参数检验（如Mann-Whitney U、Kruskal-Wallis）不需要正态性假设
    \item 样本量极大（$n>5000$）时，中心极限定理保证均值近似正态，可不必纠结
\end{itemize}

\textbf{常见坑}
\begin{itemize}
    \item \textbf{坑1}：只看p值不看图。大样本下p值几乎必然显著，应结合Q-Q图判断
    \item \textbf{坑2}：对回归模型检验原始变量的正态性。回归要求的是\textbf{残差}正态，不是自变量或因变量本身正态
    \item \textbf{坑3}：数据非正态就放弃。可以尝试对数变换、平方根变换等，变换后可能就正态了
\end{itemize}

\textbf{数据非正态怎么办？}
\begin{enumerate}
    \item 尝试变量变换（对数log、平方根、Box-Cox变换）
    \item 改用非参数检验方法（见第3章）
    \item 使用稳健统计方法（如稳健回归）
\end{enumerate}
\end{algonotes}

%% ========== 频数分析 ==========
\basicalgo{频数分析}{freq_analysis}

\begin{algointro}
频数分析就是数一数每个类别出现了多少次、占百分之多少。比如调查了100个人，其中男52人、女48人，这就是频数分析。它是分类变量的"描述性统计"。
\end{algointro}

\begin{algoscene}
\begin{itemize}
    \item 分析问卷中的单选题、多选题的选项分布
    \item 美赛中对分类数据（如地区、行业、等级）的分布描述
    \item 论文中"样本结构"部分的统计
\end{itemize}
\end{algoscene}

\begin{algoprinciple}
频数分析的核心就是两个数：
\begin{itemize}
    \item \textbf{频数（Frequency）}：每个类别出现的次数
    \item \textbf{百分比（Percent）}：每个类别占总数的比例
\end{itemize}

还可以计算\textbf{累计百分比}，比如"前两个类别合计占75\%"，帮助你快速了解主要类别。

如果有分组变量（如不同性别的满意度分布），还可以做\textbf{分组结构对比}，看不同群体的选择是否有差异。
\end{algoprinciple}

\begin{algosposs}
\begin{enumerate}
    \item 点击菜单 \textbf{分析 → 描述统计 → 频率}
    \item 将分类变量选入\textbf{变量}框
    \item 点击\textbf{统计量}按钮，可勾选百分位数、集中趋势等（分类变量一般不需要）
    \item 点击\textbf{图表}按钮，选择\textbf{条形图}或\textbf{饼图}
    \item 点击\textbf{格式}按钮，可选择按频数升序/降序排列
    \item 点击\textbf{确定}
\end{enumerate}

\textbf{默认参数}：显示频数、百分比、有效百分比、累计百分比；图表可选条形图+饼图。
\end{algosposs}

\begin{algoresult}
输出"频率"表格，包含四列：
\begin{itemize}
    \item \textbf{频数}：该类别的实际数量
    \item \textbf{百分比}：占全部样本（含缺失）的比例
    \item \textbf{有效百分比}：占有效样本（剔除缺失）的比例——通常看这个
    \item \textbf{累计百分比}：从上到下累加的比例
\end{itemize}

\textbf{怎么看}：找到有效百分比最高的类别，那就是"主流"。如果某个类别占比超过50\%，说明数据高度集中；如果各类别比较均匀，说明分布分散。
\end{algoresult}

\begin{algonotes}
\textbf{什么时候用？}分类变量的分布描述，尤其是问卷数据。

\textbf{什么时候不用？}数值变量不要用频数分析（除非已经分箱），应该用描述性统计。

\textbf{常见坑}：混淆"百分比"和"有效百分比"，有缺失值时两者不同，论文中应明确说明。
\end{algonotes}

%% ========== 数据概览 ==========
\basicalgo{数据概览}{data_overview}

\begin{algointro}
数据概览是描述性统计的"加强版"——它不仅算统计量，还自动检测数据类型、缺失值、异常值、常量列等问题，给你一份完整的"数据体检报告"。
\end{algointro}

\begin{algoscene}
\begin{itemize}
    \item 拿到一个新数据集，想快速了解整体情况
    \item 美赛C题大数据分析的第一步数据清洗
    \item 需要向队友或评委展示数据质量时
\end{itemize}
\end{algoscene}

\begin{algoprinciple}
数据概览从三个层面看数据：
\begin{enumerate}
    \item \textbf{结构层}：有多少行、多少列？每列是什么类型（数值/分类/时间/标识）？
    \item \textbf{质量层}：缺失值有多少？有没有重复行？有没有常量列（所有值都一样）？有没有异常值？
    \item \textbf{分布层}：数值变量的均值、中位数、标准差；分类变量的各类别占比
\end{enumerate}

它会综合这些信息给出一个\textbf{数据质量评分}（0-100分），并标出最严重的问题，帮你决定清洗优先级。
\end{algoprinciple}

\begin{algosposs}
\begin{enumerate}
    \item 在Dabbit AI SPSS工具箱中选择"数据概览"
    \item 上传数据文件
    \item 设置参数：概览深度选"标准"，异常值检测倍数1.5，缺失率告警阈值0.2
    \item 点击运行，获取完整的数据画像报告
\end{enumerate}

\textbf{SPSS中对应操作}：可结合 \textbf{分析 → 描述统计 → 描述} 和 \textbf{分析 → 报告 → 代码本} 实现类似功能。
\end{algosposs}

\begin{algoresult}
报告包含：
\begin{itemize}
    \item \textbf{数据集摘要}：行数、列数、文件大小
    \item \textbf{变量类型分布}：数值型几个、分类型几个、时间型几个
    \item \textbf{缺失值热力图}：直观展示哪些变量缺失严重
    \item \textbf{异常值检测}：用IQR法（超过$Q_3+1.5\times IQR$或低于$Q_1-1.5\times IQR$）标记异常值
    \item \textbf{数据质量评分}：综合评分，低于60分需要重点清洗
\end{itemize}
\end{algoresult}

\begin{algonotes}
\textbf{和描述性统计的区别}：描述性统计只算指标，数据概览还做质量诊断和异常检测，功能更全面。

\textbf{建议}：大型数据集先做数据概览，再对关键变量做详细的描述性统计。
\end{algonotes}

%% ========== 数据探索（数据质量体检） ==========
\basicalgo{数据探索（数据质量体检）}{data_exploration}

\begin{algointro}
数据探索是更深入的"数据质量体检"——它不仅发现问题，还会分析缺失值的模式、异常值的类型、数据的分布特征，给你一份可操作的清洗建议清单。
\end{algointro}

\begin{algoscene}
\begin{itemize}
    \item 数据质量较差、需要系统性清洗时
    \item 美赛中需要在论文中详细说明数据预处理过程时
    \item 建模前的最后一道数据质量检查
\end{itemize}
\end{algoscene}

\begin{algoprinciple}
数据探索分两层检查：
\begin{enumerate}
    \item \textbf{数据层}：检查文件数量、记录数、变量数，核对多表合并关系
    \item \textbf{字段层}：逐列检查缺失率、类型解析、取值集中度、异常值
\end{enumerate}

异常值检测有两种常用方法：
\begin{itemize}
    \item \textbf{IQR法}：超过$Q_3+1.5\times IQR$或低于$Q_1-1.5\times IQR$的值（默认1.5倍）
    \item \textbf{Z分数法}：Z分数绝对值超过阈值（通常为3）的值
\end{itemize}

最终合成0-100的质量评分，划分优/良/中/差四个等级。
\end{algoprinciple}

\begin{algosposs}
\begin{enumerate}
    \item 在工具箱中选择"数据探索（数据质量体检）"
    \item 上传数据
    \item 设置参数：异常值方法选"IQR法"，IQR倍数1.5，缺失率告警阈值0.30
    \item 勾选"全行重复检查"和"标识字段唯一性检查"
    \item 运行并查看报告
\end{enumerate}
\end{algosposs}

\begin{algoresult}
报告重点关注：
\begin{itemize}
    \item \textbf{缺失率超过30\%的变量}：考虑删除或用模型插补
    \item \textbf{异常值数量和位置}：判断是录入错误还是真实极端值
    \item \textbf{重复记录}：需要去重
    \item \textbf{常量列}：所有值相同的变量对建模无意义，可删除
    \item \textbf{质量评分}：低于60分必须清洗后再建模
\end{itemize}
\end{algoresult}

\begin{algonotes}
\textbf{和数据概览的区别}：数据概览偏"画像"，数据探索偏"诊断+建议"，后者更深入。

\textbf{注意}：异常值不一定是错误，可能是真实的极端情况（如收入数据中的亿万富翁），不要盲目删除，要结合业务判断。
\end{algonotes}

%% ========== 列联（交叉）分析 ==========
\basicalgo{列联（交叉）分析}{crosstab}

\begin{algointro}
列联分析就是把两个分类变量交叉起来看关系——比如"不同会员等级的续订意向有没有差异？"，它会生成一个交叉表，并通过卡方检验判断两个变量是否相关。
\end{algointro}

\begin{algoscene}
\begin{itemize}
    \item 分析两个分类变量之间的关系（如性别 vs 购买偏好）
    \item 美赛中问卷数据的交叉分析
    \item 需要判断分类变量之间是否独立时
\end{itemize}
\end{algoscene}

\begin{algoprinciple}
列联分析的核心是\textbf{交叉表}：行是一个变量的类别，列是另一个变量的类别，每个格子里是同时满足两个条件的样本数。

然后用\textbf{卡方独立性检验}判断：如果两个变量无关，那么每个格子的"期望频数"应该是多少？实际频数和期望频数的差异越大，说明两个变量越可能相关。

\textbf{卡方统计量}：$\chi^2 = \sum\frac{(O-E)^2}{E}$，其中O是观测频数，E是期望频数。$\chi^2$越大，p值越小，越可能相关。

\textbf{效应量}：Cramér's V衡量关联强度，0-1之间，越大关联越强（一般0.1以下弱，0.3中等，0.5以上强）。
\end{algoprinciple}

\begin{algosposs}
\begin{enumerate}
    \item 点击菜单 \textbf{分析 → 描述统计 → 交叉表}
    \item 将一个变量选入\textbf{行}，另一个选入\textbf{列}
    \item 点击\textbf{统计量}按钮，勾选\textbf{卡方}和\textbf{Phi和Cramér系数}
    \item 点击\textbf{单元格}按钮，勾选\textbf{观察值}、\textbf{期望值}、\textbf{行百分比}、\textbf{列百分比}
    \item 点击\textbf{确定}
\end{enumerate}

\textbf{默认参数}：显著性水平0.05，效应量用Cramér's V，期望频数诊断阈值5（低于5的格子超过20\%时需用Fisher精确检验）。
\end{algosposs}

\begin{algoresult}
\begin{itemize}
    \item \textbf{交叉表}：每个格子显示观测数、期望数、行百分比、列百分比
    \item \textbf{卡方检验表}：看Pearson卡方的\textbf{渐近显著性（p值）}
    \begin{itemize}
        \item $p > 0.05$：两变量独立（无关联）
        \item $p \leq 0.05$：两变量显著相关
    \end{itemize}
    \item \textbf{Cramér's V}：关联强度，0.1以下弱、0.1-0.3中等偏弱、0.3-0.5中等、0.5以上强
    \item \textbf{调整后标准化残差}：绝对值大于2的格子表示该类别组合显著偏离期望
\end{itemize}
\end{algoresult}

\begin{algonotes}
\textbf{前提条件}：期望频数低于5的格子不能超过20\%，否则应用Fisher精确检验。

\textbf{和卡方检验的关系}：列联分析包含交叉表+卡方检验，第3章的卡方检验侧重检验本身，这里侧重交叉表的解读。

\textbf{注意}：相关不等于因果！列联分析只能说明两个变量有关联，不能说明谁导致谁。
\end{algonotes}

\section{进阶级算法速览}

%% ========== 个体分布标识 ==========
\advancedalgo{个体分布标识}{dist_id}

\begin{algobrief}
\textbf{一句话}：判断一个连续变量最符合哪种概率分布（正态、对数正态、指数、伽玛、威布尔等）。

\textbf{美赛场景区}：可靠性分析中判断设备寿命服从什么分布；金融数据中判断收益率分布。

\textbf{核心原理}：对多种候选分布分别做极大似然估计，用AIC（赤池信息准则）比较拟合优度，AIC最小的分布最优。再用Anderson-Darling检验确认拟合效果。

\textbf{操作要点}：工具箱中选择"个体分布标识"，选分析变量，候选分布集选"常用六种"，显著性水平0.05。

\textbf{结果关键}：看AIC值最小的分布，以及Anderson-Darling检验的p值（$>0.05$表示拟合可接受）。

\textbf{注意}：这是比正态性检验更一般的方法——正态性检验只判断"是不是正态"，个体分布标识会在6种常见分布中帮你找最合适的。
\end{algobrief}

%% ========== 泊松分布检验 ==========
\advancedalgo{泊松分布检验}{poisson_test}

\begin{algobrief}
\textbf{一句话}：判断计数数据（如每小时来电数、每天事故数）是否服从泊松分布。

\textbf{美赛场景区}：排队论问题中判断到达过程是否为泊松过程；运营优化中分析事件发生规律。

\textbf{核心原理}：泊松分布的核心特征是\textbf{均值=方差}。先用离差检验比较方差与均值，再用卡方拟合优度检验比较观测频数与泊松期望频数。

\textbf{操作要点}：选择"泊松分布检验"，检验方法选"离差检验+卡方拟合优度"，显著性水平0.05。

\textbf{结果关键}：离差统计量接近1且p值$>0.05$，卡方检验p值$>0.05$，则服从泊松分布。

\textbf{注意}：如果方差远大于均值（过度离散），应考虑负二项分布；如果方差远小于均值，考虑零膨胀模型。
\end{algobrief}

%% ========== 分类汇总 ==========
\advancedalgo{分类汇总}{categorical_summary}

\begin{algobrief}
\textbf{一句话}：按一个或多个分类变量分组，对数值指标计算汇总统计量（均值、总和、中位数等）。

\textbf{美赛场景区}：按地区/年份汇总经济指标；按实验组别计算因变量的均值和标准差。

\textbf{核心原理}：就是SQL中的GROUP BY操作——先按分组变量拆分数据，再对每组计算统计量。

\textbf{操作要点}：SPSS中通过 \textbf{数据 → 分类汇总} 或 \textbf{分析 → 比较均值 → 均值} 实现。选择分组变量和汇总指标，可选合计行和分组占比。

\textbf{结果关键}：输出分组汇总表，可直接用于论文中的描述性统计表。

\textbf{注意}：分类汇总只做描述，不做组间差异检验。要判断组间差异需用t检验或方差分析（第3章）。
\end{algobrief}

% 第2章 相关分析
\chapter{相关分析}

\begin{chapteroverview}
\textbf{这个分类是干什么的？}

相关分析研究的是\textbf{变量之间"动不动一起变"}的关系。你关心两个数值变量（比如客流量和销售额）是不是同涨同跌、方向如何、关系有多强；或者两个有序变量（比如评分等级）的排序是否一致。它只回答"有没有关系、关系多强"，\textbf{不回答谁导致谁，也不做预测}——预测是回归分析（第4章）的活。

\textbf{美赛什么题型常用？}

几乎所有需要"找关系"的题目都用得上：C题（大数据题）里分析哪些因素和销量/满意度相关；A题/B题里两个物理或经济变量是否联动。论文"变量探索""相关性分析"小节的标配就是Pearson相关系数矩阵。

\textbf{学习路径建议}

先学\textbf{Pearson相关}（最常用，要求两变量都近似正态且线性）和\textbf{偏相关}（控制掉干扰因素后看净关系）；再按数据形态认识\textbf{Spearman/Kendall}（非正态或有序数据）；最后了解\textbf{典型相关、互相关、Kendall协调系数、Cochran's Q}等专门场景工具。
\end{chapteroverview}

\section{基础级算法}

%% ========== 相关分析 Pearson/Spearman/Kendall ==========
\basicalgo{{\starmark}相关分析（Pearson/Spearman/Kendall）}{pearson_corr}

\begin{algointro}
相关分析就是算一个"相关系数"，用一个介于$-1$到$1$之间的数，告诉你两个变量\textbf{方向是同涨同跌还是此消彼长}、\textbf{关系有多紧}。本案例考察门店日均客流量与当月销售总额的联动。
\end{algointro}

\begin{algoscene}
\begin{itemize}
    \item 拿到两个连续变量，想知道它们是否相关、相关多强
    \item 美赛C题中做变量间相关性扫描，为回归选变量做准备
    \item 论文"相关性分析"部分，输出相关系数矩阵和散点图
    \item 判断该用哪种相关系数：Pearson（正态+线性）、Spearman（非正态/单调）、Kendall（有序数据）
\end{itemize}
\end{algoscene}

\begin{algoprinciple}
想象身高和体重：个子高的人通常也更重。把每对$(身高,体重)$画成散点，点越贴近一条从左下到右上的直线，说明两者关系越强。相关系数就是把这种"贴线程度"压缩成一个数。

\begin{itemize}
    \item \textbf{Pearson积差相关 $r$}：最常用，要求两变量都近似正态、关系是直线。本案例中日均客流量与销售额做Shapiro检验，$W$分别为0.994（$p=0.815$）和0.992（$p=0.527$），都远大于0.05，因此用Pearson。
    \item \textbf{Spearman秩相关 $\rho$}：把数据先排序再算相关，不要求正态，只要关系单调（一起变大或一起变小）就行。
    \item \textbf{Kendall tau-b}：专门处理有序分类变量（如评级1、2、3）或结很多的数据。
\end{itemize}

\textbf{怎么读这个数}：
\begin{itemize}
    \item 符号：正号=同涨同跌，负号=此消彼长
    \item 绝对值：0.3以下弱相关，0.3--0.5中等，0.5以上较强，0.8以上极强
    \item $p$值：$p<0.05$说明这个相关不太可能是抽样巧合
\end{itemize}
\end{algoprinciple}

\begin{algosposs}
\begin{enumerate}
    \item 点击菜单 \textbf{分析 → 相关 → 双变量}
    \item 将两个变量（本案例为\textbf{日均客流量}、\textbf{当月销售总额}）选入\textbf{变量}框
    \item 勾选相关系数：默认\textbf{Pearson}；若数据非正态，改勾\textbf{Spearman}
    \item 显著性检验选\textbf{双侧检验}，勾选\textbf{标记显著性相关}
    \item 点击\textbf{选项}，可勾选"均值、标准差"和"叉积偏差、协方差"
    \item 建议先画散点图：\textbf{图形 → 旧对话框 → 散点/点状}，确认关系是直线而非曲线
    \item 点击\textbf{确定}
\end{enumerate}

\textbf{默认参数参考}：显著性水平$\alpha=0.05$，置信水平95\%，正态性检验用Shapiro，缺失值按成对删除（pairwise），方法自动选择（method=auto）。
\end{algosposs}

\begin{algoresult}
输出"相关性"矩阵，本案例关键结果：
\begin{itemize}
    \item \textbf{Pearson相关系数 $r=0.895$}：极强正相关，客流量越大销售额越高
    \item \textbf{显著性 $p<0.001$}：远小于0.05，关系高度可靠
    \item \textbf{有效对 $n=150$}，95\%置信区间为$[0.857,\ 0.923]$，区间不含0，进一步印证显著
    \item 散点图：点沿一条上升直线紧密分布，与$r=0.895$一致
\end{itemize}

\textbf{判断}：$p<0.05$且置信区间不跨0才算"显著相关"；显著只说明关系存在，\textbf{不代表因果}——客流量高可能因为门店位置好，而位置好同时抬高了销售额。
\end{algoresult}

\begin{algonotes}
\textbf{什么时候用？}两个连续（或有序）变量，想知道是否同向变动时。

\textbf{什么时候不用？}
\begin{itemize}
    \item 关系明显是曲线（如先升后降）时，Pearson会严重低估，应先画散点图
    \item 变量是无序分类（如血型）时，相关系数无意义，改用卡方检验（第3章）
\end{itemize}

\textbf{常见坑}
\begin{itemize}
    \item \textbf{坑1}：把相关当因果。冰淇淋销量和溺水人数强相关，但其实是"夏天"这个第三变量在起作用
    \item \textbf{坑2}：忽略散点图。存在一个极端点就能把$r$从0拉到0.8
    \item \textbf{坑3}：大样本下$r$很小（如0.1）也会$p<0.05$，要看$|r|$本身判断实际强弱
\end{itemize}

\textbf{和相似算法的区别}
\begin{itemize}
    \item vs \textbf{偏相关}：这里是直接相关，可能混入第三变量；偏相关会把干扰变量控制掉
    \item vs \textbf{回归分析}：相关只衡量共变强弱，回归还要给出"客流每涨1人，销售额涨多少"的方程
\end{itemize}
\end{algonotes}

%% ========== 偏相关分析 ==========
\basicalgo{偏相关分析}{partial_corr}

\begin{algointro}
偏相关就是\textbf{把干扰变量"按住不动"之后}，再看两个变量的纯净相关。本案例控制营业面积、竞争门店数、促销投入后，看客流量和销售额到底还有多强的真实关系。
\end{algointro}

\begin{algoscene}
\begin{itemize}
    \item 两变量的直接相关里混进了第三变量（如门店规模同时抬高客流和销售额），想剥离后看"净效应"
    \item 美赛中判断某因素是否在控制其他因素后仍独立起作用
    \item 论文中区分"表面相关"和"净相关"，让结论更可信
\end{itemize}
\end{algoscene}

\begin{algoprinciple}
直接算客流和销售额的相关是$r=0.902$，但这里面有水分——大门店本来客流多、销售也多，是"面积"在背后同时推着两边。

偏相关的思路像\textbf{做实验时控制变量}：先用营业面积、竞争门店数、促销投入分别去"预测"客流和销售额，把这两部分能被控制变量解释的波动\textbf{减掉（取残差）}，再对剩下的残差算相关。本案例控制3个变量后净相关$r=0.541$，比表面的0.902降了一大截，说明原来的强相关里有相当一部分是规模和投入造成的假象，但控制后依然显著。

\textbf{自由度}会随控制变量个数减少：$df=n-2-k$（$k$为控制变量数），本案例$n=135$、$k=3$，故$df=130$。
\end{algoprinciple}

\begin{algosposs}
\begin{enumerate}
    \item 点击菜单 \textbf{分析 → 相关 → 偏相关}
    \item 将要研究的两个变量（\textbf{客流量}、\textbf{销售额}）选入\textbf{变量}框
    \item 将需要控制的变量（\textbf{营业面积}、\textbf{竞争门店数}、\textbf{促销投入}）选入\textbf{控制}框
    \item 显著性选\textbf{双侧}
    \item 点击\textbf{选项}，勾选"零阶相关"——这样能同时看到控制前和控制后的相关，方便对比
    \item 点击\textbf{确定}
\end{enumerate}

\textbf{默认参数参考}：方法用Pearson，缺失按整行删除（listwise），双侧检验，控制变量手动指定。
\end{algosposs}

\begin{algoresult}
输出"偏相关系数"表，关键看两行：
\begin{itemize}
    \item \textbf{零阶相关（未控制）}：$r=0.902$，$p<0.001$——表面看极强
    \item \textbf{偏相关（控制3个变量后）}：$r=0.541$，$df=130$，$p<0.001$，95\%置信区间$[0.408,\ 0.652]$
    \item 解读：相关系数从0.902降到0.541，下降0.361，说明原关系约三分之一是被控制变量解释掉的；但控制后仍显著为正，说明客流对销售的带动是\textbf{稳健存在}的
\end{itemize}

\textbf{辅助检查}：报告还给出控制变量间VIF（均约1.0，远小于10，不存在新的共线性问题）和残差正态性（Shapiro $p=0.464$和$0.360$），这些都说明结果可靠。
\end{algoresult}

\begin{algonotes}
\textbf{什么时候用？}怀疑两变量相关被第三/第四变量污染时，偏相关是最直观的净化工具。

\textbf{什么时候不用？}
\begin{itemize}
    \item 没有明确的、可测量的干扰变量时，硬加控制变量只会消耗自由度
    \item 控制变量之间本身严重共线时，偏相关结果不稳定（本案例VIF≈1.0才安全）
\end{itemize}

\textbf{常见坑}
\begin{itemize}
    \item \textbf{坑1}：控制变量加太多。每加一个控制变量就吃掉一份自由度，样本小时检验功效大降
    \item \textbf{坑2}：把偏相关当因果。它只是"固定其他变量水平后的共变"，依然是观察性证据
\end{itemize}

\textbf{和相似算法的区别}
\begin{itemize}
    \item vs \textbf{Pearson相关}：Pearson是"裸相关"，偏相关是"扣除干扰后的净相关"
    \item vs \textbf{回归分析}：回归能同时控制多个变量并给出每变量的斜率，偏相关只给一个净相关系数，更轻量、更适合做探索
\end{itemize}
\end{algonotes}

\section{进阶级算法速览}

%% ========== Cochran's Q ==========
\advancedalgo{Cochran's Q 检验}{cochran_q}

\begin{algobrief}
\textbf{一句话}：检验\textbf{同一批人}在多个（3个及以上）二分类条件下的"通过率/成功率"是否一致。

\textbf{美赛场景区}：同一组受试者在多种方案下的成败率比较（如四种考核模块通过率是否相同）；问卷中同一批人对多个政策的支持率差异。

\textbf{核心原理}：普通卡方检验要求样本相互独立，但这里同一个人被重复观测，不独立。Cochran's Q把每一行（每个人）的成功次数汇总，构造$Q$统计量，与自由度$k-1$的卡方分布比较。本案例$Q=81.90$，$p<0.001$，四个模块通过率（88.7\%/64.7\%/54.7\%/44.0\%）差异显著。

\textbf{操作要点}：在Dabbit AI SPSS工具箱选择"Cochran's Q检验"，数据按"一人一行、每个条件一列0/1"排列；$\alpha=0.05$，置换次数5000；整体显著后用McNemar做两两比较并用Bonferroni校正。

\textbf{结果关键}：$p<0.05$说明各条件成功率不全相等；再看事后两两比较中校正后$p$仍显著的条件对，定位差异来源（本案例"设备操作"显著弱于"理论考核"）。

\textbf{注意}：它是\textbf{配对版卡方}。两个条件时退化为McNemar检验；三个及以上二分类配对条件才用它。若结果变量是连续的，改用Friedman检验。
\end{algobrief}

%% ========== 典型相关分析 ==========
\advancedalgo{典型相关分析}{cca}

\begin{algobrief}
\textbf{一句话}：研究\textbf{两组变量整体之间}是否相关——不是一对一，而是找两组各自的线性组合，使组合之间的相关最大。

\textbf{美赛场景区}：一组工艺参数（炉温、保温时间等）和一组质量指标（强度、硬度等）之间的整体联动；一组企业财务指标和一组市场表现指标的关联。

\textbf{核心原理}：分别从X组和Y组各造一个"综合得分"（典型变量），让这两个得分的相关$\rho$达到最大；用完一对再在正交方向上找下一对。本案例最强第一对$\rho=0.881$（共享方差77.7\%），共2对显著。载荷绝对值$\geq 0.3$的变量是该典型对的主要贡献者。

\textbf{操作要点}：在工具箱选择"典型相关分析"，X组选工艺变量、Y组选质量变量；数据做z-score标准化；载荷解释阈值0.30；按Wilks检验逐对判断哪些典型对显著。

\textbf{结果关键}：看典型相关系数$\rho$及其$p$值（本案例第1对$p=8.56\times10^{-58}$），再看载荷方向（同号=同向联动）和冗余度（典型相关高但冗余度低则解释价值有限）。

\textbf{注意}：它是Pearson相关在"两组多变量"上的推广。只有两个变量时它就退化成普通相关；样本量要求较高，建议样本数远大于变量总数。
\end{algobrief}

%% ========== 互相关性分析 ==========
\advancedalgo{互相关性分析}{ccf}

\begin{algobrief}
\textbf{一句话}：分析\textbf{两条时间序列}之间的\textbf{滞后相关}——A序列变动后，过几天B序列才跟着动，以及滞后几天时相关最强。

\textbf{美赛场景区}：广告投放量与后续销量的滞后关系；政策出台与经济指标变化的时间差；判断哪个序列是"领先指标"。

\textbf{核心原理}：计算$r(k)=\mathrm{corr}(X_t,\ Y_{t+k})$在不同滞后阶数$k$下的值。本案例广告曝光量与下单笔数同期相关$r(0)=0.079$几乎为0，但滞后$k=3$天处$r=0.724$（$p=1.72\times10^{-18}$）最强——即曝光领先下单约3天。

\textbf{操作要点}：工具箱选"互相关性分析"，指定X序列（曝光量）和Y序列（下单笔数）；最大滞后阶数max\_lag=10；先做ADF平稳性检验，非平稳则差分或预白化；$\alpha=0.05$。

\textbf{结果关键}：互相关函数图中超出蓝色显著带的红点是可靠滞后；最强峰所在的$k^*$就是领先期数。$k>0$表示X领先Y，$k<0$表示Y领先X。

\textbf{注意}：它只描述统计上的先后联动，\textbf{不等于因果}；要证明因果方向需再做Granger检验。两序列不平稳时直接算互相关会得到假相关。
\end{algobrief}

%% ========== Kendall 协调系数 W ==========
\advancedalgo{Kendall协调系数（W）}{kendall_w}

\begin{algobrief}
\textbf{一句话}：衡量\textbf{多个评分者}对\textbf{同一批对象}打分的\textbf{排序一致性}——大家是不是把同一家门店排在前面。

\textbf{美赛场景区}：多位评委/专家对多个方案的评分一致性检验；问卷评分者信度（inter-rater reliability）；综合排名是否可信。

\textbf{核心原理}：把每位评分者的打分转成排名，看所有人给同一对象的排名是否一致。$W$介于0到1，越接近1排序越一致。本案例10名稽核员对15家门店，$W=0.849$（$p=1.06\times10^{-18}$），一致性"很强"。

\textbf{操作要点}：工具箱选"Kendall协调系数"，数据按"评分者为行/列、对象为列/行"排列；结（同分）用平均秩处理；$p$值用卡方近似；分数方向ascending（越高越好）。

\textbf{结果关键}：$W<0.2$很弱、$0.2$--$0.4$弱、$0.4$--$0.6$中等、$0.6$--$0.8$强、$>0.8$很强；同时输出对象平均秩综合名次和每位评分者的方向一致性（识别打反方向的异常评分者）。

\textbf{注意}：$W$只反映\textbf{排序}一致性，不管评分者宽严造成的水平差异（有人普遍打高分）。它和Spearman相关的关系：两两Spearman的平均值约为$(k\cdot W-1)/(k-1)$。
\end{algobrief}

% 第3章 差异比较
\chapter{差异比较}

\begin{chapteroverview}
\textbf{这个分类是干什么的？}

差异比较回答的是\textbf{"两组或多组之间真的有差别，还是只是随机波动？"}比如：新版包装是不是真的卖得更好？A班和B班平均分差那几分是实力差距还是运气？它是所有假设检验的核心——用一个$p$值告诉你"观察到的差异，有多大可能是碰巧"。

\textbf{美赛什么题型常用？}

美赛中只要出现"实验组vs对照组""不同地区/方案/人群"的对比，几乎都要做差异检验。C题里比较不同政策效果、不同客群表现；A/B题里验证参数变化前后系统输出是否不同。论文"结果分析""模型验证"部分常出现t检验、方差分析、卡方检验。

\textbf{学习路径建议}
先搞清楚\textbf{数据是啥类型}再选方法：连续变量比较两组用\textbf{t检验}，比较三组及以上用\textbf{方差分析}；分类变量用\textbf{卡方检验}；比较前先看\textbf{方差齐性}和正态性——不满足就换\textbf{非参数检验}（Wilcoxon、Mann-Whitney、Kruskal-Wallis、Friedman）。最后再了解事后多重比较、协方差分析等进阶工具。
\end{chapteroverview}

\section{基础级算法}

%% ========== 方差齐性检验 ==========
\basicalgo{方差齐性检验}{levene_homog}

\begin{algointro}
方差齐性检验就是看\textbf{几组数据的"波动幅度"是不是一样大}。本案例比较4条灌装生产线净含量的波动：是不是某条线特别不稳定。它本身是个"体检项"，却决定了你后面能不能用普通t检验和方差分析。
\end{algointro}

\begin{algoscene}
\begin{itemize}
    \item 做独立样本t检验、方差分析之前，必须先过这一关
    \item 质量管理中比较几条生产线/几台设备的稳定性
    \item 美赛中判断各组数据离散程度是否一致，决定后续用哪种检验
\end{itemize}
\end{algoscene}

\begin{algoprinciple}
两个班平均分都是80，一个班分数挤在75--85，另一个班从60到100都有——虽然均值一样，但\textbf{稳定性}天差地别。方差齐性检验就是比"挤不挤"。

\begin{itemize}
    \item \textbf{原假设}：各组方差相等（波动一致）
    \item \textbf{Levene检验}：把每个数据点到本组中位数的距离拿出来，再对这些"距离"做一次方差分析。它对非正态比较稳健，是默认推荐
    \item \textbf{Bartlett检验}：要求各组都近似正态，检验更灵敏但扛不住偏态
\end{itemize}

本案例4条线标准差分别为1.84、1.88、2.02、6.40——第4条线波动是其他线的3倍多，一眼就不齐。
\end{algoprinciple}

\begin{algosposs}
\begin{enumerate}
    \item 点击菜单 \textbf{分析 → 描述统计 → 探索}
    \item 将因变量（\textbf{净含量}）选入\textbf{因变量列表}，分组变量（\textbf{生产线}）选入\textbf{因子列表}
    \item 点击右侧\textbf{绘制}按钮，勾选\textbf{含莱文检验的分布水平图}
    \item 点击\textbf{继续}，再\textbf{确定}
\end{enumerate}

\textbf{另一种方式}：在独立样本t检验、单因素ANOVA对话框里直接点\textbf{选项}，勾选"方差齐性Levene检验"，会自动附带输出。

\textbf{默认参数参考}：显著性水平$\alpha=0.05$，方法自动选择（各组正态用Bartlett，否则用Brown-Forsythe/Levene）。本案例各组正态、自动选用Bartlett，统计量91.31。
\end{algosposs}

\begin{algoresult}
输出"方差齐性检验"表，看Levene（或Bartlett）那一行的\textbf{显著性}：
\begin{itemize}
    \item $p>0.05$：方差齐，可以放心用普通t检验/方差分析
    \item $p<0.05$：方差不齐。本案例Bartlett $p=1.15\times10^{-19}\ll0.05$，方差不齐，灌装4线（标准差6.40）明显是"问题线"
\end{itemize}

\textbf{方差不齐怎么办？}
\begin{itemize}
    \item 两组比较：用\textbf{Welch校正t检验}（SPSS独立样本t检验表里自动给出"不假定等方差"那一行）
    \item 多组比较：用\textbf{Welch校正方差分析}或\textbf{Games-Howell}事后比较
    \item 或直接改用非参数方法（Kruskal-Wallis）
\end{itemize}
\end{algoresult}

\begin{algonotes}
\textbf{什么时候用？}任何要比较多组均值的参数检验（t检验、ANOVA）之前。

\textbf{什么时候不用？}非参数检验（Mann-Whitney、Kruskal-Wallis）不要求方差齐性，可以跳过。

\textbf{常见坑}
\begin{itemize}
    \item \textbf{坑1}：方差不齐还硬用普通ANOVA，导致$p$值失真
    \item \textbf{坑2}：样本量悬殊时方差不齐后果更严重（如一组45个、一组30个），更要警惕
\end{itemize}

\textbf{和相似算法的区别}：它和正态性检验是参数检验前的\textbf{两道并列安检}——一个查波动齐不齐，一个查分布正不正态，两道都过才能上普通t检验/ANOVA。
\end{algonotes}

%% ========== 卡方检验 ==========
\basicalgo{{\starmark}卡方检验}{chi_square}

\begin{algointro}
卡方检验用来判断\textbf{两个分类变量是否有关系}。本案例看"门店业态（社区店/超市/大卖场）"和"支付方式（现金/移动/银行卡）"是不是相互独立——不同店的人付钱习惯真的不一样吗？
\end{algointro}

\begin{algoscene}
\begin{itemize}
    \item 两个都是分类变量，想看它们是否相关
    \item 问卷数据：性别与偏好、地区与选择是否有关
    \item 美赛C题中分析离散类别之间的关联结构
\end{itemize}
\end{algoscene}

\begin{algoprinciple}
先做一个\textbf{交叉表}：行是门店业态，列是支付方式，格子里是人数。如果两变量无关，那么"各店用移动支付的比例"应该\textbf{差不多}。

卡方统计量$\chi^2=\sum\frac{(O-E)^2}{E}$里：
\begin{itemize}
    \item $O$是格子里\textbf{实际}人数
    \item $E$是"假如无关"时\textbf{期望}的人数
    \item 实际和期望差得越远，$\chi^2$越大，$p$越小，越说明两变量\textbf{不独立}
\end{itemize}

本案例$\chi^2=27.76$（自由度$df=4$），$p<0.001$，说明业态和支付方式确实有关。
\end{algoprinciple}

\begin{algosposs}
\begin{enumerate}
    \item 点击菜单 \textbf{分析 → 描述统计 → 交叉表}
    \item 行变量选\textbf{门店业态}，列变量选\textbf{支付方式}
    \item 点击\textbf{统计量}，勾选\textbf{卡方}和\textbf{Phi和Cramér's V}
    \item 点击\textbf{单元格}，勾选\textbf{实测}、\textbf{期望}、\textbf{实测/期望/行}百分比
    \item 点击\textbf{确定}
\end{enumerate}

\textbf{默认参数参考}：独立性检验，效应量用Cramér's V，不做连续性校正；期望频数不足时自动处理（小表用Fisher精确检验）；事后两两比较用Bonferroni校正。
\end{algosposs}

\begin{algoresult}
\begin{itemize}
    \item \textbf{卡方检验表}：看Pearson卡方那行的$p$值。本案例$\chi^2=27.76$，$df=4$，$p<0.001$，两变量显著相关
    \item \textbf{Cramér's V}：衡量关联强度，本案例0.304，属\textbf{中等效应}（0.1以下弱、0.3中等、0.5以上强）。注意：$p$只说明相关存不存在，效应量才说明关系大不大
    \item \textbf{调整后残差}：绝对值$\geq 2$的格子是差异主要来源。本案例大卖场用银行卡特别多（残差+4.08）、用移动支付特别少（残差$-4.74$）
    \item \textbf{行百分比}：大卖场52.6\%用银行卡，社区店54.5\%用移动支付——差异一目了然
\end{itemize}
\end{algoresult}

\begin{algonotes}
\textbf{什么时候用？}两个无序分类变量，想知道是否独立。

\textbf{什么时候不用？}
\begin{itemize}
    \item 期望频数低于5的格子超过20\%时，卡方近似不准，要改用Fisher精确检验
    \item 同一批人被重复测量（配对设计）时，普通卡方不适用，用McNemar/Cochran's Q
\end{itemize}

\textbf{常见坑}
\begin{itemize}
    \item \textbf{坑1}：把$p<0.05$当成"关系很强"。大样本下很弱的关系也显著，必须看Cramér's V
    \item \textbf{坑2}：相关不等于因果，也不代表方向
\end{itemize}

\textbf{和相似算法的区别}
\begin{itemize}
    \item vs \textbf{列联分析}：列联分析侧重交叉表描述，这里侧重卡方检验本身
    \item vs \textbf{相关分析}：相关分析用于连续变量，卡方用于分类变量
\end{itemize}
\end{algonotes}

%% ========== 独立样本T检验 ==========
\basicalgo{{\starmark}独立样本T检验}{indep_ttest}

\begin{algointro}
独立样本t检验比较\textbf{两个互不搭界的群体}的均值是否真有差别。本案例比较白班和夜班拣货员的人均拣货效率，看差那1.8件/小时是真差距还是运气。
\end{algointro}

\begin{algoscene}
\begin{itemize}
    \item 两组独立样本（如实验组/对照组、男生/女生、两个城市）比较均值
    \item 美赛A/B验证题：新方案组和旧方案组的关键指标差异
    \item 论文"组间差异检验"的标配
\end{itemize}
\end{algoscene}

\begin{algoprinciple}
夜班均值19.80，白班21.65，差1.84。但两组人本身就有波动，这点差距会不会只是抽样碰巧？

t统计量本质是：\textbf{均值差 ÷ 抽样误差}。
\begin{itemize}
    \item 这个比值越大（$|t|$越大），说明差异相对于随机波动越突出
    \item 本案例$t=-5.18$，$df=148$，$p<0.0001$，说明白班确实显著更快
\end{itemize}

\textbf{效应量Cohen's $d$}：$d=-0.85$。$|d|\approx 0.2$小效应、0.5中、0.8大。$p$受样本量影响，效应量才告诉你"差距值不值得关心"。
\end{algoprinciple}

\begin{algosposs}
\begin{enumerate}
    \item 点击菜单 \textbf{分析 → 比较均值 → 独立样本T检验}
    \item 检验变量选\textbf{拣货效率}，分组变量选\textbf{班组}
    \item 点击\textbf{定义组}，输入组1、组2的取值（如1=夜班，2=白班）
    \item 点击\textbf{选项}，置信区间填95\%
    \item 点击\textbf{确定}
\end{enumerate}

\textbf{默认参数参考}：$\alpha=0.05$，双侧检验，效应量Cohen's d，正态性Shapiro，方差齐性Levene，离群值按IQR报告不删除。
\end{algosposs}

\begin{algoresult}
输出两张表：
\begin{itemize}
    \item \textbf{组统计量表}：两组各$n=75$，夜班均值19.80、白班21.65
    \item \textbf{独立样本检验表}是核心，分两行：
    \begin{itemize}
        \item 先看\textbf{Levene方差齐性}：本案例$p=0.575>0.05$，方差齐，所以看\textbf{第一行"假设方差相等"}
        \item 若Levene $p<0.05$，改看第二行"不假设方差相等"（Welch校正）
        \item 看$t$那行的\textbf{显著性（双侧）}：本案例$p<0.0001$，差异显著
        \item \textbf{均值差}：$-1.84$（夜班$-$白班），95\%置信区间$[-2.43,\ -1.25]$不含0，与显著结论一致
    \end{itemize}
\end{itemize}
\end{algoresult}

\begin{algonotes}
\textbf{什么时候用？}两个独立组、一个连续因变量，比均值。

\textbf{什么时候不用？}
\begin{itemize}
    \item 三组及以上比较：用方差分析，别反复做两两t检验（会放大假阳性）
    \item 同一组人前后两次测量：用配对样本t检验
    \item 数据严重非正态：用Mann-Whitney U检验
\end{itemize}

\textbf{常见坑}
\begin{itemize}
    \item \textbf{坑1}：不看Levene就直接读第一行结果。方差不齐时第一行是错的
    \item \textbf{坑2}：多组两两分别做t检验，$p$值失真——必须做多重比较校正
\end{itemize}

\textbf{和相似算法的区别}
\begin{itemize}
    \item vs \textbf{配对样本t检验}：独立=两组不同人；配对=同一批人前后两次
    \item vs \textbf{方差分析}：两组用t检验，三组及以上用ANOVA
\end{itemize}
\end{algonotes}

%% ========== 单样本T检验 ==========
\basicalgo{单样本T检验}{one_sample_ttest}

\begin{algointro}
单样本t检验把\textbf{一组数据的均值}和\textbf{一个已知标准值}比。本案例：瓶身标500mL，当天抽测均值497.3mL，到底是灌装线偏了，还是抽测碰巧？
\end{algointro}

\begin{algoscene}
\begin{itemize}
    \item 产品实际指标 vs 标称值/国家标准，判断是否达标
    \item 某群体均值是否等于某个理论值（如全国平均水平）
    \item 美赛中验证模型输出是否与设定目标一致
\end{itemize}
\end{algoscene}

\begin{algoprinciple}
抽测150瓶均值497.3mL，比标称500少2.7mL。但抽样总有波动，这2.7mL是"真偏了"还是"运气"？

t统计量=$(\bar{x}-\mu_0)/$标准误，即\textbf{样本均值离标准值有多远，用抽样误差做单位来量}。
\begin{itemize}
    \item 本案例$t=-8.60$，$df=149$，$p<0.001$，远小于0.05
    \item 95\%置信区间$[496.7,\ 497.9]$，标准值500\textbf{没落在区间里}，再次印证偏了
\end{itemize}

Cohen's $d=-0.70$，中等效应——偏差在统计上显著，实际也值得调校。
\end{algoprinciple}

\begin{algosposs}
\begin{enumerate}
    \item 点击菜单 \textbf{分析 → 比较均值 → 单样本T检验}
    \item 检验变量选\textbf{净含量 net\_content\_ml}
    \item 在\textbf{检验值}框输入标准值\textbf{500}
    \item 点击\textbf{选项}，置信区间95\%
    \item 点击\textbf{确定}
\end{enumerate}

\textbf{默认参数参考}：双侧检验（均值$\neq$标准值），正态性自动检测（Shapiro），离群值先核查不删。本案例$n=150$，样本均值497.3，标准差3.82。
\end{algosposs}

\begin{algoresult}
\begin{itemize}
    \item \textbf{单样本统计量表}：$n=150$，均值497.3，标准差3.82，标准误0.31
    \item \textbf{单样本检验表}：看\textbf{检验值=500}那行
    \begin{itemize}
        \item $t=-8.60$，$df=149$，双侧$p<0.001$——显著不等于500
        \item \textbf{均值差值}$=-2.68$
        \item 95\%置信区间$[496.7,\ 497.9]$：因为不含500，所以拒绝"均值=500"
    \end{itemize}
\end{itemize}

\textbf{判断}：$p<0.05$且置信区间不含检验值，就说明样本均值和标准值有真实差距。
\end{algoresult}

\begin{algonotes}
\textbf{什么时候用？}手里一组数据，想跟一个既定标准/目标/历史均值比。

\textbf{什么时候不用？}
\begin{itemize}
    \item 小样本且严重非正态时，改用单样本Wilcoxon符号秩检验
    \item 要比的是两个组（不是一个标准值），用独立样本t检验
\end{itemize}

\textbf{常见坑}
\begin{itemize}
    \item \textbf{坑1}：只看$p$不看效应量。大样本下差0.1mL也$p<0.05$，但实际无所谓
    \item \textbf{坑2}：忘记填检验值，默认填成0会得到莫名其妙的结果
\end{itemize}

\textbf{和相似算法的区别}：单样本=一组比标准值；独立样本t=两组互相比；配对t=同一组前后比。
\end{algonotes}

%% ========== 配对样本T检验 ==========
\basicalgo{配对样本T检验}{paired_ttest}

\begin{algointro}
配对样本t检验比的是\textbf{同一批人（或同一对象）前后两次测量}的差值是否不为零。本案例：同一批前台员工培训前70.9分、培训后75.2分，培训到底有没有用？
\end{algointro}

\begin{algoscene}
\begin{itemize}
    \item 同一组对象"前测--后测"，看干预（培训/用药/新流程）是否有效
    \item 配对设计实验：同一个体左右眼、前后两次测量
    \item 美赛中评估政策/方案实施前后指标变化
\end{itemize}
\end{algoscene}

\begin{algoprinciple}
直接拿后测均值75.2和前测70.9比，混进了"人与人本来就不同"的噪声。配对t的聪明之处是\textbf{先给每个人算差值}（后$-$前），再看这堆差值的均值是不是显著不为0。

\begin{itemize}
    \item 本案例差值均值$=4.37$，差值标准差6.81
    \item $t=7.87$，$df=149$，$p<0.001$，培训确实有效
    \item 95\%置信区间$[3.28,\ 5.47]$不含0；Cohen's $d_z=0.64$，中等效应
\end{itemize}

\textbf{注意}：要检验的是\textbf{差值}是否正态，不是前后两次单独正态。本案例差值Shapiro $p=0.985$，通过。前后测相关0.781——相关越高，配对设计越划算。
\end{algoprinciple}

\begin{algosposs}
\begin{enumerate}
    \item 数据整理成一行一对：一列\textbf{培训前成绩}，一列\textbf{培训后成绩}
    \item 点击菜单 \textbf{分析 → 比较均值 → 配对样本T检验}
    \item 同时选中两列，作为\textbf{成对变量}Pair 1
    \item 点击\textbf{选项}，置信区间95\%
    \item 点击\textbf{确定}
\end{enumerate}

\textbf{默认参数参考}：双侧检验，效应量Cohen's $d_z$，差值正态性Shapiro；差值非正态且样本小时自动切换Wilcoxon符号秩。
\end{algosposs}

\begin{algoresult}
\begin{itemize}
    \item \textbf{配对样本统计量}：前测均值70.87，后测75.24
    \item \textbf{配对样本相关系数}：$r=0.781$，前后强相关（合理，因为是同一批人）
    \item \textbf{配对样本检验}是核心：
    \begin{itemize}
        \item 差值均值$=4.37$，$t=7.87$，$p<0.001$——后测显著更高
        \item 差值95\%置信区间$[3.28,\ 5.47]$不含0
    \end{itemize}
\end{itemize}
\end{algoresult}

\begin{algonotes}
\textbf{什么时候用？}同一对象前后两次（或配对对象）测量，看差值。

\textbf{什么时候不用？}
\begin{itemize}
    \item 两组是不同的人、互不配对：用独立样本t检验
    \item 差值严重非正态、样本又小：改用配对Wilcoxon符号秩检验
\end{itemize}

\textbf{常见坑}
\begin{itemize}
    \item \textbf{坑1}：把配对数据当成独立两组做独立样本t检验，浪费了配对信息、降低检验功效
    \item \textbf{坑2}：前后测之间隔了太久，混入其他干扰因素（如同一批人还接受了别的培训）
\end{itemize}

\textbf{和相似算法的区别}：它是"控制了个体差异"的t检验——因为每个人都和自己比，自动消除了人本身的底子差异。
\end{algonotes}

%% ========== 方差分析 ==========
\basicalgo{{\starmark}方差分析（单因素/多因素）}{anova}

\begin{algointro}
方差分析（ANOVA）比的是\textbf{三个及以上组}的均值是否全相等，还能看多个因素\textbf{各自的影响}以及\textbf{因素之间的搭配（交互作用）}。本案例研究商圈类型和促销方案如何影响门店日均营业额。
\end{algointro}

\begin{algoscene}
\begin{itemize}
    \item 三组及以上组间均值比较（如三种促销、三个地区、四种配方）
    \item 想看两个因素各自及搭配效果（如"促销效果是否因商圈而异"）
    \item 美赛C题中多因素对结果的影响分解
\end{itemize}
\end{algoscene}

\begin{algoprinciple}
方差分析名字叫"方差"，比的却是均值——它的思路是：\textbf{把总的波动拆成"组间差异"和"组内随机波动"两部分}。

\begin{itemize}
    \item $F$统计量 = 组间差异 ÷ 组内差异
    \item $F$越大（$p$越小），越说明各组均值真不一样，而不是随机波动
    \item 本案例：商圈类型$F=83.7$，促销方案$F=27.5$，两者交互$F=10.35$，都$p<0.001$
\end{itemize}

\textbf{交互作用}是关键概念：促销方案的效果\textbf{因商圈不同而不同}就叫有交互。比如"满减"在购物中心最赚钱、在社区店一般——这时就不能笼统说"哪种促销最好"，得说"在哪种商圈配哪种促销"。
\end{algoprinciple}

\begin{algosposs}
\begin{enumerate}
    \item 单因素ANOVA：\textbf{分析 → 比较均值 → 单因素ANOVA}，因变量进"因变量列表"，分组进"因子"；点\textbf{事后多重比较}勾Tukey
    \item 多因素/含交互：\textbf{分析 → 一般线性模型 → 单变量}
    \item 因变量选\textbf{日均营业额}，固定因子选\textbf{商圈类型}、\textbf{促销方案}
    \item 点\textbf{模型}选"全因子"（自动含主效应和交互项）
    \item 点\textbf{事后比较}选商圈、促销，勾\textbf{Tukey HSD}
    \item 点\textbf{选项}勾"描述统计""方差齐性检验""效应量估计"
    \item 点击\textbf{确定}
\end{enumerate}

\textbf{默认参数参考}：$\alpha=0.05$，II型平方和（非平衡设计稳健），效应量eta$^2$与偏eta$^2$，方差齐性用Levene，非参数灵敏度用Kruskal-Wallis。
\end{algosposs}

\begin{algoresult}
\begin{itemize}
    \item \textbf{方差齐性Levene}：本案例$p=0.202>0.05$，方差齐，可用普通ANOVA
    \item \textbf{主体间效应检验表}是核心：
    \begin{itemize}
        \item 商圈类型$F=83.7$，$p<0.001$，偏$\eta^2=0.543$（商圈解释了54\%的营业额变异，影响很大）
        \item 促销方案$F=27.5$，$p<0.001$，偏$\eta^2=0.280$
        \item \textbf{交互项}$F=10.35$，$p<0.001$——存在交互，必须结合单元格均值解读
    \end{itemize}
    \item \textbf{事后Tukey HSD}：具体哪两组不同。本案例"写字楼+会员专享"均值13.45最高，显著高于多数组合
    \item \textbf{怎么看}：主效应$p<0.05$说明该因素重要；交互$p<0.05$说明不能孤立看主效应，要画交互图看"哪组搭配最好"
\end{itemize}
\end{algoresult}

\begin{algonotes}
\textbf{什么时候用？}三组及以上均值比较，或多因素影响分解。

\textbf{什么时候不用？}
\begin{itemize}
    \item 只有两组：用t检验即可
    \item 方差不齐且样本悬殊：用Welch-ANOVA或Kruskal-Wallis
    \item 数据严重非正态：改用Kruskal-Wallis（非参数版ANOVA）
\end{itemize}

\textbf{常见坑}
\begin{itemize}
    \item \textbf{坑1}：ANOVA整体显著就完事。必须做事后多重比较，才知道\textbf{到底哪两组不同}
    \item \textbf{坑2}：忽略交互作用。交互显著时直接解释主效应会闹笑话
    \item \textbf{坑3}：多组反复做两两t检验，人为放大假阳性
\end{itemize}

\textbf{和相似算法的区别}：t检验只能比两组；ANOVA能一次比多组还能拆因素和交互。非正态的"ANOVA替代"是Kruskal-Wallis。
\end{algonotes}

\section{进阶级算法速览}

%% ========== Kruskal-Wallis ==========
\advancedalgo{Kruskal-Wallis检验（多样本秩和）}{kruskal_wallis}

\begin{algobrief}
\textbf{一句话}：方差分析的\textbf{非参数版}——比较三组及以上独立样本的\textbf{分布/中位数}是否不同，不要求正态。

\textbf{美赛场景区}：数据偏态严重（如耗时、收入、评分）时多组比较；如不同分拨中心的单件分拣耗时差异。

\textbf{核心原理}：把所有数据混在一起排序编秩，比较各组的平均秩是否差很多。$H$统计量近似卡方分布。本案例分组变量为分拨中心、指标为单件分拣耗时秒数；$p<0.05$说明各组耗时分布不全相同。

\textbf{操作要点}：\textbf{分析 → 非参数检验 → 独立样本}，选Kruskal-Wallis；效应量用$\varepsilon^2$；显著后用Bonferroni校正做两两比较。

\textbf{结果关键}：看Kruskal-Wallis的$p$值；显著后看各组平均秩，秩越高组数值越大。

\textbf{注意}：它是"升级版Mann-Whitney"。两组时它就等价于Mann-Whitney U；正态且方差齐时，普通ANOVA更有检验功效。
\end{algobrief}

%% ========== 配对Wilcoxon ==========
\advancedalgo{配对样本Wilcoxon符号秩检验}{wilcoxon_paired}

\begin{algobrief}
\textbf{一句话}：配对样本t检验的\textbf{非参数版}——同一批人前后两次测量，但差值严重非正态时用它。

\textbf{美赛场景区}：前后测差值明显偏态（如治疗前后的康复评分）；小样本+非正态配对数据。

\textbf{核心原理}：不算差值本身，而是给差值的绝对值编秩，再按差值正负号分组，比较正负秩和是否平衡。本案例判断干预后指标是否显著变化。

\textbf{操作要点}：\textbf{分析 → 非参数检验 → 相关样本}，选Wilcoxon符号秩；双侧检验，效应量$r$。

\textbf{结果关键}：$p<0.05$说明前后测量有显著的系统性升降；结合差值中位数判断方向。

\textbf{注意}：当配对差值近似正态时，配对t检验更灵敏；它是差值非正态时的稳健替代，不是首选。
\end{algobrief}

%% ========== Z检验 ==========
\advancedalgo{Z检验（均值/比率）}{z_test}

\begin{algobrief}
\textbf{一句话}：大样本下用正态分布做的均值或比率检验——样本量大到$n\geq30$时，t分布近似正态，可用更简单的$Z$检验。

\textbf{美赛场景区}：大样本下比较两总体比率（如两城市移动支付使用率）；比例类指标的假设检验。

\textbf{核心原理}：大样本时统计量近似标准正态$N(0,1)$，$Z=\frac{\text{样本统计量}-\text{参数}}{\text{标准误}}$。$|Z|>1.96$对应$\alpha=0.05$双侧显著。

\textbf{操作要点}：工具箱选"Z检验（均值/比率）"，明确是单样本还是两样本、均值还是比率；$\alpha=0.05$。

\textbf{结果关键}：$|Z|>1.96$（双侧0.05）或$>2.58$（0.01）即显著；同时报告置信区间。

\textbf{注意}：小样本仍用t检验。Z检验和t检验结论在大样本时几乎一致，区别主要在小样本。
\end{algobrief}

%% ========== Friedman ==========
\advancedalgo{Friedman检验}{friedman}

\begin{algobrief}
\textbf{一句话}：随机区组/重复测量设计下的\textbf{非参数ANOVA}——同一批对象接受\textbf{三个及以上}条件，比较是否有差异。

\textbf{美赛场景区}：同一组被试在多种方案/多个时间点下的评分比较；如五位评委对四款产品的排序一致性检验。

\textbf{核心原理}：在\textbf{每个对象内部}给各条件编秩，再汇总各条件的秩和。本案例为三水平及以上的重复测量设计，$p<0.05$说明各条件分布不全相同。

\textbf{操作要点}：\textbf{分析 → 非参数检验 → 相关样本}，选Friedman；显著后做事后两两比较。

\textbf{结果关键}：Friedman卡方统计量的$p$值；显著后用Wilcoxon两两比较并Bonferroni校正。

\textbf{注意}：它是\textbf{配对版Kruskal-Wallis}。两个配对条件时用Wilcoxon符号秩；三个及以上配对条件用Friedman。
\end{algobrief}

%% ========== MANOVA ==========
\advancedalgo{多元方差分析（MANOVA）}{manova}

\begin{algobrief}
\textbf{一句话}：同时考察一个（多个）分组因素对\textbf{多个相关因变量}的整体影响。

\textbf{美赛场景区}：比较不同组在"销量+满意度+复购率"等多个相关指标上的综合差异；避免多个单变量ANOVA放大假阳性。

\textbf{核心原理}：把多个因变量当一个整体向量，用Wilks' Lambda等统计量检验组间质心是否不同。本案例输出Wilks' $\Lambda$、$F$近似和$p$值。

\textbf{操作要点}：\textbf{分析 → 一般线性模型 → 多变量}，多个因变量同时入框，固定因子放分组变量。

\textbf{结果关键}：先看多变量检验表的Wilks' $\Lambda$及其$p$；显著后再回头看各单变量ANOVA表，定位差异来自哪个因变量。

\textbf{注意}：要求因变量多元正态、协方差阵齐性（Box's M检验）。它是"一揽子"检验，显著了还得拆到单变量层面解释。
\end{algobrief}

%% ========== Mood中位数 ==========
\advancedalgo{Mood中位数检验}{mood_median}

\begin{algobrief}
\textbf{一句话}：最稳健的多组中位数比较——只看每组落在总中位数之上/之下的个数，用卡方判断各组中位数是否相同。

\textbf{美赛场景区}：数据有极端值、严重偏态，连Kruskal-Wallis都嫌不稳时的兜底方法。

\textbf{核心原理}：求所有数据的共同中位数，数每组有多少个点在它上面/下面，列成列联表做卡方检验。

\textbf{操作要点}：工具箱选"Mood中位数检验"，指定分组变量和数值变量；$\alpha=0.05$。

\textbf{结果关键}：卡方$p<0.05$说明各组中位数不全相等；它检验的是中位数位置，不是分布形状。

\textbf{注意}：检验功效低（浪费了秩信息），是\textbf{最保守}的方法——只有数据特别脏时才用，优先Kruskal-Wallis。
\end{algobrief}

%% ========== 协方差分析 ==========
\advancedalgo{协方差分析（ANCOVA）}{ancova}

\begin{algobrief}
\textbf{一句话}：方差分析+回归的混血——比较组间均值差异时，\textbf{顺手扣掉一个连续干扰变量（协变量）}的影响。

\textbf{美赛场景区}：比较各班成绩时扣掉"入学基础分"；比较各方案效果时扣掉"初始水平"。

\textbf{核心原理}：先把因变量对协变量做回归，用\textbf{校正后}的组均值再做ANOVA。本案例在比较组别差异前，先消除协变量造成的基线不平衡。

\textbf{操作要点}：\textbf{分析 → 一般线性模型 → 单变量}，固定因子放分组，\textbf{协变量}框放干扰变量。

\textbf{结果关键}：先看"组别"效应校正后是否仍显著；若校正后差异消失，说明原差异主要来自协变量而非分组本身。

\textbf{注意}：前提是\textbf{各组与协变量回归斜率平行}（无交互），否则不能直接用。
\end{algobrief}

%% ========== 游程检验 ==========
\advancedalgo{游程检验}{runs_test}

\begin{algobrief}
\textbf{一句话}：判断一串按顺序排列的数据是不是\textbf{随机出现}的——有没有人为规律或趋势。

\textbf{美赛场景区}：检验抽样是否随机；检验残差序列是否随机（回归假设）；判断比赛/股票涨跌是否随机游走。

\textbf{核心原理}：以中位数（或某阈值）为界把数据分正负，数连续同号的"游程"个数。游程太少=有聚类/趋势，太多=有周期振荡，都说明不随机。

\textbf{操作要点}：\textbf{分析 → 非参数检验 → 游程检验}，选检验变量，临界值选"中位数"。

\textbf{结果关键}：$p>0.05$说明序列随机；$p<0.05$说明存在系统性模式。

\textbf{注意}：它检验的是\textbf{随机性}，不是均值差异——别和单样本t检验搞混。
\end{algobrief}

%% ========== 事后多重比较 ==========
\advancedalgo{事后多重比较（Post Hoc）}{posthoc}

\begin{algobrief}
\textbf{一句话}：ANOVA告诉你"至少两组不同"之后，再\textbf{逐对比较}找出到底是哪几对不同，同时控制假阳性。

\textbf{美赛场景区}：三四种促销/配方/地区之间两两找差异；是方差分析的"后半程"必备。

\textbf{核心原理}：两两比较会累积假阳性（3组就3次比较、假阳性率涨到14\%），事后比较用Tukey HSD、Bonferroni等方法给$p$值"打折"，把总体错误率压回0.05。

\textbf{操作要点}：单因素ANOVA对话框点\textbf{事后多重比较}，方差齐选\textbf{Tukey}，不齐选\textbf{Games-Howell}。

\textbf{结果关键}："多重比较表"中标注*或校正$p<0.05$的组别对，就是真正有差异的对；同列字母法（a/b/c）更直观。

\textbf{注意}：\textbf{必须在ANOVA整体显著之后再做}。整体不显著却硬做两两比较是"钓鱼"。
\end{algobrief}

%% ========== Mann-Whitney U ==========
\advancedalgo{独立样本Mann-Whitney U检验}{mannwhitney_u}

\begin{algobrief}
\textbf{一句话}：独立样本t检验的\textbf{非参数版}——两组独立样本，但数据非正态时，比两组分布/位置是否不同。

\textbf{美赛场景区}：两组偏态数据（如金额、耗时、严重度评分）的比较；不满足t检验正态前提时。

\textbf{核心原理}：把两组数据混在一起编秩，比较两组的秩和。$U$统计量对应的$p<0.05$说明两组分布位置不同。

\textbf{操作要点}：\textbf{分析 → 非参数检验 → 独立样本}，选Mann-Whitney U；两组按分组变量定义。

\textbf{结果关键}：看Mann-Whitney U（或Z）的渐近显著性；结合两组平均秩判断哪组整体更高。

\textbf{注意}：两组独立+非正态时用它；三组及以上用Kruskal-Wallis。数据正态时t检验更灵敏。
\end{algobrief}

%% ========== 等价性检验 ==========
\advancedalgo{等价性检验（TOST）}{equivalence_test}

\begin{algobrief}
\textbf{一句话}：反过来证明"两组\textbf{没有实质差别}"——普通t检验不显著只能说"没检出差异"，等价检验才能说"确实差不多"。

\textbf{美赛场景区}：证明新方案不比老方案差（等效）；证明两测量方法结果一致；环保/医药的"非劣效"论证。

\textbf{核心原理}：预先设定一个"可接受差异边界$\pm\Delta$"，做两次单侧检验（TOST）：只有当均值差的置信区间\textbf{整个落在$[-\Delta,\Delta]$内}，才判定等价。

\textbf{操作要点}：工具箱选"等价性检验"，设定等价界值$\Delta$（结合业务意义，不能随便定）。

\textbf{结果关键}：均值差的90\%置信区间完全落在等价带内，才结论"等价"；若区间跨出边界，则不能宣称等效。

\textbf{注意}：它和普通检验逻辑相反——\textbf{显著才算数}。$p<0.05$才说明等价，不是$p>0.05$。
\end{algobrief}

%% ========== 单样本Wilcoxon ==========
\advancedalgo{单样本Wilcoxon符号秩检验}{wilcoxon_onesample}

\begin{algobrief}
\textbf{一句话}：单样本t检验的\textbf{非参数版}——一组数据的中位数是否等于某个标准值，但数据非正态时用。

\textbf{美赛场景区}：小样本+偏态数据（如收入、反应时间）检验中位数是否达标/等于目标。

\textbf{核心原理}：计算每个观测与标准值之差，给差值绝对值编秩，按正负号分两组比较秩和。

\textbf{操作要点}：\textbf{分析 → 非参数检验 → 单样本}，选Wilcoxon符号秩，输入检验值。

\textbf{结果关键}：$p<0.05$说明样本中位数与标准值有显著差异；结合正/负秩方向判断偏高还是偏低。

\textbf{注意}：它检验的是\textbf{中位数/分布位置}，不是均值。数据近似正态时，单样本t检验更有功效。
\end{algobrief}

% 第4章 回归分析
\chapter{回归分析}

\begin{chapteroverview}
\textbf{这个分类是干什么的？}

回归分析是\textbf{建模和预测}的主力工具。它要回答的是：给定一堆自变量（客流、面积、促销天数……），因变量（销售额）怎么变？每个自变量\textbf{影响多大、方向是正是负}？能不能用这个方程去\textbf{预测}新店的销售额？相关分析只问"动不动一起变"，回归则进一步给出"变多少"的定量方程。

\textbf{美赛什么题型常用？}

回归是美赛\textbf{出镜率最高}的建模方法。C题几乎必做：用回归找关键影响因素、做预测；A/B题建立变量间的定量关系模型。论文"模型建立""结果分析"部分，回归系数表是核心证据。

\textbf{学习路径建议}
先掌握\textbf{线性回归}（一切的基础）和它的体检项\textbf{共线性VIF}；再学\textbf{曲线回归}（关系不是直线时）、\textbf{分层回归}（控制背景后看增量贡献）、\textbf{逐步回归}（变量太多时自动筛选）；因变量是\textbf{分类}时用\textbf{逻辑回归}。最后再根据数据特点认识岭回归、Lasso、Poisson、Logistic族等进阶模型。
\end{chapteroverview}

\section{基础级算法}

%% ========== 线性回归 ==========
\basicalgo{{\starmark}线性回归（最小二乘法）}{linear_reg}

\begin{algointro}
线性回归就是\textbf{找一条最佳直线（其实是多维平面）}，让它尽量穿过所有数据点，从而用自变量预测因变量。本案例用日均客流量、营业面积、月促销天数、竞争门店数预测月销售额。
\end{algointro}

\begin{algoscene}
\begin{itemize}
    \item 因变量是连续数值（销售额、温度、成绩），想找多个自变量对它的线性影响
    \item 美赛C题建立预测模型、量化各因素边际效应
    \item 论文中"影响因素分析"的核心模型
\end{itemize}
\end{algoscene}

\begin{algoprinciple}
想象散点图上一堆点，你画一条直线去"贴近"它们。怎么算最佳？让每个点到直线的\textbf{纵向距离平方和}最小——这就是"最小二乘法"。

回归方程长这样：
$$y = b_0 + b_1 x_1 + b_2 x_2 + \cdots + b_k x_k$$
\begin{itemize}
    \item $b_0$是截距（所有自变量为0时的基线值）
    \item $b_j$是$x_j$的系数：$x_j$每涨1单位，$y$平均涨$b_j$单位（控制其他变量不变）
    \item 本案例：竞争门店数系数$-2.025$，即周边每多一家竞争店，月销售额平均降2.025万元；月促销天数系数$+1.149$，每多促销1天销售额涨1.149万元
\end{itemize}

\textbf{$R^2$}是拟合优度：本案例0.788，说明四个自变量一起解释了78.8\%的销售额波动。
\end{algoprinciple}

\begin{algosposs}
\begin{enumerate}
    \item 点击菜单 \textbf{分析 → 回归 → 线性}
    \item 因变量选\textbf{月销售额}，自变量选\textbf{日均客流量、营业面积、月促销天数、竞争门店数}
    \item 方法选\textbf{输入}（Enter，全部强行进）
    \item 点击\textbf{统计量}，勾选"估计""置信区间""模型拟合度""共线性诊断""Durbin-Watson"
    \item 点击\textbf{绘制}，把\verb|*ZRESID|作$y$、\verb|*ZPRED|作$x$，画残差图查异方差
    \item 点击\textbf{确定}
\end{enumerate}

\textbf{默认参数参考}：缺失按删除（drop），含截距，协方差类型自动，VIF阈值10。本案例有效样本147行，$R^2=0.788$。
\end{algosposs}

\begin{algoresult}
\begin{itemize}
    \item \textbf{模型汇总表}：$R^2=0.788$、调整$R^2=0.782$；Durbin-Watson=1.84（接近2，残差无自相关）
    \item \textbf{ANOVA表}：$F=131.9$，$p=8.55\times10^{-47}$，模型整体显著
    \item \textbf{系数表}是核心：每个自变量一行
    \begin{itemize}
        \item \textbf{系数$B$}：方向和大小。客流$+0.02$、面积$+0.07$、促销天数$+1.15$、竞争店$-2.03$
        \item \textbf{$t$和$p$}：本案例4个变量全部$p<0.001$，都显著
        \item \textbf{95\%置信区间}：不含0即显著
    \end{itemize}
    \item \textbf{诊断前提}：本案例VIF最大仅1.03（无共线性），Breusch-Pagan $p=0.983$（无异方差），残差Shapiro $p=0.49$（正态）——前提全满足，结果可信
\end{itemize}
\end{algoresult}

\begin{algonotes}
\textbf{什么时候用？}连续因变量+多个连续/分类自变量，关系近似直线。

\textbf{什么时候不用？}
\begin{itemize}
    \item 因变量是0/1（是否逾期）：用逻辑回归
    \item 因变量是计数（次数）：用Poisson/负二项回归
    \item 散点明显弯曲：用曲线回归或加二次项
\end{itemize}

\textbf{常见坑}
\begin{itemize}
    \item \textbf{坑1}：不做前提诊断就报系数。残差异方差、自相关、共线性都会让$p$值失真
    \item \textbf{坑2}：把回归关系当因果。控制了变量也只能说"关联"，不一定是因果
    \item \textbf{坑3}：用回归方程外推到数据范围之外，结果不可靠
\end{itemize}

\textbf{和相似算法的区别}：它是回归家族的"母版"——后面所有回归都是它在不同因变量/不同数据形态下的变体。
\end{algonotes}

%% ========== 共线性 VIF ==========
\basicalgo{共线性分析（VIF诊断）}{vif_diag}

\begin{algointro}
共线性检查就是看\textbf{自变量之间是不是太像了}——如果两个自变量几乎在说同一件事，回归就分不清功劳该记给谁，系数会乱跳。本案例在做房价回归前，先查面积类变量是否互相重复。
\end{algointro}

\begin{algoscene}
\begin{itemize}
    \item 多元线性/逻辑回归\textbf{建模前}的必做体检
    \item 自变量多且可能相关（如建筑面积/套内面积/使用面积）
    \item 论文中报告VIF以证明模型稳定
\end{itemize}
\end{algoscene}

\begin{algoprinciple}
想象让两个特别像的人（套内面积和建筑面积）一起解释房价。既然它俩几乎同步变大，回归就搞不清"到底是谁把房价抬上去的"，于是把系数在两者间胡乱分配——标准误被放大、显著性忽高忽低、符号甚至反过来。

\begin{itemize}
    \item \textbf{VIF（方差膨胀因子）}：$VIF_j=\frac{1}{1-R_j^2}$，其中$R_j^2$是"用其他所有自变量预测$x_j$"得到的
    \item $VIF=1$：完全不共线；$VIF>5$：要警惕；$VIF>10$：严重共线性，必须处理
    \item 本案例建筑面积$VIF=150$、套内面积$VIF=356$——爆表！因为三者$r=1.00$几乎完全重复
\end{itemize}

\textbf{条件指数}$>30$也提示整个自变量系统病态（本案例44.1）。
\end{algoprinciple}

\begin{algosposs}
\begin{enumerate}
    \item 线性回归对话框（\textbf{分析 → 回归 → 线性}）中放好因变量和全部候选自变量
    \item 点击\textbf{统计量}，勾选\textbf{共线性诊断}
    \item 运行后查看系数表中的"容差"和"VIF"列
    \item 也可在Dabbit AI SPSS工具箱单独跑"共线性分析"，自动逐步剔除
\end{enumerate}

\textbf{默认参数参考}：预警阈值$VIF=5$，严重阈值$VIF=10$；处理方式stepwise\_remove（逐步剔除最大VIF）；缺失整行删除。
\end{algosposs}

\begin{algoresult}
\begin{itemize}
    \item \textbf{初始诊断表}：建筑面积VIF=150、套内面积=356、使用面积=252，全是"严重"
    \item \textbf{处置}：逐步剔除——先踢套内面积（VIF=356），再踢建筑面积（VIF=104）
    \item \textbf{复检}：保留变量最大VIF降到7.92（建筑层数），低于严重阈值10，共线性基本消除
    \item \textbf{决策}：三选一即可（使用面积/建筑面积/套内面积），保留业务上最易解释的那个，不要全塞进去
\end{itemize}
\end{algoresult}

\begin{algonotes}
\textbf{什么时候用？}只要做含多个自变量的回归，建模前必查。

\textbf{什么时候不用？}简单一元回归（只有一个自变量）不存在共线性问题。

\textbf{常见坑}
\begin{itemize}
    \item \textbf{坑1}：VIF高了就删重要变量。优先删\textbf{业务上重复}的，保留解释力强的；或合并成一个综合指标
    \item \textbf{坑2}：VIF高但$R^2$也高、预测效果好时，对预测用途影响不大，主要影响系数解释
\end{itemize}

\textbf{和相似算法的区别}：共线性是\textbf{自变量之间}的问题（VIF查它）；异方差、自相关是\textbf{残差}的问题（各自有专门检验），别混。
\end{algonotes}

%% ========== 曲线回归 ==========
\basicalgo{曲线回归}{curve_reg}

\begin{algointro}
散点图不是直线而是\textbf{弯的}怎么办？曲线回归就是在一堆候选曲线（对数、指数、幂函数、多项式）里挑最贴合的那条。本案例：推广预算越高，销售额涨得越慢——明显是边际递减的弯形。
\end{algointro}

\begin{algoscene}
\begin{itemize}
    \item 散点图显示非线性关系（饱和、递减、加速）
    \item 边际效用递减、S形增长、复利增长等典型曲线现象
    \item 美赛中经济/生物/增长类过程建模
\end{itemize}
\end{algoscene}

\begin{algoprinciple}
直线画不平的地方，换条曲线试试。常见几种：
\begin{itemize}
    \item \textbf{对数曲线}$y=a+b\ln x$：增长越来越慢（边际递减）——本案例正是这种
    \item \textbf{指数曲线}$y=ab^x$：按固定比例爆炸式增长
    \item \textbf{幂函数}$y=ax^b$：齐次增长
    \item \textbf{多项式（二次/三次）}：先升后降或有拐点
\end{itemize}

\textbf{怎么选？}用\textbf{AIC}做裁判：同样本上谁的AIC最小且模型简单，选谁。本案例对数曲线调整$R^2=0.991$、$AIC=141.5$，比直线（$R^2=0.83$、$AIC=584.7$）好一大截，方程$y=17.83+16.11\ln x$。
\end{algoprinciple}

\begin{algosposs}
\begin{enumerate}
    \item 点击菜单 \textbf{分析 → 回归 → 曲线估计}
    \item 因变量选\textbf{销售额}，自变量选\textbf{推广投入}
    \item 勾选候选模型：线性、对数、二次、三次、指数、幂（"全部候选"）
    \item 勾选\textbf{显示ANOVA表格}和"绘制模型图"
    \item 点击\textbf{确定}
\end{enumerate}

\textbf{默认参数参考}：模型选择准则AIC，多项式最高3次，残差诊断开启，置信水平95\%。注意对数/指数/幂函数要求$x>0$、$y>0$。
\end{algosposs}

\begin{algoresult}
\begin{itemize}
    \item \textbf{模型比较表}：每种曲线的$R^2$、$AIC$。本案例对数曲线$R^2=0.991$、$\Delta AIC=-443$（相对直线大幅为负），最优
    \item \textbf{选定方程}：$y=17.83+16.11\ln(x)$，截距和系数$p<0.001$都显著
    \item \textbf{残差诊断}：残差正态$p=0.988$、游程$p=0.101$、异方差$\rho=0.052$——都正常，曲线选得合适
    \item \textbf{解释}：$x$均值处边际效应$dy/dx\approx1.20$，即投入越大、新增一块钱带来的销售增量越小
\end{itemize}
\end{algoresult}

\begin{algonotes}
\textbf{什么时候用？}散点明显弯曲且能说出具体形态时。

\textbf{什么时候不用？}
\begin{itemize}
    \item 散点近似直线：别为了高$R^2$硬上高阶多项式，会过拟合
    \item 关系是多个自变量共同造成的：单变量曲线回归不够，要做多变量非线性回归
\end{itemize}

\textbf{常见坑}
\begin{itemize}
    \item \textbf{坑1}：为追求$R^2$盲目上三次、四次多项式，曲线在两端剧烈摆动（龙格现象）
    \item \textbf{坑2}：超出数据范围外推。本案例结论只在$x\in[1.2,\ 42.85]$内成立
\end{itemize}

\textbf{和相似算法的区别}：它是\textbf{单自变量}的曲线建模；多变量非线性关系用非线性回归或广义可加模型。
\end{algonotes}

%% ========== 分层回归 ==========
\basicalgo{分层回归}{hierarchical_reg}

\begin{algointro}
分层回归就是\textbf{按理论顺序分批次}往模型里加自变量，专门看"新加入的这批变量，能在已有基础上\textbf{额外多解释多少}"。本案例先放员工背景，再放工作资源，最后放团队氛围，逐层看增量。
\end{algointro}

\begin{algoscene}
\begin{itemize}
    \item 想证明"某组变量在控制了其他变量后仍有独立贡献"
    \item 心理学/组织行为学问卷分析：先控制人口学变量，再看核心变量
    \item 美赛中区分"背景因素"和"核心干预因素"谁更重要
\end{itemize}
\end{algoscene}

\begin{algoprinciple}
普通回归一次把所有变量丢进去，分不清谁是谁的功劳。分层回归像\textbf{盖楼}：
\begin{itemize}
    \item \textbf{第1层（基线）}：先放必须控制的背景变量（年龄、司龄、工作量），$R^2=0.414$
    \item \textbf{第2层}：加入工作资源（自主性、上级支持、发展机会），$R^2$涨到0.756，\textbf{增量}$\Delta R^2=0.342$，$\Delta F=66.7$，$p<0.001$——工作资源贡献巨大
    \item \textbf{第3层}：再加入心理安全感，$R^2=0.793$，$\Delta R^2=0.038$，$\Delta F=25.9$，$p<0.001$——仍有显著增量
\end{itemize}

\textbf{关键}：分层顺序是\textbf{按理论定的}，不是按数据$p$值。控制变量先进、核心变量后进，才能干净地看到核心变量的"净增量"。
\end{algoprinciple}

\begin{algosposs}
\begin{enumerate}
    \item \textbf{分析 → 回归 → 线性}，因变量选\textbf{敬业度 job\_satisfaction}
    \item \textbf{Block 1 of 1}放第1层变量：年龄age、司龄tenure、工作量workload，点"下一个"
    \item \textbf{Block 2 of 2}放第2层：工作自主性、上级支持、职业发展
    \item \textbf{Block 3 of 3}放第3层：心理安全感
    \item 点\textbf{统计量}勾"R平方变化""共线性诊断""Durbin-Watson"
    \item 点击\textbf{确定}
\end{enumerate}

\textbf{默认参数参考}：稳健标准误HC3，$\alpha=0.05$，VIF阈值10。
\end{algosposs}

\begin{algoresult}
\begin{itemize}
    \item \textbf{模型汇总表}：重点看\textbf{R平方变化$\Delta R^2$}和对应$\Delta F$的$p$值
    \begin{itemize}
        \item 第1层$R^2=0.414$（基线模型显著）
        \item 第2层$\Delta R^2=0.342$，$\Delta F=66.7$，$p<0.001$——工作资源显著增量
        \item 第3层$\Delta R^2=0.038$，$\Delta F=25.9$，$p<0.001$——心理安全感仍有额外贡献
    \end{itemize}
    \item \textbf{最终系数表}：用HC3稳健标准误。本案例7个变量全部显著，工作自主性$\beta=0.313$效应最大
    \item \textbf{前提诊断}：VIF最大4.31（无共线），Durbin-Watson=1.92（无自相关），残差正态$p=0.753$
\end{itemize}
\end{algoresult}

\begin{algonotes}
\textbf{什么时候用？}要回答"在控制了背景变量后，核心变量是否仍有独立解释力"。

\textbf{什么时候不用？}只是想找最优预测方程、没有明确理论层级时，用逐步回归更合适。

\textbf{常见坑}
\begin{itemize}
    \item \textbf{坑1}：分层顺序乱定。必须按"先外生后内生、先控制后核心"的理论逻辑
    \item \textbf{坑2}：只看最终$R^2$，不看每一层的$\Delta R^2$和$\Delta F$——那就失去了分层的意义
\end{itemize}

\textbf{和相似算法的区别}：它和逐步回归都在选变量，但\textbf{方向相反}——逐步是数据驱动自动选，分层是理论驱动手动分层。
\end{algonotes}

%% ========== 逐步回归 ==========
\basicalgo{逐步回归}{stepwise_reg}

\begin{algointro}
候选自变量太多（十几个甚至几十个），手动选不过来？逐步回归让计算机\textbf{按$p$值自动挑变量}：最有用的先进，没用的踢出去。本案例从12个经营指标里自动挑出5个预测月销售额。
\end{algointro}

\begin{algoscene}
\begin{itemize}
    \item 自变量很多、缺乏明确理论，想自动找出最重要的少数几个
    \item 美赛C题特征筛选、变量降维
    \item 探索性分析：先粗筛出关键变量，再深入建模
\end{itemize}
\end{algoscene}

\begin{algoprinciple}
像\textbf{选秀海选}：
\begin{itemize}
    \item \textbf{向前引入}：一个个试，把$p$最小（最显著）的变量拉进模型
    \item \textbf{向后剔除}：把模型里变得不显著的踢出去
    \item \textbf{逐步（both）}：边进边出，每加一个新变量都重新检查老变量要不要留
\end{itemize}

本案例用$p$值准则、both方向，从12个候选里最终留下5个：活跃会员数、商圈客流、客单价、促销费用、门店面积，调整$R^2=0.838$，整体$F=154.8$，$p=4.3\times10^{-56}$。
\end{algoprinciple}

\begin{algosposs}
\begin{enumerate}
    \item \textbf{分析 → 回归 → 线性}，因变量选\textbf{月销售额}，把\textbf{全部候选自变量}放入自变量框
    \item "方法"下拉从"输入"改为\textbf{逐步}
    \item 点\textbf{选项}：进入$F$的概率0.05，剔除$F$的概率0.10
    \item 点\textbf{统计量}勾"共线性诊断""模型拟合度"
    \item 点击\textbf{确定}
\end{enumerate}

\textbf{默认参数参考}：筛选准则pvalue，方向both，入选$p<0.05$、剔除$p>0.1$，最大入选20个。本案例5折交叉验证测试$R^2=0.814$。
\end{algosposs}

\begin{algoresult}
\begin{itemize}
    \item \textbf{输入/除去的变量表}：一步步显示哪些变量被纳入、哪些被踢。本案例纳入5次、剔除0次
    \item \textbf{最终模型系数表}：活跃会员数系数$0.563$（$p=3.6\times10^{-50}$）、商圈客流$0.804$、客单价$0.373$——这三个关联最强
    \item \textbf{整体}：$R^2=0.843$、调整$R^2=0.838$、$F$检验$p<0.001$
    \item \textbf{交叉验证}：测试$R^2=0.814$与训练$R^2$接近，说明模型稳定不过拟合
    \item 入选变量最大VIF=1.06，共线性可控
\end{itemize}
\end{algoresult}

\begin{algonotes}
\textbf{什么时候用？}变量多、理论不清，需要自动粗筛。

\textbf{什么时候不用？}
\begin{itemize}
    \item 有明确理论模型时：用分层/全变量回归，别让数据乱选
    \item 样本量小（事件数不足）时：逐步会过拟合
\end{itemize}

\textbf{常见坑}
\begin{itemize}
    \item \textbf{坑1}：把逐步选出的变量当"因果结论"。它只是统计相关，且路径依赖——换个样本选出来的变量可能完全不同
    \item \textbf{坑2}：不做交叉验证。训练$R^2$很高但新数据崩了，就是过拟合
\end{itemize}

\textbf{和相似算法的区别}：逐步回归=自动版变量筛选；共线性VIF是它的\textbf{前置安检}；Lasso是它的正则化升级版（见进阶部分）。
\end{algonotes}

%% ========== 逻辑回归 ==========
\basicalgo{{\starmark}逻辑回归（Logistic）}{logistic_reg}

\begin{algointro}
当因变量是\textbf{是/否}这种二分类（如是否逾期、是否购买、是否患病）时，线性回归就不灵了——逻辑回归应运而生。本案例预测客户未来12个月是否逾期。
\end{algointro}

\begin{algoscene}
\begin{itemize}
    \item 因变量是0/1二分类，想找影响因素并算概率
    \item 信用评分、医疗诊断、营销响应预测
    \item 美赛中"是否发生某事件"的概率建模与分类
\end{itemize}
\end{algoscene}

\begin{algoprinciple}
线性回归输出的是"销售额=多少万"，可能算出负数、也可能算出超出范围的值。但我们想知道的是\textbf{"逾期的概率"}，它必须在0到1之间。

逻辑回归用一个叫\textbf{logistic函数（S形曲线）}把任意实数压到0--1之间：
$$\ln\frac{p}{1-p}=b_0+b_1x_1+\cdots+b_kx_k$$
左边是"对数发生比"。
\begin{itemize}
    \item 系数$b$解释成\textbf{优势比$OR=e^b$}：本案例"近12月逾期次数"系数0.685，$OR=1.983$——每多一次历史逾期，未来逾期几率约翻倍
    \item $OR>1$风险上升，$OR<1$风险下降
\end{itemize}

本案例模型AUC=0.832，区分能力不错。
\end{algoprinciple}

\begin{algosposs}
\begin{enumerate}
    \item \textbf{分析 → 回归 → 二元Logistic}
    \item 因变量选\textbf{是否逾期}（取值0/1，事件类别=1）
    \item 协变量选\textbf{年龄、年收入、负债收入比、额度使用率、近12月逾期次数、性别}
    \item 分类变量点"分类"，把性别设为哑变量（首类为参照）
    \item 点\textbf{选项}勾"Hosmer-Lemeshow""分类图""95\%置信区间"
    \item 点击\textbf{确定}
\end{enumerate}

\textbf{默认参数参考}：事件类别1，参考类别0，$\alpha=0.05$，最大迭代100，普通MLE。本案例事件95例、非事件55例，EPV=15.8（达标）。
\end{algosposs}

\begin{algoresult}
\begin{itemize}
    \item \textbf{模型整体检验}：似然比$\chi^2=51.1$，$p<0.001$，模型整体显著；McFadden $R^2=0.259$
    \item \textbf{Hosmer-Lemeshow}：$\chi^2=8.95$，$p=0.346>0.05$——预测概率和实际吻合，\textbf{校准良好}
    \item \textbf{系数表}看$p$和$OR$：
    \begin{itemize}
        \item 显著变量：负债收入比$OR=1.035$、额度使用率$OR=1.048$、近12月逾期次数$OR=1.983$
        \item 不显著：年龄、年收入、性别（$p>0.05$）
    \end{itemize}
    \item \textbf{AUC=0.832}：接近0.5是瞎猜，接近1是完美区分；0.83属于良好
    \item Youden最优阈值0.74时灵敏度0.66、特异度0.89
\end{itemize}
\end{algoresult}

\begin{algonotes}
\textbf{什么时候用？}因变量是二分类（0/1），要找影响因素或算概率。

\textbf{什么时候不用？}
\begin{itemize}
    \item 因变量是连续值：用线性回归
    \item 因变量是多分类无序：用多项Logistic；有序：用有序Logistic
    \item 事件数太少（EPV$<10$）：结果不稳，需精简变量或用精确方法
\end{itemize}

\textbf{常见坑}
\begin{itemize}
    \item \textbf{坑1}：把$OR$当概率。$OR=2$是"几率翻倍"，不是"概率翻倍"
    \item \textbf{坑2}：只看系数$p$不看AUC。$p$显著不代表模型分类能力强
\end{itemize}

\textbf{和相似算法的区别}：线性回归预测连续值，逻辑回归预测\textbf{分类概率}；Probit是它的姊妹版（用正态分布而非logistic）。
\end{algonotes}

\section{进阶级算法速览}

%% ========== 分位数回归 ==========
\advancedalgo{分位数回归}{quantile_reg}

\begin{algobrief}
\textbf{一句话}：不只看"平均"，而是拟合\textbf{条件中位数（及任意分位数）}的回归——看自变量对高分位、低分位人群的不同影响。

\textbf{美赛场景区}：收入研究中，教育对高收入者和低收入者的回报是否不同；极端值敏感的政策效应分析。

\textbf{核心原理}：普通最小二乘 minimizes 残差平方和，只估条件均值；分位数回归最小化不对称加权绝对残差，可估$\tau=0.1/0.5/0.9$等任意分位。

\textbf{操作要点}：在Dabbit AI SPSS工具箱选"分位数回归"，指定因变量、自变量和目标分位数（如0.25、0.5、0.75）。

\textbf{结果关键}：比较不同分位点上系数的大小变化——若高分位系数远大于低分位，说明效应在分布高端更强。

\textbf{注意}：它对异方差、厚尾数据更稳健，能揭示均值回归看不到的"分布异质性"。
\end{algobrief}

%% ========== Deming回归 ==========
\advancedalgo{Deming回归（正交回归）}{deming_reg}

\begin{algobrief}
\textbf{一句话}：自变量和因变量\textbf{都有测量误差}时用的回归——普通回归假设$x$完全精确，这在仪器比对时不成立。

\textbf{美赛场景区}：两种测量方法/仪器的结果一致性比较；方法学比对实验。

\textbf{核心原理}：普通OLS只让$y$方向残差最小，Deming让\textbf{垂直于回归线方向}的总误差最小（正交距离），$x$和$y$都被当作有误差。需预先知道两者测量误差方差比。

\textbf{操作要点}：工具箱选"Deming回归"，输入两列测量值，设定误差方差比。

\textbf{结果关键}：斜率和截距；若斜率接近1、截距接近0，说明两种方法一致性好。

\textbf{注意}：它和Passing-Bablok常用于方法比对。当$x$无误差时，退化为普通OLS。
\end{algobrief}

%% ========== Probit回归 ==========
\advancedalgo{Probit回归}{probit_reg}

\begin{algobrief}
\textbf{一句话}：和逻辑回归目的一样（二分类因变量），只是把S形曲线从logistic换成\textbf{正态累积分布}。

\textbf{美赛场景区}：离散选择、阈值模型（如"努力超过某阈值才成功"）；经济学probit模型。

\textbf{核心原理}：把$p$通过标准正态的反函数（probit链接）线性化，$\Phi^{-1}(p)=b_0+b_1x_1+\cdots$。系数解释不如OR直观，通常报告边际效应。

\textbf{操作要点}：\textbf{分析 → 回归 → 有序/Probit}，或在工具箱选"Probit回归"，指定响应频率和观测值。

\textbf{结果关键}：系数符号和显著性；效应大小用边际概率（$x$每变1单位，成功概率变多少）。

\textbf{注意}：Probit和Logistic结果通常非常接近，选哪个主要看习惯和理论假设；逻辑回归的OR更好解释。
\end{algobrief}

%% ========== 非线性回归 ==========
\advancedalgo{非线性回归}{nonlinear_reg}

\begin{algobrief}
\textbf{一句话}：当关系是\textbf{没法通过变量变换变成直线}的曲线（如Logistic增长、Michaelis-Menten酶动力）时，直接拟合自定义非线性方程。

\textbf{美赛场景区}：S形生长曲线、饱和动力学、传染病传播模型（logistic增长）、物理经验公式。

\textbf{核心原理}：不做线性化，直接用迭代算法（如Levenberg-Marquardt）最小化残差平方和。需要\textbf{人为指定方程形式和初始参数猜测}。

\textbf{操作要点}：\textbf{分析 → 回归 → 非线性估计}，输入自定义模型表达式和参数初值。

\textbf{结果关键}：参数的估计值、标准误、$p$值；收敛是否成功；$R^2$和残差图。

\textbf{注意}：初值给得不好容易不收敛或收敛到局部最优——这是它和曲线回归最大的门槛。
\end{algobrief}

%% ========== 对数线性模型 ==========
\advancedalgo{对数线性模型}{loglinear_model}

\begin{algobrief}
\textbf{一句话}：用回归的思路分析\textbf{多个分类变量}的列联表——把期望频数取对数后建模，拆解主效应和交互效应。

\textbf{美赛场景区}：高维列联表（三个以上分类变量）的结构分析；问卷中多类别变量的关联分解。

\textbf{核心原理}：$\log(\text{期望频数})$可分解为各变量主效应+交互项。通过比较嵌套模型的似然比$\chi^2$，判断哪些交互项重要。

\textbf{操作要点}：\textbf{分析 → 对数线性 → 模型选择}，选分类变量和设计类型。

\textbf{结果关键}：模型拟合的似然比$\chi^2$（$p>0.05$说明模型拟合好）；显著的高阶交互项说明变量间复杂关联。

\textbf{注意}：它是卡方检验在多维表上的推广。二维表时它就回到卡方独立性检验。
\end{algobrief}

%% ========== HLM ==========
\advancedalgo{多层线性模型（HLM/混合效应）}{hlm}

\begin{algobrief}
\textbf{一句话}：数据天然\textbf{嵌套分层}时用它——学生嵌套在班级里、班级嵌套在学校里，不能把它们当独立样本。

\textbf{美赛场景区}：分层抽样数据、纵向重复测量、组织内部嵌套结构；组内相关问题。

\textbf{核心原理}：把回归系数拆成"固定效应"（大家通用）和"随机效应"（各组围绕它浮动的部分），同时处理层1（个体）和层2（群体）变量。

\textbf{操作要点}：在Dabbit AI SPSS工具箱选"多层线性模型"，指定分层结构（如学生$\to$班级）和各层变量。

\textbf{结果关键}：组内相关系数ICC（高说明必须用多层模型）；固定效应系数；随机效应方差占比。

\textbf{注意}：数据有明显嵌套/聚类时，用普通回归会低估标准误、夸大显著性——这是它和普通回归最大的区别。
\end{algobrief}

%% ========== Heckman ==========
\advancedalgo{Heckman两阶段模型}{heckman}

\begin{algobrief}
\textbf{一句话}：解决\textbf{样本选择性偏差}——你能观测到的数据本身不是随机样本（如只观测到"参加了工作"的人的工资）。

\textbf{美赛场景区}：工资决定模型（只有就业者才有工资）；自选择样本下的效应估计。

\textbf{核心原理}：第一阶段用Probit估计"入选概率"，算出逆米尔斯比$\lambda$；第二阶段把$\lambda$作为额外自变量放回回归，校正选择性偏差。

\textbf{操作要点}：工具箱选"Heckman两阶段"，指定选择方程（是否入选）和结果方程（被解释的连续变量）。

\textbf{结果关键}：$\lambda$的系数显著=存在选择性偏差；校正后核心变量系数与原OLS相比变化多少。

\textbf{注意}：它要求有工具变量（只影响是否入选、不直接影响结果变量），否则识别困难。
\end{algobrief}

%% ========== 稳健回归 ==========
\advancedalgo{稳健回归（RANSAC/M估计）}{robust_reg}

\begin{algobrief}
\textbf{一句话}：普通最小二乘会被\textbf{几个极端点带偏}，稳健回归给异常点"降权"，让直线不被 outliers 绑架。

\textbf{美赛场景区}：数据有明显异常值（长尾分布）；需要回归结果抗干扰。

\textbf{核心原理}：M估计用更温和的损失函数（如Huber），对大残差不那么敏感；RANSAC反复随机抽小点子集拟合，找"最多点支持"的模型。

\textbf{操作要点}：工具箱选"稳健回归（RANSAC/M估计）"，指定因变量和自变量，选择方法。

\textbf{结果关键}：与OLS系数对比——若稳健系数和OLS差很多，说明OLS被异常点影响了；同时看被标记为异常的样本。

\textbf{注意}：它不是删异常值，而是\textbf{降权}。数据干净时，OLS更有效。
\end{algobrief}

%% ========== Poisson回归 ==========
\advancedalgo{Poisson回归}{poisson_reg}

\begin{algobrief}
\textbf{一句话}：因变量是\textbf{计数数据}（发生次数，如事故数、来电数）时用的回归。

\textbf{美赛场景区}：一天内门店客诉次数、单位面积内事件数；计数类因变量的影响因素分析。

\textbf{核心原理}：假设因变量服从泊松分布，用对数链接$\ln(\mu)=b_0+b_1x_1+\cdots$。系数解释为"发生率比IRR"（$e^b$）。

\textbf{操作要点}：\textbf{分析 → 广义线性模型}，计数分布选Poisson；或工具箱选Poisson回归。

\textbf{结果关键}：系数/$IRR$的$p$值；重点检查\textbf{过度离散}（方差远大于均值）——若过度离散，改用负二项回归。

\textbf{注意}：泊松要求\textbf{均值=方差}。现实计数数据常过度离散，此时Poisson的标准误会被低估，负二项回归更合适。
\end{algobrief}

%% ========== PLS ==========
\advancedalgo{偏最小二乘回归（PLS）}{pls_reg}

\begin{algobrief}
\textbf{一句话}：自变量多、又严重共线、样本还不多时，PLS一边降维一边回归，两全其美。

\textbf{美赛场景区}：光谱数据、高维经济指标下的预测；共线性严重且样本少的回归。

\textbf{核心原理}：同时对$X$和$Y$提取潜变量（主成分），并保证潜变量间协方差最大——既考虑$X$内部结构，又紧扣$Y$。

\textbf{操作要点}：工具箱选"PLS偏最小二乘"，指定因变量和全部自变量，用交叉验证选最佳成分数。

\textbf{结果关键}：成分数选择（交叉验证$Q^2$）；回归系数；$R^2$和预测$R^2$。

\textbf{注意}：它是\textbf{PCA+回归的结合体}。自变量严重共线、样本量$<$变量数时，它比普通回归稳。
\end{algobrief}

%% ========== Tobit ==========
\advancedalgo{Tobit回归}{tobit_reg}

\begin{algobrief}
\textbf{一句话}：因变量被\textbf{截断/删失}时用——比如一堆人"支出为0"（其实是没花），普通回归会偏。

\textbf{美赛场景区}：消费支出（大量0值）、贷款额度、持续时间等有下界且堆积在边界的数据。

\textbf{核心原理}：假设存在一个潜变量（真实支出），观测值在0处被"左删失"。用极大似然同时估计回归和删失概率。

\textbf{操作要点}：工具箱选"Tobit回归"，指定因变量、删失方向（左删失于0）和自变量。

\textbf{结果关键}：系数符号方向；删失比例；与普通OLS对比看校正幅度。

\textbf{注意}：因变量在0处堆积很多时才用；若只是大量0但不是删失概念，考虑零膨胀模型。
\end{algobrief}

%% ========== 岭回归 ==========
\advancedalgo{岭回归（Ridge）}{ridge_reg}

\begin{algobrief}
\textbf{一句话}：共线性严重时，给系数估计"加一点惩罚"，牺牲一点点无偏性换稳定——岭回归。

\textbf{美赛场景区}：自变量高度相关、VIF爆表，又不想删变量时的回归稳定化。

\textbf{核心原理}：在最小二乘目标里加$\lambda\sum b_j^2$惩罚项，把系数往0收缩。$\lambda$越大收缩越强，用交叉验证选。

\textbf{操作要点}：工具箱选"岭回归"，指定因变量和自变量，交叉验证选正则化参数$\lambda$。

\textbf{结果关键}：岭迹图（各系数随$\lambda$的变化）；交叉验证误差最小的$\lambda$；收缩后更稳定的系数。

\textbf{注意}：岭回归\textbf{不做变量选择}（系数只会缩小不会变0）。要选变量用Lasso。
\end{algobrief}

%% ========== 负二项 ==========
\advancedalgo{负二项回归}{negbin_reg}

\begin{algobrief}
\textbf{一句话}：计数数据\textbf{过度离散}（方差远大于均值）时，它是Poisson回归的修正版。

\textbf{美赛场景区}：日就诊人数、事故次数、客户投诉数——这类数据往往方差$>$均值。

\textbf{核心原理}：在泊松基础上加一个离散参数，允许方差$=\mu+\alpha\mu^2$。$\alpha=0$就退化为Poisson。

\textbf{操作要点}：工具箱选"负二项回归"，指定计数因变量和自变量。

\textbf{结果关键}：离散参数$\alpha$显著$>0$=确实过度离散、用负二项对；发生率比IRR及其$p$值。

\textbf{注意}：先做Poisson的离差检验，确认过度离散了再上负二项；若0值特别多，改用零膨胀模型。
\end{algobrief}

%% ========== 零膨胀泊松 ==========
\advancedalgo{零膨胀泊松回归（ZIP）}{zip_reg}

\begin{algobrief}
\textbf{一句话}：计数数据里\textbf{0特别多}（既有"本来就不会发生"的 structural zero，又有"可能发生但这次没发生"的）时用它。

\textbf{美赛场景区}：一年就医次数（健康人大量0）、消费频次、事故数——0堆积异常。

\textbf{核心原理}：混合模型——一部分概率用Logistic决定"是不是0"，剩下的用Poisson决定"发生几次"。

\textbf{操作要点}：工具箱选"零膨胀泊松"，指定计数因变量、零膨胀模型的自变量和计数模型自变量。

\textbf{结果关键}：Vuong检验判断是否需要零膨胀；两部分各自的系数解读。

\textbf{注意}：先做ZIP vs Poisson的Vuong检验；如果不仅0多还过度离散，升级为零膨胀负二项（ZINB）。
\end{algobrief}

%% ========== GEE ==========
\advancedalgo{广义估计方程（GEE）}{gee}

\begin{algobrief}
\textbf{一句话}：处理\textbf{重复测量/聚类数据}的回归——同一对象被反复测多次，观测不独立时用它。

\textbf{美赛场景区}：面板数据、纵向追踪数据、嵌套在群体内的重复观测。

\textbf{核心原理}：不指定完整的随机效应分布，只指定"工作相关矩阵"（如可交换、AR-1），用准似然估计边际效应。

\textbf{操作要点}：工具箱选"广义估计方程"，指定聚类ID、因变量分布族、链接函数和工作相关结构。

\textbf{结果关键}：群体平均边际效应系数；比较不同工作相关结构下结果的稳健性。

\textbf{注意}：它和HLM都处理重复测量——GEE估\textbf{总体平均效应}，HLM估个体层面的随机效应。
\end{algobrief}

%% ========== GLM ==========
\advancedalgo{广义线性模型（GLM）}{glm}

\begin{algobrief}
\textbf{一句话}：一个\textbf{总框架}，把线性回归、逻辑回归、Poisson回归统一起来——换个"分布+链接函数"就能适应不同因变量。

\textbf{美赛场景区}：因变量既不是连续正态、也不是简单二分类时的统一建模入口。

\textbf{核心原理}：三个件——指数族分布（正态/二项/泊松/伽玛）+线性预测$\eta=Xb$+链接函数$g(\mu)=\eta$。选对组合就得到各种专用模型。

\textbf{操作要点}：\textbf{分析 → 广义线性模型}，在对话框里选分布族和链接函数。

\textbf{结果关键}：模型整体似然比检验；各系数$p$值；拟合优度（偏差/自由度）。

\textbf{注意}：逻辑回归是"二项分布+logit链接"，Poisson是"泊松+对数链接"——它们都是GLM的特例。
\end{algobrief}

%% ========== 零膨胀负二项 ==========
\advancedalgo{零膨胀负二项回归（ZINB）}{zinb_reg}

\begin{algobrief}
\textbf{一句话}：计数数据\textbf{既0特别多、又过度离散}时的终极版——ZIP的离散修正。

\textbf{美赛场景区}：大量0且方差远大于均值的计数（如慢病就诊次数、极端消费行为）。

\textbf{核心原理}：两部分混合——Logistic部分区分"本来不会发生"，负二项部分处理"发生了几次"，同时吸收过度离散。

\textbf{操作要点}：工具箱选"零膨胀负二项"，分别指定零膨胀模型和计数模型的自变量。

\textbf{结果关键}：Vuong检验支持ZINB优于负二项；两部分系数；离散参数。

\textbf{注意}：模型复杂、参数多，需要较大样本。先用ZIP/ZINB/负二项三者比较（AIC最小）再定稿。
\end{algobrief}

%% ========== Lasso ==========
\advancedalgo{Lasso回归}{lasso_reg}

\begin{algobrief}
\textbf{一句话}：岭回归的"会选变量"版本——惩罚不仅把系数缩小，还能\textbf{直接把不重要变量压到0}，实现自动变量筛选。

\textbf{美赛场景区}：自变量极多（基因、词袋、海量指标）的高维建模；需要稀疏模型。

\textbf{核心原理}：把惩罚项从$\lambda\sum b_j^2$（岭）换成$\lambda\sum|b_j|$（L1）。L1惩罚的几何形状让部分系数恰好变0。

\textbf{操作要点}：工具箱选"Lasso回归"，交叉验证选正则化参数$\lambda$。

\textbf{结果关键}：交叉验证误差最小的$\lambda$；非零系数对应的变量即"被选中"的关键变量；系数路径图。

\textbf{注意}：它和逐步回归都选变量，但Lasso是\textbf{正则化}方法，更稳定、不易过拟合，适合高维数据。
\end{algobrief}

% 第5章 聚类与降维
\chapter{聚类与降维}

\begin{chapteroverview}
\textbf{这个分类是干什么的？}

前面的方法都有"标准答案"（因变量、分组标签）。这一章是\textbf{无监督学习}——\textbf{没有标准答案}。聚类是把"长得像"的对象自动归成一群（会员分群、客户分层）；降维是把"十几个互相纠缠的指标"压缩成"两三个综合指标"（主成分、因子），让数据从啰嗦变简洁。

\textbf{美赛什么题型常用？}

C题（大数据题）高频出现：客户分群、城市分类、指标体系构建；A/B题中把高维观测投影到二维做可视化；论文中"综合评价"前常用PCA先把指标浓缩。

\textbf{学习路径建议}
先学\textbf{K-Means聚类}（最经典）和\textbf{主成分分析PCA}（降维入门）；再认识\textbf{二阶聚类}（能同时处理连续+分类变量）和\textbf{探索性因子分析}（问卷量表维度提炼）；进阶了解分层聚类、DBSCAN、判别分析LDA、MDS和RDA。
\end{chapteroverview}

\section{基础级算法}

%% ========== K-Means ==========
\basicalgo{{\starmark}聚类分析（K-Means）}{kmeans}

\begin{algointro}
K-Means就是\textbf{让计算机自动把样本分成$K$堆}，堆内的人像、堆间的人不像。本案例按会员的消费频次、客单价、活跃天数，把150名会员自动分成3个群体。
\end{algointro}

\begin{algoscene}
\begin{itemize}
    \item 没有标签，想把样本自动分成若干"人群包/客户群"
    \item 美赛C题客户分群、城市/地区画像、行为模式识别
    \item 论文中"分群运营""分层管理"的依据
\end{itemize}
\end{algoscene}

\begin{algoprinciple}
想象在桌面上撒一把珠子，要把它们分成$K$堆。
\begin{enumerate}
    \item 先随机选$K$个点当"堆心"
    \item 每个珠子归到离它最近的堆心
    \item 重新计算每堆的中心（均值），移动堆心
    \item 重复直到堆心不再移动
\end{enumerate}

\textbf{关键问题：$K$取几？}本案例在$k=2$到10间试，看\textbf{轮廓系数}——越接近1说明堆内紧、堆间分。结果$k=3$时轮廓系数0.77最高，故选3类。

三类画像：簇0（40\%）活跃高频但客单低；簇1（22.7\%）三项都低（沉默客户）；簇2（37.3\%）客单价高但频次低（高价值低频）。
\end{algoprinciple}

\begin{algosposs}
\begin{enumerate}
    \item 点击菜单 \textbf{分析 → 分类 → K-平均值聚类}
    \item 把参与分群的变量（\textbf{月消费频次、平均客单价、月活跃天数}）选入\textbf{变量}
    \item "聚类数"填3（也可先用工具箱自动选$K$）
    \item 点击\textbf{选项}，勾选"初始聚类中心""最终聚类中心""ANOVA表"
    \item 方法选\textbf{k-means++}初始化，固定随机种子保证可复现
    \item 点击\textbf{确定}
\end{enumerate}

\textbf{默认参数参考}：先把变量做\textbf{标准化}（消除量纲，否则客单价数值大会主导距离）；$n$\_init多跑几次取最优；$k=3$。
\end{algosposs}

\begin{algoresult}
\begin{itemize}
    \item \textbf{初始/最终聚类中心}：每堆在各变量上的均值，是"画像"核心
    \item \textbf{簇数诊断表}：$k=3$时轮廓系数0.77、Calinski-Harabasz=897（比$k=2$的0.63、$k=4$的0.59都高），故定3类
    \item \textbf{簇规模}：簇0=60（40\%）、簇1=34（22.7\%）、簇2=56（37.3\%）
    \item \textbf{轮廓系数}$0.77>0.5$：分群结构清晰；若$<0.25$说明边界模糊、分群失败
    \item 变量贡献eta$^2$：客单价0.93、消费频次0.92、活跃天数0.92——三者都很能区分人群
\end{itemize}
\end{algoresult}

\begin{algonotes}
\textbf{什么时候用？}连续变量、想要球形簇、样本量大时。

\textbf{什么时候不用？}
\begin{itemize}
    \item 不知道$K$且不想猜：用二阶聚类或分层聚类
    \item 簇形状不规则（月牙形）或有噪声点：K-Means不行，用DBSCAN
    \item 变量量纲差异大：必须先标准化，否则大数值变量主导
\end{itemize}

\textbf{常见坑}
\begin{itemize}
    \item \textbf{坑1}：不标准化直接聚类，客单价几百块把频次几次的距离全盖住
    \item \textbf{坑2}：把簇均值差当"显著差异"。K-Means是描述性分群，不是假设检验
\end{itemize}

\textbf{和相似算法的区别}：K-Means要事先定$K$、是球形簇；二阶聚类自动定$K$还能吃分类变量；DBSCAN能识别噪声和任意形状簇。
\end{algonotes}

%% ========== 二阶聚类 ==========
\basicalgo{二阶聚类（两步聚类）}{two_step}

\begin{algointro}
二阶聚类是\textbf{全自动}的K-Means升级版——\textbf{连续变量和分类变量都能吃}，还能\textbf{自动帮你决定分几类}。本案例混合年消费金额（连续）和性别、会员等级（分类）给健身会员分群。
\end{algointro}

\begin{algoscene}
\begin{itemize}
    \item 聚类变量里\textbf{既有数值又有分类}（K-Means处理不了分类）
    \item 不想手动试$K$，想让算法自动定最优簇数
    \item 美赛中混合属性的客户/对象分群
\end{itemize}
\end{algoscene}

\begin{algoprinciple}
它分两步，这也是"两步聚类"名字的由来：
\begin{enumerate}
    \item \textbf{阶段一}：逐个扫描样本，把相近的先攒成小"子簇"（像先把散珠子拢成小堆）
    \item \textbf{阶段二}：用对数似然距离把这些小簇逐层合并，自动试$k=2$到10，用\textbf{BIC/AIC}和\textbf{轮廓系数}选最佳
\end{enumerate}

本案例试$k=2$--10，轮廓系数在$k=4$处最高0.271，故定4簇。连续变量用欧式距离、分类变量用对数似然距离，所以它能混合两类变量。
\end{algoprinciple}

\begin{algosposs}
\begin{enumerate}
    \item 点击菜单 \textbf{分析 → 分类 → 两步聚类}
    \item 连续变量（\textbf{年消费金额、月均到店次数、单次时长、满意度}）放入"连续变量"
    \item 分类变量（\textbf{性别、会员等级、到店时段、运动偏好}）放入"分类变量"
    \item "聚类数"选\textbf{自动确定}，范围2--10
    \item 点\textbf{输出}勾"聚类质量轮廓图"
    \item 点击\textbf{确定}
\end{enumerate}

\textbf{默认参数参考}：选择准则auto（轮廓系数优先），连续变量自动标准化，缺失按阈值处理。本案例$n=150$，4连续+4分类变量。
\end{algosposs}

\begin{algoresult}
\begin{itemize}
    \item \textbf{自动定簇}：$k=4$时轮廓系数0.271（中等质量），BIC曲线佐证
    \item \textbf{簇规模}：簇0=39、簇1=35、簇2=36、簇3=40，较均衡
    \item \textbf{连续变量画像}：簇3年消费17063元、到店17.7次——"高价值晚间器械客"；簇1年消费仅1717元、到店2.6次——"低活跃午后团操客"
    \item \textbf{分类变量画像}：簇3以铂金会员+晚间+器械为主；簇1以普通会员+午后+团体操课为主
    \item \textbf{变量区分度}：8个变量区分度检验$p$都$<0.01$（性别$p=0.008$），说明它们都在有效分群
\end{itemize}
\end{algoresult}

\begin{algonotes}
\textbf{什么时候用？}混合连续+分类变量、想要自动定$K$、样本较大。

\textbf{什么时候不用？}全是连续变量且形状规则，K-Means更直观；簇形状不规则需要噪声识别，用DBSCAN。

\textbf{常见坑}
\begin{itemize}
    \item \textbf{坑1}：轮廓系数0.27其实只是"中等"，不要硬说"分群非常好"——按0.25以下弱、0.5以上强的标准诚实报告
    \item \textbf{坑2}：把常量列或高缺失列塞进聚类，污染距离计算
\end{itemize}

\textbf{和相似算法的区别}：它=K-Means的自动版+能吃分类变量。K-Means要手定$K$、只吃连续；分层聚类画树状图但不吃混合变量。
\end{algonotes}

%% ========== PCA ==========
\basicalgo{{\starmark}主成分分析（PCA）}{pca}

\begin{algointro}
PCA把\textbf{十几个互相关的指标}压缩成\textbf{两三个互不相关的综合指标}（主成分），用少量维度代表大部分信息。本案例从9个门店经营指标压缩到3个主成分。
\end{algointro}

\begin{algoscene}
\begin{itemize}
    \item 指标太多且互相纠缠，想降维、去共线性
    \item 综合评价前先浓缩指标；后续聚类/回归前先做特征压缩
    \item 美赛高维数据可视化、指标体系构建
\end{itemize}
\end{algoscene}

\begin{algoprinciple}
9个指标（客流、销售额、客单价、复购率……）很多在说同一件事。PCA的思路像\textbf{找最佳观察角度}：
\begin{itemize}
    \item 第一主成分PC1：数据"波动最大"的方向，本它解释了47.4\%的方差
    \item 第二主成分PC2：和PC1垂直、又能解释最多剩余波动（20.6\%）
    \item 第三PC3：再解释18.5\%
\end{itemize}

前3个累计解释86.6\%——几乎用3个综合指标就装下了原来9个指标的信息。\textbf{适用性前提}：先做Bartlett球度检验（本案例$\chi^2=1035.7$，$p=5\times10^{-194}$，适合降维）和KMO（本案例0.811$>0.6$，适合）。
\end{algoprinciple}

\begin{algosposs}
\begin{enumerate}
    \item 点击菜单 \textbf{分析 → 降维 → 因子}
    \item 把9个数值变量全部选入"变量"框
    \item 点\textbf{描述}，勾选\textbf{KMO和Bartlett球形检验}
    \item 点\textbf{抽取}，选"基于特征值$>$1"或固定提取个数
    \item 点\textbf{确定}
\end{enumerate}

\textbf{默认参数参考}：先做z-score标准化（等效于基于相关矩阵）；选择模式cumulative（累计贡献率到80\%停）；SVD求解器auto。本案例保留3个主成分。
\end{algosposs}

\begin{algoresult}
\begin{itemize}
    \item \textbf{适用性}：KMO=0.811$>0.6$；Bartlett $p<0.001$——数据适合做PCA
    \item \textbf{总方差解释表}：看特征值和累计贡献率
    \begin{itemize}
        \item PC1特征值4.298，解释47.4\%
        \item PC2特征值1.865，累计68.0\%
        \item PC3特征值1.680，累计\textbf{86.6\%}——保留到此为止
    \end{itemize}
    \item \textbf{成分矩阵（载荷）}：看每个主成分主要由哪些变量贡献。PC1主要是会员占比(+0.78)、线上占比(+0.74)、客单价(+0.73)——可命名为"会员价值维度"
    \item \textbf{碎石图}：肘部位置就是该停的地方
\end{itemize}
\end{algoresult}

\begin{algonotes}
\textbf{什么时候用？}变量多且相关，想压缩成少数综合指标、消除共线性。

\textbf{什么时候不用？}
\begin{itemize}
    \item 变量间几乎不相关（KMO$<0.5$）：没有可提取的公共结构
    \item 需要每个指标都保持可解释性：PCA合成的主成分有时难命名
\end{itemize}

\textbf{常见坑}
\begin{itemize}
    \item \textbf{坑1}：不标准化直接做。量纲大的变量会主导主成分
    \item \textbf{坑2}：强行解释主成分含义。主成分是数学方向，业务命名要结合载荷判断
\end{itemize}

\textbf{和相似算法的区别}：PCA是\textbf{无监督}的纯降维；因子分析更进一步找"潜在维度"并做旋转；LDA是有监督的降维（用到类别标签）。
\end{algonotes}

%% ========== 因子分析 ==========
\basicalgo{因子分析（探索性）}{exploratory_fa}

\begin{algointro}
问卷问了15道题，但背后其实只有"菜品、服务、环境"3个\textbf{看不见的维度}。探索性因子分析就是把这些表面题目归并成背后的潜在维度。
\end{algointro}

\begin{algoscene}
\begin{itemize}
    \item 问卷量表开发：一堆题项背后有几个维度
    \item 检验问卷结构效度、为算维度总分做准备
    \item 美赛中量表/评价体系的结构提炼
\end{itemize}
\end{algoscene}

\begin{algoprinciple}
顾客给"口味""分量""温度""新鲜度""出品"都打了高分——这5题其实都在测"菜品品质"这个\textbf{潜变量}。因子分析把它们归成一簇。

\begin{itemize}
    \item \textbf{KMO=0.762}$>0.6$，Bartlett $\chi^2=564.99$（$df=105$，$p=3.8\times10^{-64}$）——适合做
    \item 按特征值$>$1提取3个公因子，用\textbf{最大方差旋转(varimax)}让载荷结构清晰
    \item 本案例因子1=服务（服务员态度、响应速度等5题），因子2=菜品（口味、新鲜等5题），因子3=环境（整洁、氛围等5题）
\end{itemize}

旋转后累计解释方差40.6\%。\textbf{载荷}$\geq0.5$算有效归属。
\end{algoprinciple}

\begin{algosposs}
\begin{enumerate}
    \item \textbf{分析 → 降维 → 因子}
    \item 把15个题项全部选入变量框
    \item 点\textbf{描述}勾KMO和Bartlett
    \item 点\textbf{抽取}：方法选主轴因子法PAF，按特征值$>$1定因子数
    \item 点\textbf{旋转}：选\textbf{最大方差法varimax}
    \item 点\textbf{确定}
\end{enumerate}

\textbf{默认参数参考}：提取方法PAF，旋转varimax，有效载荷阈值0.50，特征值$>$1自动定因子数。本案例$n=150$（题项的10倍）。
\end{algosposs}

\begin{algoresult}
\begin{itemize}
    \item \textbf{适用性}：KMO=0.762，Bartlett $p<0.001$——通过
    \item \textbf{总方差解释}：3个因子旋转后累计解释40.6\%
    \item \textbf{旋转后载荷矩阵}是核心：每道题在哪个因子上载荷最高就归哪个因子。本案例15题干净地分进3簇，无交叉载荷
    \item \textbf{共同度}：每题被公因子解释的比例，偏低（$<0.3$）的题考虑删改
    \item \textbf{碎石图}：肘部确认3个因子合理
\end{itemize}
\end{algoresult}

\begin{algonotes}
\textbf{什么时候用？}问卷/量表题项多，想提炼背后维度。

\textbf{什么时候不用？}
\begin{itemize}
    \item KMO$<0.6$或Bartlett不显著：题项间没公共结构
    \item 只是想压缩维度做回归/聚类：PCA更直接
\end{itemize}

\textbf{常见坑}
\begin{itemize}
    \item \textbf{坑1}：不旋转就硬读载荷。未旋转的因子含义模糊，旋转后才能干净归类
    \item \textbf{坑2}：样本量太小。经验要求$n$至少是题项数的5--10倍
\end{itemize}

\textbf{和相似算法的区别}：PCA主成分是变量的线性组合（重在降维）；因子分析假设背后有\textbf{潜变量}（重在找结构）。两者常被混用，但理论出发点不同。
\end{algonotes}

\section{进阶级算法速览}

%% ========== DBSCAN ==========
\advancedalgo{DBSCAN密度聚类}{dbscan}

\begin{algobrief}
\textbf{一句话}：不靠预设簇数，按\textbf{样本密集程度}自动成团，还能把"离群点"标成噪声。

\textbf{美赛场景区}：形状不规则的簇、有异常点的地理/轨迹数据（本案例打车订单的出行模式+异常订单识别）。

\textbf{核心原理}：若某点周围半径eps内有至少min\_samples个点，就把这片稠密区域发展成一个簇；稀疏处的点标记为噪声。本案例eps=0.5、min\_samples=5，分出3簇+16个噪声点，非噪声轮廓系数0.795。

\textbf{操作要点}：工具箱选"DBSCAN密度聚类"，标准化特征；用k-距离图拐点选eps；欧氏距离。

\textbf{结果关键}：簇数自动确定；噪声占比（本案例10.7\%）；非噪声样本轮廓系数；各簇特征均值剖面。

\textbf{注意}：它不需要预设$K$、能识别任意形状和噪声，但对eps和min\_samples敏感，且密度不均时效果打折。
\end{algobrief}

%% ========== RDA ==========
\advancedalgo{冗余分析（RDA）}{rda}

\begin{algobrief}
\textbf{一句话}：研究\textbf{多个响应变量}如何被\textbf{多个解释变量}共同解释的约束排序方法——生态学常用。

\textbf{美赛场景区}：环境因子驱动群落结构；多响应变量受多环境变量影响的分析（本案例水质指标如何驱动藻类丰度）。

\textbf{核心原理}：先对响应矩阵做多元回归得拟合值，再对拟合值做主成分分解得"约束排序轴"。全局置换检验$p=0.001$，$R^2=89.24\%$说明环境变量显著解释了藻类变异。

\textbf{操作要点}：工具箱选"冗余分析RDA"，指定响应变量组和解释变量组；分类解释变量自动哑变量化；置换999次。

\textbf{结果关键}：全局置换检验$p$；各解释变量单独解释率；排序图箭头方向与长度。

\textbf{注意}：它是\textbf{有约束}的PCA。若响应沿环境梯度呈单峰分布，应改用典范对应分析CCA。
\end{algobrief}

%% ========== LDA ==========
\advancedalgo{判别分析（LDA）}{lda}

\begin{algobrief}
\textbf{一句话}：已知历史分组标签，找一条"最佳分界线"把各组分开，并给新样本归类——有监督的分类+降维。

\textbf{美赛场景区}：信用评级、疾病判别、模式识别；用已知类别的指标建立判别规则给新对象分类（本案例企业信用等级判别）。

\textbf{核心原理}：找线性判别函数，使\textbf{组间差异相对组内差异}最大。本案例5个财务指标建判别函数，留一法准确率80.6\%（Kappa=0.708），远高于猜最大类的33.3\%基线。

\textbf{操作要点}：工具箱选"判别分析LDA"，指定分组列和特征列；Box's M检验协方差齐性；留一法交叉验证。

\textbf{结果关键}：Wilks' Lambda（整体判别力）；标准化判别系数（谁最能区分）；混淆矩阵和各类命中率。

\textbf{注意}：它和PCA都降维，但\textbf{LDA用到类别标签}、目标是"分开组"；PCA无标签、目标是"保留方差"。协方差不齐时用二次判别QDA。
\end{algobrief}

%% ========== MDS ==========
\advancedalgo{多维尺度分析（MDS）}{mds}

\begin{algobrief}
\textbf{一句话}：把"谁和谁像、差多远"的抽象关系，画成\textbf{一张二维地图}——点越近越相似。

\textbf{美赛场景区}：品牌/门店感知图；对象间相似性的可视化；无法直接坐标化但知道两两距离时的定位。

\textbf{核心原理}：先算两两距离，再在二维空间摆点，让"图上距离"尽量贴近"原始距离"。本案例150家门店在7个感知维度上，标准化应力0.019（优），前两维解释99\%。

\textbf{操作要点}：工具箱选"多维尺度MDS"，标准化特征，度量MDS，降维到2维，随机种子42。

\textbf{结果关键}：标准化应力$<0.05$为优；Shepard图点贴近对角线；配置图上的点群和孤立点。

\textbf{注意}：它只保留\textbf{远近关系}，不保留绝对坐标和方向。非度量MDS用于只知道排序相似性的数据。
\end{algobrief}

%% ========== 分层聚类 ==========
\advancedalgo{分层聚类（系统聚类）}{hierarchical_cluster}

\begin{algobrief}
\textbf{一句话}：不预先定簇数，像\textbf{族谱}一样把样本一层层两两合并，最后画成一棵树（树状图/谱系图）。

\textbf{美赛场景区}：小样本下想看聚类层次结构；变量少、想直观展示"谁和谁先聚在一起"。

\textbf{核心原理}：先每个样本自成一簇，用ward连接把最像的两簇反复合并。本案例$n=150$、6指标，ward+欧氏距离，轮廓系数在$k=3$处最高0.683，cophenetic系数0.974。

\textbf{操作要点}：\textbf{分析 → 分类 → 系统聚类}，选变量、 ward法、欧氏距离、z-score标准化，输出树状图。

\textbf{结果关键}：树状图在合适高度横切得簇数；轮廓系数曲线选$k$；各簇规模和均值画像。

\textbf{注意}：它\textbf{不需要预跑$K$}，但样本量大时树状图会非常臃肿，大样本更适合K-Means。
\end{algobrief}

% 第6章 时间序列与计量经济
\chapter{时间序列与计量经济}

\begin{chapteroverview}
\textbf{这个分类是干什么的？}

这一章分两块。\textbf{时间序列}部分处理"按时间排好队"的数据（如月度销售额、日客流），关心序列平不平稳、有没有自相关、变量之间谁领先谁；\textbf{计量经济}部分则更雄心勃勃——它要回答"\textbf{因果}"问题：上了这个系统、投了这笔钱，到底带来了多少真实效果？它通过基准回归、面板模型、双重差分、工具变量等方法，努力从数据里把因果关系和"相关关系"区分开。

\textbf{美赛什么题型常用？}

计量方法在美赛中越来越受重视：B题、C题常要求评估一项政策/措施的效果（DID、PSM、RDD、工具变量都是得分利器）；时间序列方法（ADF、格兰杰、VAR、协整）则用于分析随时间演化的变量关系，是论文"建模严谨性"的体现。凡是题目问"$X$ 对 $Y$ 有没有影响、影响多大、是不是因果"，基本都落在这一章。

\textbf{学习路径建议}

先打地基：\textbf{基准回归}（带星标，一切计量的起点）和\textbf{ADF单位根检验}（时间序列的入场券）。然后学\textbf{面板模型}（带星标，处理多年多主体数据）和\textbf{双重差分DID}（带星标，政策评估王牌）。再补\textbf{ACF/PACF}和\textbf{格兰杰因果检验}这两个时间序列诊断工具。进阶部分则是DID、PSM、工具变量、RDD、VAR等"因果识别工具箱"——注意其中不少方法SPSS原生支持有限，需要通过Dabbit AI SPSS工具箱或SPSS的R/Python扩展实现。
\end{chapteroverview}

\section{基础级算法}

%% ========== 基准回归 ==========
\basicalgo{{\starmark}基准回归}{ols_baseline}

\begin{algointro}
基准回归就是用普通最小二乘法（OLS）拟合一条"最佳拟合直线"，量化核心解释变量对因变量的平均影响，并通过一系列诊断检验保证这个系数是可信的。它是几乎所有计量分析的第一步。
\end{algointro}

\begin{algoscene}
\begin{itemize}
    \item 美赛中估计"某因素对某结果的平均影响"，如数字化营销投入对门店月均销售额的影响
    \item 在加入控制变量前后比较核心系数是否稳定，为后续复杂模型打基础
    \item 论文中"基准模型/基准回归表"的标准做法
    \item 需要判断多重共线性、异方差、异常值等是否威胁结论时
\end{itemize}
\end{algoscene}

\begin{algoprinciple}
想象你要画一条直线穿过一堆散点，让所有点到直线的\textbf{垂直距离平方之和最小}——这就是OLS。

\begin{itemize}
    \item 模型写作 $y=\beta_0+\beta_1 x_1+\beta_2 x_2+\cdots+\varepsilon$，其中 $y$ 是因变量（月均销售额），$x_1$ 是核心解释变量（数字化营销投入），$x_2,\ldots$ 是控制变量（门店面积、营业年限、商圈客流、是否核心商圈），$\varepsilon$ 是误差。
    \item \textbf{核心看 $\beta_1$}：本案例为3.357，意思是在控制其他变量后，数字化营销投入每增加1个单位，月均销售额平均增加3.357个单位。
    \item \textbf{为什么要加控制变量}：如果只看营销投入和销售额，可能是"面积大的门店既投得多又卖得多"。加入面积等控制变量后，才能把这种混杂因素剥离，看到营销投入的"净效应"。
\end{itemize}

\textbf{诊断三板斧}：VIF检验多重共线性（$VIF>10$ 危险）、Breusch-Pagan检验异方差、RESET检验模型设定是否正确。本案例采用HC1稳健标准误应对异方差。
\end{algoprinciple}

\begin{algosposs}
\begin{enumerate}
    \item 点击菜单 \textbf{分析 → 回归 → 线性}
    \item 把因变量"月均销售额"选入\textbf{因变量}
    \item 把核心解释变量"数字化营销投入"和控制变量（门店面积、营业年限、商圈日均客流、是否核心商圈）选入\textbf{自变量}
    \item 点击\textbf{统计量}，勾选：估计系数、置信区间、\textbf{共线性诊断}、\textbf{异方差稳健标准误（HC1）}
    \item 在"保存"中可输出残差与库克距离，用于检查强影响点
    \item 依次估计三个规格：只有核心变量、只有控制变量、核心+全部控制变量，对比核心系数变化
    \item 点击\textbf{确定}
\end{enumerate}

\textbf{默认参数参考}：显著性水平0.05，VIF阈值10，稳健标准误HC1，包含常数项，不缩尾、不标准化。
\end{algosposs}

\begin{algoresult}
\begin{itemize}
    \item \textbf{系数（B）}：核心解释变量数字化营销投入为3.357，方向为正，说明投入越多销售额越高。
    \item \textbf{标准误与p值（Sig.）}：稳健标准误0.274，$p<0.001$，高度显著。$p<0.05$ 才认为系数可信。
    \item \item \textbf{95\%置信区间}：不跨过0就对应显著，区间越窄估计越精确。
    \item \textbf{加入控制变量后的变化}：核心系数方向保持一致，相对仅含核心变量时变化约22\%——说明效应稳健，但也提示控制变量确实吸收了一部分混杂。
    \item \textbf{VIF}：各自变量VIF均应小于10，否则存在多重共线性。
\end{itemize}

\textbf{判断好坏}：核心系数方向符合预期、$p<0.05$、VIF合格、加入控制后系数不发生方向性反转，即为稳健的基准结果。
\end{algoresult}

\begin{algonotes}
\textbf{什么时候用？}任何"X对Y平均影响"的定量分析第一步，几乎无例外。

\textbf{什么时候不用？}
\begin{itemize}
    \item X本身是"被Y影响"的（互为因果）或有遗漏变量时，OLS系数有偏，需用工具变量/DID等因果方法
    \item 因变量是0/1二分类时，应用逻辑回归（见第4章）
\end{itemize}

\textbf{常见坑}
\begin{itemize}
    \item \textbf{坑1}：把"显著"当"重要"。$p<0.05$ 只说明效应不为零，不代表效应大；还要看系数大小和 $R^2$。
    \item \textbf{坑2}：忽略异方差，用普通标准误导致 $p$ 值不可信。异方差存在时务必用HC稳健标准误。
    \item \textbf{坑3}：不做共线性诊断，变量高度相关时系数符号甚至会反常。
\end{itemize}

\textbf{和相似算法的区别}：vs \textbf{面板模型}——基准回归是横截面混合OLS，面板模型能控制个体固定效应；vs \textbf{PSM/IV}——基准回归只给"条件相关"，后两者追求更干净的因果解释。
\end{algonotes}

%% ========== ADF单位根检验 ==========
\basicalgo{ADF单位根检验}{adf_test}

\begin{algointro}
ADF单位根检验判断一条时间序列是不是"平稳"的——即它的均值和波动是否长期稳定。如果序列带"单位根"（不平稳），直接拿来回归会产生"伪回归"，所以它是时间序列分析前的必做体检。
\end{algointro}

\begin{algoscene}
\begin{itemize}
    \item 任何时间序列建模、协整分析之前的必做前提检验
    \item 本案例：评估销售额与营销费用是否平稳、要差分几阶才能预测
    \item 判断两个序列能否直接回归，还是必须进入协整框架
    \item 论文中报告"序列为I(0)或I(1)"的依据
\end{itemize}
\end{algoscene}

\begin{algoprinciple}
想象一个醉汉走路：他每一步的方向完全随机，走着走着就漂走了，没有回到出发点的趋势——这就是带单位根的\textbf{不平稳序列}，比如"醉汉漫步"。而一个挂在弹簧下的重物，无论怎么扰动都会回到平衡位置——这是\textbf{平稳序列}。

\begin{itemize}
    \item \textbf{原假设 $H_0$}：序列存在单位根（不平稳）。我们希望拒绝它。
    \item \textbf{判断逻辑}：把ADF统计量和临界值比较，或直接看 $p$ 值。$p<0.05$ 拒绝原假设，即序列平稳；$p\geq0.05$ 不能拒绝，即不平稳。
    \item \textbf{单整阶数}：水平序列不平稳就做一阶差分再检验。一阶差分后平稳，称为 $I(1)$；本身平稳则为 $I(0)$。本案例sales水平 $p=0.2651$（不平稳），一阶差分后 $p=0.0000$（平稳），故为 $I(1)$。
\end{itemize}

\textbf{为什么重要}：两个各自随机漂移（不平稳）的序列做回归，即使毫不相干也可能得到很高的 $R^2$——这就是"伪回归"。必须先确认平稳性或同阶单整。
\end{algoprinciple}

\begin{algosposs}
\begin{enumerate}
    \item 在Dabbit AI SPSS工具箱中选择"ADF单位根检验"
    \item 上传含时间列与数值列的月度序列（本案例销售额sales\_amount、营销费用）
    \item 设置参数：
    \begin{itemize}
        \item 回归形式：默认含常数项（regression='c'）
        \item 滞后阶数：AIC自动选择（autolag='AIC'）
        \item 是否检验一阶差分：是
        \item 显著性水平0.05；缺失值线性插值；KPSS交叉验证可选
    \end{itemize}
    \item 运行，依次查看水平序列与一阶差分序列的ADF结果
\end{enumerate}

\textbf{SPSS中对应操作}：SPSS无原生ADF菜单，通过 \textbf{分析 → 时间序列} 或R/Python扩展（statsmodels的adfuller）实现。
\end{algosposs}

\begin{algoresult}
\begin{itemize}
    \item \textbf{ADF统计量与p值}：水平序列 $p=0.2651\geq0.05$，不能拒绝单位根，序列不平稳。
    \item \textbf{一阶差分结果}：$p=0.0000<0.05$，差分后平稳。
    \item \textbf{结论判定}：sales为 $I(1)$ 序列（需一阶差分）。若营销费用同为 $I(1)$，两者才可能做协整。
    \item \textbf{临界值对照}：除p值外，也可把ADF统计量与1\%、5\%、10\%临界值比较，统计量小于临界值才平稳。
\end{itemize}

\textbf{判断好坏}：$p<0.05$ 即平稳。报告时要写明"原序列几阶单整、差分后是否平稳"，这决定后续建模路径。
\end{algoresult}

\begin{algonotes}
\textbf{什么时候用？}任何时间序列回归/预测前，判断平稳性。

\textbf{什么时候不用？}横截面数据（无时间顺序）不需要做单位根检验。

\textbf{常见坑}
\begin{itemize}
    \item \textbf{坑1}：对两个 $I(1)$ 序列直接回归而不做协整检验，得到伪回归。
    \item \textbf{坑2}：过度差分。序列本就平稳却反复差分，会人为引入噪声，可用KPSS检验交叉验证。
    \item \textbf{坑3}：回归形式选错（如该加趋势项却没加），会导致平稳性误判。
\end{itemize}

\textbf{和相似算法的区别}：vs \textbf{差分分析}——ADF是检验"是否平稳、几阶单整"，差分分析是实际执行差分并看变动；vs \textbf{协整检验}——ADF判断单整阶数，协整检验在此基础上判断多个 $I(1)$ 序列有无长期均衡关系。
\end{algonotes}

%% ========== 双重差分（DID） ==========
\basicalgo{{\starmark}双重差分（DID）}{did_model}

\begin{algointro}
DID是评估政策效果的"王牌方法"。它同时用了"处理前后的差"和"处理组与对照组的差"，把两个差再相减，剥离掉时间趋势和固有差异，干净地估计出政策本身带来的效果。
\end{algointro}

\begin{algoscene}
\begin{itemize}
    \item 美赛中评估一项政策/措施的效果，如本案例2019年试点智能补货系统是否提升销售额
    \item 存在"部分对象接受处理、部分不处理"的准实验情境
    \item 有政策前后多年的面板数据，且处理组与对照组事前走势可比
    \item 需要回答"政策带来了多大净效应、能否推广"
\end{itemize}
\end{algoscene}

\begin{algoprinciple}
想象你想知道一门网课有没有用。直接比较"上过课的人"和"没上课的人"不公平（也许本来就更爱学习）；只看"上课前后"也不公平（也许这段时间整体都在进步）。DID的妙招是\textbf{两个差一起算}：

\begin{itemize}
    \item \textbf{第一个差}：处理组试点后减试点前——但这里面混了"随时间自然增长"。
    \item \textbf{第二个差}：对照组同期后减前——它只反映"随时间自然增长"。
    \item \textbf{双重差分}：把第一个差减去第二个差，自然趋势被抵消，剩下的就是政策的净效应（平均处理效应ATT）。
\end{itemize}

\textbf{核心前提——平行趋势}：政策实施前，处理组和对照组的结果走势必须大致平行。本案例用事件研究联合F检验验证这一点；前提成立，结论才可信。估计用双向固定效应回归，标准误在个体层面聚类。本案例政策使处理组销售额平均提高162.827（$p<0.001$）。
\end{algoprinciple}

\begin{algosposs}
\begin{enumerate}
    \item 在Dabbit AI SPSS工具箱中选择"双重差分DID"
    \item 准备面板数据：个体标识unit、时间year、结果变量sales、处理组标识treat
    \item 设置参数：
    \begin{itemize}
        \item 估计方法：双向固定效应回归
        \item 政策时间：2019年；聚类变量unit、个体层面聚类稳健标准误
        \item 平行趋势检验：事件研究联合F检验，参考期取政策前一期
        \item 安慰剂检验：100次置换
    \end{itemize}
    \item 运行，依次查看平行趋势图、ATT系数与安慰剂检验结果
\end{enumerate}

\textbf{SPSS实现}：DID需用面板/固定效应回归，SPSS原生面板支持有限，通过Dabbit AI SPSS工具箱或SPSS的R扩展（fixest/plm包）实现。
\end{algosposs}

\begin{algoresult}
\begin{itemize}
    \item \textbf{ATT（政策净效应）}：本案例为162.827，$p<0.001$，95\%置信区间[109.46, 216.19]——试点系统显著提升了销售额。
    \item \textbf{平行趋势检验}：政策前各期事件研究系数联合不显著（$p>0.05$），说明事前两组走势平行，前提成立。
    \item \textbf{安慰剂检验}：随机分配"伪处理组"重复100次，真实估计应明显落在伪估计分布之外，说明结论不是随机巧合。
    \item \textbf{动态效应}：事件研究图显示政策后处理组系数持续抬升，可观察效果随时间的变化。
\end{itemize}

\textbf{判断好坏}：ATT显著、平行趋势通过、安慰剂检验过关，三者齐备才是可信的政策评估。
\end{algoresult}

\begin{algonotes}
\textbf{什么时候用？}有处理组/对照组、有政策前后面板、且事前走势平行的政策评估。

\textbf{什么时候不用？}
\begin{itemize}
    \item 处理组和对照组事前走势明显不平行时，DID有偏
    \item 政策"一刀切"全员实施、没有对照组时，无法用DID
\end{itemize}

\textbf{常见坑}
\begin{itemize}
    \item \textbf{坑1}：不做平行趋势检验就报ATT，这是审稿人必挑的硬伤。
    \item \textbf{坑2}：标准误不聚类。面板数据误差在个体内相关，必须用个体层面聚类稳健标准误。
    \item \textbf{坑3}：把DID当"相关"，它识别的是因果，但前提是平行趋势成立。
\end{itemize}

\textbf{和相似算法的区别}：vs \textbf{多期DID}——处理在同一时点用标准DID，分批错时上线用多期DID；vs \textbf{PSM}——PSM解决"可观测特征不齐"，DID解决"随时间变化的混杂"，二者常结合（PSM-DID）。
\end{algonotes}

%% ========== 面板模型 ==========
\basicalgo{{\starmark}面板模型}{panel_model}

\begin{algointro}
面板数据同时有"横截面（多个门店）"和"时间（多年）"两个维度。面板模型能把门店天生的、不随时间变的差异（区位、铺面）和全行业共同的年度波动分离出去，得到各项经营投入对销售额的净影响。
\end{algointro}

\begin{algoscene}
\begin{itemize}
    \item 美赛中分析"多个主体、多个时期"的数据，如多家门店连续多年的经营记录
    \item 需要控制个体固有差异和共同时间冲击后估计各因素的净效应
    \item 决定该用固定效应、随机效应还是混合OLS时
    \item 本案例：客流、员工、促销对年销售额各有多大影响
\end{itemize}
\end{algoscene}

\begin{algoprinciple}
你有一张"门店 $\times$ 年份"的大表。问题在于：每家门店天生不同（有的在旺角），每年又有共同行情（某年整体生意好做）。面板模型就是把这两部分固定住：

\begin{itemize}
    \item \textbf{混合OLS}：把所有数据堆一起回归，假装门店之间没区别——最简单但最粗糙。
    \item \textbf{个体固定效应（FE）}：给每家门店一个"截距项"，吸收掉所有不随时间变的个体特征（区位、铺面）。这样比较的是"同一家门店自身变化"。
    \item \textbf{随机效应（RE）}：把个体差异当作随机扰动，更省参数但要求个体与解释变量不相关。
\end{itemize}

\textbf{怎么选}：用F检验、BP-LM检验和Hausman检验按 $\alpha=0.05$ 自动判定。本案例检验结果支持\textbf{个体固定效应}。结果：客流强度每+1单位，年销售额平均+4.769（$p<0.001$）；员工人数+1单位，+25.866；促销投入+1单位，+3.224（$p=0.019$）。
\end{algoprinciple}

\begin{algosposs}
\begin{enumerate}
    \item 在Dabbit AI SPSS工具箱中选择"面板模型"
    \item 设定：实体列门店编号、时间列年份、因变量年销售额、解释变量客流强度/员工人数/促销投入
    \item 设置参数：
    \begin{itemize}
        \item 模型选择：自动诊断选择（自动做F、LM、Hausman检验）
        \item 时间效应：自动判断；实体聚类稳健标准误
        \item 显著性水平0.05
    \end{itemize}
    \item 运行，查看模型选择结论与最终固定效应回归系数表
\end{enumerate}

\textbf{SPSS实现}：SPSS原生面板/固定效应支持有限，通过Dabbit AI SPSS工具箱或R扩展（plm包）实现。
\end{algosposs}

\begin{algoresult}
\begin{itemize}
    \item \textbf{模型选择结论}：报告先给出F检验、BP-LM、Hausman的结果，据此选定固定效应。
    \item \textbf{系数与p值}：客流4.769、员工25.866、促销3.224，均在0.05水平显著。
    \item \textbf{解释}：员工人数系数最大，说明增派人手对销售拉动最明显；客流、促销次之。
    \item \textbf{拟合与诊断}：关注 $R^2$、聚类稳健标准误与系数置信区间。
\end{itemize}

\textbf{判断好坏}：诊断检验明确选了哪种效应形式、核心系数显著且方向合理，即为合格的面板分析。
\end{algoresult}

\begin{algonotes}
\textbf{什么时候用？}既有横截面又有时间维度的"面板"数据，想控制个体与时间效应。

\textbf{什么时候不用？}纯横截面或纯时间序列数据，用不上面板模型。

\textbf{常见坑}
\begin{itemize}
    \item \textbf{坑1}：跳过Hausman检验随便选固定/随机效应。选错会导致系数不一致。
    \item \textbf{坑2}：忽略异方差和序列相关，普通标准误偏小，过度自信。应使用聚类稳健标准误。
\end{itemize}

\textbf{和相似算法的区别}：vs \textbf{动态面板模型}——本面板不含滞后因变量；若上一年销售本身影响今年，需用动态面板（GMM）。vs \textbf{基准回归}——面板控制了个体/时间固定效应，比混合OLS干净。
\end{algonotes}

%% ========== 自相关/偏自相关分析（ACF/PACF） ==========
\basicalgo{自相关/偏自相关分析（ACF/PACF）}{acf_pacf}

\begin{algointro}
ACF/PACF图是时间序列的"身份证"——它们画出当前值与过去各期值的相关程度，帮你判断这条序列内部有没有惯性、该用几阶的AR或MA模型来描述它。
\end{algointro}

\begin{algoscene}
\begin{itemize}
    \item 建立ARIMA/ARMA模型前的\textbf{定阶依据}（决定 $p,q$）
    \item 本案例：门店日销售额相邻日期为何连续同涨同跌，该用AR几阶
    \item 判断序列是否为白噪声（完全随机）
    \item 论文中说明模型阶数选择过程的必备图
\end{itemize}
\end{algoscene}

\begin{algoprinciple}
你想知道"今天的销售额和昨天、前天、大前天……到底有多像"。

\begin{itemize}
    \item \textbf{ACF（自相关函数）}：直接算 $y_t$ 与 $y_{t-k}$ 的相关系数。它把"中间各期的间接影响"也算进去了。
    \item \textbf{PACF（偏自相关函数）}：在剔除了中间各期影响后，单独看 $y_t$ 与 $y_{t-k}$ 的直接相关——像"控制其他变量后的净相关"。
    \item \textbf{看形状定阶}：PACF在 $p$ 阶后\textbf{截断}（立刻落入置信带）、ACF\textbf{拖尾}，提示用AR($p$)；反之ACF截断、PACF拖尾提示用MA($q$)；两者都拖尾则用ARMA。
\end{itemize}

本案例150个日销售额序列：ACF拖尾、PACF在2阶后截断，且ADF判定 $d=0$，故建议以 \textbf{AR(2)}（即ARIMA(2,0,0)）为候选。超出蓝色置信带的柱子才算显著相关。
\end{algoprinciple}

\begin{algosposs}
\begin{enumerate}
    \item 在Dabbit AI SPSS工具箱中选择"ACF/PACF分析"
    \item 上传按时间有序的单变量序列（本案例门店150个自然日日销售额）
    \item 设置参数：
    \begin{itemize}
        \item 时间列date、分析序列默认首个数值列
        \item 最大滞后阶数：自动（约min(n/3,40)，本案例滞后1--21）
        \item 显著性水平0.05；差分阶数自动（先ADF，本案例 $d=0$）
        \item 季节周期自动检测；输出AR/MA/ARMA定阶建议
    \end{itemize}
    \item 运行，查看ACF图、PACF图与自动定阶提示
\end{enumerate}

\textbf{SPSS中对应操作}：\textbf{分析 → 预测 → 自相关}，可输出ACF和PACF图。
\end{algosposs}

\begin{algoresult}
\begin{itemize}
    \item \textbf{ACF图}：柱子逐阶衰减拖尾，说明存在短期惯性。
    \item \textbf{PACF图}：前两阶显著（超出置信带），第3阶起落入带内——在2阶后截断。
    \item \textbf{Ljung-Box检验}：若 $p<0.05$，序列非白噪声、存在可利用的自相关结构。
    \item \textbf{自动定阶建议}：本案例推荐AR(2)。
\end{itemize}

\textbf{判断好坏}：图形形状规则、能明确判断截断/拖尾位置，定阶就有依据；若所有柱子都在带内，则序列是白噪声，无需建模。
\end{algoresult}

\begin{algonotes}
\textbf{什么时候用？}ARIMA/SARIMA建模前的定阶，以及判断序列自相关结构。

\textbf{什么时候不用？}横截面数据、或只关心因果而不做时间预测时。

\textbf{常见坑}
\begin{itemize}
    \item \textbf{坑1}：对不平稳序列直接看ACF。不平稳序列ACF会缓慢衰减、假装拖尾，必须先差分再看图。
    \item \textbf{坑2}：把ACF和PACF的截断/拖尾看反。记住"PACF截断→AR，ACF截断→MA"。
\end{itemize}

\textbf{和相似算法的区别}：vs \textbf{ADF检验}——ADF判断平不平稳，ACF/PACF判断自相关结构并定阶；vs \textbf{格兰杰因果}——ACF/PACF看单个序列的自相关，格兰杰看两个序列间的领先滞后。
\end{algonotes}

%% ========== 格兰杰因果检验 ==========
\basicalgo{格兰杰因果检验}{granger_test}

\begin{algointro}
格兰杰因果检验回答："A的历史数据，能不能帮助更好地预测B？"如果能，就说"A在预测意义上是B的格兰杰原因"。注意它是\textbf{预测因果}，不是哲学意义上的真因果。
\end{algointro}

\begin{algoscene}
\begin{itemize}
    \item 本案例：广告点击量的历史是否提升对成交额的预测，能否据此提前备货
    \item 判断两个时间序列谁领先谁、谁是先行指标
    \item 建模VAR前判断变量间的领先—滞后关系
    \item 区分"同期相关"与"可用的先行关系"
\end{itemize}
\end{algoscene}

\begin{algoprinciple}
直觉是：如果A真的先动、B后动，那么\textbf{知道A的过去}应该让你对B的预测更准。

\begin{itemize}
    \item \textbf{受限模型}：只用B自己的历史预测B。
    \item \textbf{非受限模型}：再加上A的历史一起预测B。
    \item \textbf{F检验比较}：如果加入A的历史后预测误差显著下降（$\Delta R^2$ 显著增大），就说A是B的格兰杰原因。
\end{itemize}

本案例：加入广告点击量(ads\_clicks)的历史后，对成交额(sales\_amount)的预测显著改善（$p=3.3\times10^{-25}$），即\textbf{点击量是成交额的格兰杰原因}——可据点击数据提前安排营销与备货。同时反向检验（成交额→点击量），并控制行业景气指数避免误判。注意：必须先ADF确认序列平稳，滞后阶数按AIC选择。
\end{algoprinciple}

\begin{algosposs}
\begin{enumerate}
    \item 在Dabbit AI SPSS工具箱中选择"格兰杰因果检验"
    \item 上传月度数据（本案例ads\_clicks、sales\_amount、行业景气指数）
    \item 设置参数：
    \begin{itemize}
        \item 原因序列cause、结果序列effect；时间列month
        \item 平稳性auto（自动ADF，不平稳则差分）；滞后准则aic
        \item 双向检验、p值校正none
    \end{itemize}
    \item 运行，查看各方向的F检验p值与 $\Delta R^2$
\end{enumerate}

\textbf{SPSS实现}：SPSS无原生格兰杰检验菜单，通过R/Python扩展（statsmodels的grangercausalitytests）实现。
\end{algosposs}

\begin{algoresult}
\begin{itemize}
    \item \textbf{F检验p值}：ads\_clicks→sales\_amount 为 $p=3.3\times10^{-25}<0.05$，存在格兰杰因果。
    \item \textbf{反向检验}：sales\_amount→ads\_clicks 不显著，则不存在反向预测关系。
    \item \textbf{$\Delta R^2$}：加入原因序列后 $R^2$ 的提升幅度，越大说明先行信息越有价值。
    \item \textbf{残差白噪声}：通过则模型设定可靠。
\end{itemize}

\textbf{判断好坏}：目标方向 $p<0.05$ 即存在预测意义上的领先关系；要结合业务解释，别把它当真正的物理因果。
\end{algoresult}

\begin{algonotes}
\textbf{什么时候用？}想找先行指标、判断两序列谁领先谁，或为VAR定关系。

\textbf{什么时候不用？}需要严格因果识别时（应配合DID/IV）；序列不平稳也不协整时直接检验会误判。

\textbf{常见坑}
\begin{itemize}
    \item \textbf{坑1}：把"格兰杰原因"当"真因果"。它只表示预测上有增量信息，可能有共同的第三方驱动。
    \item \textbf{坑2}：滞后阶数随便定。阶数选错结论会变，应按AIC选择并做敏感性。
    \item \textbf{坑3}：忽略不平稳。必须先差分/协整，否则是伪检验。
\end{itemize}

\textbf{和相似算法的区别}：vs \textbf{VAR}——格兰杰是VAR的一个后续诊断环节，VAR同时建模多变量并做脉冲响应；vs \textbf{ACF/PACF}——后者看单序列自相关，格兰杰看两序列间的跨序列领先关系。
\end{algonotes}

\section{进阶级算法速览}

%% ========== 异质性检验 ==========
\advancedalgo{异质性检验}{heterogeneity_test}

\begin{algobrief}
\textbf{一句话}：检验核心解释变量的效应是否因群体（如不同区域）而异，回答"该不该一刀切推广"。

\textbf{美赛场景区}：政策/措施在不同地区、人群、门店间效果是否一致，指导差异化投放。

\textbf{核心原理}：先分组子样本回归看各组系数，再用"组哑元 $\times$ 核心解释变量"交互项做Wald F检验，判断组间系数差异是否整体显著，辅以Chow检验交叉验证。本案例交互项Wald $F=9.402$（$p=10^{-5}$），"华东"组效应最大（$\beta=4.290$）。

\textbf{操作要点}：工具箱选"异质性检验"，指定结果变量、核心解释变量、分组变量，标准误HC1，参考组自动。

\textbf{结果关键}：交互项联合F检验 $p<0.05$ 即存在显著组间异质性；再看哪组系数最大。

\textbf{注意}：组内样本不能太小（否则系数不稳）；异质性检验显著后要进一步分组回归定位差异来源。
\end{algobrief}

%% ========== PSM ==========
\advancedalgo{倾向得分匹配（PSM）}{psm}

\begin{algobrief}
\textbf{一句话}：解决"处理组和对照组本来就不可比"的问题——给每个处理对象找一个各方面特征几乎一样的对照对象，再比较结果。

\textbf{美赛场景区}：评估培训、干预等非随机政策对工资/绩效的效果，消除自选择偏差。

\textbf{核心原理}：先用Logit模型估计每个对象接受处理的倾向得分，再按得分做最近邻匹配，用标准化均值差（SMD）检验匹配后协变量是否平衡，最后在平衡样本上估计ATT。本案例匹配后最大$|SMD|=0.071$（已平衡），ATT=742.7（95\%CI [463,1110]，$p<0.001$）。

\textbf{操作要点}：工具箱选"PSM"，指定处理列、结果列、协变量，方法最近邻匹配，estimand=ATT，bootstrap给置信区间。

\textbf{结果关键}：匹配后所有协变量$|SMD|<0.1$才算平衡；ATT显著为正。

\textbf{注意}：PSM只能平衡"可观测"特征，无法纠正不可观测混杂；若选不到共同支撑域内的对照，匹配失效。SPSS原生不支持，通过R/Python扩展实现。
\end{algobrief}

%% ========== Dagum基尼系数 ==========
\advancedalgo{Dagum基尼系数}{dagum_gini}

\begin{algobrief}
\textbf{一句话}：把总体收入差异分解成"组内差异、组间净差距、组间分布交叠"三部分，回答差异到底来自哪。

\textbf{美赛场景区}：区域协调发展评估中，衡量居民收入差异并定位主导来源。

\textbf{核心原理}：先算总体基尼系数，再沿三条路径分解——各组内部差异、各组均值水平差（组间净贡献）、组间分布交叠（超变密度），三者之和等于总体基尼。本案例总体基尼0.2051，组间净差距贡献率最高（71.06\%），差异主要体现为各组平均水平分层。

\textbf{操作要点}：工具箱选"Dagum基尼系数"，指定分组列、数值列，组均值降序排列，Bootstrap 200次给95\%置信区间。

\textbf{结果关键}：看三类贡献率谁最大——组间净差距主导说明是"分层"，组内主导说明是"组内分化"。

\textbf{注意}：需要足够的组内样本量；它是普通基尼系数的升级，能告诉你差异的结构来源。
\end{algobrief}

%% ========== 动态面板模型 ==========
\advancedalgo{动态面板模型}{dynamic_panel}

\begin{algobrief}
\textbf{一句话}：在面板模型里加入"上一期结果"作为解释变量，捕捉经济行为的惯性，并解决由此带来的内生性。

\textbf{美赛场景区}：销售额等指标存在跨期延续时，估计滞后效应与解释变量的短期净效应。

\textbf{核心原理}：滞后因变量与个体固定效应相关，OLS/固定效应都有偏。采用一阶差分GMM：先差分消去个体效应，再用滞后两期及更早的水平值作内部工具变量。本案例滞后因变量系数0.544（强惯性），traffic、promo短期效应均显著；AR(2)与Hansen过度识别检验支持工具设定。

\textbf{操作要点}：工具箱选"动态面板"，指定个体列、时间列、因变量、解释变量，估计方法two-step，缺失值strict。

\textbf{结果关键}：AR(2)检验$p>0.05$（无二阶序列相关）、Hansen J检验$p>0.1$（工具外生）才可采信。

\textbf{注意}：这是高级计量方法，SPSS原生不支持，需通过R扩展（panelvar/pdm）或Dabbit AI工具箱实现；工具不能过多（工具数应少于个体数）。
\end{algobrief}

%% ========== SFA ==========
\advancedalgo{随机前沿分析（SFA）}{sfa_model}

\begin{algobrief}
\textbf{一句话}：把实际产出与"理论最优产出前沿"的差距拆成"运气等随机波动"和"真实技术无效率"两部分，测算每个主体的效率。

\textbf{美赛场景区}：评估企业/合作社/医院的生产效率，识别低效主体作为帮扶对象。

\textbf{核心原理}：把误差写成随机噪声 $v$（正态）减技术无效率 $u$（半正态、非负），用极大似然估计前沿函数与两方差；$\gamma$ 衡量无效率占总误差的比重，LR检验无效率是否显著。本案例平均技术效率约0.899，$\gamma=0.855$，无效率显著。

\textbf{操作要点}：工具箱选"随机前沿分析"，指定产出、投入、主体编号、年份，函数形式log-log，优化器L-BFGS-B多起点。

\textbf{结果关键}：$\gamma$ 接近1说明差距主要来自真实效率损失；报告各主体技术效率并排序找低效者。

\textbf{注意}：需面板或截面数据且投入产出方向一致；与DEA的区别是SFA区分了随机噪声、对异常值稳健。
\end{algobrief}

%% ========== 门槛回归 ==========
\advancedalgo{门槛回归}{threshold_reg}

\begin{algobrief}
\textbf{一句话}：当解释变量对结果的影响跨过某个"临界值"后发生结构性突变时，用门槛回归自动找到这个临界值并分段回归。

\textbf{美赛场景区}：如客流量超过某阈值后，营销投入的销售拉动效应才显著增强。

\textbf{核心原理}：以门槛变量的实测值为候选网格逐点拟合分区模型，用残差平方和最小定位门槛值；再用自助法supF检验判断"无门槛/一个门槛/多个门槛"。本案例客流量门槛49.73（自助$p=0.003$），跨过它后促销系数由0.28跃升到1.89。

\textbf{操作要点}：工具箱选"门槛回归"，指定结果变量、门槛变量、协变量，自助法300次，残差重抽。

\textbf{结果关键}：门槛自助$p<0.05$确认存在门槛；报告门槛值置信区间与两侧系数差异。

\textbf{注意}：门槛两侧都要有足够样本；本质是分段线性，适合"阈值效应"类问题，SPSS原生不支持，需R/Python扩展。
\end{algobrief}

%% ========== 泰尔指数 ==========
\advancedalgo{泰尔指数}{theil_index}

\begin{algobrief}
\textbf{一句话}：从信息熵角度量化组间/组内差异，并能把总差异分解为组间贡献和组内贡献。

\textbf{美赛场景区}：评估各城区个体经营收入的均衡性，判断帮扶应侧重低水平城区还是城区内分化。

\textbf{核心原理}：泰尔T指数GE(1)，越大越不平等；存在分组时分解为组间差异与组内差异并算贡献率。本案例总泰尔指数0.2066（95\%CI [0.166,0.248]），组间贡献59.8\%、组内40.2\%。

\textbf{操作要点}：工具箱选"泰尔指数"，指定分组列、指标列，默认GE(1)，剔除非正值，自助1000次给置信区间。

\textbf{结果关键}：看组间/组内贡献率谁大——组间大说明是"区域水平差距"，组内大说明是"区域内分化"。

\textbf{注意}：指标必须为正；它与基尼系数同属不平等度量，但泰尔的可分解性更好。
\end{algobrief}

%% ========== Mann-Kendall趋势检验 ==========
\advancedalgo{Mann-Kendall趋势检验}{mk_test}

\begin{algobrief}
\textbf{一句话}：一种不要求数据正态、对异常值稳健的非参数趋势检验，判断序列是否显著上升/下降，并估计趋势斜率。

\textbf{美赛场景区}：水文水质、气象、环境监测序列的长期趋势诊断，如各断面COD是否显著下降。

\textbf{核心原理}：计算秩序列的Mann-Kendall统计量S与p值判断趋势显著性，再用Sen斜率估计趋势幅度；序列有自相关或周期时用方差修正/TFPW预白化。本案例3条断面序列中2条检出显著下降趋势。

\textbf{操作要点}：工具箱选"Mann-Kendall趋势检验"，分组列站点、双侧检验$\alpha=0.05$，自相关自动诊断、连续性修正。

\textbf{结果关键}：$p<0.05$ 且Sen斜率为负/正对应下降/上升趋势；报告斜率及其95\%置信区间。

\textbf{注意}：适用于偏态、有异常值的环境时间序列；它只判断单调趋势，不拟合具体曲线。
\end{algobrief}

%% ========== 频谱分析 ==========
\advancedalgo{频谱分析}{spectrum_analysis}

\begin{algobrief}
\textbf{一句话}：把时间信号从"时域"变换到"频域"，看能量集中在哪些频率、主周期是多少。

\textbf{美赛场景区}：设备振动故障诊断、周期信号识别，判断能量是否集中在转子工频及其倍频。

\textbf{核心原理}：用Welch平均法估计功率谱密度（hann窗），以背景中位功率为基准检测显著谱峰，峰值频率换算成周期（周期=1/频率）。本案例约8Hz主频占主导（主周期约0.125秒），该频段功率占总功率81.5\%。

\textbf{操作要点}：工具箱选"频谱分析"，指定时间列与信号列，Welch法，窗hann，重叠50\%，峰突出度阈值3倍背景。

\textbf{结果关键}：主峰频率与周期、各峰能量占比；主峰能量占比越高周期性越强。

\textbf{注意}：要求等间隔采样；它回答"周期成分多强"，与ACF/PACF的时域视角互补。
\end{algobrief}

%% ========== SUR ==========
\advancedalgo{似不相关回归（SUR）}{sur_model}

\begin{algobrief}
\textbf{一句话}：当同一批对象同时有多个相关的结果方程时，把它们联合估计，利用方程间的误差相关提高效率。

\textbf{美赛场景区}：同时估计促销对销售额、促销对成交笔数两个方程，且两误差可能同期相关。

\textbf{核心原理}：先逐方程OLS得残差，估计跨方程误差协方差矩阵，再用可行广义最小二乘（FGLS）堆叠联合求解；用Breusch-Pagan LM检验跨方程误差是否相关，相关才有联合估计的必要。本案例促销支出对销售额系数0.706（$p<0.001$）。

\textbf{操作要点}：工具箱选"似不相关回归"，逐方程设定因变量与自变量，FGLS两步，VIF阈值10。

\textbf{结果关键}：BP-LM检验$p<0.05$说明跨方程误差相关、SUR有效；报告各方程系数。

\textbf{注意}：若跨方程误差不相关，SUR与逐方程OLS几乎无差别；SPSS无原生菜单，需R扩展。
\end{algobrief}

%% ========== ECM ==========
\advancedalgo{误差修正模型（ECM）}{ecm_model}

\begin{algobrief}
\textbf{一句话}：在确认两序列协整后，量化短期偏离长期均衡后会以多快速度被"拉回"。

\textbf{美赛场景区}：如集装箱吞吐量与外贸规模的长期均衡关系、短期回调速度。

\textbf{核心原理}：先确认变量同为I(1)且协整，再以长期回归残差的滞后项作误差修正项，连同解释变量差分项一起回归；误差修正系数为负，表示偏离后向均衡回调。本案例短期偏离每期回调约31\%，约2.23期回调一半。

\textbf{操作要点}：工具箱选"误差修正模型"，指定x、y、时间列，最大滞后1阶，AIC选阶，含截距。

\textbf{结果关键}：误差修正项系数显著为负即机制成立；其大小决定回调速度与半衰期。

\textbf{注意}：必须先通过协整检验，否则ECM无意义；它把长期均衡和短期动态统一在一个模型里。
\end{algobrief}

%% ========== 信号滤波 ==========
\advancedalgo{信号滤波}{signal_filter}

\begin{algobrief}
\textbf{一句话}：在保留有用特征频率成分的同时，滤掉高频噪声和特定工频干扰，得到干净波形。

\textbf{美赛场景区}：轴承振动在线监测中保留轴转频特征、抑制电子噪声和50Hz工频干扰。

\textbf{核心原理}：先做频谱诊断确定主频与噪声频带，再选Butterworth低通/带通滤波器并定截止频率，采用零相位处理避免相位偏移，用保形相关系数、信噪比增益等指标验证。本案例Butterworth低通（order=4，截止30Hz）保形相关0.970，信噪比从11.6dB提升到21.0dB。

\textbf{操作要点}：工具箱选"信号滤波"，指定时间列与信号列，方法butter\_lowpass，必要时设截止频率。

\textbf{结果关键}：保形相关高、上频带抑制充分、信噪比增益大；滤后波形用于趋势判断。

\textbf{注意}：截止频率要根据频谱诊断设定，误删有用频带会丢信息；滤波应零相位处理。
\end{algobrief}

%% ========== 基尼系数 ==========
\advancedalgo{基尼系数}{gini_coef}

\begin{algobrief}
\textbf{一句话}：衡量收入/资源不平等程度的经典指标，0为完全平等、1为极端不平等，0.4是国际警戒线。

\textbf{美赛场景区}：居民家庭收入分配年度监测，判断总体不平等水平及与警戒线的关系。

\textbf{核心原理}：把家庭收入升序排列，以人口份额对收入份额作累积画出洛伦兹曲线，基尼系数=洛伦兹曲线与绝对平均线之间面积的两倍；用Bootstrap给置信区间。本案例基尼0.468（95\%CI [0.40,0.52]），已越过0.4警戒线。

\textbf{操作要点}：工具箱选"基尼系数"，指定数值字段income、权重字段weight，直接公式法，自举1000次，10组区间。

\textbf{结果关键}：与0.4警戒线比较；可分组算组内基尼定位差异来源。

\textbf{注意}：它只给总体不平等一个数，不告诉你差异结构；要分解结构用Dagum基尼或泰尔指数。
\end{algobrief}

%% ========== 分组回归 ==========
\advancedalgo{分组回归}{group_reg}

\begin{algobrief}
\textbf{一句话}：按某分类（如东中西部）分别回归，比较自变量对因变量的影响在各组间是否有结构性差异。

\textbf{美赛场景区}：评估营销投入对营业额的影响是否因区域而异，决定是否差异化投放。

\textbf{核心原理}：先分组建OLS，再用Chow检验对整体组间差异做F检验、用"组 $\times$ 自变量"交互项联合Wald检验定位差异来源；异方差时用稳健标准误。本案例Chow检验$p<10^{-63}$，最突出差异为"西部$\times$营销投入"（系数差-1.79）。

\textbf{操作要点}：工具箱选"分组回归"，指定因变量、自变量、分组变量，参考组按取值第一组，组间差异用Chow+交互Wald。

\textbf{结果关键}：Chow/交互Wald检验$p<0.05$即组间结构差异显著；报告各组系数对照。

\textbf{注意}：与异质性检验思路一致，但分组回归更侧重"逐组建模+整体结构差异检验"；每组样本要够。
\end{algobrief}

%% ========== 协整检验 ==========
\advancedalgo{协整检验}{coint_test}

\begin{algobrief}
\textbf{一句话}：判断多个各自不平稳（随机游走）的序列，是否存在一个长期稳定的均衡关系把它们"绑在一起"。

\textbf{美赛场景区}：如港口集装箱吞吐量与腹地外贸额是否长期联动、短期背离后能否回归均衡。

\textbf{核心原理}：先确认各序列同阶单整（通常都I(1)），再用Engle-Granger检验两两配对（看长期回归残差是否平稳）或Johansen迹检验确定协整个数。本案例两序列存在协整，短期偏离会向均衡回归。

\textbf{操作要点}：工具箱选"协整检验"，指定x、y列，方法eg（Engle-Granger），autolag=aic，trend=c，迹检验。

\textbf{结果关键}：$p<0.05$ 或迹统计量大于临界值即存在协整；之后可接ECM估计回调速度。

\textbf{注意}：前提是各序列同阶单整；不协整的I(1)序列直接回归就是伪回归。
\end{algobrief}

%% ========== VAR ==========
\advancedalgo{VAR向量自回归}{var_model}

\begin{algobrief}
\textbf{一句话}：把多个相互影响的时间序列放在一起，用彼此的滞后项联立建模，刻画它们的动态联动和冲击传导。

\textbf{美赛场景区}：如促销投入、客流、销售三者如何相互传导，一次投放多久后达到峰值、贡献多大。

\textbf{核心原理}：先ADF确认平稳（不平稳则差分），按AIC选滞后阶，估计VAR方程组并检验稳定性（特征根在单位圆内）与残差白噪声；再做格兰杰因果、脉冲响应、方差分解。本案例促销投入对客流有显著领先影响（$p\approx6\times10^{-8}$），冲击约第3期达峰值18.43，对客流贡献约40\%。

\textbf{操作要点}：工具箱选"VAR"，指定多个变量列、时间列，ic=aic，trend=c，设脉冲响应与方差分解期数。

\textbf{结果关键}：稳定性与残差白噪声通过；脉冲响应图看传导路径，方差分解看贡献占比。

\textbf{注意}：变量不宜过多、样本要长；它是"方程组"而非单一因果，结果用于刻画动态关系。
\end{algobrief}

%% ========== TSLS/2SLS ==========
\advancedalgo{两阶段最小二乘（TSLS/2SLS）}{tsls_2sls}

\begin{algobrief}
\textbf{一句话}：当解释变量"内生"（与误差相关）导致OLS有偏时，找一个只通过该解释变量影响结果的"工具变量"，分两步把因果效应干净地识别出来。

\textbf{美赛场景区}：如广告投放因"市场好的地方既多投又多卖"而内生，用广告点击成本作工具识别真实广告效应。

\textbf{核心原理}：第一阶段用工具变量预测内生变量，第二阶段用预测值回归因变量；用弱工具F检验（$F\geq10$）、Wu-Hausman内生性检验、Sargan过度识别检验把关。本案例广告投放系数2.569，第一阶段F=29.67（无弱工具），工具外生性通过，内生性显著。

\textbf{操作要点}：工具箱选"两阶段最小二乘"，指定因变量、内生变量、工具变量、控制变量，稳健标准误，输出OLS对照。

\textbf{结果关键}：第一阶段$F\geq10$、Sargan $p>0.05$、Hausman确认内生；以2SLS系数为准并与OLS比偏差方向。

\textbf{注意}：找工具变量是难点——必须满足"相关性+外生性（排他）"；SPSS原生不支持，通过R扩展实现。
\end{algobrief}

%% ========== 多期DID ==========
\advancedalgo{多期DID}{multi_period_did}

\begin{algobrief}
\textbf{一句话}：当处理对象在不同时点分批上线政策时（交错处理），用它估计平均处理效应，并校正传统双向固定效应可能的偏误。

\textbf{美赛场景区}：如2018、2020两批门店先后上线智能补货系统，部分门店始终不参与作对照。

\textbf{核心原理}：整理个体—时间面板，用事件研究检验平行趋势，再用双向固定效应估ATT，并补充Sun-Abraham队列加权估计作稳健性，全程个体聚类标准误。本案例政策使结果平均提高3.572（$p<0.001$，95\%CI [2.87,4.27]）。

\textbf{操作要点}：工具箱选"多期DID"，指定个体列、时间列、处理列、处理时点列，事件窗口[-4,+1]，个体聚类。

\textbf{结果关键}：政策前事件研究系数不显著（平行趋势）；ATT显著；Sun-Abraham结果与TWFE一致。

\textbf{注意}：交错处理下传统TWFE可能因"坏对照"有偏，务必做事件研究与稳健性对照；高级计量，需R扩展实现。
\end{algobrief}

%% ========== 方差分解 ==========
\advancedalgo{方差分解}{variance_decomp}

\begin{algobrief}
\textbf{一句话}：把结果变量的总波动按各因素（运营模式、面积档位、商圈类型等）分解，看谁解释了大部分差异。

\textbf{美赛场景区}：弄清门店间销售额差异主要由哪个经营因素造成，决定优先抓哪个环节。

\textbf{核心原理}：用方差分析把总平方和序贯拆成各因素与误差，各来源占比之和为100\%，再用F检验判断因素显著。本案例"商圈类型"贡献最大（54.92\%，显著）。

\textbf{操作要点}：工具箱选"方差分解"，指定因变量与多个因素变量，不纳入交互项，$\alpha=0.05$。

\textbf{结果关键}：按贡献率排序找主控因素；注意顺序（序贯分解）会影响各因素占比。

\textbf{注意}：它是描述+检验"谁解释波动"，不直接给因果；因素水平数和样本量要够。
\end{algobrief}

%% ========== RDD ==========
\advancedalgo{断点回归（RDD）}{rdd_model}

\begin{algobrief}
\textbf{一句话}：当政策资格由某"评分是否跨过一个明确临界值"决定时，比较断点两侧紧邻对象的结果，识别政策的局部因果效应。

\textbf{美赛场景区}：如技改评估达60分即获专项资助，比较60分上下企业次年人均产出差异。

\textbf{核心原理}：利用分配变量在断点处的确定性规则构造拟随机分组，用局部线性加权回归估计结果在断点处的跳跃；做密度连续性检验（有无操纵）、协变量平衡、安慰剂断点和带宽敏感性。本案例h=17.19下效应6.214（95\%CI [3.85,8.58]，$p=0.000$）。

\textbf{操作要点}：工具箱选"断点回归"，指定结果变量、分配变量、断点、处理变量，核triangular，带宽法cv。

\textbf{结果关键}：断点处跳跃显著、密度连续（无操纵）、不同带宽结论稳定。

\textbf{注意}：结论只在断点附近局部有效；精确断点（断点处处理完全切换）用sharp RDD，否则模糊RDD；需R扩展实现。
\end{algobrief}

%% ========== 差分分析 ==========
\advancedalgo{差分分析}{diff_analysis}

\begin{algobrief}
\textbf{一句话}：对持续上升/波动的时间序列逐期相减，判断要几阶差分才平稳，并量化逐期增减。

\textbf{美赛场景区}：如月度销售额持续攀升，判断差分阶数并刻画逐期变动，为预测和库存做准备。

\textbf{核心原理}：以ADF为主、KPSS为辅，从d=0逐阶检验，不平稳就差分一次直到平稳或达上限；再对差分序列做Ljung-Box白噪声检验。本案例sales需1阶差分后平稳（$p<0.0001$），差分后每期平均变动约3.15。

\textbf{操作要点}：工具箱选"差分分析"，指定时间列、值列，普通差分，阶数auto，ADF为主KPSS辅助。

\textbf{结果关键}：报告使序列平稳所需差分阶数d，以及差分后均值变动幅度。

\textbf{注意}：差分是ADF/ARIMA的前置步骤；差分丢失水平信息后不再直接代表原值，外推需累减还原。
\end{algobrief}

%% ========== GMM估计 ==========
\advancedalgo{GMM估计}{gmm_model}

\begin{algobrief}
\textbf{一句话}：当内生解释变量有多个工具变量（过度识别）时，用广义矩估计得到比2SLS更渐近有效的系数。

\textbf{美赛场景区}：如广告投放内生，用媒介成本指数、竞品广告密度等多个排他工具估计广告对销售的真实边际贡献。

\textbf{核心原理}：写成线性矩条件，两步GMM——第一步用普通工具权重（等价2SLS）得初值与残差，第二步用残差构造最优权重矩阵重估；工具数多于内生变量时用Hansen J检验外生性。本案例广告投放每+1单位带来销售+1.746，门店数+1单位、售价+1单位均显著。

\textbf{操作要点}：工具箱选"GMM估计"，指定被解释变量、内生解释变量、外生解释变量、工具变量，权重矩阵two\_step。

\textbf{结果关键}：一阶段F达弱工具阈值、Hansen J $p>0.1$（工具外生）、系数显著。

\textbf{注意}：GMM是2SLS在过度识别下的推广；工具要外生且相关，否则Hansen J会拒绝；需R/Python扩展实现。
\end{algobrief}

% 第7章 预测模型
\chapter{预测模型}

\begin{chapteroverview}
\textbf{这个分类是干什么的？}

预测模型回答的是一个向前看的问题：\textbf{基于历史数据，未来会变成什么样？}它把按时间顺序排列的观测（如月度销售额、周订单量、日收益率）当作一个"会自己延续下去的故事"，从中提炼出趋势、季节波动和随机惯性，再外推出未来若干期的数值。与第6章关注"变量之间是什么关系"不同，这一章更关注"这条时间序列本身接下来怎么走"。

\textbf{美赛什么题型常用？}

预测是美赛的高频任务：C题（数据驱动题）几乎每年都要对销量、客流、需求做短期预测；A题连续型问题里也常需要先预测人口、负荷、污染物浓度再做优化。只要题目出现"forecast、predict、estimate future demand、未来趋势"等字样，基本就落到这一章。常见产出包括未来6--12个月的点预测、预测区间和精度评估。

\textbf{学习路径建议}

先学\textbf{灰色预测GM(1,1)}和\textbf{指数平滑}——它们原理直观、对数据量要求低，小样本、趋势型序列很好用；再啃\textbf{ARIMA}（带星标，是时间序列预测的"主力模型"）。进阶部分了解\textbf{SARIMA}（处理明显季节周期）、\textbf{GARCH}（专门预测波动率而非水平值）和\textbf{马尔可夫预测}（预测状态而非数值）。选型时记住一句话：\textbf{先看数据有没有趋势、有没有季节、样本够不够多，再决定用哪个模型。}
\end{chapteroverview}

\section{基础级算法}

%% ========== 灰色预测GM(1,1) ==========
\basicalgo{灰色预测GM(1,1)}{grey_gm11}

\begin{algointro}
灰色预测GM(1,1)是一种专门对付\textbf{小样本、贫信息}的预测方法——只要你有四五个以上连续的、近似单调增长（或下降）的数据点，它就能帮你外推未来。它不要求数据服从什么分布，特别适合数据量不够、用不了复杂统计模型的场景。
\end{algointro}

\begin{algoscene}
\begin{itemize}
    \item 历史数据点很少（比如只有5--15个年度数据），却要做短期外推
    \item 美赛中刚起步的新兴业务量、新产品销量、某指标逐年爬升的早期预测
    \item 序列整体近似指数增长/衰减，随机波动不大
    \item 需要一个计算简单、能快速给出趋势外推参考的方法
\end{itemize}
\end{algoscene}

\begin{algoprinciple}
想象你手里只有几个零散的点，但它们大体在往一个方向走。灰色预测的巧思是：\textbf{不去直接预测那些抖动的原始数，而是先把它们"累加"成一条平滑上升的曲线}。

\begin{itemize}
    \item \textbf{第一步：累加生成（AGO）}。把原始序列 $x^{(0)}(1),x^{(0)}(2),\ldots$ 逐项累加，得到 $x^{(1)}(k)=\sum_{i=1}^{k}x^{(0)}(i)$。累加会把随机波动"抹平"，让新序列接近一条光滑的指数曲线——就像把起伏的海浪累加后变成缓慢抬升的海平面。
    \item \textbf{第二步：建一阶微分方程}。假设累加后的曲线服从 $\frac{dx^{(1)}}{dt}+a x^{(1)}=b$，其中 $a$ 叫发展系数（决定增长速度），$b$ 叫灰作用量（内生驱动项），用最小二乘把 $a,b$ 估出来。
    \item \textbf{第三步：还原预测}。解出微分方程的时间响应式得到 $x^{(1)}$ 的预测，再通过相邻两项相减（累减）还原回原始序列尺度的预测值。
\end{itemize}

\textbf{精度怎么看}：用平均相对误差、后验差比值 $C$ 和小误差概率 $p$ 评定精度等级（一级最好，二级合格，三级勉强，四级不合格）。$C$ 越小、$p$ 越接近1越好。
\end{algoprinciple}

\begin{algosposs}
\begin{enumerate}
    \item 在Dabbit AI SPSS工具箱中选择"灰色预测GM(1,1)"
    \item 上传等间隔的单变量序列（本案例为某生鲜渠道连续150期周订单量）
    \item 设置参数：
    \begin{itemize}
        \item 预测期数 \texttt{n forecast}：5（外推未来5周）
        \item 级比检验 \texttt{ratio test}：on（自动判断序列是否满足建模条件）
        \item 精度接受等级 \texttt{accept grade}：二级合格
        \item 残差修正 \texttt{residual correction}：off（必要时可开GM(1,1)残差修正提高精度）
    \end{itemize}
    \item 点击运行，查看拟合值、预测值与精度等级
\end{enumerate}

\textbf{默认参数参考}：变换方式 auto（自动），平移值 auto，预测5期。若级比落在可覆盖区间外，工具箱会自动做平移处理使序列满足建模条件。
\end{algosposs}

\begin{algoresult}
报告重点看三块：

\begin{itemize}
    \item \textbf{级比检验}：判断原始序列是否落在GM(1,1)的可用区间内。通过则可直接建模；不通过需平移变换。
    \item \textbf{发展系数 $a$}：$a<0$ 表示增长，$|a|$ 越大增长越快；一般 $-a<0.3$ 适合中短期预测，$0.3<-a<0.5$ 只宜短期预测。
    \item \textbf{预测值序列与精度}：本案例未来5期周订单量预测约为 14.24、14.38、14.52、14.67、14.81（万件），期末较首期上升约4.1\%。精度达到"二级合格"即可作为短期参考。
\end{itemize}

\textbf{判断好坏}：平均相对误差小于10\%为一级（好），10\%--20\%为二级（合格），大于20\%就要慎用。外推期数不要太多，GM(1,1)一般只适合预测未来3--5期。
\end{algoresult}

\begin{algonotes}
\textbf{什么时候用？}数据点少（4--20个）、序列单调近似指数、波动小的短期预测。

\textbf{什么时候不用？}
\begin{itemize}
    \item 序列剧烈震荡、无明显单调趋势时，GM(1,1)会给出误导性的平滑外推
    \item 样本量足够大（几十个以上）时，ARIMA、指数平滑等统计模型通常更准
    \item 有明显季节周期时不要单用GM(1,1)
\end{itemize}

\textbf{常见坑}
\begin{itemize}
    \item \textbf{坑1}：跳过级比检验直接建模。级比不合格会让模型系统性失真。
    \item \textbf{坑2}：把GM(1,1)当长期预测工具。它本质是指数外推，外推期数一多就严重脱离实际。
\end{itemize}

\textbf{和相似算法的区别}：vs \textbf{指数平滑}——两者都适合短期预测，但灰色预测专为小样本设计、能处理贫信息；指数平滑需要更长的序列且更擅长捕捉季节。vs \textbf{ARIMA}——ARIMA需要较大样本并定阶，灰色预测则"给点数据就能跑"。
\end{algonotes}

%% ========== 指数平滑 ==========
\basicalgo{指数平滑}{exp_smoothing}

\begin{algointro}
指数平滑是一种"越近的数据越重要"的预测方法——它给最近的观测最大权重，越久远的数据权重按指数衰减，从而自动跟踪序列的水平、趋势和季节变化。你只需要把数据交给它，它会自动帮你选最合适的平滑模型。
\end{algointro}

\begin{algoscene}
\begin{itemize}
    \item 有趋势和/或季节波动的月度、季度销售预测（如本案例2012--2024共150个月全渠道销售额）
    \item 美赛中需要快速给出未来若干期点预测和预测区间时
    \item 数据量中等（几十个以上）、希望模型自动拟合、不想手动定阶
    \item 备货、预算编制、排班等需要近期数值参考的场景
\end{itemize}
\end{algoscene}

\begin{algoprinciple}
想象你每天估明天的销售额：\textbf{昨天的实际值比一个月前的数据更有参考价值}。指数平滑就是按这个直觉设计的。

\begin{itemize}
    \item \textbf{简单指数平滑（一次）}：只适合没有趋势的平稳序列。预测值 $S_t=\alpha y_t+(1-\alpha)S_{t-1}$，其中 $\alpha$ 是平滑系数（0--1），$\alpha$ 越大越信新数据、越小越平滑。
    \item \textbf{Holt线性趋势（二次）}：在水平之外再加一个趋势分量 $\beta$，适合有上升/下降趋势但没有季节的序列。
    \item \textbf{Holt-Winters三次}：再加季节分量 $\gamma$，同时刻画趋势和年度季节周期，分加法季节（季节波动幅度恒定）和乘法季节（幅度随水平放大）。
\end{itemize}

三个参数 $\alpha,\beta,\gamma$ 都由计算机在历史数据上自动优化（最小化预留段预测误差RMSE），你不用手调。本案例序列有明显趋势和年度季节波动，因此自动选中\textbf{三次Holt-Winters加法季节}，得到 $\alpha=0.100,\beta=0.021,\gamma=0.570$。
\end{algoprinciple}

\begin{algosposs}
\begin{enumerate}
    \item 在Dabbit AI SPSS工具箱中选择"指数平滑"
    \item 上传含时间列和数值列的月度序列（本案例时间跨度2012年1月--2024年6月，共150个月）
    \item 设置参数：
    \begin{itemize}
        \item 平滑方法 \texttt{平滑方法}：auto（在一次/二次/三次模型间按预留误差自动择优）
        \item 趋势设定、季节设定：auto（自动诊断）
        \item 季节周期：auto（月度自动识别为12）
        \item 参数估计：自动优化；误差指标 RMSE
        \item 预测步数：12（预测未来12个月）
    \end{itemize}
    \item 运行并查看自动选中的模型类型、参数与预测区间
\end{enumerate}

\textbf{SPSS中对应操作}：\textbf{分析 → 预测 → 创建传统模型}，方法选"指数平滑"，并在"条件"中指定简单/Holt/Winters。
\end{algosposs}

\begin{algoresult}
\begin{itemize}
    \item \textbf{自动选中的模型}：报告开头会告诉你最终用了哪种平滑（本案例为三次Holt-Winters加法季节），以及 $\alpha,\beta,\gamma$ 三个参数值。
    \item \textbf{未来各期点预测}：本案例下一期点预测约4082.19万元。
    \item \textbf{95\%预测区间}：下一期约为 [3989.5, 4176.5]。区间越往后越宽，反映不确定性随预测期数增大。
    \item \textbf{残差诊断}：若残差自相关检验通过（接近白噪声），说明模型已把可预测的规律都捕捉到了。
\end{itemize}

\textbf{判断好坏}：比较预留段上的RMSE/MAE，残差无明显自相关即为合格。预测区间覆盖住真实值比例约95\%即说明区间合理。
\end{algoresult}

\begin{algonotes}
\textbf{什么时候用？}有趋势/季节的中短期序列预测，希望自动拟合、少调参。

\textbf{什么时候不用？}
\begin{itemize}
    \item 需要解释"哪个因素驱动了预测"时——指数平滑是纯外推，不解释因果
    \item 序列存在突变点或外部冲击（如疫情）时，平滑会反应滞后
\end{itemize}

\textbf{常见坑}
\begin{itemize}
    \item \textbf{坑1}：对有强趋势的序列用简单指数平滑，结果永远追不上实际值。
    \item \textbf{坑2}：忽略季节周期。月度数据一定要让模型自动识别季节（$m=12$），否则节假日、大促造成的波峰会被误判为噪声。
\end{itemize}

\textbf{和相似算法的区别}：vs \textbf{ARIMA}——指数平滑"上手快、自动选模型"，但理论框架不如ARIMA完整；ARIMA通过ACF/PACF定阶、更灵活，适合对精度要求更高的场景。两者在很多时候预测精度接近，可对比后择优。
\end{algonotes}

%% ========== ARIMA模型 ==========
\basicalgo{{\starmark}ARIMA模型}{arima_model}

\begin{algointro}
ARIMA是时间序列预测里最经典、最强大的"主力模型"。它把序列拆成三部分：差分让序列平稳（I）、用历史值回归现在（AR）、用历史误差修正现在（MA），三者组合起来捕捉序列的惯性，从而做出有统计依据的预测。
\end{algointro}

\begin{algoscene}
\begin{itemize}
    \item 美赛C题最常用的预测模型之一，用于月度/季度销售额、需求量预测
    \item 序列有趋势但无明显季节（有季节请用SARIMA，见进阶部分）
    \item 样本量足够（一般50个观测以上），需要给出带置信区间的短期预测
    \item 论文中需要严谨地报告定阶过程、残差白噪声检验和回测误差
\end{itemize}
\end{algoscene}

\begin{algoprinciple}
ARIMA写作 $\mathrm{ARIMA}(p,d,q)$，三个数字各管一件事：

\begin{itemize}
    \item \textbf{$d$（差分阶数，Integrated）}：先对序列做 $d$ 次差分把它变平稳。本案例月销售额持续上升，做1阶差分后平稳，故 $d=1$。
    \item \textbf{$p$（自回归阶数，AR）}：用\textbf{自身过去 $p$ 期}的值预测现在，即 $y_t$ 受 $y_{t-1},\ldots,y_{t-p}$ 影响。PACF图在 $p$ 阶后截断提示 $p$。
    \item \textbf{$q$（移动平均阶数，MA）}：用\textbf{过去 $q$ 期的预测误差}修正现在，即 $y_t$ 受 $\varepsilon_{t-1},\ldots,\varepsilon_{t-q}$ 影响。ACF图在 $q$ 阶后截断提示 $q$。
\end{itemize}

\textbf{定阶直觉}：先靠ACF/PACF图给出候选，再用AIC信息准则在 $(p,q)$ 网格里选最小者——AIC同时奖励拟合好、惩罚参数多，防止过拟合。本案例自动定阶得到 $\mathrm{ARIMA}(0,1,1)$，即一个"一阶差分+一阶移动平均"的简洁模型。
\end{algoprinciple}

\begin{algosposs}
\begin{enumerate}
    \item 在Dabbit AI SPSS工具箱中选择"ARIMA模型"
    \item 上传等间隔月度序列（本案例2013年1月--2025年6月共150期月销售额）
    \item 设置参数：
    \begin{itemize}
        \item 阶数 \texttt{order}：auto（在 $(p_{\max},d_{\max},q_{\max})$ 网格内按AIC自动择优）
        \item 信息准则 \texttt{ic}：aic；趋势项 auto
        \item 回测期 \texttt{test periods}：auto（留尾滚动回测评估真实误差）
        \item 预测步数：6（预测未来半年）
    \end{itemize}
    \item 运行并查看自动定阶结果、残差诊断与预测区间
\end{enumerate}

\textbf{SPSS中对应操作}：\textbf{分析 → 预测 → 创建传统模型}，方法选"ARIMA"，在"条件"中手动或自动指定 $p,d,q$。
\end{algosposs}

\begin{algoresult}
\begin{itemize}
    \item \textbf{最终阶数}：本案例为 ARIMA(0,1,1)。若工具箱给出多个候选，选AIC最小且残差白噪声的。
    \item \textbf{预测值与区间}：未来6期预测从6876.16变化到6999.92，95\%置信区间由 [6794.75, 6957.57] 逐步放宽到 [6709.97, 7289.87]——越远越不确定。
    \item \textbf{回测误差}：留尾滚动12步的 RMSE=45.67、MAE=34.09、MAPE=0.50\%。MAPE小于5\%即为精度良好。
    \item \textbf{残差白噪声检验}：Ljung-Box检验 $p=0.9857>0.05$，说明残差已是随机噪声，模型已无遗漏结构。
\end{itemize}

\textbf{判断好坏}：残差必须通过白噪声检验（否则说明还有规律没被模型抓住）；MAPE越小越好；预测区间应基本覆盖真实值。
\end{algoresult}

\begin{algonotes}
\textbf{什么时候用？}无明显季节、样本量充足的单变量序列短期预测，是美赛预测的"标配"。

\textbf{什么时候不用？}
\begin{itemize}
    \item 有强烈年度季节周期时，普通ARIMA抓不住，要用SARIMA
    \item 数据点太少（少于20）时定阶不稳，改用灰色预测或指数平滑
    \item 序列非平稳又无法差分平稳时，考虑协整/向量误差修正（见第6章）
\end{itemize}

\textbf{常见坑}
\begin{itemize}
    \item \textbf{坑1}：对非平稳序列直接建ARIMA。必须先做ADF检验，差分平稳后再定 $p,q$，否则是伪模型。
    \item \textbf{坑2}：只看AIC选了复杂高阶模型却不做残差白噪声检验。过拟合的模型在样本内很好、样本外很差。
    \item \textbf{坑3}：外推太多期。ARIMA是短期模型，一般预测3--6期，期数一多区间宽到失去意义。
\end{itemize}

\textbf{和相似算法的区别}：vs \textbf{指数平滑}——ARIMA有严格统计框架、可做假设检验，指数平滑更"开箱即用"；vs \textbf{SARIMA}——后者多了季节项 $(P,D,Q)_m$，处理月度/季度周期数据。
\end{algonotes}

\section{进阶级算法速览}

%% ========== GARCH模型 ==========
\advancedalgo{GARCH模型}{garch_model}

\begin{algobrief}
\textbf{一句话}：GARCH不预测序列的"数值水平"，而是专门预测它的\textbf{波动率（波动剧烈程度）}，刻画"大波动后面跟着大波动"的聚集现象。

\textbf{美赛场景区}：金融风控中预测股票/收益率的未来波动、计算VaR风险价值；可靠性与信号分析中估计噪声强度。

\textbf{核心原理}：先检验收益率是否存在ARCH效应（波动聚集），再用 $\mathrm{GARCH}(p,q)$ 同时估计均值方程和条件方差方程；波动持续性用 $\alpha+\beta$ 衡量，接近1表示波动会长期延续。本案例 $\alpha+\beta=0.9723$，下一期波动率约0.7218。

\textbf{操作要点}：工具箱选"GARCH"，指定收益率列与日期列，分布默认normal，均值模型常数，置信水平95\%/99\%，自动或手动指定 $p,q$。

\textbf{结果关键}：ARCH检验 $p<0.05$ 确认波动聚集显著；$\alpha,\beta$ 系数显著；标准化残差无残余ARCH效应；输出未来多期波动率与VaR分位数。

\textbf{注意}：GARCH只适用于围绕某均值上下波动的收益率类序列，不要拿去预测销售额水平；SPSS原生支持有限，建议通过SPSS的R扩展（arch包）或Dabbit AI工具箱实现。
\end{algobrief}

%% ========== 季节性ARIMA（SARIMA） ==========
\advancedalgo{季节性ARIMA（SARIMA）}{sarima_model}

\begin{algobrief}
\textbf{一句话}：SARIMA在普通ARIMA基础上增加一组"季节阶数"，专门处理以固定周期（如12个月、4个季度）重复波动的序列。

\textbf{美赛场景区}：月度客运量、月度销量等既含趋势又含年度季节周期的预测，如本案例城轨线路月客运量预测未来12个月。

\textbf{核心原理}：写作 $\mathrm{SARIMA}(p,d,q)(P,D,Q)_m$，其中 $m$ 是季节周期（月度=12），后三个大写阶数管"季节层"的自回归/差分/移动平均。先用季节分解和季节强度判断季节是否存在，再定 $(p,d,q)$ 与 $(P,D,Q)$。本案例为 $\mathrm{SARIMA}(1,1,1)(1,1,1)_{12}$，残差白噪声检验 $p=0.919$。

\textbf{操作要点}：工具箱选"SARIMA"，指定季节周期 $m$，非季节与季节阶数自动或手动设定，外推12期并输出95\%预测区间。

\textbf{结果关键}：残差通过Ljung-Box白噪声检验；首期预测1735.72（万人次），区间[1733.05, 1738.40]，较最近观测上升约2.8\%。

\textbf{注意}：必须有足够多的完整季节周期（至少2--3年月度数据）才能估计季节项；数据太少时季节阶数定不稳，应退回普通指数平滑。
\end{algobrief}

%% ========== 马尔可夫预测 ==========
\advancedalgo{马尔可夫预测}{markov_forecast}

\begin{algobrief}
\textbf{一句话}：马尔可夫预测面对的是"只有有限几种状态"的序列（如设备正常/异常/停机），它通过统计状态之间的转移概率，预测未来各状态出现的可能性。

\textbf{美赛场景区}：设备健康状态、天气等级、市场景气状态、信用等级等离散状态的演进预测；为检修安排、排产预案提供概率依据。

\textbf{核心原理}：先统计相邻时刻的状态转移频次，用最大似然估计一步转移概率矩阵；核心假设是"下一状态只依赖当前状态"（一阶马尔可夫性），用卡方检验验证。之后从当前状态出发逐期乘转移矩阵做多步预测，并求解稳态分布。本案例预测5步后"正常运行"概率约50.8\%，长期稳态约正常52.3\%、轻微异常22.7\%、明显劣化15.9\%。

\textbf{操作要点}：工具箱选"马尔可夫预测"，指定状态列（与可选实体列），初始分布默认取当前状态，预测步数5。

\textbf{结果关键}：马尔可夫性检验 $p<0.05$（转移确实依赖当前状态）；输出多步状态概率与长期稳态分布。

\textbf{注意}：若检验发现转移规律随时间变化（非齐次）或下一状态还依赖更早的历史，则一阶马尔可夫假设不成立，需要用时变或高阶马尔可夫模型；SPSS原生不支持，通过R/Python扩展实现。
\end{algobrief}

% 第8章 综合评价
\chapter{综合评价}

\begin{chapteroverview}
\textbf{这个分类是干什么的？}

综合评价解决的是"把多个指标合成一个总得分来排名"的问题。现实中评价一个对象往往要同时看好几个维度——评价门店要看坪效、毛利率、复购率、损耗率……每个指标量纲不同、方向不一，还可能有的"越大越好"有的"越小越好"。综合评价就是一套系统方法：先\textbf{给每个指标定权重}（谁更重要），再\textbf{把它们合成一个可比的总分}，最后排出优劣次序。

\textbf{美赛什么题型常用？}

综合评价是美赛\textbf{最高频}的建模题型之一。凡是题目出现"评价、排名、排序、优选、选最佳方案、确定指标权重"等字样，几乎都用这一章。其中\textbf{熵值法、AHP、TOPSIS}是三大核心方法，几乎每场比赛都至少用到一个；DEA用于效率评价，耦合协调度、障碍度模型用于发展质量诊断。

\textbf{学习路径建议}
先掌握\textbf{赋权}：客观赋权学\textbf{熵值法}（带星标，完全靠数据说话），主观赋权学\textbf{AHP层次分析法}（带星标，靠专家成对比较）；再学\textbf{CRITIC权重法}。然后学\textbf{TOPSIS}（带星标，与理想方案比距离排名）这个最常用的"打分器"，以及它和熵值法结合的\textbf{熵权TOPSIS}。\textbf{模糊综合评价}用于指标模糊、难以精确量化的场景。进阶部分了解DEA/SBM效率评价、BWM、耦合协调、障碍度、DEMATEL、ISM等专门工具。
\end{chapteroverview}

\section{基础级算法}

%% ========== 熵值法 ==========
\basicalgo{{\starmark}熵值法}{entropy_weight}

\begin{algointro}
熵值法是一种\textbf{完全客观}的定权方法：它不请专家打分，而是让数据自己说话——哪个指标在各评价对象之间差异越大、越能"拉开档次"，就给它越高的权重。最后按权重加权求和得到每个对象的综合得分并排名。
\end{algointro}

\begin{algoscene}
\begin{itemize}
    \item 美赛中需要\textbf{客观赋权}、避免主观偏好争议的综合评价
    \item 本案例：150家门店在坪效、毛利率、复购率、周转率、损耗率、投诉率六维下的经营质量排名
    \item 有较多样本（几十个以上）、各指标方向不一需要先同向化
    \item 论文中"客观权重确定"环节的标准方法
\end{itemize}
\end{algoscene}

\begin{algoprinciple}
"熵"本来是信息论里衡量\textbf{混乱程度}的量。熵值法的反直觉逻辑是：\textbf{一个指标在所有对象上取值都差不多，说明它区分不了谁好谁坏，信息量少、权重就低；反之取值越参差，越能拉开差距，权重就高}。

\begin{itemize}
    \item \textbf{第一步：方向同向化+归一化}。负向指标（如损耗率、投诉率，越低越好）先按"最大值平移"转成正向，再做min-max归一化到$[0,1]$。
    \item \textbf{第二步：算信息熵 $e_j$}。对第 $j$ 个指标，把各对象归一化取值转成比重 $p_{ij}$，算 $e_j=-k\sum p_{ij}\ln p_{ij}$。$e_j$ 越接近1，说明取值越均匀。
    \item \textbf{第三步：定权}。差异系数 $d_j=1-e_j$，再归一化 $w_j=d_j/\sum d_j$ 就是该指标权重。
    \item \textbf{第四步：合成得分}。$S_i=\sum_j w_j\cdot x'_{ij}$，把150家门店按 $S_i$ 降序排。
\end{itemize}

本案例权重：坪效0.284最高（门店间差距最大），其次库存周转率0.234、毛利率0.166；损耗率、投诉率因门店间差异小而权重低。
\end{algoprinciple}

\begin{algosposs}
\begin{enumerate}
    \item 在Dabbit AI SPSS工具箱中选择"熵值法"
    \item 上传评价数据：对象标识列选"门店编号"，纳入6个指标
    \item 设置参数：
    \begin{itemize}
        \item 指标方向：坪效/毛利率/会员复购率/库存周转率为正向；损耗率/顾客投诉率为负向
        \item 方向化方式：最大值平移；缺失处理整行删除
        \item 得分量纲：百分制换算；排序输出降序
        \item 常量指标权重记零并保留
    \end{itemize}
    \item 运行，查看权重表与综合得分排名
\end{enumerate}

\textbf{SPSS中实现}：SPSS无原生熵值法菜单，通过 \textbf{分析 → 描述统计} 配合变换计算，或用Dabbit AI SPSS工具箱一键完成。
\end{algosposs}

\begin{algoresult}
\begin{itemize}
    \item \textbf{权重表}：报告列出每个指标的熵值 $e_j$、差异系数 $d_j$、权重 $w_j$。本案例坪效权重0.284最高。
    \item \textbf{权重怎么解读}：权重高$\neq$指标业务上最重要，而是"它在样本间最能拉开差距"。这是熵值法的核心特点——客观但可能与直觉重要性不符。
    \item \textbf{综合得分排名}：本案例S115得分100分居首，S148得0分垫底，中间连续分布。
    \item \textbf{领先原因}：冠军店的优势主要落在权重最高（差异最大）的坪效等指标上。
\end{itemize}

\textbf{判断好坏}：权重之和=1，无指标权重为负；若某指标取值全相同（常量），其权重记0，应考虑剔除。
\end{algoresult}

\begin{algonotes}
\textbf{什么时候用？}需要客观、可复现、不掺杂专家主观的赋权排名；样本量较大。

\textbf{什么时候不用？}
\begin{itemize}
    \item 样本很少（几个对象）时，"差异"本身不稳定，权重不可靠
    \item 业务上某些指标虽差异小但极端重要（如安全性），纯数据赋权会低估它
\end{itemize}

\textbf{常见坑}
\begin{itemize}
    \item \textbf{坑1}：忘记先把负向指标同向化，直接拿"投诉率"算熵，方向就反了。
    \item \textbf{坑2}：归一化后出现0，导致 $\ln 0$ 报错——需加平移量。
    \item \textbf{坑3}：把"客观权重"当成"正确权重"。它反映区分度，不反映重要性，重要场合应与AHP等主客观方法结合。
\end{itemize}

\textbf{和相似算法的区别}：vs \textbf{AHP}——熵值法客观、靠数据，AHP主观、靠专家；vs \textbf{变异系数法/CRITIC}——都客观，但熵值法用信息熵，变异系数用标准差/均值，CRITIC兼顾对比强度与指标间冲突性。
\end{algonotes}

%% ========== AHP ==========
\basicalgo{{\starmark}层次分析法（AHP）}{ahp_method}

\begin{algointro}
AHP层次分析法是最经典的\textbf{主观赋权}方法。当指标间谁更重要难以直接量化时，它让专家只做"两两比较"（A和B哪个更重要、重要几倍），再把这些两两判断整理成一套权重，并检查专家判断是否自相矛盾。
\end{algointro}

\begin{algoscene}
\begin{itemize}
    \item 美赛中指标重要性无法从数据直接得出、需要专家经验判断时（选址、方案优选）
    \item 本案例：三家供应商在技术、价格、交付、服务、数据安全五准则下的择优
    \item 多准则决策、需要把定性经验转成定量权重
    \item 论文中"权重确定"环节，常与熵值法组合成主客观组合赋权
\end{itemize}
\end{algoscene}

\begin{algoprinciple}
让你直接给5个指标各打一个权重很难，但问你"技术和价格哪个更重要、重要多少倍"就容易多了。AHP就是把这些两两比较拼成一张判断矩阵，再反推出权重。

\begin{itemize}
    \item \textbf{构造判断矩阵}：用Saaty的1--9标度打分。$a_{ij}=1$ 表示两指标同等重要，$3$ 表示前者稍重要，$5$ 明显重要，$7$ 强烈重要，$9$ 极端重要；反之取倒数。
    \item \textbf{求权重}：对判断矩阵求最大特征根 $\lambda_{\max}$ 对应的特征向量，归一化即权重。
    \item \textbf{一致性检验（关键！）}：算一致性指标 $CI=(\lambda_{\max}-n)/(n-1)$，一致性比率 $CR=CI/RI$（$RI$ 是随机一致性指标，随阶数查表）。\textbf{$CR<0.1$ 才算专家判断前后一致、权重可用}；$CR\geq0.1$ 要回去改判断。
\end{itemize}

本案例准则层：技术能力权重0.317最高，价格0.260次之；综合后供应商A权重0.389最优。判断矩阵 $\lambda_{\max}=5.002$，$CR=0.001<0.1$，一致性很好。
\end{algoprinciple}

\begin{algosposs}
\begin{enumerate}
    \item 在Dabbit AI SPSS工具箱中选择"层次分析法AHP"
    \item 录入准则层两两判断（如"技术能力 vs 价格竞争力"的1--9标度）
    \item 若为双层结构，再逐准则录入"各方案 vs 该准则"的局部判断
    \item 设置参数：一致性检验阈值CR=0.10；专家判断聚合方式geometric（几何平均）
    \item 运行，查看权重向量、$\lambda_{\max}$、CI、CR一致性检验结果
\end{enumerate}

\textbf{SPSS中实现}：SPSS无原生AHP，通过Dabbit AI SPSS工具箱或R/Python扩展实现。
\end{algosposs}

\begin{algoresult}
\begin{itemize}
    \item \textbf{权重表}：每个准则的权重及排序。本案例技术能力0.317居首。
    \item \textbf{一致性比率CR}：本案例 $CR=0.001<0.1$，通过。若 $CR\geq0.1$，说明专家判断自相矛盾（如说A比B重要、B比C重要、C又比A重要），结果不可用。
    \item \textbf{综合权重}：方案综合权重 $=\sum$（准则权重 $\times$ 方案在该准则下局部权重）。本案例供应商A 0.389 $>$ B 0.319 $>$ C 0.292。
    \item \textbf{敏感性}：可对判断加小扰动看权重是否稳定；若两方案权重差小于0.03，视为排序接近。
\end{itemize}

\textbf{判断好坏}：CR$<$0.1是硬门槛；权重符合业务直觉、方案排序清晰即可采信。
\end{algoresult}

\begin{algonotes}
\textbf{什么时候用？}指标重要性主要靠经验判断、缺乏客观数据时；多方案比选。

\textbf{什么时候不用？}
\begin{itemize}
    \item 能拿到大量样本数据时，优先用熵值法等客观方法减少主观性
    \item 专家之间意见严重分歧时，AHP权重会不稳定
\end{itemize}

\textbf{常见坑}
\begin{itemize}
    \item \textbf{坑1}：不做一致性检验就报权重。CR超0.1的判断矩阵是无效的。
    \item \textbf{坑2}：指标过多（超过9个）时两两比较次数爆炸、一致性难保证。
    \item \textbf{坑3}：把AHP当客观方法。它本质是主观经验的量化，应说明专家构成。
\end{itemize}

\textbf{和相似算法的区别}：vs \textbf{熵值法}——AHP主观、熵值法客观，常组合赋权；vs \textbf{模糊层次分析FAHP}——AHP用精确数，FAHP用三角模糊数表达"大致重要"，更贴合专家模糊判断；vs \textbf{BWM}——BWM只需比较"最优/最劣"，比较次数更少、一致性更好。
\end{algonotes}

%% ========== CRITIC权重法 ==========
\basicalgo{CRITIC权重法}{critic_weight}

\begin{algointro}
CRITIC是一种比熵值法考虑更全面的客观赋权法：它不仅看指标自身的对比强度（标准差），还看指标之间的冲突性（相关程度）——既差异大、又与其他指标不重复的指标，权重最高。
\end{algointro}

\begin{algoscene}
\begin{itemize}
    \item 多个指标间可能信息重复，希望定权时兼顾"区分度"和"独立性"
    \item 本案例：150家供应商在报价、质量、交付、服务四维下的客观赋权
    \item 论文中想说明权重既考虑波动又考虑指标冗余
    \item 作为熵值法的改进版对照使用
\end{itemize}
\end{algoscene}

\begin{algoprinciple}
CRITIC考虑两个因素：

\begin{itemize}
    \item \textbf{对比强度}：指标标准差越大，对象间差异越大，信息量越大。
    \item \textbf{冲突性}：两个指标若高度相关（都在反映同一件事），信息重叠、贡献应打折；与其他指标相关越低，独立性越强、贡献越大。
    \item \textbf{信息量 $C_j$}$=\sigma_j\sum_k(1-r_{jk})$：标准差乘上"与其他指标不相关程度"之和。最后把 $C_j$ 归一化得到权重 $w_j$。
\end{itemize}

直觉：一个指标又能拉开差距、又不跟别人重复，才是真正"独当一面"的指标。
\end{algoprinciple}

\begin{algosposs}
\begin{enumerate}
    \item 在Dabbit AI SPSS工具箱中选择"CRITIC权重法"
    \item 上传数据，指定对象标识与各指标列、各指标方向
    \item 设置参数：标准化方式默认min-max；缺失值均值填补
    \item 运行，查看各指标标准差、相关阵、信息量 $C_j$ 与最终权重
\end{enumerate}
\end{algosposs}

\begin{algoresult}
\begin{itemize}
    \item \textbf{权重表}：列出每个指标的标准差、与其他指标的平均相关、信息量 $C_j$、权重 $w_j$。
    \item \textbf{解读}：标准差大且与其他指标相关低的指标权重最高；若某指标与别人高度相关，其权重被压下来。
    \item \textbf{后续}：把CRITIC权重代入TOPSIS或加权综合得分即可排名。
\end{itemize}
\end{algoresult}

\begin{algonotes}
\textbf{什么时候用？}指标间存在信息重叠、希望客观赋权同时去冗余。

\textbf{什么时候不用？}样本量小、相关阵估计不稳时慎用。

\textbf{常见坑}
\begin{itemize}
    \item \textbf{坑1}：不做指标方向同向化，负向指标混入会扭曲结果。
    \item \textbf{坑2}：混淆CRITIC与熵值法——前者兼顾相关性，后者只看离散度。
\end{itemize}

\textbf{和相似算法的区别}：vs \textbf{熵值法}——熵值法只用信息熵（离散度），CRITIC额外惩罚指标间相关；vs \textbf{独立性权系数法}——后者从复相关角度压冗余指标，思路相近。
\end{algonotes}

%% ========== 模糊综合评价 ==========
\basicalgo{模糊综合评价}{fuzzy_eval}

\begin{algointro}
很多评价是模糊的——"服务好不好"没法用一个精确数字衡量，只能说"很好/较好/一般/较差"。模糊综合评价用\textbf{隶属度}把这种"亦此亦彼"的模糊判断量化，再综合成一个评价等级。
\end{algointro}

\begin{algoscene}
\begin{itemize}
    \item 本案例：养老服务质量、满意度等难以精确打分的定性/半定性评价
    \item 指标评语是"优/良/中/差"等模糊等级
    \item 美赛中评价服务、生态、舒适度等带有主观模糊性的对象
    \item 需要给出"该对象综合属于哪一档"的结论
\end{itemize}
\end{algoscene}

\begin{algoprinciple}
想象让你给"服务水平"打分，你很难说它到底是8.3还是8.7分，但你可以说：\textbf{它"隶属"于"优"的程度是0.2、"良"0.5、"中"0.3}——这就是隶属度。

\begin{itemize}
    \item \textbf{第一步：定因素集和评语集}。因素集如\{可及性、专业性、规范性……\}，评语集如\{优、良、中、差\}。
    \item \textbf{第二步：建隶属度矩阵 $R$}。对每个因素，专家/数据给出它对各评语的隶属度（一行和为1）。
    \item \textbf{第三步：定权重向量 $W$}。可用AHP或熵值法得到各因素权重。
    \item \textbf{第四步：模糊合成}。$B=W\circ R$（矩阵乘法或模糊算子），得到对象对各评语的综合隶属度。
    \item \textbf{第五步：最大隶属度原则}。$B$ 中最大值对应的评语，就是综合评价等级。
\end{itemize}
\end{algoprinciple}

\begin{algosposs}
\begin{enumerate}
    \item 在Dabbit AI SPSS工具箱中选择"模糊综合评价"
    \item 录入评价因素、评语等级，以及各因素对各等级的隶属度
    \item 设定因素权重（可对接AHP/熵值法结果）
    \item 运行，查看综合隶属度向量与最终判定等级
\end{enumerate}
\end{algosposs}

\begin{algoresult}
\begin{itemize}
    \item \textbf{综合隶属度向量 $B$}：如 (0.2, 0.5, 0.2, 0.1)，分别对应优/良/中/差。
    \item \textbf{判定等级}：取最大分量对应等级——本例最大0.5在"良"，故综合评为"良"。
    \item \textbf{注意}：若最大两个分量很接近（如0.45 vs 0.42），说明对象处于两档之间，不要过度解读。
\end{itemize}
\end{algoresult}

\begin{algonotes}
\textbf{什么时候用？}评价标准模糊、定性指标多、需要给出等级判定时。

\textbf{什么时候不用？}指标全部精确可测、直接算分更清晰时，不必套模糊。

\textbf{常见坑}
\begin{itemize}
    \item \textbf{坑1}：隶属度矩阵每一行不归一（和不为1），合成结果失真。
    \item \textbf{坑2}：最大隶属度过低（如只占0.3）却硬判一档，应说明"接近临界"。
\end{itemize}

\textbf{和相似算法的区别}：vs \textbf{AHP}——AHP解决"权重"，模糊综合解决"模糊评价本身"，二者常组合（FAHP是把AHP模糊化）。
\end{algonotes}

%% ========== TOPSIS ==========
\basicalgo{{\starmark}TOPSIS法}{topsis_method}

\begin{algointro}
TOPSIS（逼近理想解排序法）是综合评价里\textbf{最常用的打分器}。它的思路很直观：在所有方案中找一个"完美方案（正理想解，每项都最好）"和一个"最差方案（负理想解，每项都最差）"，然后看每个对象\textbf{离谁更近}——离正理想越近、离负理想越远，就越好。
\end{algointro}

\begin{algoscene}
\begin{itemize}
    \item 美赛中\textbf{排名、优选、方案排序}的头号方法，几乎万能
    \item 本案例：150家门店在六项经营指标下的综合排序
    \item 多指标、多对象，需要一个连续可比的综合得分
    \item 常与熵值法组合成"熵权TOPSIS"，是论文高频组合
\end{itemize}
\end{algoscene}

\begin{algoprinciple}
想象一场选秀：评委心里有个"全能冠军模板"（每项都是全场最佳）和一个"全能垫底模板"（每项都最差）。每个选手看他\textbf{离冠军模板近不近、离垫底模板远不远}。

\begin{itemize}
    \item \textbf{第一步：向量归一化+加权}。把各指标归一化，乘上权重（等权或熵权）。
    \item \textbf{第二步：确定正负理想解}。正理想解 $Z^+$ = 各指标最优值，负理想解 $Z^-$ = 各指标最差值（先统一方向，负向指标越小越好）。
    \item \textbf{第三步：算距离}。每个对象到正理想解的欧氏距离 $D^+$，到负理想解的距离 $D^-$。
    \item \textbf{第四步：算贴近度} $C_i=\dfrac{D_i^-}{D_i^+ + D_i^-}$。$C_i$ 越接近1越接近正理想、越好；越接近0越差。
\end{itemize}

本案例在加权归一化空间中，对象M0067贴近度1.0000（即它本身就是正理想解）排名第一。
\end{algoprinciple}

\begin{algosposs}
\begin{enumerate}
    \item 在Dabbit AI SPSS工具箱中选择"TOPSIS法"
    \item 上传数据，指定对象标识列与各指标列
    \item 设置参数：
    \begin{itemize}
        \item 指标方向：正向（越大越好）/负向（越小越好，先同向化）
        \item 权重：等权，或接入熵权/CRITIC权重
        \item 归一化方式：向量归一化
    \end{itemize}
    \item 运行，查看正/负理想解、各对象距离 $D^+,D^-$ 与贴近度 $C$ 排名
\end{enumerate}

\textbf{SPSS中实现}：SPSS无原生TOPSIS，通过Dabbit AI SPSS工具箱或R/Python扩展实现。
\end{algosposs}

\begin{algoresult}
\begin{itemize}
    \item \textbf{正/负理想解}：列出各指标的最优、最差值。
    \item \textbf{$D^+,D^-$}：每个对象到两个理想解的距离。
    \item \textbf{贴近度 $C$}：核心结果，$C\in[0,1]$，按 $C$ 降序排名。本案例M0067的 $C=1.0000$（与正理想重合）居首。
    \item \textbf{稳健性}：可做留一敏感性（去掉某指标重排），看名次是否稳定。
\end{itemize}

\textbf{判断好坏}：$C$ 是相对值，只在本样本内可比；$C=1$ 即该对象在所有指标上都最优。
\end{algoresult}

\begin{algonotes}
\textbf{什么时候用？}多对象多指标的排名优选，是综合评价的"通用底座"。

\textbf{什么时候不用？}
\begin{itemize}
    \item 指标间存在不可公度、或有缺失严重时，先清洗
    \item 需要判断"属于哪一档"而非连续排名时，模糊综合或RSR更合适
\end{itemize}

\textbf{常见坑}
\begin{itemize}
    \item \textbf{坑1}：负向指标不同向化就直接算距离，方向错误。
    \item \textbf{坑2}：用等权却声称"客观"。等权也是一种主观假设，建议配合熵权。
    \item \textbf{坑3}：只看贴近度绝对值。$C$ 是相对排名工具，不同批次样本的 $C$ 不能直接比大小。
\end{itemize}

\textbf{和相似算法的区别}：vs \textbf{熵权TOPSIS}——本TOPSIS权重可等权可主观，熵权TOPSIS用熵值法客观定权；vs \textbf{VIKOR}——TOPSIS看距离理想解的贴近度，VIKOR是折中排序，引入群体效用与个体遗憾。
\end{algonotes}

%% ========== 熵权TOPSIS ==========
\basicalgo{熵权TOPSIS}{entropy_topsis}

\begin{algointro}
熵权TOPSIS是\textbf{熵值法定权 + TOPSIS排序}的强强组合：先用熵值法从数据客观算出各指标权重，再把这套权重喂给TOPSIS算贴近度排名。它兼顾了客观赋权和距离排序，是美赛综合评价的"黄金组合"。
\end{algointro}

\begin{algoscene}
\begin{itemize}
    \item 美赛中最常被推荐的综合评价"标准答案"之一
    \item 本案例：150家连锁门店在销售额、坪效、复购率、满意度、成本、损耗六维下的运营绩效排名
    \item 既不想主观定权、又想要连续可比的贴近度得分
    \item 论文中权重客观、排序清晰、可做稳健性检验
\end{itemize}
\end{algoscene}

\begin{algoprinciple}
两步走：

\begin{itemize}
    \item \textbf{第一步：熵法定权}。对各指标算信息熵 $e_j$、效用值 $d_j=1-e_j$、熵权 $w_j=d_j/\sum d$。某指标对象间差异越大、熵越小、权重越高。本案例商品损耗率权重0.275最高、顾客满意度0.005最低。
    \item \textbf{第二步：TOPSIS排序}。把归一化矩阵乘熵权，找正/负理想解，算 $D^+,D^-$，贴近度 $C=D^-/(D^++D^-)$。
\end{itemize}

负向指标（运营成本、商品损耗率）先按最大值差法 $max-x$ 同向化。本案例M0067贴近度1.0000居首，M0100次之0.810。
\end{algoprinciple}

\begin{algosposs}
\begin{enumerate}
    \item 在Dabbit AI SPSS工具箱中选择"熵权TOPSIS"
    \item 上传数据，指定对象标识列（门店编号）
    \item 设置参数：
    \begin{itemize}
        \item 方向映射：商品损耗率、运营成本为负向（negative），其余正向
        \item 负向变换：max\_minus；归一化：vector向量归一化
        \item 缺失处理：median\_impute中位数填补
        \item 敏感性：leave\_one\_out留一法
    \end{itemize}
    \item 运行，查看熵权表、贴近度排名与留一稳健性
\end{enumerate}
\end{algosposs}

\begin{algoresult}
\begin{itemize}
    \item \textbf{熵权表}：本案例商品损耗率0.275、运营成本0.243权重最高（最能区分门店），顾客满意度0.005最低。
    \item \textbf{贴近度排名}：M0067为1.0000、M0100为0.8100……按 $C$ 降序。
    \item \textbf{留一稳健性}：逐一去掉指标重排，若头部名次稳定则结论可信。
\end{itemize}

\textbf{判断好坏}：权重反映区分度（非重要性）；排名对留一指标不敏感即稳健。
\end{algoresult}

\begin{algonotes}
\textbf{什么时候用？}希望权重完全客观、又要连续贴近度排名时的首选。

\textbf{什么时候不用？}样本很少、或业务上必须体现某些指标重要性时，纯熵权会失真。

\textbf{常见坑}
\begin{itemize}
    \item \textbf{坑1}：忽略负向指标同向化，成本/损耗方向算反。
    \item \textbf{坑2}：把熵权当重要性权重。它只代表区分度，低权重的满意度可能业务上极重要。
\end{itemize}

\textbf{和相似算法的区别}：它是熵值法与TOPSIS的串联；若想权重更贴合业务，可换成AHP权重再代入TOPSIS。
\end{algonotes}

\section{进阶级算法速览}

%% ========== SBM ==========
\advancedalgo{SBM模型}{sbm_model}

\begin{algobrief}
\textbf{一句话}：考虑了"松弛变量"和"非期望产出"的DEA效率评价模型，能同时处理少投入、多产出、少排污。

\textbf{美赛场景区}：绿色生产效率评价，如制造业企业在投入人力能耗、产出产值的同时排放二氧化碳。

\textbf{核心原理}：用Tone的非径向SBM模型，把投入冗余、期望产出不足、非期望产出过多的松弛量直接计入目标函数，经线性规划逐单元求效率。本案例150家企业平均效率0.731，21家有效，冗余主要在碳排放（松弛率26.2\%）。

\textbf{操作要点}：工具箱选"SBM"，指定投入、期望产出、非期望产出，VRS，非期望产出弱可处置。

\textbf{结果关键}：效率得分0--1，越接近1越好；看松弛率定位冗余来源（投入侧还是排放侧）。

\textbf{注意}：效率是组内相对值，跨期跨组不可直接比；需线性规划求解，SPSS原生不支持，用R/Python扩展。
\end{algobrief}

%% ========== 变异系数法 ==========
\advancedalgo{变异系数法（信息量权重法）}{cv_weight}

\begin{algobrief}
\textbf{一句话}：用"标准差/均值"（变异系数）衡量指标在样本间的相对波动，波动越大权重越高的客观赋权法。

\textbf{美赛场景区}：完全不引入主观打分，依据数据自身区分信息为门店经营指标定权。

\textbf{核心原理}：逐列算均值与标准差，变异系数$=$标准差/均值，归一化为权重，再对极差归一化数据加权求和。本案例"线上销售占比"权重最高（0.2805），综合得分最高对象M1142（0.7068）。

\textbf{操作要点}：工具箱选"变异系数法"，指标列自动识别，负向指标倒数变换，缺失均值填补。

\textbf{结果关键}：变异系数大的指标权重高；注意要求指标取值为正才能算CV。

\textbf{注意}：与熵值法同为客观赋权，思路相近但CV用离散系数、熵值用信息熵；量纲不同时要先同向化。
\end{algobrief}

%% ========== 模糊层次分析法 ==========
\advancedalgo{模糊层次分析法}{fahp_method}

\begin{algobrief}
\textbf{一句话}：把AHP中专家的"点估计"判断升级为"三角模糊数(l,m,u)"，表达"大概在这区间"的模糊判断，更贴近真实决策心理。

\textbf{美赛场景区}：多准则供应商择优，专家难以用精确数比较、判断带模糊性时。

\textbf{核心原理}：成对比较表示为三角模糊数$(l,m,u)$，用中值矩阵做一致性检验（CR$<$0.1），再用Buckley几何平均法算模糊权重、去模糊化归一化排序，并做对数扰动敏感性。本案例供应商A综合权重0.389最优，技术能力准则权重0.317。

\textbf{操作要点}：工具箱选"模糊层次分析"，去模糊化centroid，权重geometric\_mean，CR阈值0.10，专家几何聚合。

\textbf{结果关键}：CR$<$0.1通过；报告模糊权重区间与去模糊权重，权重差小于0.03视为排序接近。

\textbf{注意}：它是AHP的模糊升级版，适合判断不确定的场景；计算更复杂，需R/Python扩展。
\end{algobrief}

%% ========== RFM ==========
\advancedalgo{RFM客户价值模型}{rfm_model}

\begin{algobrief}
\textbf{一句话}：用最近消费时间(R)、消费频率(F)、消费金额(M)三个维度给客户打分分层，识别高价值客户。

\textbf{美赛场景区}：电商/零售客户分群与差异化营销，如找出贡献大部分营收的核心会员。

\textbf{核心原理}：从交易流水按客户聚合出R/F/M，分位数各打1--5分，以每维中位数为界分高低，组合成八类客群（重要价值/发展/保持/挽留等）。本案例重要价值客户占29.8\%却贡献68.3\%金额。

\textbf{操作要点}：工具箱选"RFM"，指定客户编号、消费日期、金额列，参考日期取数据最近日，等权，中位数分界。

\textbf{结果关键}：各客群的数量占比、金额占比、价值强度；帕累托曲线看集中度。

\textbf{注意}：R越小越好（越近越活跃）、F和M越大越好；权重默认等权，重要业务可调整。
\end{algobrief}

%% ========== VIKOR ==========
\advancedalgo{VIKOR法}{vikor_method}

\begin{algobrief}
\textbf{一句话}：一种"折中"排序法，在群体效用最优和个体遗憾最小之间权衡，给出妥协最优解。

\textbf{美赛场景区}：多指标供应商/方案选优，指标相互冲突、需要可接受的折中方案时。

\textbf{核心原理}：把各方案到理想点的加权距离拆成群体效用$S$和个体遗憾$R$，以折中系数$v=0.5$合成$Q$，再用可接受优势和稳定性两个条件裁决唯一最优。本案例S090的$Q=0$为唯一折中最优。

\textbf{操作要点}：工具箱选"VIKOR"，等权或自定义权重，极差归一化，$v=0.5$，方向自动识别。

\textbf{结果关键}：$Q$越小越好；需同时通过可接受优势（$Q$差达阈值）与稳定性（$S$或$R$也最优）检验。

\textbf{注意}：与TOPSIS同属多属性决策，但VIKOR显式引入"遗憾"和妥协条件，适合冲突型方案选择。
\end{algobrief}

%% ========== 功效系数法 ==========
\advancedalgo{功效系数法}{efficacy_coef}

\begin{algobrief}
\textbf{一句话}：把每个指标按"不允许值→满意值"线性映射到60--100分，再加权汇总成百分制综合得分。

\textbf{美赛场景区}：门店服务质量、政府绩效等需要百分制打分和分等级（优/良/中/及格）的评价。

\textbf{核心原理}：为每个指标定满意值（95分位）和不允许值（5分位），单项功效系数$=60+40\times$(实际-不允许)/(满意-不允许)，越界截断，加权求和。本案例门店056得分98.47（优），门店140为62.94（及格）。

\textbf{操作要点}：工具箱选"功效系数法"，指定正向/逆向指标，满意/不允许分位点0.95/0.05，范围60--100，越界截断。

\textbf{结果关键}：综合得分+等级+短板指标；平均功效系数低的指标是共性短板。

\textbf{注意}：基准值（满意/不允许）设定主观，要说明依据；它输出百分制便于管理沟通。
\end{algobrief}

%% ========== BWM ==========
\advancedalgo{最好最差法（BWM）}{bwm_method}

\begin{algobrief}
\textbf{一句话}：只需专家指定"最好"和"最差"准则，分别比较它们与其他准则，比较次数比AHP少、一致性更好的主观赋权法。

\textbf{美赛场景区}：多专家参与的指标权重确定，如30人评审组确定供应商评价准则权重。

\textbf{核心原理}：专家给出最优准则到各准则的BO向量、各准则到最劣准则的OW向量，用最小化最大偏差的线性规划求权重，以$CR=\xi^*/CI$检验一致性。本案例产品质量权重0.421最高，30组判断一致性均值0.0083全部通过。

\textbf{操作要点}：工具箱选"BWM"，指定best/worst准则，几何聚合，一致性阈值0.10。

\textbf{结果关键}：$CR<0.1$通过；权重和为1；报告专家间分歧幅度。

\textbf{注意}：是AHP的改进，比较次数从$n(n-1)/2$降到$2n-3$；仍属主观赋权，需R/Python扩展。
\end{algobrief}

%% ========== 障碍度模型 ==========
\advancedalgo{障碍度模型}{obstacle_model}

\begin{algobrief}
\textbf{一句话}：不只算综合得分，更专门诊断"到底是哪些指标在拖后腿、阻碍系统提升"。

\textbf{美赛场景区}：县域发展、城市宜居等综合评价的短板诊断，回答"该优先抓哪个指标/子系统"。

\textbf{核心原理}：极差标准化后用"1-标准化值"表示偏离最优的差距，障碍度$=$偏离度$\times$权重占全部加权偏离的比例；累计占比首次达60\%的指标为主要障碍因子。本案例"乡村旅游收入占比"障碍度最高（17.12\%），产业融合子系统合计35.4\%。

\textbf{操作要点}：工具箱选"障碍度模型"，指定对象列、指标列、准则层归属，权重entropy，缺失均值填补。

\textbf{结果关键}：按平均障碍度降序找主要障碍因子；子系统障碍度找短板模块。

\textbf{注意}：它与综合得分互补——得分告诉你"谁差"，障碍度告诉你"差在哪、先改哪"。
\end{algobrief}

%% ========== 德尔菲法 ==========
\advancedalgo{德尔菲法（Delphi）}{delphi_method}

\begin{algobrief}
\textbf{一句话}：通过多轮匿名征求专家意见、逐轮反馈统计结果，直到意见收敛，从而筛选出高共识的评价条目。

\textbf{美赛场景区}：构建评价指标体系之初，筛选哪些指标真正重要、专家意见一致。

\textbf{核心原理}：多轮匿名打分（1--5分），每轮算均值、满分频率、变异系数、四分位距，用Kendall协调系数$W$衡量专家整体协调，按均值和变异系数双阈值筛选条目，直至收敛。本案例3轮后$W=0.874$，保留3/5条目。

\textbf{操作要点}：工具箱选"德尔菲法"，录入多轮专家评分，均值阈值3.5、CV阈值0.25，Kendall $W$检验。

\textbf{结果关键}：末轮$W$接近1且显著、均值位移小说明收敛；保留达标的条目。

\textbf{注意}：它是指标\textbf{筛选}工具，本身不算综合得分；常作为AHP/模糊综合的前置步骤。
\end{algobrief}

%% ========== 灰色关联分析 ==========
\advancedalgo{灰色关联分析}{grey_relational}

\begin{algobrief}
\textbf{一句话}：以一个参考序列为基准，算各比较序列与它的"形状相似度"，判断哪些因素与结果最同步。

\textbf{美赛场景区}：小样本下识别哪些经营因素与销售额走势最贴近、谁是主导驱动。

\textbf{核心原理}：先极差标准化，指定参考列（如日销售额），逐点算绝对差和关联系数，加权平均成灰色关联度并排序，分辨系数$\rho=0.5$。本案例日客流量关联度最高（$r=0.917$）。

\textbf{操作要点}：工具箱选"灰色关联分析"，指定参考列，分辨系数0.5，等权，极差标准化。

\textbf{结果关键}：关联度越大越贴近参考；$\rho$在0.1--0.9变化排序稳定即可靠。

\textbf{注意}：它衡量"形状同步"而非因果，适合样本少、信息不完全的小样本系统分析。
\end{algobrief}

%% ========== 独立性权系数法 ==========
\advancedalgo{独立性权系数法}{independence_weight}

\begin{algobrief}
\textbf{一句话}：认为与其他指标重叠（相关）越多的指标冗余越重、权重应越小，按信息独立性客观定权。

\textbf{美赛场景区}：指标间高度相关（如客流、销售额、坪效互相牵连）时，给门店评价定一组不重复信息的权重。

\textbf{核心原理}：算指标相关阵，以复相关系数衡量该指标被其余指标线性解释的程度，按$1/$复相关归一化得权重。本案例顾客满意度复相关最低（0.766）获最高权重18.6\%。

\textbf{操作要点}：工具箱选"独立性权系数法"，标准化minmax，pearson相关，复相关口径。

\textbf{结果关键}：复相关越低越独立、权重越高；权重反映信息独立性而非业务重要性。

\textbf{注意}：与CRITIC思路相近（都惩罚相关），但独立性权直接用复相关，CRITIC用标准差$\times$(1-相关)。
\end{algobrief}

%% ========== ISM ==========
\advancedalgo{解释结构模型（ISM）}{ism_model}

\begin{algobrief}
\textbf{一句话}：把一堆相互影响的因素，通过邻接矩阵和可达矩阵，梳理成"表层结果—中层传导—深层根源"的层级结构。

\textbf{美赛场景区}：分析影响工程质量/系统性能的众多因素，找出哪些是根源、哪些是表象。

\textbf{核心原理}：由专家两两判断建0/1邻接矩阵，加单位阵做布尔传递闭包得可达矩阵，用"可达集$\subseteq$先行集"逐层剥离分层，再做MICMAC驱动-依赖分类。本案例20因素分5层，高层领导重视、制度建设是根源驱动。

\textbf{操作要点}：工具箱选"ISM"，录入因素间二元关系，一致率阈值0.5，含MICMAC分类。

\textbf{结果关键}：层级图（表层L1到根源L5）；驱动型因素优先治理，依赖型只是结果。

\textbf{注意}：它梳理因素\textbf{结构关系}，不打分排名；判断矩阵来自专家，需多人共识。
\end{algobrief}

%% ========== 耦合协调度模型 ==========
\advancedalgo{耦合协调度模型}{coupling_coord}

\begin{algobrief}
\textbf{一句话}：衡量两个（或多个）子系统之间"相互配合、协调发展"的程度，输出协调等级和短板系统。

\textbf{美赛场景区}：如城镇化与生态环境、经济与社会两系统的协同发展评估，面板时序数据。

\textbf{核心原理}：子系统内用熵权合成综合指数$U$，算耦合度$C$、综合发展指数$T$、耦合协调度$D=\sqrt{C\times T}$，按0.1间隔分十级（失调—协调），按指数差判断滞后系统。本案例$D$均值0.693（中级协调），逐期+0.014。

\textbf{操作要点}：工具箱选"耦合协调度"，指定时间/个体列、各子系统指标组、负向指标，功效系数标准化+熵权。

\textbf{结果关键}：$D$的协调等级、随时间趋势斜率、哪个子系统滞后。

\textbf{注意}：耦合度高不代表发展好（低水平也可耦合），必须结合$T$看$D$；需R/Python扩展。
\end{algobrief}

%% ========== Malmquist指数 ==========
\advancedalgo{Malmquist指数}{malmquist_index}

\begin{algobrief}
\textbf{一句话}：基于DEA前沿，衡量决策单元全要素生产率的跨期变化，并拆成"效率追赶"和"技术进步"两部分。

\textbf{美赛场景区}：评价多家门店/企业连续几年生产率是进步还是退步、靠管理追赶还是技术升级。

\textbf{核心原理}：用DEA求各期前沿距离函数，按Färe分解算Malmquist指数$M=$技术效率变化$\times$技术进步，再细分纯技术效率与规模效率变化。本案例平均每期提升6.5\%，主要靠技术进步（1.089）。

\textbf{操作要点}：工具箱选"Malmquist"，指定DMU、时期、投入列、产出列，输入导向，几何平均汇总。

\textbf{结果关键}：$M>1$进步、$<1$退步；分解看是EFFCH还是TECHCH主导。

\textbf{注意}：需平衡投入产出面板；它是DEA的动态版，静态效率用DEA/SBM。
\end{algobrief}

%% ========== DEA ==========
\advancedalgo{数据包络分析（DEA）}{dea_method}

\begin{algobrief}
\textbf{一句话}：不用预设权重和生产函数，靠线性规划构造"生产前沿"，评价多投入多产出决策单元的相对效率。

\textbf{美赛场景区}：门店/银行/医院等单位投入资源转化为产出的效率评价，找标杆和低效改进空间。

\textbf{核心原理}：用包络线性规划构造由最有效单元围成的前沿面，测各单元到前沿的径向距离得效率分（0--1），再分解纯技术效率与规模效率。本案例150店中34家前沿有效，平均效率0.888，低效店平均有14.4\%投入压缩空间。

\textbf{操作要点}：工具箱选"DEA"，指定DMU列、投入列、产出列，投入导向、VRS，样本量$\geq 3(m+s)$。

\textbf{结果关键}：效率=1为有效；报告松弛量、规模报酬类型与参照标杆单元。

\textbf{注意}：效率是组内相对值，跨期不可直接比；不区分随机噪声（SBM/SFA改进）；需R/Python扩展。
\end{algobrief}

%% ========== 综合指数 ==========
\advancedalgo{综合指数}{composite_index}

\begin{algobrief}
\textbf{一句话}：把多个量纲方向不一的指标同向化、标准化后加权平均成一个0--1的综合指数，简单直接地排名分档。

\textbf{美赛场景区}：多区域门店经营绩效的统一打分、五档分档和区域比较。

\textbf{核心原理}：标定各指标方向（正向/逆向），统一变换后min-max标准化，按权重加权平均得0--1指数，降序排名并按分位数分五档。本案例ST0056指数0.959最高。

\textbf{操作要点}：工具箱选"综合指数"，指定id列、指标列与方向，加权平均，等权，分位数分档，minmax。

\textbf{结果关键}：综合指数均值/区间、五档分布、领先与落后对象。

\textbf{注意}：这是最朴素的加权综合，权重主观；追求客观可换熵权，追求距离可比可用TOPSIS。
\end{algobrief}

%% ========== DEMATEL ==========
\advancedalgo{DEMATEL}{dematel_method}

\begin{algobrief}
\textbf{一句话}：从专家两两影响打分中，算每个因素的"影响度、被影响度、中心度、原因度"，区分原因因素和结果因素。

\textbf{美赛场景区}：供应链/系统改进中找"牵一发而动全身"的关键源头因素。

\textbf{核心原理}：由专家0--4打分建直接影响矩阵，规范化后求综合影响矩阵，中心度$M=D+R$（谁关键）、原因度$N=D-R$（正为原因、负为结果）。本案例信息共享水平中心度最高，是核心原因因素。

\textbf{操作要点}：工具箱选"DEMATEL"，0--4级直接打分，最大行和归一化，阈值取矩阵均值。

\textbf{结果关键}：中心度排序找关键因素；原因度正负区分原因/结果，优先改原因因素。

\textbf{注意}：原因/结果是打分相对属性，不等于经过检验的因果链；与ISM互补——ISM分层级、DEMATEL判因果属性。
\end{algobrief}

%% ========== RSR ==========
\advancedalgo{秩和比法（RSR）}{rsr_method}

\begin{algobrief}
\textbf{一句话}：把各指标先编秩（排名次），再把秩次加权平均成一个0--1的综合秩和比，最后分档。

\textbf{美赛场景区}：医院科室、多指标评价的综合排序与相对档位划分，尤其适合指标量纲差异大时。

\textbf{核心原理}：高优指标从小到大编秩、低优指标反向，秩次按权重相加并归一得RSR，用概率单位回归检验其近似正态，再按分档概率划档。本案例科室011的RSR=0.989最高，分较差/中/良/优四档。

\textbf{操作要点}：工具箱选"秩和比法"，指定对象列、低优列，并列取平均秩，等权，正态性检验开启。

\textbf{结果关键}：RSR越大越好；probit回归$R^2$高说明近似正态、分档可信。

\textbf{注意}：编秩会损失原始数值信息（只保留顺序）；它简单稳健，适合数据质量一般的综合评价。
\end{algobrief}

% 第9章 问卷研究
\chapter{问卷研究}

\begin{chapteroverview}
\textbf{这个分类是干什么的？}

问卷研究是通过发放量表、收集被试主观评分来测量"态度、满意度、意愿"这类看不见摸不着的构念（construct）的一整套方法。它要回答两个核心问题：你的问卷\textbf{靠不靠谱}（信度）、测的是不是你想测的东西（效度）？以及问卷数据之间的\textbf{关系结构}是怎样的（谁影响谁、有没有中介和调节）。从发问卷前的项目分析、到回收后的信效度检验、再到用结构方程模型验证理论模型，是一条完整的链条。

\textbf{美赛什么题型常用？}

美赛中凡是涉及"人"的主观评价的题目几乎都要用到本章方法：C 题（大数据题）里若带用户评分、满意度、调查数据；B 题中需要做公众偏好、政策接受度分析；以及任何需要设计问卷或解释问卷数据的赛题。尤其\textbf{信度分析、效度分析、结构方程模型（SEM）}是问卷类论文的"标配三件套"，评审看到问卷数据却没有报告信效度会直接质疑结论的可靠性。

\textbf{学习路径建议}

按真实做问卷的顺序学：先学\textbf{项目分析（区分度）}筛掉烂题，再学\textbf{信度分析}确认题项内部一致、\textbf{效度分析}（探索性因子分析）确认结构合理；然后学\textbf{中介效应}和\textbf{调节效应}看懂变量间的作用机制；最后用\textbf{结构方程模型（SEM）}把测量和关系统一在一个模型里验证理论假设。其余如 KANO、PSM、联合分析、NPS 等是面向特定商业问题的进阶工具，遇到对应题型再查即可。
\end{chapteroverview}

\section{基础级算法}

%% ========== 项目分析（区分度分析） ==========
\basicalgo{项目分析（区分度分析）}{item_analysis}

\begin{algointro}
项目分析又叫区分度分析，是正式发问卷前的"预试体检"——先小范围找一批人填一遍，看每道题能不能把"总分高的人"和"总分低的人"区分开。区分不开的题就是烂题，应该删掉或改掉。
\end{algointro}

\begin{algoscene}
\begin{itemize}
    \item 自编问卷正式大规模发放\textbf{之前}的预试环节，必备步骤
    \item 美赛中若题目要求设计问卷或量表，可在方法部分说明"经项目分析删去区分度不足的题项"
    \item 检验李克特量表（Likert）中每一道题的鉴别能力
\end{itemize}
\end{algoscene}

\begin{algoprinciple}
想象你出了一套考卷，想知道每道题有没有水平区分度。做法很朴素：

\begin{itemize}
    \item 先把所有人按总分从高到低排队，取前 27\% 当"高分组"、后 27\% 当"低分组"（这是教育测量里的经典切法）
    \item 对每一道题，看高分组和低分组在这道题上的平均分差多少
    \item 如果一道题好学生和差学生得分差不多，说明这道题\textbf{没有区分度}，谁答都一样，是废题
\end{itemize}

\textbf{两个核心指标：}
\begin{itemize}
    \item \textbf{决断值 CR（Critical Ratio）}：就是高分组与低分组在该题上得分差的独立样本 t 值。CR $\geq 3.0$ 且 $p<0.05$ 表示该题区分力足够
    \item \textbf{题总相关 r}：该题得分与"剔除本题后总分"的相关。$r \geq 0.3$ 且 $p<0.05$ 表示该题和整个量表方向一致
\end{itemize}

两个指标都达标就保留；都不达标建议删除；介于中间建议修改后保留。
\end{algoprinciple}

\begin{algosposs}
\begin{enumerate}
    \item 先处理反向计分题（如本题 1=非常同意、5=非常不同意，需在分析前反转成统一方向）
    \item 在 Dabbit AI SPSS 工具箱中选择"项目分析（区分度分析）"
    \item 选择题项所在列，指定计分类型为 likert，填写反向题编号（如 q6、q11、q16）
    \item 高/低分组比例默认各 27\%，CR 阈值默认 3.0，题总相关阈值默认 0.3
    \item 若已有现成总分列可指定，否则由工具自动把所有题项求和构造总分
    \item 点击运行，得到逐题的 CR、p 值、题总相关与"保留/修改/删除"判定
\end{enumerate}

\textbf{默认参数参考}：分组比例各 27\%，CR 阈值 3.0，题总相关下限 0.3，相关类型自动选择。
\end{algosposs}

\begin{algoresult}
输出一张逐题判定表，关键列：
\begin{itemize}
    \item \textbf{高分组均值 / 低分组均值}：两组在该题上的平均得分
    \item \textbf{决断值 CR}：$\geq 3.0$ 且 $p<0.05$ 为达标
    \item \textbf{题总相关 r}：$\geq 0.3$ 且 $p<0.05$ 为达标；若为负说明该题方向可能反了
    \item \textbf{判定}：工具自动给出"保留 / 建议修改 / 建议删除"
\end{itemize}

\textbf{怎么判断好坏}：大多数题 CR 显著、题总相关达标，说明问卷质量整体不错；若某道题 CR 不显著且题总相关很低（如案例中 q20 的 CR=0.21、r=-0.07），说明这道题谁答都一样，应列入删题清单。
\end{algoresult}

\begin{algonotes}
\textbf{什么时候用？} 自编量表预试阶段、正式施测前的删题决策。

\textbf{什么时候不用？} 直接使用成熟量表（已有信效度证据）时不必再做项目分析；数据已正式收集完毕、无法再改题时也只能事后报告。

\textbf{常见坑}
\begin{itemize}
    \item \textbf{坑1}：忘记反转反向题。反向题不反转，题总相关会变成负的，导致好题被误删
    \item \textbf{坑2}：只看 CR 不看题总相关。CR 高但题总相关低，说明这道题虽然能分组但和量表其他题不同向
    \item \textbf{坑3}：预试样本太小（如 $n<50$），高低分组各只有十几人，CR 很不稳定
\end{itemize}

\textbf{和相似算法的区别}：项目分析看"单题区分度"，信度分析看"整组题的内部一致性"；前者筛题、后者验表，二者常配合使用。
\end{algonotes}

%% ========== 信度分析 ==========
\basicalgo{{\starmark}信度分析}{reliability_analysis}

\begin{algointro}
信度分析回答一个问题：你的量表\textbf{稳不稳定、可不可靠}？如果同一批人今天填一遍、下周再填一遍结果差很远，或者同一维度下的题项各说各话，那这个量表就是不可靠的，后面基于它做的任何分析都站不住脚。
\end{algointro}

\begin{algoscene}
\begin{itemize}
    \item 凡是使用李克特量表收集问卷数据的论文，\textbf{必须}报告信度系数，否则评审会质疑数据质量
    \item 美赛 C 题中涉及满意度、态度、心理量表数据时，先做信度再做后续建模
    \item 判断量表中哪些题项拖了后腿、需要删除
\end{itemize}
\end{algoscene}

\begin{algoprinciple}
最常用的信度指标是 \textbf{Cronbach's $\alpha$（克隆巴赫系数）}。

打个比方：一个"服务满意度"量表有 10 道题，如果它们真的都在测同一个东西，那顾客在第 1 题打高分、其他题也应该倾向打高分——题目之间应该高度相关。$\alpha$ 衡量的就是这种"内部一致性"：

\begin{itemize}
    \item $\alpha$ 的大小取决于\textbf{题项间平均相关}和\textbf{题项数量}
    \item $\alpha$ 越接近 1，说明题项越是在同测一个构念
\end{itemize}

\textbf{判断标准（很重要，必须背）：}
\begin{itemize}
    \item $\alpha \geq 0.9$：优秀
    \item $0.8 \leq \alpha < 0.9$：良好
    \item $0.7 \leq \alpha < 0.8$：可接受（探索性研究可放宽到 0.6）
    \item $\alpha < 0.6$：不佳，需要修订量表
\end{itemize}

除 $\alpha$ 外，还会算\textbf{折半信度}（把题项劈成两半看两半是否一致）和\textbf{组合信度 CR}（结构方程框架下的信度指标）。

另一个关键诊断是\textbf{"删除该题项后的 $\alpha$"}：如果删掉某题后整表 $\alpha$ 反而上升很多（$\Delta\alpha \geq 0.02$），说明这道题在拖后腿；同时看\textbf{修正题总相关}（CITC），低于 0.3 说明该题与其他题关联太弱。
\end{algoprinciple}

\begin{algosposs}
\begin{enumerate}
    \item 在 SPSS 中点击 \textbf{分析 → 标度 → 可靠性分析}
    \item 将同一维度下的所有题项选入\textbf{项}框，模型选择"Alpha"
    \item 点击\textbf{统计}按钮，勾选"项"、"标度"、"删除项后的标度"、"相似矩阵"
    \item 多个维度需\textbf{分别}计算各自的 $\alpha$（不要把所有题混在一起算）
    \item 在 Dabbit AI 工具箱中可一键同时输出 $\alpha$、折半信度、CR，并用 Bootstrap 1000 次给出 $\alpha$ 的 95\% 置信区间
    \item 缺失值默认成对删除；$\alpha$ 可接受阈值设 0.7，题总相关下限 0.3，删除提升阈值 $\Delta\alpha=0.02$
\end{enumerate}

\textbf{默认参数参考}：同时计算 Cronbach's $\alpha$、折半信度、组合信度 CR；Bootstrap 1000 次；$\alpha$ 阈值 0.7。
\end{algosposs}

\begin{algoresult}
\begin{itemize}
    \item \textbf{Cronbach's $\alpha$}：整表或该维度的核心信度系数。案例中整表 10 题 $\alpha=0.865$，判级"良好"
    \item \textbf{Bootstrap 95\% 置信区间}：如 $[0.830, 0.894]$，区间不跨过 0.7 才能说信度可靠
    \item \textbf{删除项后的 $\alpha$}：逐题列出删掉该题后的 $\alpha$。若某行删除后 $\alpha$ 明显高于整表 $\alpha$，且该题修正题总相关 $<0.3$，则建议考虑删除（案例中"会员优惠力度"删除后 $\alpha$ 从 0.865 升到 0.889，CITC 仅 0.148）
    \item \textbf{折半信度}：Spearman-Brown 校正后的值，应与 $\alpha$ 接近
\end{itemize}

\textbf{判断好坏}：各维度 $\alpha$ 都 $\geq 0.7$ 即为合格量表；论文中通常报告"各维度 Cronbach's $\alpha$ 介于 0.xx--0.xx 之间，均大于 0.7，量表信度良好"。
\end{algoresult}

\begin{algonotes}
\textbf{什么时候用？} 任何使用多题项量表的数据，做进一步分析之前必须先报告。

\textbf{什么时候不用？} 单题测量（只有一道题）无法算 $\alpha$；分类单选题也不需要信度。

\textbf{常见坑}
\begin{itemize}
    \item \textbf{坑1}：把不同维度的题混在一起算一个 $\alpha$。不同构念的题混算会虚高或虚低，必须\textbf{分维度}算
    \item \textbf{坑2}：$\alpha$ 越高越好于是疯狂加题。题越多 $\alpha$ 自然越高，但题目冗余无意义；$\alpha$ 在 0.8--0.9 最佳
    \item \textbf{坑3}：看到某题删除后 $\alpha$ 升高就盲目删题。是否删题必须结合题项内容和理论，不能只看统计量
    \item \textbf{坑4}：$\alpha$ 只反映内部一致性，\textbf{不代表效度高}——大家一致地答错也是一致地错
\end{itemize}

\textbf{和相似算法的区别}：信度分析看"题项是否一致"，效度分析看"题项是否测对了东西"；信度是效度的必要不充分条件。
\end{algonotes}

%% ========== 效度分析 ==========
\basicalgo{{\starmark}效度分析}{validity_analysis}

\begin{algointro}
效度分析回答：你的量表\textbf{测准了没有}？你设计量表时以为它在测"工作投入"，数据是否真的支持这一点？效度分析（这里指结构效度，用探索性因子分析 EFA 实现）就是拿数据来验证你预设的维度结构。
\end{algointro}

\begin{algoscene}
\begin{itemize}
    \item 自编或改编量表在正式分析前，\textbf{必须}做结构效度检验，与信度分析一起报告
    \item 美赛问卷类论文中，需证明"我们的问卷确实测到了想测的构念"
    \item 探索量表潜在结构，发现哪些题项应该归在同一因子下
\end{itemize}
\end{algoscene}

\begin{algoprinciple}
\textbf{探索性因子分析（EFA）}的思想：你设计了 15 道题，预设它们分属"工作投入、职业认同、离职意愿"三个维度。因子分析会把这些题项的相关关系提炼成少数几个"潜在因子"——

\begin{itemize}
    \item 如果数据真的如你所料，那"工作投入"的 5 道题会\textbf{扎堆}落在同一个因子上，其他题不沾边
    \item 这 5 道题在该因子上的载荷（loading）应该都 $\geq 0.5$
    \item 如果某道题在两个因子上载荷都差不多（交叉载荷），说明这道题归属不清
\end{itemize}

\textbf{两个前提检验（必须先通过才能做因子分析）：}
\begin{itemize}
    \item \textbf{KMO 检验}：取值 0--1。$>0.9$ 极好，$0.8$--$0.9$ 很好，$0.7$ 以上尚可，$<0.5$ 不宜做因子分析
    \item \textbf{Bartlett 球形检验}：$p<0.05$ 表示题项间存在相关，适合提取因子
\end{itemize}

\textbf{关键指标：}
\begin{itemize}
    \item \textbf{因子载荷}：题项与所属因子的相关，$\geq 0.5$ 可接受，$\geq 0.7$ 很好
    \item \textbf{累计方差解释率}：提取的因子加起来能解释多少题项总变异，$>50\%$ 可接受
    \item \textbf{共同度（communalities）}：每题被提取因子解释的比例，$>0.4$ 为宜
\end{itemize}
\end{algoprinciple}

\begin{algosposs}
\begin{enumerate}
    \item 在 SPSS 中点击 \textbf{分析 → 降维 → 因子}
    \item 将所有题项选入\textbf{变量}框
    \item 点击\textbf{描述}：勾选"KMO 和 Bartlett 球形检验"
    \item 点击\textbf{提取}：方法选"主轴因子"或"主成分"，勾选"碎石图"，按特征值 $>1$ 或平行分析确定因子数
    \item 点击\textbf{旋转}：选\textbf{最大方差法（Varimax）}，使载荷矩阵更易读
    \item 点击\textbf{选项}：勾选"按大小排序"、"隐藏小系数"（绝对值 $<0.4$ 不显示）
    \item 在 Dabbit AI 工具箱中可直接按预设维度（如"工作投入：投入1--投入5"）校验载荷归属，自动判定交叉载荷
\end{enumerate}

\textbf{默认参数参考}：提取方法主轴因子法，旋转 Varimax，因子数 auto（结合特征值 $>1$ 与平行分析），主载荷阈值 0.5，次高载荷与主载荷差 $\geq 0.2$。
\end{algosposs}

\begin{algoresult}
\begin{itemize}
    \item \textbf{KMO 与 Bartlett 表}：KMO 如案例 0.903（"适合做因子分析"），Bartlett $\chi^2=1801.88$、$p<0.001$，两前提都通过
    \item \textbf{方差解释表}：各因子特征值与累计解释率。案例中 3 个因子累计解释 73.8\%，远高于 50\% 参考线
    \item \textbf{旋转后载荷矩阵}：每行一题，看载荷最大的那一列是否落在\textbf{预设维度}上。案例中投入1--5 全部落在"工作投入"列（载荷 0.82--0.87），无交叉载荷，15/15 题全部达标
    \item \textbf{共同度}：每行最后一列，案例均在 0.7 左右，说明每题被因子解释得不错
\end{itemize}

\textbf{判断好坏}：题项按预期归到对应因子、主载荷 $\geq 0.5$、无严重交叉载荷、累计方差解释率 $>50\%$，则结构效度良好。
\end{algoresult}

\begin{algonotes}
\textbf{什么时候用？} 预试问卷结构效度检验；新量表开发。

\textbf{什么时候不用？} 已有成熟量表且样本为验证性用途时，应用\textbf{验证性因子分析（CFA）}而非 EFA。

\textbf{常见坑}
\begin{itemize}
    \item \textbf{坑1}：KMO 或 Bartlett 不通过就强行做因子分析，结果无意义
    \item \textbf{坑2}：因子数靠"特征值 $>1$"硬定，不结合碎石图和平行分析，常常多提或少提因子
    \item \textbf{坑3}：题项落在了\textbf{没预料到的因子}上却不报告，假装一切正常——这正是效度分析要发现的问题
    \item \textbf{坑4}：把 EFA 当 CFA 用。EFA 是"让数据说话找结构"，CFA 是"验证预设结构"，预试用 EFA、正式验证用 CFA
\end{itemize}

\textbf{和相似算法的区别}：效度分析（EFA，探索性）与验证性因子分析（CFA，验证性）互补；信度（$\alpha$）回答"一致吗"，效度（因子分析）回答"测对了吗"。
\end{algonotes}

%% ========== 中介效应 ==========
\basicalgo{中介效应（含平行/链式中介）}{mediation_analysis}

\begin{algointro}
中介效应研究"A 是怎么影响 C 的"——不是直接影响，而是 A 先影响 B、B 再影响 C，B 就是中间的"传声筒"（中介变量）。比如"组织支持感"不直接提升"工作绩效"，而是先提升"工作满意度"，满意度高了绩效才好——满意度就是中介。
\end{algointro}

\begin{algoscene}
\begin{itemize}
    \item 美赛中需要解释变量间\textbf{作用机制/传导路径}时（A 通过什么中间环节影响 C）
    \item 问卷类、管理类、社科类建模中常见，比单纯报告"显著相关"更深入
    \item 可分解总效应为直接效应和间接效应，量化各条路径的贡献比例
\end{itemize}
\end{algoscene}

\begin{algoprinciple}
用生活化的例子理解：你发现"多学习（X）"和"成绩好（Y）"相关，但这是为什么？可能是因为"学习提升了理解深度（M）"，理解深了成绩才好。M 就是中介。

\begin{itemize}
    \item \textbf{总效应 c}：X 对 Y 的总影响（不考虑中间环节）
    \item \textbf{路径 a}：X 对 M 的影响；\textbf{路径 b}：M 对 Y 的影响
    \item \textbf{间接效应} $= a \times b$：X 经由 M 传导到 Y 的部分
    \item \textbf{直接效应 c'}：去掉 M 的传导后，X 还剩多少直接影响 Y
    \item 关系：$c = c' + a \times b$
\end{itemize}

\textbf{平行中介}：有多个中介并列（如满意度 M1、工作投入 M2 同时传导），各算各的 $a_i \times b_i$。\textbf{链式中介}：M1 先影响 M2、M2 再影响 Y，形成链条。

\textbf{显著性怎么判断}：传统 Sobel 检验功效低，现在主流用 \textbf{Bootstrap}——有放回抽样 5000 次，每次重估间接效应，得到 95\% 置信区间。\textbf{区间不含 0 就说明间接效应显著}。

若直接效应 $c'$ 显著、间接效应也显著 = 部分中介；若 $c'$ 不显著只有间接效应 = 完全中介。
\end{algoprinciple}

\begin{algosposs}
\begin{enumerate}
    \item 在 Dabbit AI SPSS 工具箱中选择"中介效应分析"
    \item 选入自变量 X、中介变量 M（可多个）、因变量 Y，控制变量（如任职年限）一并放入
    \item 模型类型选"平行中介"（多个中介并列）或"链式中介"（按顺序传导）
    \item Bootstrap 方法选百分位法，抽样 5000 次，置信水平 95\%
    \item 运行后查看各路径系数与间接效应的 Bootstrap 置信区间
\end{enumerate}

\textbf{SPSS 原生对应}：需安装 PROCESS 插件（Hayes），通过 \textbf{回归}或语法实现；SPSS 菜单本身不直接提供中介效应检验。

\textbf{默认参数参考}：Bootstrap 百分位法，模型 parallel（平行），可纳入控制变量。
\end{algosposs}

\begin{algoresult}
\begin{itemize}
    \item \textbf{总效应 c}：如案例 $0.602$、$p<0.001$，X 总体上显著影响 Y
    \item \textbf{直接效应 c'}：如 $0.160$、$p=0.069$（不显著），说明去掉中介后 X 不再直接影响 Y
    \item \textbf{各条间接效应}：X$\to$M1$\to$Y = $0.342$（95\%CI $0.206$--$0.494$，不含 0，显著）；X$\to$M2$\to$Y = $0.099$（CI $0.023$--$0.185$，显著）
    \item \textbf{总间接效应}：$0.441$，占总效应约 73\%，即大部分影响是通过中介传导的
    \item 每条间接效应的 CI\textbf{不跨 0} 即显著
\end{itemize}

\textbf{怎么判断}：直接效应不显著 + 间接效应显著 = 完全中介（案例即此情形）；两者都显著 = 部分中介。
\end{algoresult}

\begin{algonotes}
\textbf{什么时候用？} 需要揭示"A 经 B 影响 C"的内在机制时。

\textbf{什么时候不用？} 只是想预测 Y 或看 X 是否影响 Y 时，普通回归即可，不必硬加中介。

\textbf{常见坑}
\begin{itemize}
    \item \textbf{坑1}：用横截面问卷数据说"因果中介"。中介分析只能说"统计上的传导路径"，不能证明因果关系
    \item \textbf{坑2}：中介变量和自变量高度共线却不报告。平行中介需检查 VIF
    \item \textbf{坑3}：只看间接效应点估计，不看 Bootstrap 置信区间是否跨 0
\end{itemize}

\textbf{和相似算法的区别}：中介效应看"X 通过 M 影响 Y"的\textbf{中间通道}；调节效应看"M 改变 X 对 Y 的影响\textbf{强度}"，一个是"传话筒"、一个是"放大器"。
\end{algonotes}

%% ========== 调节作用 ==========
\basicalgo{调节作用（调节效应）分析}{moderation_analysis}

\begin{algointro}
调节效应研究"X 对 Y 的影响，会不会因为另一个变量 M 而变强或变弱"？M 是个"放大器/削弱器"。比如"工作自主性提升工作投入"这个关系，在"上级支持感高"的员工身上特别明显、在上级不支持时几乎没用——上级支持感就是调节变量。
\end{algointro}

\begin{algoscene}
\begin{itemize}
    \item 美赛中分析"某因素的作用在不同人群/情境下是否不同"时
    \item 问卷研究中探究边界条件："在什么情况下，A 对 B 的作用更强/更弱"
    \item 比主效应分析更精细，能给出差异化建议（如"对高支持感员工应多授权"）
\end{itemize}
\end{algoscene}

\begin{algoprinciple}
调节效应的核心是\textbf{交互项} $X \times M$。

打个比方：自主性（X）对投入（Y）有正向作用，但这个作用的斜率会随上级支持（M）变化。做法：

\begin{itemize}
    \item 先把连续的 X、M 各自\textbf{中心化}（减去均值），降低交互项与主效应的共线性
    \item 构造交互项 $X \times M$
    \item 回归 $Y = b_0 + b_1 X + b_2 M + b_3 (X \times M) + $ 控制变量
    \item 看 $b_3$（交互项系数）是否显著：$b_3$ 显著就说明存在调节效应
\end{itemize}

\begin{itemize}
    \item $b_3 > 0$：M 越高，X 对 Y 的正向作用\textbf{越强}（加强型调节）
    \item $b_3 < 0$：M 越高，X 对 Y 的作用\textbf{越弱}（削弱型调节）
\end{itemize}

交互项显著后，再做\textbf{简单斜率检验}：把 M 分成"低（均值 $-1$ 个标准差）、中、高（均值 $+1$ 个标准差）"三组，分别看每组中 X 对 Y 的斜率，画成两条交互线，直观展示调节方向。
\end{algoprinciple}

\begin{algosposs}
\begin{enumerate}
    \item 在 Dabbit AI SPSS 工具箱中选择"调节效应分析"
    \item 选入自变量 X（工作自主性）、调节变量 M（上级支持感）、因变量 Y（工作投入），控制变量（在职年限）
    \item 中心化方式选"均值中心化"，自动构造交互项
    \item 勾选"简单斜率检验"，分层水平设为均值 $\pm 1$ 个标准差
    \item 比较含交互项模型（模型2）与不含交互项模型（模型1）的 $\Delta R^2$
\end{enumerate}

\textbf{SPSS 原生对应}：通过 \textbf{分析 → 回归 → 线性}，手动将 X、M 中心化后相乘生成交互项，再放入回归；简单斜率需用手动语法或 PROCESS 插件。

\textbf{默认参数参考}：均值中心化，普通标准误，简单斜率按均值 $\pm 1$SD 分组。
\end{algosposs}

\begin{algoresult}
\begin{itemize}
    \item \textbf{$\Delta R^2$ 与 F 变化检验}：案例中加入交互项后 $R^2$ 从 0.332 升到 0.493，$\Delta R^2=0.162$，$p<0.001$，说明交互项显著提升解释力
    \item \textbf{交互项系数 $b_3$}：案例 $0.387$、$p<0.001$，显著且为正，即上级支持感越高，自主性对投入的正向带动越强
    \item \textbf{简单斜率表}：低支持时 X$\to$Y 斜率仅 $0.057$（不显著）；高支持时斜率 $0.861$（显著）——直观说明"只有上级支持到位，授权才有效"
    \item \textbf{VIF}：案例最大 1.037，远低于 10，说明中心化后共线性可控
\end{itemize}

\textbf{怎么判断}：交互项 $p<0.05$ 即存在调节；再结合简单斜率方向判断是加强还是削弱。
\end{algoresult}

\begin{algonotes}
\textbf{什么时候用？} 探究变量关系的边界条件、异质性效应时。

\textbf{什么时候不用？} 只想知道 X 是否影响 Y 时不必加调节；样本量小（如 $n<100$）时交互项检验功效低。

\textbf{常见坑}
\begin{itemize}
    \item \textbf{坑1}：不中心化就构造交互项，导致主效应与交互项严重共线、系数不稳定
    \item \textbf{坑2}：只报告交互项显著，不做简单斜率、不画交互图，读者看不懂调节方向
    \item \textbf{坑3}：把调节当因果。调节效应同样是观察性关联
\end{itemize}

\textbf{和相似算法的区别}：调节效应（$X \times M$，M 是放大器）与中介效应（$X \to M \to Y$，M 是传话筒）方向完全不同，可结合成"有调节的中介"。
\end{algonotes}

%% ========== 结构方程模型 SEM ==========
\basicalgo{{\starmark}结构方程模型（SEM）}{sem_model}

\begin{algointro}
结构方程模型（Structural Equation Model, SEM）是问卷研究的"集大成者"。它一次性同时做两件事：一是\textbf{测量模型}（CFA，验证你的量表题项能否准确测出潜变量），二是\textbf{结构模型}（检验潜变量之间谁影响谁的理论路径）。简单说，它把信效度检验和路径分析合并成一个大模型来验证。
\end{algointro}

\begin{algoscene}
\begin{itemize}
    \item 问卷研究论文的\textbf{核心方法}，尤其当你有多个潜变量、要验证一套理论假设时
    \item 美赛中涉及"组织支持感$\to$工作投入$\to$离职意向"这类多变量理论模型时
    \item 同时回答"量表好不好"（测量）和"假设成不成立"（结构）两个问题
\end{itemize}
\end{algoscene}

\begin{algoprinciple}
现实中很多概念（如"工作投入""满意度"）是\textbf{看不见的潜变量}，只能靠几道题间接测量。SEM 的巧妙之处在于它把这两件事放在一个图里：

\begin{itemize}
    \item \textbf{测量模型（外圈）}：每个潜变量由 3--4 道题指标测量（如工作投入 $\leftarrow$ we1、we2、we3），这部分本质是 CFA
    \item \textbf{结构模型（内圈）}：潜变量之间画箭头表示假设的因果路径（如组织支持感 $\to$ 工作投入 $\to$ 离职意向）
\end{itemize}

模型用\textbf{极大似然估计}反复调整，直到模型隐含的协方差矩阵尽可能接近真实数据的协方差矩阵。

\textbf{拟合优度指标（必须背的判断标准）：}
\begin{itemize}
    \item $\chi^2/df$（卡方自由度比）：$\leq 3$ 良好，越小越好
    \item CFI、TLI（比较拟合指数）：$\geq 0.90$ 可接受，$\geq 0.95$ 良好
    \item RMSEA（近似误差均方根）：$\leq 0.08$ 可接受，$\leq 0.05$ 良好
    \item SRMR（标准化残差均方根）：$\leq 0.08$ 良好
\end{itemize}

\textbf{测量质量指标}：标准化载荷 $\geq 0.5$，组合信度 CR $\geq 0.7$，平均方差抽取 AVE $\geq 0.5$，且 $\sqrt{AVE}$ 大于该因子与其他因子的相关（区分效度）。
\end{algoprinciple}

\begin{algosposs}
\begin{enumerate}
    \item SEM 在 SPSS 中需借助 AMOS 模块（画图）或在 Dabbit AI SPSS 工具箱中用代码指定模型
    \item 在工具箱中填写模型设定：测量部分写明"潜变量：题项列表"，结构部分写明"因变量 $\sim$ 自变量"
    \item 估计方法选\textbf{极大似然（ML）}，适合连续指标
    \item 缺失值处理默认成列删除（listwise），最大迭代 600
    \item 拟合标准选"常规"（$\chi^2/df \leq 3$、CFI $\geq 0.90$、RMSEA $\leq 0.08$、SRMR $\leq 0.08$）
    \item 因子载荷阈值 0.5，信度阈值 0.7，最小样本量建议 200
\end{enumerate}

\textbf{SPSS 路径}：\textbf{分析 → 降维}只做 EFA；SEM 本身在 SPSS Statistics 原生菜单中不直接提供，需用 AMOS 或扩展模块。
\end{algosposs}

\begin{algoresult}
\begin{itemize}
    \item \textbf{测量模型拟合}：案例 $\chi^2/df=1.10$、CFI=0.992、RMSEA=0.026、SRMR=0.047，全部达标
    \item \textbf{标准化载荷}：各题载荷 0.70--0.86，均 $>0.5$，说明题项很好地测量了对应潜变量
    \item \textbf{CR 与 AVE}：各因子 CR $>0.8$、AVE $>0.57$，区分效度（$\sqrt{AVE}$ 大于因子相关）全部通过
    \item \textbf{结构路径系数}：组织支持感$\to$工作投入 $\beta=0.307$（$p<0.001$）；职业成长$\to$工作投入 $\beta=0.493$；工作投入$\to$离职意向 $\beta=-0.421$（负向，投入越高越不想走）
    \item \textbf{$R^2$}：工作投入被解释了 45.5\% 的方差，离职意向被解释了 17.8\%
\end{itemize}

\textbf{怎么判断}：五项拟合指标全部达标 + 假设路径显著，就说明理论模型得到数据支持。
\end{algoresult}

\begin{algonotes}
\textbf{什么时候用？} 有多个潜变量、要同时验证测量质量和理论关系时，是问卷研究的高级形态。

\textbf{什么时候不用？} 只有单一因变量、几个观测变量的简单关系，用回归即可；样本量太小（$n<100$）时 SEM 不稳定，建议至少 $n=200$。

\textbf{常见坑}
\begin{itemize}
    \item \textbf{坑1}：只看 $\chi^2$ 不看其他指标。$\chi^2$ 对样本量敏感，大样本几乎都显著，必须综合 $\chi^2/df$、CFI、RMSEA、SRMR 一起判断
    \item \textbf{坑2}：拟合不好就疯狂删题、加相关残差，直到"凑"出好结果——这是过度拟合，失去验证意义
    \item \textbf{坑3}：把横截面数据的路径系数说成因果。SEM 验证的是"模型与数据拟合"，不是因果证明
    \item \textbf{坑4}：样本量不够就跑 SEM。案例 $n=150$ 已被提示"低于建议值 200"，需谨慎解释
\end{itemize}

\textbf{和相似算法的区别}：SEM = CFA（测量）+ 路径分析（结构）；路径分析只用观测变量、不带测量模型，CFA 只验证测量结构、不检验潜变量间关系。
\end{algonotes}

\section{进阶级算法速览}

%% ========== 内容效度 ==========
\advancedalgo{内容效度}{content_validity}

\begin{algobrief}
\textbf{一句话}：邀请领域专家逐条判断"题项是否测到了想测的内容"，用 I-CVI/S-CVI 量化条目内容代表性。

\textbf{美赛场景区}：自编量表正式使用前，由专家评议题项覆盖面；美赛中设计问卷时可借鉴专家评分思路。

\textbf{核心原理}：专家按 1--4 级评题项相关性，评分 $\geq 3$ 算"相关"。条目级 I-CVI = 评相关的专家数/专家总数；$\geq 6$ 名专家时 I-CVI $\geq 0.78$ 达标，$3$--$5$ 名时要求 $=1.00$。量表级 S-CVI/Ave $\geq 0.90$、S-CVI/UA $\geq 0.80$。

\textbf{操作要点}：工具箱选"内容效度"，评分尺度 1--4，缺失策略 exclude，6 名专家对 25 条条目评分。

\textbf{结果关键}：逐题 I-CVI 对照阈值，未达标条目（如案例 Q5、Q17）需修订；S-CVI/Ave 与 S-CVI/UA 判断整体。

\textbf{注意}：内容效度是"专家共识"层面的效度，与后续统计导向的结构效度（EFA/CFA）互补，不能互相替代。
\end{algobrief}

%% ========== 共同方法偏差 Harman ==========
\advancedalgo{共同方法偏差检验（Harman）}{harman_cm_bias}

\begin{algobrief}
\textbf{一句话}：同一批人、同一时间、自评方式收集的问卷，可能因"共同方法"让变量间相关被人为夸大，Harman 单因子检验筛查这一风险。

\textbf{美赛场景区}：问卷全部来自同源自评数据、要做相关/回归分析前的必做筛查。

\textbf{核心原理}：对全部题项做无旋转主成分分析，若第一个公因子解释的方差 $<50\%$，说明数据没被单一共同因子主导，共同方法偏差不严重。

\textbf{操作要点}：工具箱选"Harman 单因子检验"，特征值 $>1$ 提取，首因子方差阈值 50\%，先看 KMO/Bartlett 是否适合。

\textbf{结果关键}：案例首因子解释率 39.0\% $<50\%$，判定低风险；若首因子 $>50\%$ 则存在严重共同方法偏差。

\textbf{注意}：Harman 只是\textbf{粗略筛查}，灵敏度低——没超标不等于没有偏差；严格验证需用标记变量或潜在方法因子法。
\end{algobrief}

%% ========== 验证性因子分析 CFA ==========
\advancedalgo{验证性因子分析（CFA）}{cfa_analysis}

\begin{algobrief}
\textbf{一句话}：与 EFA"让数据找结构"相反，CFA 是"拿预设的因子结构去验证数据支不支持"，是正式量表验证的标准方法。

\textbf{美赛场景区}：使用成熟量表或已有理论维度时，用 CFA 验证测量结构；常作为 SEM 的第一步。

\textbf{核心原理}：用极大似然估计题项在预设因子上的载荷，用 $\chi^2/df$、CFI、TLI、RMSEA、SRMR 判断模型是否复现数据相关结构；再用 CR $\geq 0.7$、AVE $\geq 0.5$、Fornell-Larcker 准则（$\sqrt{AVE}>$ 因子相关）检验聚合效度与区分效度。

\textbf{操作要点}：工具箱选"验证性因子分析"，写明"因子：题项列表"，缺失 listwise，标准档位判定。

\textbf{结果关键}：案例 $\chi^2/df=0.78$、CFI=1.000、RMSEA=0.000，全部载荷显著且 $>0.5$，三维结构获数据支持。

\textbf{注意}：CFA 检验的是"预设结构与数据吻合度"，预设错了再完美的拟合也无意义；样本建议 $n \geq 200$。
\end{algobrief}

%% ========== 路径分析 ==========
\advancedalgo{路径分析}{path_analysis}

\begin{algobrief}
\textbf{一句话}：用可观测变量（不带潜变量测量模型）画有向路径图，估计各条直接、间接、总效应。

\textbf{美赛场景区}：研究多变量间传导链条（如产品品质 $\to$ 感知价值 $\to$ 满意度 $\to$ 复购意愿）时。

\textbf{核心原理}：对每个内生变量做一次回归，按路径追踪法则把影响分解为直接效应、间接效应（路径系数乘积之和）与总效应；间接效应显著性用 Bootstrap 区间检验；整体用 $\chi^2$、CFI、RMSEA、SRMR 评价。

\textbf{操作要点}：工具箱选"路径分析"，画定有向无环路径图，标准化系数输出，Bootstrap 2000 次，百分位区间。

\textbf{结果关键}：案例价格感知$\to$复购意愿总效应 0.404，其中直接 0.247、间接 0.157（CI $[0.094,0.228]$ 不含 0）；CFI=0.997、RMSEA=0.050。

\textbf{注意}：路径分析是 SEM 的特例（只用观测变量）；间接效应必须看 Bootstrap 区间是否跨 0，不能只看点估计。
\end{algobrief}

%% ========== 调节中介 ==========
\advancedalgo{调节中介作用（有调节的中介）}{moderated_mediation}

\begin{algobrief}
\textbf{一句话}：中介效应的传导强度本身又被另一个变量 W 调节——"X 经 M 影响 Y"的间接效应，在 W 高低时不一样。

\textbf{美赛场景区}：研究更精细的"在哪类人群身上，传导机制更有效"，如高心理资本员工的组织支持转化效率更高。

\textbf{核心原理}：先检验无条件中介（$a \times b$ 的 Bootstrap 区间），再在被调节路径（如第一阶段 X$\to$M）加入 W 交互项，算\textbf{调节中介指数 $\omega$}；$\omega$ 的 CI 不含 0 即成立，并报告 W 低/中/高三个水平上的条件间接效应。

\textbf{操作要点}：工具箱选"有调节的中介"，指定 X/M/Y/W 及被调节路径（第一阶段），均值中心化，Bootstrap 5000 次，探针水平均值 $\pm 1$SD。

\textbf{结果关键}：案例 $\omega=0.169$（CI $[0.094,0.262]$ 不含 0）；低心理资本时间接效应 0.199、高时 0.538，即随 W 升高而增强。

\textbf{注意}：观察性数据只能说统计关系，不构成因果；需先有理论依据再指定被调节路径，不能数据挖掘式乱试。
\end{algobrief}

%% ========== 组内评分者信度 rwg ==========
\advancedalgo{组内评分者信度（rwg）}{rwg_agreement}

\begin{algobrief}
\textbf{一句话}：当同一组（如同一家门店）有多名评分者对同一对象打分时，检验组内评分是否足够一致、能否聚合为组级得分。

\textbf{美赛场景区}：多评委/多观察员打分数据（如门店服务评分、专家打分）聚合成团队/组织层面指标前。

\textbf{核心原理}：$rwg = 1 - S^2 / \sigma^2_{EU}$，其中 $S^2$ 是组内实际评分方差，$\sigma^2_{EU}$ 是"完全随机打分"的期望方差（均匀分布下为 $(A^2-1)/12$）。$rwg$ 越接近 1 表示评分越一致；$\geq 0.70$ 可聚合，负值表示分歧比随机还大。

\textbf{操作要点}：工具箱选"组内评分者信度 rwg"，指定分组字段、题项列，量表 1--5 档，阈值 0.70，每组至少 3 人。

\textbf{结果关键}：案例 20 组 rwg(j) 均值 0.83、达标率 90\%，支持聚合；门店 A19 为负值（-0.81），存在极端分歧不宜聚合。

\textbf{注意}：rwg 回答"组内一致性"，与 ICC（组间信度）互补；负值组要回头核查评分标准或数据。
\end{algobrief}

%% ========== 对应分析 ==========
\advancedalgo{对应分析}{correspondence_analysis}

\begin{algobrief}
\textbf{一句话}：把两个分类变量的交叉表投影到二维图上，直观展示"哪类人偏好哪类东西"，相似的类别会在图上聚在一起。

\textbf{美赛场景区}：分析"会员等级 $\times$ 商品偏好""地区 $\times$ 品牌选择"等双类别关联，比交叉表更直观。

\textbf{核心原理}：对列联表做奇异值分解，把行列类别投影到低维空间，总惯量等于 $\chi^2/n$。同一象限的行列类别关联更强。第一维解释率常过半，前两维累计解释率可达 80\% 以上。

\textbf{操作要点}：SPSS 中 \textbf{分析 → 降维 → 对应分析}，指定行变量与列变量，坐标类型选 principal，维度数 auto。

\textbf{结果关键}：案例 $\chi^2=67.5$、$p<0.001$，第一维解释 59.2\%；铂金会员与美妆个护同侧、普通会员与生鲜果蔬同侧。

\textbf{注意}：对应分析只展示关联结构，不做因果；期望频数小于 5 的格子过多时结论不稳。
\end{algobrief}

%% ========== 多重响应交叉分析 ==========
\advancedalgo{多重响应交叉分析}{multi_response_crosstab}

\begin{algobrief}
\textbf{一句话}：专门处理问卷\textbf{多选题}——统计每个选项的普及率，并按分组看不同群体的多选偏好差异。

\textbf{美赛场景区}：问卷中"您经常参与哪些项目（可多选）"这类数据的分析，普通频数分析处理不了。

\textbf{核心原理}：把多选答案转成"每选项一列 0/1"，算\textbf{响应率}（该选项占总选择次数）与\textbf{普及率}（多少比例人选了它，可超 100\%）；再对每个选项做"组别 $\times$ 是否选中"的卡方检验，并用 FDR 校正多重比较。

\textbf{操作要点}：SPSS 中 \textbf{分析 → 多重响应}定义变量集，再做交叉表；工具箱中直接选响应变量与分组变量。

\textbf{结果关键}：案例跑步机有氧普及率 48\% 最高；经校正后瑜伽普拉提、游泳等选项组间差异显著，Cramér's V 衡量强度。

\textbf{注意}：各选项普及率之和可超过 100\%（因为一人可多选），不要当成普通百分比解读。
\end{algobrief}

%% ========== KANO 模型 ==========
\advancedalgo{KANO模型}{kano_model}

\begin{algobrief}
\textbf{一句话}：把产品功能分成五类——必备型、期望型、魅力型、无差异型、反向型，告诉你哪些功能是"必须有"、哪些是"惊喜点"。

\textbf{美赛场景区}：产品/服务功能优先级排序，资源有限时决定先做什么；美赛中涉及用户需求分类、功能规划的题目。

\textbf{核心原理}：对每个功能问两道题——"如果有这个功能你感觉如何""如果没有你感觉如何"，按标准 KANO 分类表映射类别。Better=(魅力+期望)/有效数（增加带来的满意提升），Worse=-(期望+必备)/有效数（缺失带来的不满）。

\textbf{操作要点}：工具箱选"KANO 模型"，正向反向题同向编码，象限分界用整体均值，平票顺序"必备$>$期望$>$魅力$>$无差异$>$反向"。

\textbf{结果关键}：案例"余额变动实时提醒"是必备属性（$|Worse|=0.667$，缺失就不满），"开屏广告"是反向属性（提供反而反感）。

\textbf{注意}：优先保障高 $|Worse|$ 的必备功能，再投入高 Better 的魅力功能；反向功能应移除。
\end{algobrief}

%% ========== MaxDiff ==========
\advancedalgo{MaxDiff最大差异测量}{maxdiff}

\begin{algobrief}
\textbf{一句话}：让受访者在每组几个属性里选出"最喜欢"和"最不喜欢"的，从而排出属性重要性的绝对顺序，避免李克特打分的趋中偏差。

\textbf{美赛场景区}：产品属性取舍（如笔记本电脑更看重性能还是续航）、卖点优先级。

\textbf{核心原理}：用层次贝叶斯（HB）估计每个属性的相对效用，偏好份额 $= \exp(\text{效用})/\sum \exp(\text{效用}) \times 100\%$，合计 100\%，份额越高越受偏好。

\textbf{操作要点}：工具箱选"MaxDiff"，数据为长表（每受访者每任务一行），估计方法 auto（小样本自动 HB），可按城市分组比较。

\textbf{结果关键}：案例处理器性能偏好份额 51.2\% 居首，键盘手感 1.1\% 最易被牺牲；$\rho^2=0.296$、任务命中率 62.7\%。

\textbf{注意}：MaxDiff 比直接打分更能拉开属性差距，但问卷设计要求属性组合呈正交，否则效用不可识别。
\end{algobrief}

%% ========== 联合分析 CBC ==========
\advancedalgo{联合分析（含CBC）}{conjoint_analysis}

\begin{algobrief}
\textbf{一句话}：把产品拆成品牌、价格、配置等属性水平，让受访者在多个完整产品方案中做选择，反推每个属性水平的效用和重要性。

\textbf{美赛场景区}：新品概念测试、属性配置与定价决策；美赛中涉及产品/服务设计权衡的题目。

\textbf{核心原理}：用条件 Logit（多项 Logit）估计各属性水平效用，被选概率经指数归一化；属性相对重要性 = 该属性水平效用极差 / 所有属性极差之和。McFadden $R^2$ 衡量模型拟合。

\textbf{操作要点}：工具箱选"联合分析/CBC"，编码方式 effect，优化器 BFGS，价格可作线性项，输出相对重要性排序。

\textbf{结果关键}：案例品牌相对重要性 32.8\% 最高，价格每涨 1000 元效用降 1.23；McFadden $R^2=0.295$。

\textbf{注意}：样本量通常需要每个属性组合有足够重复；某水平不显著时不能宣称有真实偏好差异。
\end{algobrief}

%% ========== NPS ==========
\advancedalgo{NPS净推荐值分析}{nps_analysis}

\begin{algobrief}
\textbf{一句话}：用一个 0--10 分的"推荐意愿"问题算净推荐值 $NPS=$ 推荐者占比 $-$ 贬损者占比，衡量客户口碑。

\textbf{美赛场景区}：客户满意度/忠诚度评价；品牌口碑追踪。

\textbf{核心原理}：9--10 分为推荐者、7--8 分为被动者、0--6 分为贬损者。$NPS = $ 推荐者\% $-$ 贬损者\%。经验阈值：$>0$ 可接受、$>50$ 良好、$>70$ 优秀。用 Bootstrap 给 CI，分组用三档 $\times$ 分组卡方 + Bonferroni 校正。

\textbf{操作要点}：工具箱选"NPS"，指定评分列与分组列，CI 用 bootstrap，比较方法 chi2。

\textbf{结果关键}：案例总体 NPS=14.0（CI $0.7$--$26.7$，不含 0），华南 59.2 显著优于西部 -18.8。

\textbf{注意}：NPS 是个单一口碑指标，不能替代满意度细节；分组差异要做多重比较校正。
\end{algobrief}

%% ========== PSM 价格敏感度 ==========
\advancedalgo{价格敏感度分析（PSM）}{psm_analysis}

\begin{algobrief}
\textbf{一句话}：用"太便宜/便宜/贵/太贵"四个价格问题，找出消费者心理上最优定价点和可接受价格区间。

\textbf{美赛场景区}：新品定价决策；美赛中涉及产品定价、需求估计的题目。

\textbf{核心原理}：Van Westendorp 模型画四条累积曲线：太便宜(P1)、便宜(P2)随价格下降，贵(P3)、太贵(P4)随价格上升。四条曲线交点给出四个关键价：PMC 可接受下限（P1$\cap$P3）、OPP 最优价（P1$\cap$P4）、IDP 无差异点（P2$\cap$P3）、PME 可接受上限（P2$\cap$P4）。

\textbf{操作要点}：工具箱选"价格敏感度 PSM"，线性插值，最小有效样本 30，剔除逻辑顺序违规（P1$>$P2$>$P3$>$P4）的答卷。

\textbf{结果关键}：案例最优价 OPP $\approx 12.4$ 元，可接受区间 $[9.67, 13.86]$ 元。

\textbf{注意}：PSM 是"心理价位"调查，与真实购买行为有差距，正式定价还需结合成本与竞品。
\end{algobrief}

%% ========== Turf ==========
\advancedalgo{Turf组合模型}{turf_model}

\begin{algobrief}
\textbf{一句话}：从一堆候选项（口味、渠道、广告位）里挑出固定数量的组合，使"至少喜欢其中一项"的去重人数最多。

\textbf{美赛场景区}：菜单组合、媒体投放组合、产品线规划——在有限名额下最大化覆盖。

\textbf{核心原理}：本质是组合优化。对 0/1 喜好数据，枚举（或贪心+邻域搜索）所有 $k$ 项组合，算去重触达人数，取最大者。比"直接选单项人气最高的 $k$ 个"更优，因为考虑了人群重叠。

\textbf{操作要点}：工具箱选"Turf"，目标 reach\_frequency，组合大小设 k，候选数小时精确枚举、过大时启发式近似。

\textbf{结果关键}：案例选 3 款口味去重触达 120 人（80\%），比"单项最高 3 款"基线多 33 人。

\textbf{注意}：Turf 优化的是\textbf{覆盖广度}，不是每个单项的人气；重叠大的项目组合后增量很小。
\end{algobrief}

%% ========== BPTO ==========
\advancedalgo{品牌价格抵补模型（BPTO）}{bpto_model}

\begin{algobrief}
\textbf{一句话}：把消费者对品牌的偏好折算成"愿意多付多少钱"，并测算各品牌提价后的份额流失（价格弹性）。

\textbf{美赛场景区}：多品牌/多产品线定价、品牌资产评估。

\textbf{核心原理}：用条件 Logit 把选择效用拆成品牌虚拟变量系数 + 价格系数。品牌相对参考品牌的货币溢价 = $-$品牌系数/价格系数；价格弹性 = 提价 10\% 后的份额变化率。

\textbf{操作要点}：工具箱选"BPTO"，指定品牌、价格、选择、任务、受访者字段，收敛容差 1e-6，加价幅度 10\%。

\textbf{结果关键}：案例以漫享为参照，云屿品牌溢价 3.42 元显著为正，谷咖需折价 6.64 元；各品牌价格弹性 $-2.9$ 至 $-3.5$（富有弹性）。

\textbf{注意}：BPTO 回答"品牌值多少钱"，参考品牌应选份额最高者；弹性绝对值 $>1$ 说明提价得不偿失。
\end{algobrief}

%% ========== 价格断裂点 ==========
\advancedalgo{价格断裂点模型}{price_breakpoint}

\begin{algobrief}
\textbf{一句话}：让不同消费者组只看一个候选价格，画出价格--需求曲线，找出购买意愿骤降的"断裂台阶"，并定收益最大化价。

\textbf{美赛场景区}：定价调研，找出消费者心理价格上限；与 PSM 互补（PSM 问主观感受，本模型看购买意愿突变）。

\textbf{核心原理}：把"意愿 $\geq 4$ 分"定义为需求率，形成价格--需求曲线；相邻价格档做两比例 Z 检验（BH 校正）+ Chow 断点检验找突变点；再用"价格 $\times$ 需求率"算各档收益，取最大者定价。

\textbf{操作要点}：工具箱选"价格断裂点"，独立组设计，购买意愿阈值 4 分，最小突变幅度 10 个百分点，BH 校正。

\textbf{结果关键}：案例 109$\to$129 元时需求骤降 36 个百分点（校正 $p=0.027$），收益在 109 元达峰，建议定价 109 元。

\textbf{注意}：购买意愿是声明意愿，与真实购买有差距；断裂点建议结合成本与竞品复核。
\end{algobrief}

% 第10章 医学统计
\chapter{医学统计}

\begin{chapteroverview}
\textbf{这个分类是干什么的？}

医学统计是为"随访数据、诊断评价、预后预测"这类问题发展出的一套专门方法。它处理两类特殊数据：一是\textbf{时间到事件数据}（从起点到发生某件事经过多久，且有人中途退出观察），二是\textbf{诊断/预测模型的性能评价}（一个打分准不准、能不能用）。这些方法虽然名字带"医学"，但核心思想在美赛里用途极广。

\textbf{美赛什么题型常用？}

千万别被"医学"两个字骗了——本章很多方法是\textbf{模型评价}和\textbf{风险预测}的通用工具：
\begin{itemize}
    \item \textbf{ROC 曲线 / AUC}：评价任何二分类预测模型（会不会流失、是否违约、是否生病）的区分能力，是美赛建模论文的必报指标
    \item \textbf{Kaplan-Meier、Cox 回归}：迁移到"客户存续时间、设备寿命、项目完成时间"等带删失的时间数据
    \item \textbf{校准曲线、DCA、NRI/IDI}：评价概率预测是否可靠、新模型比旧模型好在哪
    \item \textbf{单因素/多因素回归、亚组分析、基线分析}：任何"找影响因素、分组比较"问题的标准流程
\end{itemize}

\textbf{学习路径建议}

先学\textbf{基线分析}和\textbf{单因素/多因素回归}这两个找因素的基础工具；再学\textbf{ROC 曲线}（评价二分类模型性能的核心）；然后学\textbf{Kaplan-Meier 生存曲线}和\textbf{Cox 回归}（处理带删失的时间数据）；最后按需要查校准曲线、DCA、竞争风险等进阶模型评价工具。
\end{chapteroverview}

\section{基础级算法}

%% ========== 基线分析 ==========
\basicalgo{基线分析}{baseline_analysis}

\begin{algointro}
基线分析就是在对比两组（如干预组 vs 对照组）之前，先检查两组在起点上是否"旗鼓相当"。如果两组一开始就不一样，后面观察到的结局差异就说不清是处理的效果还是本来就不同。
\end{algointro}

\begin{algoscene}
\begin{itemize}
    \item 美赛中比较两组/多组对象前，\textbf{先做基线均衡性检验}，是论文规范性的体现
    \item 任何"实验组 vs 对照组"的研究设计，表格 1 通常就是基线特征表
    \item 定位两组在哪些变量上本来就不均衡，后续需在模型中校正
\end{itemize}
\end{algoscene}

\begin{algoprinciple}
思路很直白：把所有基线变量逐一在两组间做比较。

\begin{itemize}
    \item 连续变量：先看正态性（Shapiro-Wilk）和方差齐性（Levene），正态方差齐用独立样本 t 检验，否则用 Mann-Whitney U 非参数检验；多组用 ANOVA 或 Kruskal-Wallis
    \item 分类变量：用卡方检验，期望频数不足时改用 Fisher 精确检验
    \item 每个变量同时报告效应量（Cohen's d、Cramér's V），不只看 p 值
\end{itemize}

$p>0.05$ 表示该基线变量两组\textbf{未发现显著差异}（均衡）；$p<0.05$ 表示两组本来就不同，需要警惕混杂。
\end{algoprinciple}

\begin{algosposs}
\begin{enumerate}
    \item 在 Dabbit AI SPSS 工具箱中选择"基线分析"
    \item 指定分组列（如 group），选入连续基线变量（年龄、BMI、血压等）和分类基线变量（性别、吸烟等）
    \item 排除 ID 列；连续变量自动按正态性/方差齐性选 t 或非参数检验
    \item 分类变量期望频数不足时自动用 Fisher 精确检验
    \item 显著性水平 $\alpha=0.05$，输出效应量，p 值校正默认 none
\end{enumerate}

\textbf{SPSS 原生对应}：连续变量用 \textbf{分析 → 比较均值 → 独立样本 T 检验}，分类变量用 \textbf{分析 → 描述统计 → 交叉表}（勾选卡方）。
\end{algosposs}

\begin{algoresult}
输出一张"基线特征表"，每个变量一行：
\begin{itemize}
    \item \textbf{组内描述}：连续变量报均值$\pm$标准差或中位数（四分位），分类变量报频数（构成比）
    \item \textbf{检验方法与 p 值}：如案例 triglyceride 用 Mann-Whitney U，$p=0.002$ 显著；smoking 卡方 $p=0.006$ 显著
    \item \textbf{效应量}：Cohen's d（连续）或 Cramér's V（分类），衡量差异大小
    \item 案例 8 个基线变量中 2 个（甘油三酯、吸烟）显著不均衡，需在后续模型中校正
\end{itemize}

\textbf{注意}：$p>0.05$ 只表示"未发现显著差异"，不等于两组完全相等（可能是样本量不够）；报告时要如实说明。
\end{algoresult}

\begin{algonotes}
\textbf{什么时候用？} 两组/多组比较研究的第一步，写"Table 1"。

\textbf{什么时候不用？} 随机对照试验理论上基线应均衡，但仍建议报告；纯预测建模（不分组）不需要。

\textbf{常见坑}
\begin{itemize}
    \item \textbf{坑1}：对所有变量都做 t 检验而不看正态性，偏态数据用 t 检验结果不可靠
    \item \textbf{坑2}：基线有差异却不校正，直接归因于处理效应
    \item \textbf{坑3}：把"基线均衡"当成研究质量的全部——均衡性差时应改用多因素回归校正
\end{itemize}

\textbf{和相似算法的区别}：基线分析是"描述+均衡性检验"的标准化流程，本身不评价处理效果；处理效果靠后续的组间比较或回归。
\end{algonotes}

%% ========== 单因素与多因素回归 ==========
\basicalgo{单因素与多因素回归}{uni_multi_regression}

\begin{algointro}
这是找影响因素的标准两步法：先\textbf{单因素回归}把每个候选变量单独和结局跑一遍，筛掉无关的；再把候选变量一起放进\textbf{多因素回归}，在互相控制后看哪些因素\textbf{独立}地影响结局。
\end{algointro}

\begin{algoscene}
\begin{itemize}
    \item 美赛中"哪些因素影响 Y"类问题的标配分析（如哪些指标影响血糖、哪些特征影响违约）
    \item 医学/健康/经济类问题中识别独立风险因素
    \item 论文中"单因素筛选 $\to$ 多因素校正"是最常见的建模叙事
\end{itemize}
\end{algoscene}

\begin{algoprinciple}
打个比方，你想找影响体重的因素。单因素就像一个一个看：只看运动、只看饮食、只看睡眠。但运动多的人往往饮食也规律，单看运动的效果可能混进了饮食的功劳。

\begin{itemize}
    \item \textbf{单因素回归}：每个自变量单独对结局回归，得到未调整的粗效应。筛选阈值常用 $p<0.10$（稍宽松，避免漏掉真正重要的变量）
    \item \textbf{多因素回归}：把单因素达标的变量同时放入一个模型，每个系数都表示"控制了其他变量后"该因素的独立效应
    \item 结局是连续变量用线性回归（OLS），结局是二分类用 Logistic 回归
\end{itemize}

关键输出是调整后的系数与 95\% 置信区间，配合 VIF 检查共线性（VIF $<10$ 可接受）。
\end{algoprinciple}

\begin{algosposs}
\begin{enumerate}
    \item 在 Dabbit AI SPSS 工具箱中选择"单因素与多因素回归"
    \item 选结局变量（如空腹血糖，连续型），选入全部候选自变量
    \item 单因素筛选 p 阈值默认 0.10，达标的全部纳入多因素模型
    \item 分类变量自动哑变量化；VIF 阈值 10；异方差自动检验后必要时改稳健标准误
    \item 多因素默认"全量保留"入选变量，不做逐步筛选（避免过拟合）
\end{enumerate}

\textbf{SPSS 原生对应}：\textbf{分析 → 回归 → 线性}（连续结局）或 \textbf{二元 Logistic}（二分类结局）。
\end{algosposs}

\begin{algoresult}
\begin{itemize}
    \item \textbf{单因素表}：逐变量报告粗效应、95\%CI、p 值。案例年龄、BMI、吸烟、家族史单因素 $p<0.10$，进入多因素
    \item \textbf{多因素表}：校正后效应。案例年龄 $\beta=0.015$（$p=0.002$）、BMI $\beta=0.045$（$p=0.003$）、吸烟整体 $p<0.001$ 仍独立显著
    \item \textbf{$R^2$ 与 F 检验}：模型整体解释力，案例 $R^2=0.226$、$p<0.001$
    \item \textbf{VIF}：最大 2.08，无严重共线性
    \item 对比单因素（圆点）与多因素（方块）效应，看校正后哪些被削弱——那就是被其他变量混淆的
\end{itemize}

\textbf{怎么判断}：多因素模型中 $p<0.05$ 且 CI 不跨无效值（线性跨 0、Logistic 跨 1）的变量，是独立影响因素。
\end{algoresult}

\begin{algonotes}
\textbf{什么时候用？} 需要识别独立影响因素、构建预测模型时。

\textbf{什么时候不用？} 纯预测任务可直接用机器学习方法；变量间有强因果链条时需谨慎解释。

\textbf{常见坑}
\begin{itemize}
    \item \textbf{坑1}：把单因素 $p<0.05$ 直接当结论。单因素是未调整的，可能被混杂
    \item \textbf{坑2}：用逐步回归自动选变量。逐步法结果不稳定、易过拟合，推荐理论驱动纳入
    \item \textbf{坑3}：不检查共线性（VIF），导致系数符号反常、不稳定
    \item \textbf{坑4}：观察性数据下把独立相关说成因果
\end{itemize}

\textbf{和相似算法的区别}：单因素/多因素回归是通用工具；Cox 回归专门处理"时间到事件+删失"结局，Logistic 处理二分类结局。
\end{algonotes}

%% ========== ROC 曲线 ==========
\basicalgo{{\starmark}ROC曲线（诊断评价）}{roc_curve}

\begin{algointro}
ROC 曲线是评价二分类模型"分得准不准"的黄金标准。不管你预测的是"会不会生病""会不会流失""会不会违约"，只要结果是"是/否"两类，都能用 ROC 和它的核心指标 AUC 来衡量模型的区分能力。
\end{algointro}

\begin{algoscene}
\begin{itemize}
    \item \textbf{美赛建模论文必报}：任何二分类预测模型（逻辑回归、随机森林、SVM 等）都要报告 AUC
    \item 比较两个预测模型哪个更好（在同一样本上比 AUC）
    \item 为诊断/预测打分选择最佳"及格线"（截断值）
    \item 医学中评价某种检测手段的诊断价值，迁移到美赛就是评价任何分类器
\end{itemize}
\end{algoscene}

\begin{algoprinciple}
二分类预测有四种结果（混淆矩阵）：

\begin{itemize}
    \item \textbf{真阳性 TP}：实际是，预测也是（找对了）
    \item \textbf{真阴性 TN}：实际否，预测也否（排对了）
    \item \textbf{假阳性 FP}：实际否，预测是（误报）
    \item \textbf{假阴性 FN}：实际是，预测否（漏报）
\end{itemize}

两个核心率：
\begin{itemize}
    \item \textbf{灵敏度（召回率）} $= TP/(TP+FN)$：真病人中被找出来的比例——不漏诊
    \item \textbf{特异度} $= TN/(TN+FP)$：真健康人中被正确排除的比例——不误诊
\end{itemize}

截断阈值越高，判阳性越严，特异度越高、灵敏度越低。把\textbf{所有可能阈值}下的（1$-$特异度，灵敏度）描出来，连成 ROC 曲线。

\textbf{AUC（曲线下面积）}就是 ROC 曲线下面的面积，含义是：\textbf{随机抽一个阳性样本和一个阴性样本，模型给阳性打分更高的概率}。
\begin{itemize}
    \item AUC $=0.5$：跟瞎猜一样
    \item $0.7$--$0.8$：可接受
    \item $0.8$--$0.9$：良好
    \item $>0.9$：优秀
\end{itemize}

\textbf{约登指数} $J=$ 灵敏度 $+$ 特异度 $-1$，最大时对应的阈值就是兼顾漏诊误诊的最佳截断值。
\end{algoprinciple}

\begin{algosposs}
\begin{enumerate}
    \item 在 SPSS 中点击 \textbf{分析 → 分类 → ROC 分析}
    \item 将模型预测得分/概率选入\textbf{检验变量}，金标准结局选入\textbf{状态变量}，指定阳性取值（如 1）
    \item 勾选"输出 ROC 曲线""对角参考线""标准误差和置信区间"
    \item 在 Dabbit AI 工具箱中还能用 Bootstrap 2000 次给 AUC 置信区间，按约登指数自动选最佳截断值
    \item 若要比较两个模型，可同时放入两个得分变量
\end{enumerate}

\textbf{默认参数参考}：最佳截断准则=约登指数最大化，Bootstrap 2000 次，随机种子 42。
\end{algosposs}

\begin{algoresult}
\begin{itemize}
    \item \textbf{AUC 及其 95\%CI}：案例 AUC=0.855（CI $0.793$--$0.907$），属"区分度良好"。AUC $>0.5$ 且 CI 不跨 0.5 才说明模型有效
    \item \textbf{ROC 曲线}：越贴近左上角越好；对角线是随机参考线
    \item \textbf{最佳截断值}：按约登指数最大，案例 $-0.319$ 对应灵敏度 0.92、特异度 0.64、平衡准确率 0.78
    \item \textbf{灵敏度/特异度随阈值变化曲线}：两条线交点附近即最佳工作点
\end{itemize}

\textbf{怎么选阈值}：用于筛查（宁可错杀不可放过）就下调阈值提高灵敏度；用于确诊（避免误诊）就上调阈值提高特异度。
\end{algoresult}

\begin{algonotes}
\textbf{什么时候用？} 任何二分类模型的性能评价，论文必备。

\textbf{什么时候不用？} 回归任务（连续预测值）用 R$^2$/RMSE；多分类任务需用多分类 AUC。

\textbf{常见坑}
\begin{itemize}
    \item \textbf{坑1}：只看准确率不看 AUC。类别不平衡时（如 95\% 阴性），全猜阴性准确率 95\% 但模型毫无用处，AUC 才揭示真相
    \item \textbf{坑2}：在训练集上算 AUC 自欺欺人。必须在独立测试集/验证集上评价
    \item \textbf{坑3}：AUC 高就以为概率准。AUC 只衡量排序区分能力，\textbf{不代表概率校准}（后者看校准曲线）
    \item \textbf{坑4}：两个模型 AUC 差很小却宣称谁更好，应看 DeLong 检验差异是否显著
\end{itemize}

\textbf{和相似算法的区别}：ROC/AUC 看\textbf{区分度}；校准曲线看\textbf{概率是否准确}；DCA 看\textbf{临床/决策价值}——三者评价模型的不同侧面。
\end{algonotes}

%% ========== Kaplan-Meier 生存曲线 ==========
\basicalgo{Kaplan-Meier生存曲线}{kaplan_meier}

\begin{algointro}
生存曲线处理"从某起点到发生某事件经过了多久"的数据，而且中途有人退出观察（删失）也不怕。它能画出随时间推移"还没发生事件"的比例如何下降，并比较两组的生存过程是否不同。
\end{algointro}

\begin{algoscene}
\begin{itemize}
    \item 美赛中凡是带\textbf{时间起点 + 事件 + 中途退出}的数据都能用：客户留存时长（何时流失）、设备故障时间、项目完工时间
    \item 比较两组（如两种治疗、两种维护策略）哪个更"耐久"
    \item 估计中位生存时间（一半人发生事件的时间点）
\end{itemize}
\end{algoscene}

\begin{algoprinciple}
"删失"是生存数据的核心难点：随访结束时还没出事件的人，你只知道他"到观察截止还没出事"，不知道他之后什么时候出事。

\begin{itemize}
    \item \textbf{Kaplan-Meier 乘积极限法}：只在真正发生事件的时间点才让生存率下降一格，删失只在曲线上画个小竖线标记、不改变曲线高度
    \item 生存率 = 到每个事件时点为止"幸存下来"的累积概率
    \item \textbf{中位生存时间}：生存率降到 50\% 对应的时间
\end{itemize}

\textbf{Log-Rank 检验}：比较两条生存曲线是否有显著差异。$p<0.05$ 说明两组生存过程确实不同（如方案 A 的中位无进展生存 22.1 个月，方案 B 仅 8.8 个月）。
\end{algoprinciple}

\begin{algosposs}
\begin{enumerate}
    \item 在 SPSS 中点击 \textbf{分析 → 生存分析 → Kaplan-Meier}
    \item 将生存时间选入\textbf{时间}，事件状态选入\textbf{状态}，定义事件发生值（如 event=1）
    \item 分组变量选入\textbf{因子}（如治疗分组）
    \item 点击\textbf{保存}可加入生存概率；点击\textbf{比较因子}选 Log-Rank 检验
    \item 在 Dabbit AI 工具箱中直接指定时间列、事件列、分组列，自动输出中位生存时间与关键时点生存率
\end{enumerate}

\textbf{默认参数参考}：事件=1，自动选取关键随访时点，输出 95\% 置信区间（Greenwood 法）。
\end{algosposs}

\begin{algoresult}
\begin{itemize}
    \item \textbf{各组中位生存期}：案例方案 A 22.1 月（CI 16.2--26.5），方案 B 8.8 月（CI 6.1--13.7）
    \item \textbf{Log-Rank 检验}：案例 $\chi^2=13.15$、$p=0.0003$，两组生存曲线差异显著
    \item \textbf{关键时点生存率}：如随访 9.45 月时方案 A 生存概率 0.743、方案 B 仅 0.473
    \item \textbf{风险集人数表}：曲线后半段风险集人数变少，尾部估计不稳，解读要谨慎
\end{itemize}

\textbf{怎么判断}：曲线分的越开、Log-Rank $p<0.05$，两组差异越明显；中位生存期越长代表越"耐久"。
\end{algoresult}

\begin{algonotes}
\textbf{什么时候用？} 带删失的时间到事件数据的组间比较和生存过程描述。

\textbf{什么时候不用？} 所有人都观察到了事件、无删失时，普通均值比较即可；想同时控制多个影响因素时应用 Cox 回归。

\textbf{常见坑}
\begin{itemize}
    \item \textbf{坑1}：把删失的人直接当"没发生事件"硬算普通均值，浪费信息且有偏
    \item \textbf{坑2}：删失与分组有关（如某组中途大量退出）却不报告，会严重偏倚结论
    \item \textbf{坑3}：曲线尾部风险集只剩几个人还过度解读
\end{itemize}

\textbf{和相似算法的区别}：Kaplan-Meier 只做描述和两组比较（单因素）；Cox 回归能同时纳入多个协变量、量化各因素的风险比。
\end{algonotes}

%% ========== 亚组分析 ==========
\basicalgo{亚组分析}{subgroup_analysis}

\begin{algointro}
亚组分析问：整体上 A 比 B 好，但这个好处\textbf{对所有人都一样吗}？会不会对年轻人效果大、老年人效果小？它按某个变量（年龄、性别、地区）把人群拆开，分别看处理效应，并用交互检验判断效应是否真的随亚组变化。
\end{algointro}

\begin{algoscene}
\begin{itemize}
    \item 美赛中分析政策/方案效果的\textbf{异质性}——"对谁更有效"
    \item 精细化、分层决策依据（如"对老年人群优先推荐该方案"）
    \item 任何整体处理效应已知后，想知道效应是否因人群而异时
\end{itemize}
\end{algoscene}

\begin{algoprinciple}
核心是\textbf{交互项检验}，而不是简单地看每个亚组内部 p 值。

\begin{itemize}
    \item 先在全人群估计总体处理效应
    \item 再按亚组水平分别估计效应及 95\%CI
    \item 关键：用"处理 $\times$ 亚组"交互项的 Wald/F 检验判断亚组间差异是否显著
\end{itemize}

\textbf{重要原则}：亚组内"不显著"不等于"该亚组无效"——可能只是该亚组样本小、检验功效低。\textbf{组间差异必须看交互检验 p 值}，不能拿各亚组单独的 p 值互相比较。
\end{algoprinciple}

\begin{algosposs}
\begin{enumerate}
    \item 在 Dabbit AI SPSS 工具箱中选择"亚组分析"
    \item 选结局变量、处理变量、亚组分组变量（如 age\_group）
    \item 结局类型选 continuous（连续），效应量自动选均值差 MD
    \item 交互检验 auto；可选入协变量校正；森林图按自然顺序排列
\end{enumerate}

\textbf{SPSS 原生对应}：在回归模型中加入"处理 $\times$ 亚组"交互项，看交互项系数显著性。
\end{algosposs}

\begin{algoresult}
\begin{itemize}
    \item \textbf{总体效应}：案例均值差 0.77（CI 0.58--0.95），整体有效
    \item \textbf{交互检验 p 值}：案例 $p<0.001$，效应确实随年龄组变化——这是\textbf{最关键}的结果
    \item \textbf{各亚组效应}：$<55$ 岁 MD=0.27、$55$--$64$ 岁 MD=0.89、$\geq 65$ 岁 MD=1.22，随年龄增大效果越强
    \item \textbf{森林图}：把总体和各亚组效应及 CI 并排，直观比较
\end{itemize}

\textbf{怎么判断}：交互检验 $p<0.05$ 才宣称"效应因亚组而异"；否则应报告整体效应、不强行分层。
\end{algoresult}

\begin{algonotes}
\textbf{什么时候用？} 政策/方案效果的人群异质性分析。

\textbf{什么时候不用？} 没有先验假设、纯粹靠数据挖掘式地试各种亚组会导致假阳性。

\textbf{常见坑}
\begin{itemize}
    \item \textbf{坑1}：把"某亚组 p>0.05"当成"该亚组无效"——功效问题
    \item \textbf{坑2}：同时试很多亚组变量却不做多重比较校正，假发现率高
    \item \textbf{坑3}：亚组样本量小还下强结论
\end{itemize}

\textbf{和相似算法的区别}：亚组分析是\textbf{描述+检验效应异质性}；调节效应分析（第 9 章）是用连续调节变量建模交互，思想相通。
\end{algonotes}

%% ========== Cox 回归 ==========
\basicalgo{{\starmark}Cox回归（比例风险回归）}{cox_regression}

\begin{algointro}
Cox 回归是生存分析的"多因素版"。它既能处理带删失的时间数据，又能同时纳入多个影响因素，输出每个因素的\textbf{风险比 HR}——告诉我们"这个因素使事件发生风险变成原来的几倍"。
\end{algointro}

\begin{algoscene}
\begin{itemize}
    \item 美赛中\textbf{风险预测/预后模型}的核心工具：客户流失风险、设备失效风险、项目失败风险
    \item 同时控制多个因素，看哪些因素独立影响"生存时间"
    \item 输出 HR 便于业务解读（"年龄每增加一岁，风险升高 5.5\%"）
    \item 医学中是预后研究金标准，迁移到美赛即"多因素时间到事件预测"
\end{itemize}
\end{algoscene}

\begin{algoprinciple}
Cox 模型的形式：风险函数 $h(t|X) = h_0(t) \exp(\beta_1 X_1 + \beta_2 X_2 + \dots)$。

\begin{itemize}
    \item $h_0(t)$ 是基准风险曲线（形状不关心，Cox 用偏似然巧妙避开了估计它）
    \item 每个因素的系数 $\beta$ 取指数就是\textbf{风险比 HR}$=\exp(\beta)$
    \item HR $>1$：该因素增加风险；HR $<1$：保护因素；HR $=1$：无影响
    \item 如案例 age HR=1.055，年龄每大 1 岁死亡风险约为原来的 1.055 倍；beta\_blocker HR=0.608，用药者风险降到约 0.61 倍（保护）
\end{itemize}

\textbf{比例风险假设}：HR 不随时间变化（各因素效应恒定）。用 Schoenfeld 残差检验，$p>0.05$ 表示假设成立；违反时需用分层 Cox 或时依协变量。

\textbf{模型区分度}用 C 指数（0.5 随机、$>0.7$ 可接受），含义和 AUC 类似。
\end{algoprinciple}

\begin{algosposs}
\begin{enumerate}
    \item 在 SPSS 中点击 \textbf{分析 → 生存分析 → Cox 回归}
    \item 将时间变量选入\textbf{时间}，事件状态选入\textbf{状态}（定义事件值=1），协变量选入\textbf{协变量}
    \item 分类协变量设置哑变量；方法选"进入"（强制纳入，不做逐步）
    \item 点击\textbf{绘图}可输出生存曲线；点击\textbf{保存}可加入预测生存概率
    \item 在 Dabbit AI 工具箱中指定时间列 followup\_time、事件列 event，基准估计 Breslow，自动做 Schoenfeld 检验
\end{enumerate}

\textbf{默认参数参考}：无分层、无权重，基准估计 Breslow，自动检验比例风险假设。
\end{algosposs}

\begin{algoresult}
\begin{itemize}
    \item \textbf{HR 与 95\%CI}：报告每个变量的 HR（CI）和 p 值。案例 lvef HR=0.960（保护）、nyha=IV HR=5.766（高危）、beta\_blocker HR=0.608（保护），均 $p<0.05$
    \item \textbf{Schoenfeld 检验}：案例各变量 $p>0.05$，比例风险假设成立
    \item \textbf{C 指数}：案例 0.731，区分能力可接受
    \item \textbf{生存曲线}：按典型患者关键因素水平画出预测生存概率随时间变化
\end{itemize}

\textbf{怎么判断}：HR 的 CI 不跨 1 且 $p<0.05$ 为独立风险因素；HR $>1$ 是危险因子，$<1$ 是保护因子。
\end{algoresult}

\begin{algonotes}
\textbf{什么时候用？} 带删失的时间到事件数据，且要同时控制多个预测因素、量化风险比。

\textbf{什么时候不用？} 只描述生存过程不找因素用 Kaplan-Meier；无删失数据用普通回归。

\textbf{常见坑}
\begin{itemize}
    \item \textbf{坑1}：不检验比例风险假设。若某因素 HR 随时间变化却硬套 Cox，结论有偏
    \item \textbf{坑2}：把 HR 当成绝对风险。HR 是"相对倍数"，基线风险高低也很重要
    \item \textbf{坑3}：事件数太少（经验 EPV：每变量至少 10 个事件）还硬塞很多变量，模型不稳定
    \item \textbf{坑4}：观察性数据下把 HR 当因果效应
\end{itemize}

\textbf{和相似算法的区别}：Cox 回归是多因素生存分析；Kaplan-Meier 是单因素描述；竞争风险模型用于存在多种互斥结局时。
\end{algonotes}

\section{进阶级算法速览}

%% ========== 森林图 ==========
\advancedalgo{森林图}{forest_plot}

\begin{algobrief}
\textbf{一句话}：把多项研究（或多个亚组/变量）的效应量与 95\%CI 画在同一张图上，逐个比较并给出合并效应。

\textbf{美赛场景区}：Meta 分析汇总多来源数据；多因素回归后把各变量效应量并排展示；亚组分析结果可视化。

\textbf{核心原理}：每项研究用方块（大小代表权重）+ 横线（95\%CI）表示，底部菱形是合并效应（随机/固定效应模型加权）。$I^2 \geq 50\%$ 或 Q 检验 $p<0.10$ 时用随机效应模型。

\textbf{操作要点}：工具箱选"森林图/Meta 分析"，效应量选 MD 或 OR/RR，合并方法 auto，输出异质性 $I^2$。

\textbf{结果关键}：案例 150 项研究随机效应合并 MD=0.451（CI 0.406--0.495），$I^2=65.7\%$ 异质性较高。

\textbf{注意}：$I^2$ 大说明研究间不一致，需分析异质来源而非只看合并值。
\end{algobrief}

%% ========== Bland-Altman ==========
\advancedalgo{Bland-Altman一致性分析}{bland_altman}

\begin{algobrief}
\textbf{一句话}：评价两种测量方法（如新仪器 vs 金标准）测同一批样品时是否一致，核心是画"差值 vs 均值"图并给出 95\% 一致性界限。

\textbf{美赛场景区}：迁移到"两个模型/两套评分系统结果是否一致"；校准新测量工具时。

\textbf{核心原理}：逐对算差值 $d=A-B$，平均差 $\bar{d}$ 是系统偏倚，95\% 一致性界限 $= \bar{d} \pm 1.96\,SD$——约 95\% 的差值落在此区间；与临床可接受界限比较判断能否互换。

\textbf{操作要点}：工具箱选"Bland-Altman"，自动判断差值正态性，比例偏倚诊断开启，置信水平 95\%。

\textbf{结果关键}：案例平均差 0.018（CI 跨 0，无系统偏倚），95\% 一致性界限 $-0.39$--$0.43$；注意相关高（r=0.997）不等于一致。

\textbf{注意}：相关系数高不代表方法一致，必须看差值的一致性界限是否落在可接受范围内。
\end{algobrief}

%% ========== 校准曲线 ==========
\advancedalgo{校准曲线}{calibration_curve}

\begin{algobrief}
\textbf{一句话}：检验模型给出的预测概率\textbf{准不准}——预测说有 30\% 概率出事，实际是不是真的约 30\% 出事。

\textbf{美赛场景区}：评价任何输出概率的分类模型（不只看 AUC，还要看概率是否可信）；美赛中做风险评分模型必备。

\textbf{核心原理}：按预测概率分箱，比较每箱平均预测概率与实际事件发生率；理想状态点落在对角线上。校准截距 $=0$、斜率 $=1$ 最佳；Brier 分数越小越好。

\textbf{操作要点}：工具箱选"校准曲线"，指定概率列与结局列，分箱策略 quantile（按分位等分箱），可做 Bootstrap 置信带。

\textbf{结果关键}：案例校准截距 0.33（整体低估风险），Brier=0.195，HL 检验 $p=0.34$；建议先做截距重校准。

\textbf{注意}：AUC 高不等于校准好；概率用于决策阈值时校准比区分度更关键。
\end{algobrief}

%% ========== DCA 决策曲线 ==========
\advancedalgo{DCA决策曲线}{dca_curve}

\begin{algobrief}
\textbf{一句话}：从"净获益"角度比较一个预测模型和"全干预/不干预"两种极端策略，看在哪个阈值区间用模型决策更划算。

\textbf{美赛场景区}：评价模型的\textbf{实际决策价值}（不只是统计指标，而是按模型行动到底赚不赚）；美赛中评估预警/筛查模型的实用性。

\textbf{核心原理}：净获益 $NB(t) = $ 真阳率 $-$ 假阳率 $\times t/(1-t)$，$t$ 是干预阈值。把各阈值下 NB 画曲线，高于"全干预"和"不干预"两条基准线的区间就是临床获益区间。

\textbf{操作要点}：工具箱选"DCA"，填入各模型概率列与结局列，LOESS 平滑，报告阈值区间 auto。

\textbf{结果关键}：案例综合模型在阈值 0.09--0.97 区间净获益均高于基准，最优决策点 $t^*=0.43$ 时每 100 例净获益 +21.9。

\textbf{注意}：DCA 把"误干预代价"和"漏干预代价"用阈值权重折算，阈值设定要结合实际业务。
\end{algobrief}

%% ========== ICC ==========
\advancedalgo{组内相关系数（ICC）}{icc_coefficient}

\begin{algobrief}
\textbf{一句话}：衡量多名评分者/多次测量对同一对象打分的\textbf{绝对一致性}，比 Kappa 更适合连续测量值。

\textbf{美赛场景区}：多评委打分一致性（如专家评审、传感器重复测量）；评价测量工具可重复性。

\textbf{核心原理}：把总方差分解为对象间、评分者间、误差三部分，ICC=对象间方差/总方差。取值 0--1。Koo \& Li 分级：$<0.5$ 差、$0.5$--$0.75$ 中等、$0.75$--$0.9$ 良好、$>0.9$ 优秀。

\textbf{操作要点}：工具箱选"ICC"，模型 two-way random、绝对一致、单次测量，置信水平 95\%。

\textbf{结果关键}：案例 ICC2=0.844（CI 0.76--0.90），良好；取多人平均后 ICC2k=0.942 达优秀（平均提高信度）。

\textbf{注意}：要选对"绝对一致 vs 相对一致""单次 vs 平均"，选错会得到不同数值。
\end{algobrief}

%% ========== RCS ==========
\advancedalgo{限制性立方样条（RCS）}{rcs_spline}

\begin{algobrief}
\textbf{一句话}：不预设线性关系，用灵活的样条曲线刻画暴露与结局之间的\textbf{非线性}剂量反应（如睡眠太短太长血压都高的 U 形关系）。

\textbf{美赛场景区}：探索变量间非线性关系（避免简单线性假设误导）；找最优点/阈值（如最佳睡眠时长、最佳温度）。

\textbf{核心原理}：把连续暴露按分位放几个结点，用一段段三次多项式在结点处光滑拼接。用样条项的 Wald 联合检验判断非线性是否显著（$p<0.05$ 即存在非线性）。

\textbf{操作要点}：工具箱选"RCS"，指定暴露列与结局列，结点数 quantile 分位（默认 4 个），参考值取 P10。

\textbf{结果关键}：案例睡眠时长与收缩压呈 U 形（非线性 $p<0.001$），最低点约 6.5 小时；结点数 3/4/5 结论一致故稳健。

\textbf{注意}：结点数不宜过多（3--5 个即可），否则曲线过拟合出假波动。
\end{algobrief}

%% ========== 寿命表 ==========
\advancedalgo{寿命表分析}{life_table}

\begin{algobrief}
\textbf{一句话}：把随访时间切成若干区间，逐段计算死亡率和累积生存率，是生存分析最基础的描述工具。

\textbf{美赛场景区}：设备/产品可靠性分析（无故障使用时间）；Kaplan-Meier 的分组汇总版。

\textbf{核心原理}：每个区间用"期初人数 $-$ 删失数/2"估计有效暴露人数，算区间死亡概率，再连乘得累积生存率；Log-Rank 检验组间差异；在生存率首次低于 50\% 处插值得中位生存时间。

\textbf{操作要点}：工具箱选"寿命表"，指定时间列、状态列、分组列，比较检验 logrank。

\textbf{结果关键}：案例改进工艺组中位无故障使用 54.4 月，原工艺组 21.6 月，Log-Rank $p<0.001$。

\textbf{注意}：区间宽度会影响结果，应做宽度敏感性复核（如 $\pm$ 一倍宽度）。
\end{algobrief}

%% ========== 条件逻辑回归 ==========
\advancedalgo{条件逻辑回归}{conditional_logistic}

\begin{algobrief}
\textbf{一句话}：用于\textbf{匹配病例对照}数据——按年龄、性别等把病例和对照一一配对后，在每对内部比较，自动消去匹配因素的干扰。

\textbf{美赛场景区}：匹配设计下的风险因素识别；迁移到"配对样本/分层设计"下的二分类影响因素分析。

\textbf{核心原理}：在每个匹配层内构造条件似然，层间恒定因素（如匹配变量年龄）自动被消去、无法估计。输出各因素 OR 及 CI。

\textbf{操作要点}：工具箱选"条件逻辑回归"，指定匹配层 ID、结局、协变量，自动剔除层内恒定变量。

\textbf{结果关键}：案例 75 个匹配对，孕前 BMI OR=1.23、家族史 OR=3.42、睡眠 OR=0.63（保护），均显著。

\textbf{注意}：匹配因素在模型中自动不可估计（被消去），这是正常现象，不是出错。
\end{algobrief}

%% ========== Ridit ==========
\advancedalgo{Ridit分析}{ridit_analysis}

\begin{algobrief}
\textbf{一句话}：处理有序等级资料（如满意度五级），以参照组为基准把等级折算成得分，比较各组整体等级是偏高还是偏低。

\textbf{美赛场景区}：有序评分（满意度、疗效等级）的多组比较；迁移到任何有序分类结局的组间比较。

\textbf{核心原理}：用参照组各等级累计比例算 Ridit 得分，使参照组平均恒为 0.5；其他组平均 Ridit $>0.5$ 表示等级整体偏高、$<0.5$ 偏低，做 z 检验并 Bonferroni 校正。

\textbf{操作要点}：工具箱选"Ridit"，指定分组列与等级列，方向"等级越高越好"，参照组取最大组。

\textbf{结果关键}：案例线上预约 Ridit=0.70（更满意）、电话预约 0.37（更不满意），均与参照组差异显著。

\textbf{注意}：Ridit 只利用等级顺序、不利用间距；样本小的组区间宽、结论不稳。
\end{algobrief}

%% ========== CMH 分层卡方 ==========
\advancedalgo{分层卡方（CMH检验）}{cmh_test}

\begin{algobrief}
\textbf{一句话}：控制一个分层变量（如不同医院/地区）后，再判断暴露与二分类结局是否关联，避免分层因素造成混杂。

\textbf{美赛场景区}：多中心/多地区数据合并分析时；控制混杂后的关联性检验。

\textbf{核心原理}：每层内算 2$\times$2 表的 OR，用 Mantel-Haenszel 法合并得合并 OR 与合并卡方；Breslow-Day 检验各层 OR 是否一致；比较粗 OR 与调整 OR 判断混杂大小。

\textbf{操作要点}：工具箱选"CMH 分层卡方"，指定暴露、结局、分层列。

\textbf{结果关键}：案例调整中心后合并 OR=0.33（CI 0.15--0.72），ERAS 组并发症优势更低；粗 OR=0.43，调整后变化 30\% 提示存在混杂。

\textbf{注意}：Breslow-Day $p<0.05$ 说明各层效应不一致，应分层解释而非只报合并值。
\end{algobrief}

%% ========== OR/RR ==========
\advancedalgo{OR值和RR值（效应量指标）}{or_rr}

\begin{algobrief}
\textbf{一句话}：从 2$\times$2 四格表出发，量化两组事件风险的相对差异——OR 是比值比、RR 是相对危险度，并给出风险差 AR。

\textbf{美赛场景区}：任何两组率的比较（干预组 vs 对照组事件率）；论文中报告效应大小。

\textbf{核心原理}：RR=暴露组事件率/对照组事件率；OR=暴露组事件比值/对照组事件比值。结局罕见时 OR 近似 RR，结局常见时 OR 会夸大关联。

\textbf{操作要点}：工具箱选"OR/RR"，指定暴露与结局字段，对照组取值，零格自动修正，主指标选 RR。

\textbf{结果关键}：案例干预组事件率 16.7\% vs 对照 37.5\%，RR=0.444（CI 0.25--0.79），风险差 $-20.8$ 个百分点。

\textbf{注意}：结局常见时不要用 OR 代替 RR 解读绝对风险；观察性数据只提示关联。
\end{algobrief}

%% ========== NRI/IDI ==========
\advancedalgo{NRI和IDI指标}{nri_idi}

\begin{algobrief}
\textbf{一句话}：在旧模型基础上加入新变量后，判断预测是否真的\textbf{变好}——即使 AUC 提升很小，NRI/IDI 也能检测出重分类改善。

\textbf{美赛场景区}：比较两个嵌套预测模型（旧模型 vs 加了新特征的新模型），量化新增量价值。

\textbf{核心原理}：NRI 看事件组风险被正确上调、非事件组被正确下调的比例之和；IDI 看事件组与非事件组平均预测风险分离度的改善。两者都给 Bootstrap CI，不含 0 即显著改善。

\textbf{操作要点}：工具箱选"NRI/IDI"，指定结局、旧模型特征、新模型新增特征，模式 both。

\textbf{结果关键}：案例加入血清乳酸后连续 NRI=0.768（$p<0.001$）、IDI=0.171（$p<0.001$），AUC 仅提升 0.065 但重分类显著改善。

\textbf{注意}：AUC 变化常不敏感，NRI/IDI 是检测增量价值的更灵敏工具。
\end{algobrief}

%% ========== Kappa ==========
\advancedalgo{Kappa一致性检验}{kappa_test}

\begin{algobrief}
\textbf{一句话}：评价两个评分者对同一批对象分类结果的一致性，\textbf{扣除了"碰巧一致"}的部分，比单纯算符合率更严格。

\textbf{美赛场景区}：两个评委/两种分类器/两次标注结果的一致性；标注质量检验。

\textbf{核心原理}：观察一致率减去机遇一致率，再除以 $1-$机遇一致率。Kappa $=1$ 完全一致、$=0$ 等于瞎猜。Landis-Koch 分级：$<0.2$ 极差、$0.4$--$0.6$ 中等、$0.6$--$0.8$ 良好、$>0.8$ 很好。有序等级可用加权 Kappa。

\textbf{操作要点}：工具箱选"Kappa"，两列成对评分，权重方案按数据选，判读标准 Altman。

\textbf{结果关键}：案例两位医师 CT 分级 Kappa=0.721（CI 0.63--0.81），良好；分歧集中在"可疑恶性"等级。

\textbf{注意}：符合率高但 Kappa 低，说明两人都倾向选多数类别（边际分布偏），一致性其实差。
\end{algobrief}

%% ========== McNemar ==========
\advancedalgo{配对卡方（McNemar检验）}{mcnemar_test}

\begin{algobrief}
\textbf{一句话}：检验\textbf{同一批对象}前后两次二分类结果是否变化（如保养前 vs 保养后设备隐患率），是配对设计下的率比较。

\textbf{美赛场景区}：同一组对象干预前后对比；配对样本的两个相关率比较。

\textbf{核心原理}：只看不一致的配对——"由阳转阴"对数 $b$ 与"由阴转阳"对数 $c$ 是否对称。若 $b$ 和 $c$ 明显不等，说明前后发生了系统改变。一致对不参与检验。

\textbf{操作要点}：工具箱选"McNemar"，指定前后两列，$b+c \geq 25$ 用渐近卡方，否则精确检验。

\textbf{结果关键}：案例 $b=38$（阳性转阴性）、$c=13$（阴性转阳性），率差 $-16.7$ 个百分点、$p=0.0008$，保养后隐患率显著下降。

\textbf{注意}：McNemar 只看率是否变化，效应大小要另外看率差及其 CI。
\end{algobrief}

%% ========== 重复测量方差分析 ==========
\advancedalgo{重复测量方差分析}{rm_anova}

\begin{algobrief}
\textbf{一句话}：同一批对象在多个时间点（如基线、3 月、6 月）被反复测量，比较均值是否随时间变化，且利用"同一个人"的相关提高检验功效。

\textbf{美赛场景区}：追踪/面板数据的时间趋势分析；同一批客户/对象多次测量的比较。

\textbf{核心原理}：把个体间差异从误差中分离，做时间主效应 F 检验。前提是球形性（Mauchly 检验），不满足时用 Greenhouse-Geisser 校正自由度。显著后做事后两两比较（Bonferroni 校正）。

\textbf{操作要点}：工具箱选"重复测量方差分析"，长格式，指定被试、时间、测量值列，球形性校正 GG。

\textbf{结果关键}：案例 50 人 3 时点收缩压时间主效应 $F=193.4$、$p<0.001$，两两比较均显著下降。

\textbf{注意}：违背球形性时必须看校正后 p 值，不能直接用原始 F 检验。
\end{algobrief}

%% ========== 竞争风险模型 ==========
\advancedalgo{竞争风险模型}{competing_risk}

\begin{algobrief}
\textbf{一句话}：当一个人可能因\textbf{多种互斥原因}退出（如腹透患者要么技术失败、要么死亡），把别的原因当删失会高估关注事件发生率，竞争风险模型正确处理这一点。

\textbf{美赛场景区}：迁移到"多种互斥结局"的风险分析（如客户可能流失也可能升级套餐）；更真实地估计累积发生率。

\textbf{核心原理}：用 Aalen-Johansen 法估计各原因累计发生函数（CIF），用 Fine-Gray 子分布风险回归估计因素对关注事件的影响（SHR）。

\textbf{操作要点}：工具箱选"竞争风险"，指定时间、多分类结局、关注事件取值、协变量，Bootstrap 1000 次给 CI。

\textbf{结果关键}：案例 12 月时技术失败 CIF=0.191，血清白蛋白 SHR=0.91（保护）；注意不要用"$1-$KM"口径（会高估）。

\textbf{注意}：存在竞争事件时，普通 Cox/KM 会有偏，必须用 CIF/Fine-Gray。
\end{algobrief}

%% ========== 卡方拟合优度 ==========
\advancedalgo{卡方拟合优度检验}{chisq_gof}

\begin{algobrief}
\textbf{一句话}：检验一组分类观测的实际频数分布是否符合某个\textbf{理论比例}（如工艺基线规定的 4:3:2:1）。

\textbf{美赛场景区}：检验实际分布是否偏离预期标准（如质量缺陷构成、民意分布是否符合预测）。

\textbf{核心原理}：$\chi^2 = \sum (O-E)^2/E$，$O$ 为观测频数、$E$ 为理论期望频数。$\chi^2$ 越大、$p<0.05$ 表示实际分布显著偏离理论。用贡献 $(O-E)^2/E$ 定位是哪类偏了。

\textbf{操作要点}：工具箱选"卡方拟合优度"，指定类别列与理论比例，期望频数 $<5$ 时用蒙特卡罗精确检验。

\textbf{结果关键}：案例 $\chi^2=11.89$（df=3）、$p=0.008$，拒绝"符合基线构成"；焊接缺陷贡献最大（偏多）。

\textbf{注意}：这是\textbf{单样本}的分布检验，与两变量独立性的卡方检验不同。
\end{algobrief}

%% ========== 剂量反应 LD50 ==========
\advancedalgo{剂量反应分析（LD50）}{ld50_analysis}

\begin{algobrief}
\textbf{一句话}：拟合剂量--反应 S 形曲线，反解出使 50\% 受试对象出现反应的剂量（LD50），量化毒性/效力强弱。

\textbf{美赛场景区}：迁移到"暴露强度--响应率"关系的阈值估计（如价格区间--购买率、污染浓度--患病率）。

\textbf{核心原理}：对剂量取对数后拟合 Probit（或 Logistic）广义线性模型，得到 S 形曲线；用 Fieller 法反解反应率 10\%/50\%/90\% 对应的剂量（LD10/LD50/LD90）。失拟检验 $p>0.05$ 说明曲线拟合良好。

\textbf{操作要点}：工具箱选"剂量反应 LD50"，link 选 probit，剂量 log10 变换，输出 LD10/50/90。

\textbf{结果关键}：案例 LD50=252.8（CI 212.5--301.1），LD50 越小毒性越强；失拟 $p=0.91$ 拟合良好。

\textbf{注意}：剂量范围必须覆盖反应率 0\%--100\%，否则外推 LD50 不可靠。
\end{algobrief}
% 第11章 机器学习
\chapter{机器学习}

\begin{chapteroverview}
\textbf{这个分类是干什么的？}

前面章节的统计模型大多是"人先假设一个关系形式（如线性），再让数据去拟合"。机器学习反过来：把大量特征喂给算法，让它\textbf{自己学}出特征与结果之间的复杂规律，再拿去预测新样本。它特别擅长非线性关系、特征交互多、样本量大的预测任务。

\textbf{美赛什么题型常用？}

机器学习是美赛 C 题（大数据题）和预测类问题的\textbf{主力方法}：
\begin{itemize}
    \item \textbf{分类问题}：会不会流失、是否违约、属于哪类——随机森林、SVM、神经网络、XGBoost
    \item \textbf{回归问题}：预测租金、销量、价格——随机森林、SVR、神经网络
    \item \textbf{关联挖掘}：哪些商品常被一起买——Apriori 关联规则
    \item 美赛中\textbf{随机森林、SVM、神经网络}是三类最常用、评审最认可的方法
\end{itemize}

\textbf{学习路径建议}

先建立"预测模型怎么评价"的共同语言（准确率、精确率、召回率、F1、AUC），再从\textbf{决策树}（最可解释）入手理解树模型；然后学\textbf{随机森林}（树的集成，稳）、\textbf{SVM}（小样本非线性分类利器）、\textbf{KNN} 和\textbf{朴素贝叶斯}（简单基线）；\textbf{神经网络}理解其万能逼近能力但要注意过拟合。进阶的 GBDT/XGBoost/LightGBM/CatBoost/AdaBoost/ExtraTrees 都是树模型集成的变体，理解了随机森林再触类旁通即可。
\end{chapteroverview}

\section{基础级算法}

%% ========== 决策树 ==========
\basicalgo{决策树}{decision_tree}

\begin{algointro}
决策树就是一棵"如果...就..."的问答树：从根节点开始，按某个特征一问一答向下分叉，最后落到一片叶子上给出分类或预测。它是所有树模型（随机森林、GBDT）的基础，规则完全透明、最容易解释。
\end{algointro}

\begin{algoscene}
\begin{itemize}
    \item 美赛中需要\textbf{可解释规则}的预测问题（如"满足什么条件订单会延误"）
    \item 给复杂黑箱模型做对照基线
    \item 特征既有数值又有类别、关系非线性时
\end{itemize}
\end{algoscene}

\begin{algoprinciple}
想象你判断一个订单会不会延误：先问"发货时段是不是高峰？"是就走左支、否走右支；左支再问"骑手历史准时率是否 $<80\%$？"……一层层问下去，直到给出结论。

\begin{itemize}
    \item 每次分裂选一个\textbf{最能把两类样本分开}的特征和阈值
    \item 分类用\textbf{基尼不纯度（Gini）}衡量节点"纯不纯"，分裂后越纯越好
    \item 叶子节点里多数类的类别就是预测结果
    \item \textbf{剪枝}（max\_depth、min\_samples\_leaf）防止树长得太细、死记训练数据
\end{itemize}

生活类比：决策树就像查字典——每一步二选一把范围缩小，最后定位答案。单棵树容易过拟合（ memorize 噪声），所以实际多用它的集成版（随机森林、GBDT）。
\end{algoprinciple}

\begin{algosposs}
\begin{enumerate}
    \item 在 SPSS 中通过 \textbf{分析 → 预测建模 → 构建决策树}（或 Tree 节点）实现
    \item 在 Dabbit AI SPSS 工具箱中选"决策树"，指定目标列（如是否延误）
    \item 任务类型 auto（分类/回归自动判断），分裂准则 gini，splitter=best
    \item 控制过拟合：max\_depth 不限制时配合 min\_samples\_leaf、ccp\_alpha 剪枝
    \item 按类别分层切分训练/测试集，预处理（缺失填补、独热编码）只在训练折内拟合，避免数据泄露
\end{enumerate}

\textbf{默认参数参考}：criterion=gini，splitter=best，max\_depth 不限，min\_samples\_split=2，随机种子 42。
\end{algosposs}

\begin{algoresult}
\begin{itemize}
    \item \textbf{混淆矩阵}：行=真实类别、列=预测类别，对角线越集中越好。案例测试集准确率 76.3\%，高于多数类基线 55.3\%
    \item \textbf{精确率/召回率/F1}：见下文统一解释
    \item \textbf{特征重要性}：发货时段、骑手准时率、交通等级各贡献约 25\%，是主要判断依据
    \item \textbf{树结构图}：可直接翻译成"如果...则..."的业务规则
\end{itemize}

\textbf{怎么判断好坏}：测试集准确率显著高于多数类基线、交叉验证各折稳定，就是可用；树太深（如案例 6 层 19 叶）要警惕过拟合。
\end{algoresult}

\begin{algonotes}
\textbf{什么时候用？} 需要可解释规则、特征类型混合、关系非线性时。

\textbf{什么时候不用？} 追求极致精度时单棵树不够，用随机森林/GBDT；特征量纲敏感问题上不如线性模型直接。

\textbf{常见坑}
\begin{itemize}
    \item \textbf{坑1}：不剪枝让树无限生长，训练集准确率 100\% 测试集很差（过拟合）
    \item \textbf{坑2}：在测试集上做缺失填补/编码，造成数据泄露
    \item \textbf{坑3}：把特征重要性当因果——它只表示预测贡献大小
\end{itemize}

\textbf{和相似算法的区别}：单棵决策树可解释但不稳；随机森林是很多棵树投票，更稳更准但牺牲解释性。
\end{algonotes}

%% ========== 随机森林 ==========
\basicalgo{{\starmark}随机森林}{random_forest}

\begin{algointro}
随机stvo森林就是"一千棵树的民主投票"——训练很多棵决策树，每棵只看部分样本、部分特征，最后让所有树投票决定结果。它是美赛里最常用、最稳、调参最少的机器学习模型。
\end{algointro}

\begin{algoscene}
\begin{itemize}
    \item \textbf{美赛预测任务的首选基线}：分类（流失/违约/患病）或回归（数值预测）都能用
    \item 特征多、关系非线性、有交互时表现稳健
    \item 同时给出特征重要性，帮助识别关键变量
\end{itemize}
\end{algoscene}

\begin{algoprinciple}
生活类比：判案时不是只听一个法官（单棵树），而是请 500 个法官各自\textbf{独立}审一部分证据、投一票，少数服从多数。每个法官都可能判错，但集体投票就稳得多。

为什么"随机"？
\begin{itemize}
    \item \textbf{样本随机}：每棵树用有放回抽样（bootstrap）得到的不同子样本训练
    \item \textbf{特征随机}：每次分裂只从随机抽的一部分特征（分类常用 sqrt(p)）里选最优
    \item 这两点保证树之间\textbf{不一样}，否则大家犯一样的错，集成就没意义
\end{itemize}

最后分类按多数票、回归按平均。特征重要性由每棵树分裂时不纯度下降的总和归一化得到。
\end{algoprinciple}

\begin{algosposs}
\begin{enumerate}
    \item 在 SPSS 中通过 \textbf{分析 → 预测建模} 或在 Dabbit AI 工具箱选"随机森林"
    \item 指定目标列（如是否流失），排除 ID 列；类别不平衡时设 class\_weight=balanced
    \item 分裂准则 gini，max\_features=sqrt，n\_estimators 通常几百棵
    \item 按类别分层切训练/测试集，训练集内做 K 折交叉验证
    \item 缺失值中位数/众数填补，类别变量独热编码，全程固定随机种子
\end{enumerate}

\textbf{默认参数参考}：criterion=gini，max\_features=sqrt，class\_weight=balanced，分层切分，5 折交叉验证。
\end{algosposs}

\begin{algoresult}
\begin{itemize}
    \item \textbf{准确率}：案例测试集 0.868，显著高于多数类基线 0.750
    \item \textbf{AUC}：0.839，区分能力良好（见 ROC 章节解释）
    \item \textbf{宏 F1}：0.792
    \item \textbf{特征重要性}：在网时长（0.437）、月均流量、月通话时长、月均消费贡献最大——前 5 个主导
    \item \textbf{交叉验证}：5 折准确率 $0.813\pm0.056$，与测试集一致说明稳定不过拟合
\end{itemize}

\textbf{怎么判断}：测试集指标明显高于基线、交叉验证波动小、训练/测试差距不大，就是好模型。
\end{algoresult}

\begin{algonotes}
\textbf{什么时候用？} 大多数表格类预测任务的默认首选，"先跑个随机森林看看"。

\textbf{什么时候不用？} 需要极强可解释性（用单棵树/逻辑回归）；需要输出精确概率（概率校准偏差，先校准）；特征极稀疏高维文本（用线性模型/神经网络）。

\textbf{常见坑}
\begin{itemize}
    \item \textbf{坑1}：在训练集上算准确率自夸。必须报独立测试集和交叉验证
    \item \textbf{坑2}：类别不平衡还不加权，模型偏向多数类。用 class\_weight=balanced 或报告平衡准确率
    \item \textbf{坑3}：把特征重要性当成因果方向
    \item \textbf{坑4}：泄露字段（如用了未来才知道的信息）进模型，虚高
\end{itemize}

\textbf{和相似算法的区别}：随机森林 bagging（树之间独立并行投票）；GBDT/XGBoost boosting（树之间串行纠错）。两者都是树集成，RF 更稳更省心。
\end{algonotes}

%% ========== SVM ==========
\basicalgo{{\starmark}支持向量机（SVM/SVR）}{svm}

\begin{algointro}
支持向量机寻找一条"最宽的分界线"把两类样本分开——不仅要分开，还要让离边界最近的样本到边界的距离（间隔）尽可能大。它在小样本、高维、非线性问题上特别强，分类版叫 SVM、回归版叫 SVR。
\end{algointro}

\begin{algoscene}
\begin{itemize}
    \item 美赛中\textbf{样本量适中、特征维度不高}的二分类问题（如客户流失预测）
    \item 数据有非线性边界时，用核函数升级
    \item 作为与随机森林并列的常用分类器对照
\end{itemize}
\end{algoscene}

\begin{algoprinciple}
生活类比：在两拨人之间画一条三八线，你会尽量画在\textbf{两拨人中间最空的地方}，让两边都离得越远越好——这样即使来了新人，也不容易站错边。离边界最近的那几个样本点就叫\textbf{支持向量}（它们"支撑"了这条边界）。

关键概念：
\begin{itemize}
    \item \textbf{软间隔参数 C}：C 越大越不容忍分错样本（可能过拟合），C 越小越宽容（可能欠拟合）
    \item \textbf{核函数 kernel}：线性核（linear）处理线性可分；RBF 核处理非线性，靠 \texttt{gamma} 控制边界弯曲程度
    \item \textbf{特征标准化}：SVM 对量纲极敏感，必须先把各特征标准化到均值 0、方差 1，否则量纲大的特征会主导距离
\end{itemize}

通过网格搜索 + 交叉验证选最优的 C、gamma、kernel 组合。
\end{algoprinciple}

\begin{algosposs}
\begin{enumerate}
    \item 在 SPSS 中通过 \textbf{分析 → 预测建模} 或 Dabbit AI 工具箱选"支持向量机"
    \item 任务类型选分类（SVC），指定目标列（如是否流失）
    \item 先做特征标准化（只在训练折内 fit，避免泄露）
    \item 网格搜索：C 候选（如 0.1, 1, 10）、gamma 候选（如 0.01, 0.1, 1）、核函数 linear 与 rbf
    \item 5 折交叉验证选最优参数，再在独立测试集评估
\end{enumerate}

\textbf{默认参数参考}：C 候选 0.1/1/10，gamma 候选 0.01/0.1，核函数 linear+rbf，评分指标 accuracy。
\end{algosposs}

\begin{algoresult}
\begin{itemize}
    \item \textbf{最优参数}：案例 C=1.0、gamma=0.01、kernel=linear，CV 均分 0.902
    \item \textbf{测试集准确率}：78.9\%，高于多数类基线 68.4\%
    \item \textbf{宏平均精确率 0.761、召回率 0.734、F1 0.744}（指标含义见下方统一说明）
    \item \textbf{支持向量占比}：案例 35.7\%，占比越高说明类别重叠越多、边界越复杂
    \item \textbf{混淆矩阵}：真实类 0 中 23 判对、3 误判；真实类 1 中 7 判对、5 漏判
\end{itemize}

\textbf{分类指标统一解释（重要）}：
\begin{itemize}
    \item \textbf{准确率 Accuracy}：所有样本中预测对的比例。类别不平衡时会骗人
    \item \textbf{精确率 Precision}：预测为正的样本里真为正的比例（"说是的里面有多少真是"）——控制误报
    \item \textbf{召回率 Recall（灵敏度）}：真为正的样本里被找出来的比例（"真病人找出来多少"）——控制漏诊
    \item \textbf{F1}：精确率和召回率的调和平均，两者兼顾时的综合分
    \item \textbf{AUC}：排序区分能力（见第 10 章 ROC）
\end{itemize}
\end{algoresult}

\begin{algonotes}
\textbf{什么时候用？} 中小样本、需要非线性边界、特征数适中的分类/回归。

\textbf{什么时候不用？} 样本量巨大（SVM 训练慢）；特征类别很多/缺失复杂（树模型更省心）；需要直接输出概率（SVM 概率需额外校准）。

\textbf{常见坑}
\begin{itemize}
    \item \textbf{坑1}：不标准化特征。量纲差异会让 SVM 失效
    \item \textbf{坑2}：C 和 gamma 不调参，用默认值碰运气。必须网格搜索 + 交叉验证
    \item \textbf{坑3}：类别不平衡不处理，用 accuracy 评模型
\end{itemize}

\textbf{和相似算法的区别}：SVM 靠间隔最大化+核函数；随机森林靠树集成。小样本非线性 SVM 常更优，大数据/多类型特征 RF/GBDT 更稳。
\end{algonotes}

%% ========== KNN ==========
\basicalgo{K近邻（KNN）}{knn}

\begin{algointro}
KNN 的思想极其朴素：\textbf{近朱者赤，近墨者黑}。要判断一个新样本属于哪类，就看它周围最近的 K 个邻居大多数是哪类，它就是哪类。它不需要训练模型，是"惰性学习"的代表。
\end{algointro}

\begin{algoscene}
\begin{itemize}
    \item 美赛中作为\textbf{简单基线模型}，快速对比其他复杂方法是否真的更优
    \item 特征已经很好地标准化、类别间有清晰局部结构的分类/回归
    \item 教学演示"相似性分类"思想
\end{itemize}
\end{algoscene}

\begin{algoprinciple}
生活类比：你搬到一个新小区，想判断这小区是"高风险"还是"低风险"，不用算什么模型，只要看看离你最近的 K 户邻居大多是什么情况，投票决定。

\begin{itemize}
    \item 计算新样本与\textbf{所有训练样本}的距离（欧氏距离）
    \item 取最近的 K 个邻居，分类按多数投票、回归按平均
    \item \textbf{K 的选择很关键}：K 太小容易被噪声带偏（过拟合），K 太大边界过度平滑（欠拟合）
    \item \textbf{必须标准化特征}：否则"月消费几千元"会压倒"在网时长几十月"这类量纲小的特征
\end{itemize}
\end{algoprinciple}

\begin{algosposs}
\begin{enumerate}
    \item 在 Dabbit AI 工具箱选"K 近邻 KNN"，指定目标列与数值特征
    \item 特征标准化（均值 0、方差 1），距离度量 euclidean
    \item 候选 K 设为奇数序列（1,3,5,...,29），用 5 折交叉验证按 accuracy 选最优 K
    \item 案例交叉验证选出 K=9，权重 uniform
    \item 在独立测试集评估并输出逐类精确率/召回率
\end{enumerate}

\textbf{默认参数参考}：n\_neighbors=5，metric=euclidean，weights=uniform，standardize=True，5 折 CV。
\end{algosposs}

\begin{algoresult}
\begin{itemize}
    \item \textbf{最优 K}：案例交叉验证 K=9 时均分最高（0.752）
    \item \textbf{测试集准确率}：0.822，宏平均 F1 0.826
    \item \textbf{逐类精确率/召回率}：高风险类精确率 0.923（判高风险的基本没错）、召回率 0.857
    \item \textbf{交叉验证 K 曲线}：K 过小波动大、K 过大得分下降，据此选平衡点
\end{itemize}

\textbf{怎么判断}：测试集准确率显著高于随机基线、各类召回率均衡即可；KNN 不给因果解释，只体现局部相似结构。
\end{algoresult}

\begin{algonotes}
\textbf{什么时候用？} 快速基线、特征已标准化、样本量不大。

\textbf{什么时候不用？} 样本量大（预测时要和所有样本算距离，慢）；特征量纲差异大又不标准化；高维稀疏数据。

\textbf{常见坑}
\begin{itemize}
    \item \textbf{坑1}：不标准化特征，量纲大的列主导距离
    \item \textbf{坑2}：K 取偶数导致投票平局；K 太大/太小
    \item \textbf{坑3}：样本量大时 KNN 预测极慢却还硬用
\end{itemize}

\textbf{和相似算法的区别}：KNN 是"懒惰"的基于实例方法，不训练参数；SVM/随机森林会训练出明确的决策边界。
\end{algonotes}

%% ========== 朴素贝叶斯 ==========
\basicalgo{朴素贝叶斯}{naive_bayes}

\begin{algointro}
朴素贝叶斯基于贝叶斯定理，假设"各个特征之间相互独立（朴素）"，通过先验概率和条件概率算出每个类别的后验概率，取最大者为预测类别。它在文本分类（垃圾邮件、情感分析）里又快又好用。
\end{algointro}

\begin{algoscene}
\begin{itemize}
    \item 美赛中文本/词频类分类（如评论情感正/中/负）
    \item 高维稀疏特征（如词袋）下的快速基线
    \item 需要输出后验概率时
\end{itemize}
\end{algoscene}

\begin{algoprinciple}
生活类比：你看到一个人"带伞、看天色"，猜他是不是要下雨了——根据以往经验，下雨天的人多带伞（条件概率），再结合下雨本身的概率（先验），算出"带伞的日子下雨"的后验概率。

\begin{itemize}
    \item \textbf{贝叶斯定理}：$P(\text{类}|\text{特征}) \propto P(\text{类}) \prod P(\text{特征}|\text{类})$
    \item "朴素"在于假设各特征条件独立（现实中往往不独立，但效果居然常常不错）
    \item 多分类时选后验概率最大的类别
    \item Multinomial 变体适合计数/词频特征
\end{itemize}
\end{algoprinciple}

\begin{algosposs}
\begin{enumerate}
    \item 在 SPSS 中通过 \textbf{分析 → 分类 → 朴素贝叶斯}（如该模块可用）
    \item 在 Dabbit AI 工具箱选"朴素贝叶斯"，指定目标列（如情感 sentiment）与特征列
    \item 模型类型 multinomial，评价指标 macro\_f1，缺失用均值填补
    \item 分层切分训练/测试集，预处理只在训练折内做
\end{enumerate}

\textbf{默认参数参考}：model type=multinomial，metric=macro\_f1，missing strategy=mean。
\end{algosposs}

\begin{algoresult}
\begin{itemize}
    \item \textbf{准确率}：案例 84.2\%，宏平均 F1 0.838，ROC-AUC 0.977
    \item \textbf{逐类精确率/召回率/F1}：负面类精确率 1.0（预测负面的全对）、正面类召回率 0.923
    \item \textbf{混淆矩阵}：错误集中在中性与正面之间（这两类特征重叠大）
    \item \textbf{特征区分度}：用 F 统计量排序，"退货""好用""失望"等词最能拉开类别
\end{itemize}

\textbf{怎么判断}：AUC 高、混淆矩阵对角线集中即可；注意朴素贝叶斯的后验概率往往不准（因独立假设），但类别排序常常可用。
\end{algoresult}

\begin{algonotes}
\textbf{什么时候用？} 高维稀疏文本/计数特征、需要快速概率分类。

\textbf{什么时候不用？} 特征间强相关且要精确概率时；特征量纲连续且关系复杂时（用树模型）。

\textbf{常见坑}
\begin{itemize}
    \item \textbf{坑1}：把输出的后验概率当真概率用——独立假设导致概率失真，只适合排序
    \item \textbf{坑2}：零概率问题（某特征在某类没出现过）——需用平滑参数 alpha
    \item \textbf{坑3}：特征强相关时还盲目用（独立假设被严重违反）
\end{itemize}

\textbf{和相似算法的区别}：朴素贝叶斯是概率生成式方法；逻辑回归是判别式方法，后者在特征相关时常更准。
\end{algonotes}

%% ========== 神经网络 ==========
\basicalgo{{\starmark}神经网络（BP/MLP）}{neural_network}

\begin{algointro}
BP 神经网络（多层感知机 MLP）模拟大脑神经元：输入层接收特征，经过一层或多层隐藏层的非线性变换，最后输出预测。它能拟合极其复杂的非线性关系，是深度学习的雏形。
\end{algointro}

\begin{algoscene}
\begin{itemize}
    \item 美赛中特征与目标关系高度复杂、样本量较大的预测问题（如房源租金预测）
    \item 作为与树模型对比的"万能拟合器"
    \item 图像、文本等非结构化数据的建模基础
\end{itemize}
\end{algoscene}

\begin{algoprinciple}
生活类比：神经网络像一条流水线，每个工人（神经元）对上一层传来的信号做加权求和再套个非线性激活函数，传给下一层。层层加工后，最后一层给出结论。训练时用\textbf{误差反向传播（BP）}：先前向算预测误差，再从后往前逐层调整每个权重，让误差越来越小。

\begin{itemize}
    \item \textbf{隐藏层 (30,30)}：两层各 30 个神经元
    \item \textbf{激活函数 ReLU}：引入非线性，否则多层叠加等价于线性回归
    \item \textbf{正则系数 alpha}：抑制过拟合
    \item \textbf{标准化必须做}：神经网络对量纲和初始值敏感
    \item \textbf{早停（early stopping）}：验证集不再改善就停，防过拟合
\end{itemize}
\end{algoprinciple}

\begin{algosposs}
\begin{enumerate}
    \item 在 SPSS 中通过 \textbf{分析 → 神经网络 → 多层感知器} 实现
    \item 在 Dabbit AI 工具箱选"神经网络 BP/MLP"，指定目标列与特征列
    \item 隐藏层结构 (30,30)，激活 relu，求解器 sgd，学习率自适应
    \item 数值特征标准化、类别特征独热编码、缺失中位数/众数填补（均在训练折内）
    \item 70/30 切分训练/测试集，观察损失曲线是否收敛、是否过拟合
\end{enumerate}

\textbf{默认参数参考}：hidden layer sizes=(30,30)，activation=relu，solver=sgd，learning rate=adaptive。
\end{algosposs}

\begin{algoresult}
\begin{itemize}
    \item \textbf{回归指标}：案例测试集 $R^2=0.882$、RMSE=5.14、MAE=3.52，远优于均值基线（$R^2\approx0$）
    \item \textbf{训练 vs 测试得分}：训练 0.924、测试 0.882，差距小说明没过拟合
    \item \textbf{损失曲线}：走平说明收敛；若测试损失后期回升则过拟合
    \item \textbf{排列重要性}：配套评分、距地铁距离、所在区域影响最大
\end{itemize}

\textbf{怎么判断}：测试集 $R^2$ 高且与训练集差距小、损失曲线收敛，就是可用的神经网络。
\end{algoresult}

\begin{algonotes}
\textbf{什么时候用？} 复杂非线性、样本充足、特征经过预处理的预测问题。

\textbf{什么时候不用？} 样本量小（神经网络需要数据喂养）；需要强可解释性；简单线性关系（杀鸡用牛刀）。

\textbf{常见坑}
\begin{itemize}
    \item \textbf{坑1}：不标准化输入，网络难收敛
    \item \textbf{坑2}：隐藏层太大/训练轮数太多，死记训练噪声（过拟合）。用正则 alpha 和早停
    \item \textbf{坑3}：把神经网络当万能最优。表格数据上 GBDT/XGBoost 常常又快又好，不必硬上深度学习
    \item \textbf{坑4}：结果不可复现却不固定随机种子
\end{itemize}

\textbf{和相似算法的区别}：神经网络是万能非线性拟合器但黑箱、需调参；树集成（随机森林/GBDT）在中等规模表格数据上往往更省心。
\end{algonotes}

\section{进阶级算法速览}

%% ========== 决策树集成家族 ==========
\advancedalgo{GBDT梯度提升树}{gbdt}

\begin{algobrief}
\textbf{一句话}：串行地逐棵训练树，每棵新树专门去\textbf{拟合前面所有树残差}（预测错的部分），一步步纠错。

\textbf{美赛场景区}：表格数据分类/回归（如客户流失预测），精度常高于随机森林。

\textbf{核心原理}：加法模型，每轮拟合损失函数的负梯度（残差近似），用学习率和子采样控制过拟合；早停防止过拟合。

\textbf{操作要点}：工具箱选 GBDT，指定目标列与特征列，分类任务对数损失，调 n\_estimators、learning rate、max depth。

\textbf{结果关键}：案例测试集准确率 0.80、AUC 0.845，优于多数类基线 0.644；前 5 特征贡献 81\%。

\textbf{注意}：比随机森林更易过拟合，需配合学习率、树深、早停调参。
\end{algobrief}

\advancedalgo{XGBoost}{xgboost}

\begin{algobrief}
\textbf{一句话}：GBDT 的工程优化版，加入正则项、二阶导数信息和缺失值原生处理，是 Kaggle 表格数据竞赛的常胜将军。

\textbf{美赛场景区}：高精度分类/回归（如信贷违约预测），特征工程到位时精度最高。

\textbf{核心原理}：在损失中加入一二阶泰勒展开和正则项（L1/L2），控制模型复杂度；支持早停和缺失值自动分裂。

\textbf{操作要点}：工具箱选 XGBoost，学习率 0.1、树深 3、最多 200 轮、早停 30 轮，分类用 logloss。

\textbf{结果关键}：案例测试集准确率 0.80、AUC 0.90；信用评分、年龄、负债收入比分裂增益最高。

\textbf{注意}：参数多需调；特征重要性是增益而非因果；防泄露用分层三分（训练/验证/测试）。
\end{algobrief}

\advancedalgo{LightGBM}{lightgbm}

\begin{algobrief}
\textbf{一句话}：XGBoost 的"快速版"，用直方图算法和 leaf-wise 生长，训练更快、内存更省，尤其适合大数据。

\textbf{美赛场景区}：样本量大、特征多的表格预测（如信贷违约）；需要快速迭代调参时。

\textbf{核心原理}：按直方图把连续特征离散化加速分裂，按 leaf-wise（优先分裂损失下降最大的叶子）生长，支持类别特征原生处理。

\textbf{操作要点}：工具箱选 LightGBM，任务 binary，指定类别特征列，早停定迭代轮数。

\textbf{结果关键}：案例测试集 AUC 0.818，早停最佳 103 轮；年收入、历史逾期次数增益最高。

\textbf{注意}：leaf-wise 易过拟合小样本，需限制 num\_leaves 和 min\_child\_samples。
\end{algobrief}

\advancedalgo{CatBoost}{catboost}

\begin{algobrief}
\textbf{一句话}：专为\textbf{类别特征多}的数据优化的梯度提升树，能原生高质量处理类别变量，少做特征工程。

\textbf{美赛场景区}：类别特征丰富的预测（如电信客户流失，套餐类型/办理渠道/区域）。

\textbf{核心原理}：用Ordered Target Statistics 编码类别特征，避免目标泄露；同样是 boosting 框架，默认参数就很稳。

\textbf{操作要点}：工具箱选 CatBoost，损失 Logloss，深度 6、学习率 0.05、最多 500 轮早停，指定类别特征列。

\textbf{结果关键}：案例测试集 AUC 0.856、准确率 0.739（基线 0.522）；合约状态重要性最高。

\textbf{注意}：类别特征很多时 CatBoost 比 XGBoost 省心，但小样本上未必占优。
\end{algobrief}

\advancedalgo{AdaBoost自适应提升}{adaboost}

\begin{algobrief}
\textbf{一句话}：每轮\textbf{提高被错分样本的权重}，让下一个弱学习器（通常是很浅的树）重点关注难样本，最后加权组合。

\textbf{美赛场景区}：二分类任务的稳健 boosting 基线；样本难度不均时。

\textbf{核心原理}：用决策树桩（深度 1）作弱学习器，按样本分类对错动态调整权重，加权投票集成。

\textbf{操作要点}：工具箱选 AdaBoost，弱学习器决策树桩，轮数 50、学习率 1.0，分层 K 折交叉验证。

\textbf{结果关键}：案例测试集 AUC 0.784、准确率 0.778（基线 0.540）；信用评分、负债收入比重要性最高。

\textbf{注意}：对噪声和异常点敏感（会过度关注错分样本）；类别极不平衡时需加权。
\end{algobrief}

\advancedalgo{极端随机树（ExtraTrees）}{extratrees}

\begin{algobrief}
\textbf{一句话}：比随机森林更"随机"的 bagging 树——分裂点也随机选（不找最优），偏差略高但方差更低、更快。

\textbf{美赛场景区}：随机森林的对照集成；特征噪声大时可能更稳。

\textbf{核心原理}：与随机森林类似但分裂阈值随机选取，进一步降低方差；同时报告不纯度重要性和置换重要性互相印证。

\textbf{操作要点}：工具箱选 ExtraTrees，max\_features=sqrt，分层切分训练/测试。

\textbf{结果关键}：案例平衡准确率 0.776、AUC 0.841（基线 0.589）；缴费方式重要性最高。

\textbf{注意}：当置换重要性明显低于不纯度重要性时，说明该特征冗余或高基数干扰，解读要谨慎。
\end{algobrief}

%% ========== 关联规则 ==========
\advancedalgo{关联规则（Apriori）}{apriori}

\begin{algobrief}
\textbf{一句话}：从大量购物篮/订单里挖出"买了 A 的人很可能也买 B"的搭配规则，用于推荐和套餐设计。

\textbf{美赛场景区}：购物篮分析、交叉销售推荐；迁移到"哪些事件/项经常共现"的规律挖掘。

\textbf{核心原理}：三步——算频繁项集（支持度）、生成 A$\to$B 规则、按置信度和提升度筛选。
\begin{itemize}
    \item \textbf{支持度 Support}：A 和 B 同时出现的交易占比（覆盖广不广）
    \item \textbf{置信度 Confidence}：买了 A 的人里买 B 的比例（条件概率）
    \item \textbf{提升度 Lift}：$P(B|A)/P(B)$，$>1$ 表示 A 对 B 有正向拉动
\end{itemize}

\textbf{操作要点}：工具箱选 Apriori，指定交易 ID 与项列，设 min\_support、min\_confidence、min\_lift，按 lift 排序。

\textbf{结果关键}：案例最强规则"鲜榨果汁$\to$田园沙拉"，支持度 4.7\%、置信度 70\%、提升度 9.55。

\textbf{注意}：高支持度不等于强规则（可能只是 B 本身畅销），必须看提升度 $>1$ 才算真关联；关联不等于因果。
\end{algobrief}

% 第12章 规划求解
\chapter{规划求解}

\begin{chapteroverview}
\textbf{这个分类是干什么的？}

规划求解就是"在一堆限制条件下找最优方案"。比如：预算有限时选哪些研发项目收益最大？货车怎么装货最划算？快递路线怎么排最短？工厂排产怎么利润最高？这类问题统称为\textbf{优化问题}——有目标（最大利润/最小成本），有约束（预算/容量/时间），要找出一组决策变量让目标最好。

\textbf{美赛什么题型常用？}

规划求解是美赛B题（离散优化/调度分配）和D题（运营管理）的\textbf{核心武器}。典型场景：车辆路径规划（TSP）、任务指派、背包选品、生产排产、资源分配。遗传算法、粒子群、模拟退火等智能优化方法在美赛中\textbf{极其常用}——当问题复杂到没有解析解时，靠它们在巨大的解空间里"搜"出好方案。

\textbf{学习路径建议}

先搞懂\textbf{动态规划（背包问题）}和\textbf{遗传算法}（组合优化的两大主力），再学\textbf{粒子群}和\textbf{模拟退火}（连续优化与路径优化的利器）。进阶了解\textbf{匈牙利算法}（指派）、\textbf{表上作业法}（运输）、\textbf{分支定界/割平面}（整数规划精确求解）等经典方法。注意：SPSS原生不支持优化求解，这些算法通过\textbf{Dabbit AI SPSS工具箱}或Python/R扩展实现。
\end{chapteroverview}

\section{基础级算法}

%% ========== 遗传算法 ==========
\basicalgo{{\starmark}遗传算法}{genetic_algorithm}

\begin{algointro}
遗传算法（GA）模拟"物竞天择、适者生存"的自然进化过程：把每一个候选方案编码成"染色体"，让一群方案在迭代中通过选择、交叉、变异不断进化，最终留下最优方案。它是美赛中\textbf{最常用的智能优化算法}，尤其适合调度、分配、路径等组合优化问题。
\end{algointro}

\begin{algoscene}
\begin{itemize}
    \item \textbf{背包问题}：预算有限时选哪些项目/物品使总价值最大（如研发项目立项）
    \item \textbf{调度与分配}：任务指派、人员排班、生产排产
    \item \textbf{路径优化}：旅行商问题（TSP）、车辆路线规划
    \item 美赛B题/D题中当问题规模大、无法用精确方法求解时
    \item 目标函数不可导、约束复杂、传统方法搞不定的问题
\end{itemize}
\end{algoscene}

\begin{algoprinciple}
想象一个"方案进化岛"：一开始随机撒下一群候选方案（种群），每个方案就是一个"生物个体"，用一串编码表示（比如0/1串表示选或不选某个项目）。

\begin{itemize}
    \item \textbf{适应度评价}：每个方案根据目标函数打分——收益高的方案"更健康"
    \item \textbf{选择（Selection）}：适应度高的方案更容易"繁殖后代"（锦标赛选择：随机抽几个，留最强的）
    \item \textbf{交叉（Crossover）}：两个优秀方案"交配"，各取一段编码拼接成新方案——好的基因组合在一起
    \item \textbf{变异（Mutation）}：随机翻转某一位编码，保持种群多样性，防止早熟收敛
    \item \textbf{精英保留（Elitism）}：每代最好的几个方案直接进入下一代，不丢失
\end{itemize}

如此迭代几百代后，种群会越来越集中到最优解附近。

\textbf{关键参数}：种群大小（默认80）、交叉概率（默认0.8）、变异率（默认0.05）、迭代代数（默认100）。0/1背包问题中用repair策略处理超容量个体。
\end{algoprinciple}

\begin{algosposs}
SPSS原生不支持遗传算法，在\textbf{Dabbit AI SPSS工具箱}中操作：

\begin{enumerate}
    \item 在工具箱中选择"遗传算法"
    \item 上传数据：价值列选"value"，成本列选"weight"，容量设为预算上限（如3200）
    \item 设置参数：种群大小80，迭代100代，交叉概率0.8，变异率0.05，精英保留数2
    \item 超容量处理选"repair"（按价值密度从低到高剔除超容项目）
    \item 稳健性运行次数设为3（多次随机种子验证稳定性）
    \item 点击运行，输出最优方案与收敛曲线
\end{enumerate}

\textbf{参考实现}：也可通过Python的DEAP或scikit-opt库实现，核心是编码方式与适应度函数的设计。
\end{algosposs}

\begin{algoresult}
\begin{itemize}
    \item \textbf{最优总价值/成本}：如选中28个项目，总价值13756，总成本3194（容量利用率99.8\%）
    \item \textbf{收敛曲线}：每代最优适应度随代数上升并趋于平稳——曲线变平说明收敛
    \item \textbf{稳健性}：多次随机种子运行结果应接近（如最优值在13756--13762之间），极差越小越稳定
    \item \textbf{改进幅度}：比随机可行方案基线提升多少（如提升1231.75），说明进化搜索有效
\end{itemize}

\textbf{判断好坏}：遗传算法是\textbf{启发式方法}，不保证全局最优。如果多次运行结果差异大（极差占最优值比例高），应增大种群或代数后重跑。
\end{algoresult}

\begin{algonotes}
\textbf{什么时候用？}
\begin{itemize}
    \item 组合优化问题规模大（如150个物品的背包、100个节点的TSP），精确方法算不动
    \item 目标函数复杂、不可导、有约束
\end{itemize}

\textbf{什么时候不用？}
\begin{itemize}
    \item 小规模问题（$n<20$）——直接用动态规划或穷举更靠谱
    \item 问题有成熟精确算法（如指派问题用匈牙利算法）——没必要用启发式
\end{itemize}

\textbf{常见坑}
\begin{itemize}
    \item \textbf{坑1}：变异率太小→种群多样性不足，早熟收敛在局部最优；太大→随机游走不收敛
    \item \textbf{坑2}：不做多次随机种子运行——单次结果可能只是运气好，不代表真最优
    \item \textbf{坑3}：超容量个体直接丢弃而非修复——会丢失大量可行解，应设计repair策略
\end{itemize}

\textbf{和相似算法的区别}
\begin{itemize}
    \item vs \textbf{动态规划}：动态规划求\textbf{精确全局最优}但只适用于结构化问题（背包、最短路）；遗传算法是\textbf{启发式近似}但适用面广
    \item vs \textbf{粒子群}：遗传算法适合\textbf{离散组合优化}（0/1选择、排列）；粒子群适合\textbf{连续参数优化}
\end{itemize}
\end{algonotes}

%% ========== 动态规划（背包资源分配） ==========
\basicalgo{动态规划（背包资源分配）}{dp_knapsack}

\begin{algointro}
动态规划（DP）是一种"把大问题拆成小问题、记住小问题答案"的精确求解方法。经典应用是\textbf{背包问题}：预算有限时选哪些项目使总收益最大。它能\textbf{证明全局最优}，不像遗传算法只能给近似解。
\end{algointro}

\begin{algoscene}
\begin{itemize}
    \item \textbf{0-1背包问题}：预算/容量有限，每个项目只能选或不选（如技改项目立项）
    \item 资源分配：有限资金在多个活动间分配使总收益最大
    \item 最短路径、设备更新、库存优化等多阶段决策问题
\end{itemize}
\end{algoscene}

\begin{algoprinciple}
想象你在整理一个背包旅行，背包容量有限，面前有一堆物品各有重量和价值，怎么选最划算？

动态规划的思路是\textbf{分阶段、记中间结果}：

\begin{itemize}
    \item 设 $dp[w]$ 表示"容量为 $w$ 时能装的最大价值"
    \item 对每个物品（重量 $c$、价值 $v$），从大到小遍历容量：$dp[w] = \max(dp[w], dp[w-c] + v)$
    \item 意思是：要么不装这个物品（$dp[w]$），要么装（腾出 $c$ 的容量装它，总价值变成 $dp[w-c]+v$），取大者
    \item 从 $w=0$ 递推到容量上限，最后 $dp[W]$ 就是最大价值，再回溯找出选了哪些物品
\end{itemize}

这就像填表：每加一个物品，就把整张表更新一遍。虽然看起来笨，但\textbf{枚举了所有可能组合}，所以一定能找到全局最优。
\end{algoprinciple}

\begin{algosposs}
SPSS原生不支持动态规划，在\textbf{Dabbit AI SPSS工具箱}中操作：

\begin{enumerate}
    \item 上传候选项目数据（含成本列和收益列）
    \item 背包类型选"0-1背包"，容量自动提取问题给定值（如2600）
    \item 缺失值处理选"剔除缺失行"，价值并列处理选"按价值密度优先"
    \item 点击运行，输出最优入选方案与容量敏感性曲线
\end{enumerate}
\end{algosposs}

\begin{algoresult}
\begin{itemize}
    \item \textbf{最大总收益}：如容量2600内最大收益6404，入选15个项目
    \item \textbf{预算利用率}：如100\%，说明钱花完了仍有可收益项目；利用率不足100\%通常是剩余额度装不下任何项目
    \item \textbf{容量敏感性曲线}：展示容量从0到上限每个档位对应的最优收益——额外预算值多少钱一目了然
    \item \textbf{vs 贪心对照}：贪心按单位价值密度排序选，通常比DP少赚（如贪心6315 vs DP 6404），差异说明需要全局权衡
\end{itemize}
\end{algoresult}

\begin{algonotes}
\textbf{什么时候用？}背包/资源分配类问题，且成本和容量可离散化（整数）。

\textbf{什么时候不用？}成本是连续值且很大时（状态空间爆炸），应改用遗传算法等启发式方法。

\textbf{常见坑}
\begin{itemize}
    \item \textbf{坑1}：成本不取整导致状态空间巨大——DP要求容量映射到整数网格
    \item \textbf{坑2}：以为贪心（按价值密度排序）就是最优——贪心无法回溯替换，DP能全局权衡
\end{itemize}

\textbf{和遗传算法的区别}：DP是\textbf{精确方法}（证明最优），但只适用于结构清晰的问题；遗传算法是\textbf{启发式方法}（近似解），适用面更广。
\end{algonotes}

%% ========== 粒子群优化算法 ==========
\basicalgo{{\starmark}粒子群优化算法}{pso}

\begin{algointro}
粒子群优化（PSO）模拟鸟群觅食：一群"粒子"在搜索空间中飞，每个粒子根据自己的历史最佳位置和群体全局最佳位置调整飞行方向，最终聚集到最优解附近。它特别适合\textbf{连续参数优化}问题，比如拟合非线性模型的参数。
\end{algointro}

\begin{algoscene}
\begin{itemize}
    \item \textbf{非线性模型参数拟合}：如衰减振动模型 $y=a+b\cdot e^{-ct}\sin(dt+e)$ 的参数标定
    \item 工程优化：调整多个连续参数使某个指标最优
    \item 神经网络超参数调优
    \item 美赛中当模型是连续的、目标函数不可导时
\end{itemize}
\end{algoscene}

\begin{algoprinciple}
想象一群鸽子在广场上找食物。每只鸽子（粒子）：

\begin{itemize}
    \item 知道自己\textbf{历史上找到过食物最多的位置}（个体极值 $pbest$）
    \item 知道\textbf{整个鸽群找到过食物最多的位置}（全局极值 $gbest$）
    \item 每一步同时被这两个位置"吸引"，调整飞行方向和速度
\end{itemize}

速度更新公式：$v_{t+1} = w\cdot v_t + c_1\cdot r_1\cdot(pbest-x_t) + c_2\cdot r_2\cdot(gbest-x_t)$

\begin{itemize}
    \item $w$（惯性权重）：保持当前飞行方向的趋势，默认从0.9线性递减到0.4（前期探索、后期精修）
    \item $c_1$（认知系数）：向自己历史最佳靠拢的权重
    \item $c_2$（社会系数）：向群体最佳靠拢的权重
    \item $r_1, r_2$：随机数，增加搜索多样性
\end{itemize}

就像鸽子既记得自己以前在哪吃到过好东西，也跟着大部队去它们发现的好地方，最后全群汇聚到食物最丰富的区域。
\end{algoprinciple}

\begin{algosposs}
SPSS原生不支持粒子群优化，在\textbf{Dabbit AI SPSS工具箱}中操作：

\begin{enumerate}
    \item 在工具箱中选择"粒子群优化"
    \item 上传数据，指定自变量列和因变量列
    \item 选择模型族（如exp\_sine衰减振动模型）
    \item 设置参数：粒子数200，最大迭代500次，惯性权重线性递减0.9→0.4，认知/社会系数默认
    \item 独立运行次数设为10（多次验证稳定性），边界处理选"反射"
    \item 点击运行，输出最优参数与收敛曲线
\end{enumerate}
\end{algosposs}

\begin{algoresult}
\begin{itemize}
    \item \textbf{最优参数估计}：如5个参数 $a=1.247, b=2.073, c=0.333, d=2.149, e=0.567$
    \item \textbf{拟合优度}：$R^2=0.9929$，RMSE=0.0464——越接近1/越小越好
    \item \textbf{收敛稳定性}：10次独立运行的最优SSE标准差应极小（如1.09e-07），变异系数小说明结果稳健
    \item \textbf{残差图}：残差围绕零线随机分布、无系统性偏差，说明模型合适
\end{itemize}
\end{algoresult}

\begin{algonotes}
\textbf{什么时候用？}连续参数的非线性优化问题，目标函数不可导或求导困难。

\textbf{什么时候不用？}离散组合优化（如0/1背包、TSP）——PSO虽可改造但不如遗传算法自然。

\textbf{常见坑}
\begin{itemize}
    \item \textbf{坑1}：不做多次独立运行——单次结果可能收敛到局部最优，应至少运行5--10次取最优
    \item \textbf{坑2}：搜索边界设得不对——最优解贴在边界上说明边界设窄了，应扩大
    \item \textbf{坑3}：惯性权重固定不衰减——前期不探索会困在局部，后期不收敛会震荡
\end{itemize}

\textbf{和遗传算法的区别}：PSO是\textbf{连续优化}的利器（鸟群飞行是连续运动）；GA是\textbf{离散组合优化}的主力（编码/交叉/变异天然适合排列和0/1选择）。
\end{algonotes}

%% ========== 模拟退火算法 ==========
\basicalgo{模拟退火算法}{simulated_annealing}

\begin{algointro}
模拟退火（SA）模仿金属退火过程：先高温加热让粒子自由运动（广泛探索），再缓慢降温让粒子逐渐稳定到最低能量状态（收敛到最优）。它的\textbf{绝招是能以一定概率接受劣解}，从而跳出局部最优陷阱。
\end{algointro}

\begin{algoscene}
\begin{itemize}
    \item \textbf{旅行商问题（TSP）}：配送路线优化，找最短环游路线
    \item 调度问题、指派问题等组合优化
    \item 当问题有很多局部最优、贪心算法容易困死时
\end{itemize}
\end{algoscene}

\begin{algoprinciple}
想象你在一个有很多山头和山谷的地形上找最低点，但你不能看地图（不知道全局），只能一步步走。

\begin{itemize}
    \item \textbf{高温阶段}：你喝醉了，连走上坡路都不怕（以概率 $e^{-\Delta E/T}$ 接受劣解），到处乱窜——这样不会困在第一个遇到的小山谷里
    \item \textbf{降温阶段}：你慢慢清醒，越来越不想走上坡路，只接受下降的移动
    \item \textbf{终止}：温度很低时，你基本只往下走，最终停在一个很深的谷底
\end{itemize}

关键是\textbf{Metropolis准则}：新方案比当前好就接受；比当前差，也以概率 $e^{-\Delta E/T}$ 接受——温度高时接受概率大（大胆探索），温度低时接受概率小（谨慎收敛）。

\textbf{参数}：初始温度自动估算，衰减系数0.94（每轮乘0.94），每温迭代500次，终止温度0.0001，重启3次。
\end{algoprinciple}

\begin{algosposs}
SPSS原生不支持模拟退火，在\textbf{Dabbit AI SPSS工具箱}中操作：

\begin{enumerate}
    \item 上传站点坐标数据（TSP问题）
    \item 问题类型自动识别为"坐标型TSP（欧氏距离）"
    \item 设置参数：初始温度自动估算，衰减系数0.94，每温迭代500次，重启3次
    \item 随机种子固定为42（可复现）
    \item 点击运行，输出最优回路与收敛曲线
\end{enumerate}
\end{algosposs}

\begin{algoresult}
\begin{itemize}
    \item \textbf{最优目标值}：如TSP回路总长220.13，比贪心基准252.96改善12.98\%
    \item \textbf{收敛曲线}：前期快速下降、后期走平——走平说明已收敛
    \item \textbf{接受率}：后期接受率低（如7.2\%）说明已进入纯改善搜索
    \item \textbf{重启稳定性}：多次重启最优值离散度小（如12.25\%以内）说明结果可靠
\end{itemize}
\end{algoresult}

\begin{algonotes}
\textbf{什么时候用？}组合优化问题，尤其是TSP路径规划，且贪心/局部搜索容易陷入局部最优时。

\textbf{什么时候不用？}问题规模极小（直接穷举即可），或有精确多项式算法（如指派用匈牙利）。

\textbf{常见坑}
\begin{itemize}
    \item \textbf{坑1}：降温太快→没探索够就收敛，困在局部最优；降温太慢→计算时间太长
    \item \textbf{坑2}：不做多次重启→单次结果受随机初始解影响大
\end{itemize}

\textbf{和遗传算法的区别}：SA是\textbf{单点搜索}（一个方案在迭代中改进），GA是\textbf{种群搜索}（一群方案同时进化）；SA实现简单，GA并行性好但调参更复杂。
\end{algonotes}

\section{进阶级算法速览}

%% ========== 割平面法 ==========
\advancedalgo{割平面法}{cutting_plane}

\begin{algobrief}
\textbf{一句话}：求解整数线性规划的精确方法——先放松整数要求解线性规划，再逐步添加"切割约束"把非整数解切掉，直到得到整数最优解。

\textbf{美赛场景区}：生产排产中产品数量必须为整数（不能生产半个产品）的线性规划问题。

\textbf{核心原理}：先解LP松弛（允许小数），如果最优解含小数变量，就构造Gomory割平面（一条新约束）把小数部分切掉，重新求解。逐轮收紧，直到所有变量取整。就像雕刻：一刀一刀切掉多余的部分，最后留下整数最优解。

\textbf{操作要点}：工具箱中选择"割平面法"，目标方向max，割平面行选择规则选"最大分数部优先"，最大迭代100次。

\textbf{结果关键}：LP松弛最优值逐轮下降并逼近整数最优值（如从45.5降到45），最终找到整数最优解。与独立MILP求解器交叉验证一致即证明最优。

\textbf{注意}：割平面法和分支定界法是整数规划两大精确方法；割平面法靠加约束，分支定界靠分分支。
\end{algobrief}

%% ========== 牛顿法（拟牛顿法） ==========
\advancedalgo{牛顿法（拟牛顿法）}{newton_method}

\begin{algobrief}
\textbf{一句话}：用二阶导数（Hessian矩阵）信息加速非线性优化收敛的方法，比一阶梯度下降快得多。

\textbf{美赛场景区}：非线性模型参数拟合，如药代动力学指数模型 $C(t)=C_0 e^{-k_e t}$ 的参数估计。

\textbf{核心原理}：梯度下降只看"往哪个方向走最陡"（一阶梯度），牛顿法还看"地形的曲率"（二阶导数），像一个会开车的人不仅知道上坡方向，还知道弯道有多急，因此能更快逼近最优点。拟牛顿法（BFGS）不直接算Hessian，用梯度差近似，降低计算量。

\textbf{操作要点}：工具箱中选"牛顿法"，模型选指数模型，方法选auto（牛顿优先、失稳自动切拟牛顿），初值策略选启发式。

\textbf{结果关键}：迭代次数少（如3次就收敛），参数估计值与标准误，RMSE和AIC衡量拟合优度。

\textbf{注意}：牛顿法要求目标函数光滑可导；如果Hessian矩阵不正定会失稳，此时自动切换拟牛顿BFGS。
\end{algobrief}

%% ========== 梯度下降法 ==========
\advancedalgo{梯度下降法}{gradient_descent}

\begin{algobrief}
\textbf{一句话}：最基础的数值优化方法——沿着梯度反方向一步步走，直到到达最低点。神经网络训练的基石。

\textbf{美赛场景区}：线性回归/逻辑回归的参数求解（虽然有解析解，但梯度下降是理解优化的基础）；大规模机器学习模型训练。

\textbf{核心原理}：站在山坡上想走到谷底，每一步都沿着最陡的下坡方向走。学习率控制每步走多大——太大跨过谷底，太小走得慢。加入动量（momentum=0.9）可以像滚球一样加速通过平坦区、抑制震荡。

\textbf{操作要点}：工具箱中选"梯度下降法"，学习率0.01，动量0.9，批量全量，特征标准化开启，迭代上限2000次。

\textbf{结果关键}：损失函数随迭代下降曲线；最终系数与最小二乘解析解的偏差（应极小，如$10^{-4}$量级）；$R^2$。

\textbf{注意}：学习率是最关键参数——需调参；特征必须标准化否则量纲差异导致收敛慢。
\end{algobrief}

%% ========== 差分进化算法 ==========
\advancedalgo{差分进化算法}{differential_evolution}

\begin{algobrief}
\textbf{一句话}：一种基于种群的连续全局优化算法，利用种群中个体间的\textbf{差分向量}进行变异，特别适合非线性模型参数估计。

\textbf{美赛场景区}：酶促反应米氏方程 $v=V\cdot S/(K+S)$ 等非线性模型的全局参数拟合。

\textbf{核心原理}：和遗传算法类似但变异方式不同——随机选三个个体 $a,b,c$，新个体 = $a + F\cdot(b-c)$，用两个体的差异方向来变异，自适应调整步长。交叉后与目标个体贪婪比较，更优才保留。

\textbf{操作要点}：工具箱中选"差分进化"，策略best1bin，变异因子0.5+1，多次不同种子重启，可选L-BFGS-B局部精化。

\textbf{结果关键}：多次重启结果应一致（重启跨度极小）；$R^2$接近1说明拟合好；残差正态性检验通过。

\textbf{注意}：DE是全局搜索，比牛顿法慢但不容易陷局部最优；通常DE粗搜+局部精化（L-BFGS-B）效果最好。
\end{algobrief}

%% ========== 匈牙利算法 ==========
\advancedalgo{匈牙利算法}{hungarian_algorithm}

\begin{algobrief}
\textbf{一句话}：求解\textbf{指派问题}的精确多项式算法（$O(n^3)$）——把 $n$ 个任务分配给 $n$ 个人，使总成本最小。

\textbf{美赛场景区}：快递订单与车辆的最优配对、人员岗位分配、教师排课。

\textbf{核心原理}：每行减去该行最小值、每列减去该列最小值（等价变换不改变最优解），然后用最少的线覆盖所有零元素——如果线数等于 $n$，就能在零元素处找到完美指派；否则继续调整。就像在成本矩阵上"划线找零"，零位置就是最优配对。

\textbf{操作要点}：工具箱中选"匈牙利算法"，目标minimize，求解器auto。

\textbf{结果关键}：总成本（如8对配对总成本72）；贪心基线通常差很多（如123，多70.8\%），证明全局最优的价值。

\textbf{注意}：要求任务数=人数（平衡指派）；不平衡时先加虚拟行/列。
\end{algobrief}

%% ========== 表上作业法（运输问题） ==========
\advancedalgo{表上作业法}{transportation_method}

\begin{algobrief}
\textbf{一句话}：求解\textbf{运输问题}的经典方法——多个仓库向多个门店发货，已知各仓供应量、各店需求量和单位运价，求总运费最低的调运方案。

\textbf{美赛场景区}：连锁零售物流配送调度、工厂到仓库的原材料运输规划。

\textbf{核心原理}：先用伏格尔法（Vogel）找一个较好的初始方案（每次优先选机会成本最大的路线），再用位势法计算每条空闲路线的检验数——如果存在负检验数说明还能降本，沿闭回路调整运量，直到所有检验数$\geq 0$即为最优。

\textbf{操作要点}：工具箱中选"表上作业法"，初始方案伏格尔法，产销不平衡自动加虚拟节点，独立求解器交叉验证。

\textbf{结果关键}：最优总运费（如287802元）；初始方案经11轮调整收敛；所有检验数$\geq 0$证明最优。

\textbf{注意}：供过于求时加虚拟销地（不真运货），供不应求时加虚拟产地。
\end{algobrief}

%% ========== 蚁群算法（ACO） ==========
\advancedalgo{蚁群算法}{aco}

\begin{algobrief}
\textbf{一句话}：模拟蚂蚁觅食时释放信息素的正反馈机制，求解\textbf{旅行商问题（TSP）}的群智能算法。

\textbf{美赛场景区}：150个配送点的最短环游路线规划、网络布线最优路径。

\textbf{核心原理}：蚂蚁走路线时留下信息素，短路线上信息素挥发少、累积快，后面的蚂蚁更倾向选信息素浓的路——正反馈让全群逐渐汇聚到最短路线。$\alpha$控制信息素重要性、$\beta$控制距离启发重要性、挥发率防止信息素无限累积。

\textbf{操作要点}：工具箱中选"蚁群算法"，蚂蚁数30，迭代100次，$\alpha=1, \beta=2$，挥发率0.5，精英权重0。

\textbf{结果关键}：近似最短环游长度（如372.765），比最近邻+2-opt基线改善4.2\%；第45代后收敛稳定。

\textbf{注意}：ACO是启发式近似，不保证全局最优；节点数越大越有优势。
\end{algobrief}

%% ========== 分支定界法 ==========
\advancedalgo{分支定界法}{branch_and_bound}

\begin{algobrief}
\textbf{一句话}：求解整数规划的精确方法——把解空间分成树枝（分支），用松弛问题算上界，上界不如当前最优就剪掉该枝（剪枝），避免穷举。

\textbf{美赛场景区}：0-1背包问题（货车装货选哪些货品收益最大）的精确求解。

\textbf{核心原理}：像查字典——先翻开中间估个上界，如果这半边的上界还不如已经找到的最好方案，就不翻了（剪枝）。按价值密度排序，贪心先找一个好下界，分数背包松弛算上界，上下界收敛到重合即证明最优。

\textbf{操作要点}：工具箱中选"分支定界法"，分支策略best\_first，容量约束，gap tolerance默认。

\textbf{结果关键}：最优总价值（如1265，证明全局最优）；探索节点数（如2057个）；与MILP求解器交叉验证一致。

\textbf{注意}：分支定界和割平面法都是整数规划精确方法；分支定界靠分树剪枝，割平面靠加约束。
\end{algobrief}

%% ========== 禁忌搜索算法 ==========
\advancedalgo{禁忌搜索算法}{tabu_search}

\begin{algobrief}
\textbf{一句话}：带"记忆"的局部搜索——用禁忌表记录最近走过的路，防止立即回退循环，配合特赦准则允许突破禁忌去找更好解。

\textbf{美赛场景区}：TSP路径优化（150个门店巡回补货顺序）。

\textbf{核心原理}：普通局部搜索走两步就可能绕圈回来（A→B→A→B）。禁忌搜索把最近走过的"边"记入禁忌表（如禁忌长度15），强制搜索去新方向；但如果某个禁忌方向能得到历史最优解（特赦准则），还是允许走。

\textbf{操作要点}：工具箱中选"禁忌搜索"，邻域two\_opt，初始解nearest\_neighbor，禁忌长度15，最大迭代600，候选数40。

\textbf{结果关键}：总成本从98590降到82045（改善16.78\%）；最优解在第585次迭代发现；距最小生成树下界差距32.27\%。

\textbf{注意}：禁忌搜索是局部搜索的增强版，依赖初始解质量；适合TSP等排列优化。
\end{algobrief}

%% ========== 内点法 ==========
\advancedalgo{内点法}{interior_point}

\begin{algobrief}
\textbf{一句话}：求解大规模线性规划的现代方法——从可行域\textbf{内部}沿中心路径收敛到最优顶点，不像单纯形法沿边界爬行。

\textbf{美赛场景区}：大规模生产排产优化（60个产品品类、29项共用资源约束），求最优产量组合和瓶颈资源影子价格。

\textbf{核心原理}：单纯形法像沿多边形围墙从一个顶点走到另一个顶点；内点法像从多边形内部直接向最优顶点"钻过去"，变量多时速度快得多。还能算出\textbf{影子价格}——每种紧缺资源放宽1单位能多赚多少钱。

\textbf{操作要点}：工具箱中选"内点法"（HiGHS-IPM），presolve开启，变量下界默认0。

\textbf{结果关键}：最优目标值（如420390）；各紧约束的影子价格（如R01资源影子价格7.233，是瓶颈资源）；约束余量为0的约束即瓶颈。

\textbf{注意}：内点法适合大规模LP；整数规划仍需分支定界。影子价格是敏感性分析的关键输出。
\end{algobrief}

% 第13章 质量控制
\chapter{质量控制}

\begin{chapteroverview}
\textbf{这个分类是干什么的？}

质量控制（QC）用统计方法监控生产/服务过程是否\textbf{稳定受控}，以及产品质量是否\textbf{满足规格要求}。核心工具是控制图（判断过程是否有异常波动）和过程能力指数（判断波动是否在可接受范围内）。它回答两个问题：过程稳不稳？产品合不合格？

\textbf{美赛什么题型常用？}

美赛工业/制造类题目（如生产流程优化、产品质量改进、供应链可靠性）中非常实用。控制图和过程能力分析是六西格玛管理的核心方法。可靠性分析（寿命分布、加速试验）在设备维护、产品寿命预测类题目中也会用到。

\textbf{学习路径建议}

先学\textbf{帕累托图}（抓关键少数问题）和\textbf{Xbar-R控制图}（最常用的计量型控制图），再学\textbf{过程能力分析}（判断过程能不能满足规格）。I-MR控制图用于单值数据，然后了解属性控制图（C图、U图）和测量系统分析（量具R\&R）。SPSS可通过"分析 → 控制图"菜单操作部分功能。
\end{chapteroverview}

\section{基础级算法}

%% ========== 帕累托图 ==========
\basicalgo{{\starmark}帕累托图}{pareto_chart}

\begin{algointro}
帕累托图基于"80/20法则"——80\%的问题是由20\%的原因造成的。它把各类问题按贡献大小从高到低排序，用柱状图加累计折线，帮你找出"关键少数"原因，把改进资源集中在最值得的地方。
\end{algointro}

\begin{algoscene}
\begin{itemize}
    \item 客户投诉分析：哪些故障原因造成了大部分损失？
    \item 产品缺陷分类：哪些缺陷类型最需要改进？
    \item 成本分析：哪些成本项占了大部分费用？
    \item 美赛中需要"抓重点"的问题分析环节
\end{itemize}
\end{algoscene}

\begin{algoprinciple}
帕累托图就是"排排队、算累计"：

\begin{itemize}
    \item 把所有原因按损失金额（或发生次数）\textbf{从大到小排序}
    \item 画柱状图：柱子越高=该原因损失越大
    \item 叠加累计百分比折线：从第一个柱子开始累加
    \item 在累计80\%处画一条线——线左边的就是"A类关键原因"
\end{itemize}

这就像医生看病：不是所有症状都要治，先看哪个病最严重、最花钱，集中力量先治它。

\textbf{ABC分级}：累计0--80\%为A类（关键少数，优先改进），80--95\%为B类（次要），95--100\%为C类（一般）。
\end{algoprinciple}

\begin{algosposs}
\begin{enumerate}
    \item 在Dabbit AI SPSS工具箱中选择"帕累托图"
    \item 上传数据：类别字段选"客诉原因"，数值字段选"损失金额（元）"
    \item 汇总方式选"求和"，累计判定线设为0.80
    \item 勾选"ABC分级"，阈值80/95
    \item 勾选"尾部合并"：占比小于2\%的小类别合并为"其他"
    \item 点击运行，输出排序表与帕累托图
\end{enumerate}

\textbf{SPSS对应操作}：也可先用 \textbf{分析 → 描述统计 → 频率} 排序，再用图形手工绘制帕累托图。
\end{algosposs}

\begin{algoresult}
\begin{itemize}
    \item \item \textbf{排序表}：每个类别的数值、占比、累计占比、ABC分级
    \item \textbf{关键少数}：累计首次达到80\%需要几个类别（如前4个类别占82.9\%）
    \item \textbf{A类类别}：应优先改进（如显示屏异常35.8\%、主板失效22.8\%）
    \item \textbf{集中度}：前20\%类别贡献了多少（如3个类别占75.1\%）——如果贡献分散，说明不能简单套用80/20法则
\end{itemize}
\end{algoresult}

\begin{algonotes}
\textbf{什么时候用？}需要从众多问题原因中找出主要矛盾时。

\textbf{什么时候不用？}各类别贡献均匀分散时（前80\%需要很多类别），说明没有"关键少数"，需全面改进。

\textbf{常见坑}
\begin{itemize}
    \item \textbf{坑1}：只看发生次数不看损失金额——小问题发生次数多但损失小，大问题次数少但损失大，应按金额排序
    \item \textbf{坑2}：把尾部小类别单独画很多小柱子——应合并为"其他"
\end{itemize}

\textbf{和柱状图的区别}：普通柱状图不排序、无累计线；帕累托图专为"抓重点"设计。
\end{algonotes}

%% ========== Xbar-R控制图 ==========
\basicalgo{Xbar-R控制图}{xbar_r_chart}

\begin{algointro}
Xbar-R控制图是\textbf{最常用的计量型控制图}——Xbar图监控子组均值的波动（过程位置是否偏移），R图监控子组内极差（过程波动是否稳定）。两者配合使用，判断生产过程是否处于统计受控状态。
\end{algointro}

\begin{algoscene}
\begin{itemize}
    \item 批量生产过程的质量监控（每2小时抽5件测外径）
    \item 判断工序是否稳定、有没有异常时段
    \item 为日常监控建立控制限基准
    \item 美赛工业类题目中的过程稳定性分析
\end{itemize}
\end{algoscene}

\begin{algoprinciple}
把连续生产按时间分成小子组（如每2小时为一组，每组5件），画两张图：

\begin{itemize}
    \item \textbf{Xbar图（均值图）}：每组算一个均值，打点在图上。如果均值突然偏移，说明过程中心漂移了（如刀具磨损、原料换批）
    \item \textbf{R图（极差图）}：每组内最大值减最小值，反映组内波动。如果R变大，说明过程波动加剧了（如设备松动）
\end{itemize}

控制限不是规格限，而是根据\textbf{过程自身波动}算出来的：中心线CL=均值均值，UCL=$\bar{\bar{X}}+A_2\bar{R}$，LCL=$\bar{\bar{X}}-A_2\bar{R}$。点在限内随机散布=受控；有点超限或连续同侧=失控信号。
\end{algoprinciple}

\begin{algosposs}
\begin{enumerate}
    \item SPSS菜单：\textbf{分析 → 质量控制 → 控制图}
    \item 选择"Xbar-R"图，观测值字段选"外径测量值"，子组标识选"批次号"
    \item 规格上下限分别填入（如USL=50.20，LSL=49.80）
    \item $\sigma$估计方法选"极差法"
    \item 判异规则：单点超限、连续9点同侧、连续6点趋势
    \item 试运行阶段勾选"剔除特殊原因点后重算"
    \item 点击确定
\end{enumerate}
\end{algosposs}

\begin{algoresult}
\begin{itemize}
    \item \textbf{中心线与控制限}：Xbar图CL≈50.0，UCL≈50.06，LCL≈49.94
    \item \textbf{判异信号}：有点超出控制限→该子组存在特殊原因，需回溯该时段
    \item \textbf{受控判定}：试运行阶段剔除特殊原因点后，两图均无判异信号→过程受控，可固化控制限
    \item \textbf{过程能力}：同时输出Cp=1.49, Cpk=1.49（见过程能力分析）
\end{itemize}

\textbf{注意}：控制限≠规格限！控制限是过程"自然波动"的范围，规格限是客户要求的范围。受控不代表满足规格。
\end{algoresult}

\begin{algonotes}
\textbf{什么时候用？}计量型数据，且能按时间合理分组（每组2--9件），子组内不应有特殊原因。

\textbf{什么时候不用？}
\begin{itemize}
    \item 每件产品只有一个测量值、无法分组→用I-MR控制图
    \item 计数型数据（缺陷数/不合格率）→用C图/U图/P图
\end{itemize}

\textbf{常见坑}
\begin{itemize}
    \item \textbf{坑1}：把子组设得太大——子组应短时间内抽样，组内只有普通原因波动
    \item \textbf{坑2}：看到超限点就盲目调整工艺——先调查原因（设备/原料/人员），确认特殊原因后再调
\end{itemize}

\textbf{和I-MR的区别}：Xbar-R需要\textbf{子组}（每组多件），I-MR用于\textbf{单值}（每批只有一个数据）。
\end{algonotes}

%% ========== I-MR控制图 ==========
\basicalgo{I-MR控制图}{imr_chart}

\begin{algointro}
I-MR控制图用于\textbf{每次只有一个测量值}的场景——比如每批药品只测一个含量值。I图（单值图）监控每个测量值，MR图（移动极差图）用相邻两个值的差来估计短期波动。
\end{algointro}

\begin{algoscene}
\begin{itemize}
    \item 每批产品只测一个关键指标（如药品含量、炉温）
    \item 单件小批量生产（无法凑子组）
    \item 连续采样的过程监控
\end{itemize}
\end{algoscene}

\begin{algoprinciple}
既然每组只有一个值，怎么估计波动？用\textbf{相邻两个值的差}（移动极差MR）：

\begin{itemize}
    \item \textbf{I图（Individual）}：把每个测量值直接打点，中心线=均值，控制限$\bar{x}\pm E_2\bar{MR}$
    \item \textbf{MR图（Moving Range）}：$MR_i = |x_i - x_{i-1}|$，相邻两点差值。MR图监控短期波动是否突然变大
\end{itemize}

$E_2$是常数（span=2时约2.66）。就像虽然不知道"组内方差"，但用"今天和昨天的差别"来估计日常波动幅度。
\end{algoprinciple}

\begin{algosposs}
\begin{enumerate}
    \item 在Dabbit AI SPSS工具箱中选择"I-MR控制图"
    \item 测量值列自动识别唯一数值列
    \item 移动极差跨度设为2（相邻两点差）
    \item $\sigma$估计选"移动极差法"
    \item 判异规则选"仅超限判异"
    \item 点击运行
\end{enumerate}
\end{algosposs}

\begin{algoresult}
\begin{itemize}
    \item \textbf{I图}：CL=100，UCL=100.3，LCL=99.68——超出限的点为失控信号
    \item \textbf{MR图}：CL=0.119，UCL=0.389——MR超限说明相邻两次测量间波动骤增
    \item \textbf{信号定位}：如第48批I图和MR图同时超限，说明该批有异常
\end{itemize}
\end{algoresult}

\begin{algonotes}
\textbf{什么时候用？}单件测量、无法分组的连续数据。

\textbf{常见坑}：数据有自相关（相邻值不独立）时，MR估计的$\sigma$会不准，控制限失真。
\end{algonotes}

%% ========== 过程能力分析 ==========
\basicalgo{{\starmark}过程能力分析}{process_capability}

\begin{algointro}
过程能力分析回答一个核心问题：\textbf{这个过程生产的产品，能不能稳定落在规格范围内？}它用Cp、Cpk等指数量化"过程波动"和"规格要求"的差距，是六西格玛管理的核心工具。
\end{algointro}

\begin{algoscene}
\begin{itemize}
    \item 产线是否能满足客户规格要求（如净含量490--510g）
    \item 新设备/新工艺投用前的能力评估
    \item 客户投诉增多时诊断是中心偏移还是波动过大
    \item 美赛工业题中"过程是否合格"的定量回答
\end{itemize}
\end{algoscene}

\begin{algoprinciple}
想象一条传送带不断产出瓶子，有的装多了有的装少了。规格要求490--510g，目标500g。

\begin{itemize}
    \item \textbf{Cp（过程潜力）}：规格宽度 $\div$ 过程波动宽度 = $(USL-LSL)/(6\sigma_w)$。只看"波动够不够小"，不管中心对不对
    \item \textbf{Cpk（过程能力）}：min($CPU, CPL$)，其中 $CPU=(USL-\mu)/(3\sigma_w)$，$CPL=(\mu-LSL)/(3\sigma_w)$。同时考虑波动大小\textbf{和}中心是否对齐
    \item \textbf{Pp/Ppk}：用长期总标准差$\sigma_o$计算，反映实际表现（含组间漂移）
\end{itemize}

\textbf{判断标准}：
\begin{itemize}
    \item $Cpk \geq 1.33$：能力充分，放心生产
    \item $1.0 \leq Cpk < 1.33$：能力尚可，需监控
    \item $Cpk < 1.0$：能力不足，必须改进
\end{itemize}
\end{algoprinciple}

\begin{algosposs}
\begin{enumerate}
    \item 在Dabbit AI SPSS工具箱中选择"过程能力分析"
    \item 上传子组数据，测量值字段选"净含量\_g"，子组字段选"批次"
    \item 填入规格上下限（LSL=490, USL=510）和目标值（500）
    \item 正态性检验选Shapiro-Wilk，$\alpha=0.05$
    \item 非正态时用百分位数法构造Pp/Ppk
    \item 先运行控制图确认过程受控，再读能力指数
    \item 点击运行
\end{enumerate}

\textbf{前提}：过程必须先\textbf{统计受控}（控制图无判异），否则Cp/Cpk没有意义！
\end{algosposs}

\begin{algoresult}
\begin{itemize}
    \item \textbf{正态性}：$p=0.48 \geq 0.05$→数据近似正态，可用Cp/Cpk
    \item \textbf{Cp=1.51, Cpk=1.49}→能力充分（1.33--1.67区间），可稳定满足规格
    \item \textbf{Pp=1.47, Ppk=1.45}→长期能力略低于短期，说明有轻微组间漂移
    \item \textbf{超差率}：约11 ppm（百万分之11），即每百万件约11件超出规格
    \item \textbf{直方图}：叠加规格线和均值线，直观看分布是否在规格内
\end{itemize}
\end{algoresult}

\begin{algonotes}
\textbf{什么时候用？}过程受控后，评估是否满足客户规格要求。

\textbf{什么时候不用？}
\begin{itemize}
    \item 过程尚未受控（控制图有判异点）→先解决特殊原因，再谈能力
    \item 数据严重非正态→需Box-Cox变换或用百分位数法
\end{itemize}

\textbf{常见坑}
\begin{itemize}
    \item \textbf{坑1}：Cp很高但Cpk低→说明波动小但中心偏了，调中心即可
    \item \textbf{坑2}：Cp和Cpk都低→波动太大，必须改进工艺减小变异
    \item \textbf{坑3}：忽略过程受控前提直接算Cpk→失控过程的Cpk不可信
\end{itemize}

\textbf{Cp vs Cpk}：Cp只看波动宽度，Cpk同时看波动和中心。Cpk永远$\leq$ Cp，两者差距大说明中心偏移。
\end{algonotes}

\section{进阶级算法速览}

%% ========== Z-MR控制图 ==========
\advancedalgo{Z-MR控制图（标准值-移动极差图）}{zmr_chart}

\begin{algobrief}
\textbf{一句话}：已知标准均值和标准差时用的单值-移动极差控制图，适合单件小批量生产。

\textbf{美赛场景区}：精密阀芯逐件加工的外径监控，无法按子组凑数据。

\textbf{核心原理}：与I-MR类似，但中心线直接用给定的标准值mu0而非样本均值，$\sigma$也用已知标准值。MR图监控相邻件波动。

\textbf{操作要点}：指定mu和sigma，MR步长2，Nelson规则1（单点超限即判异）。

\textbf{结果关键}：X图超限2点、MR图超限2点→过程存在特殊原因需调查。

\textbf{注意}：Z-MR与I-MR的区别在于中心线用标准值还是样本估计值。
\end{algobrief}

%% ========== 参数分布分析（可靠性分析） ==========
\advancedalgo{参数分布分析（可靠性分析）}{reliability_distribution}

\begin{algobrief}
\textbf{一句话}：判断设备寿命服从什么分布（威布尔/指数/对数正态），估计平均寿命和B10可靠寿命。

\textbf{美赛场景区}：水泵、电容器等设备的预防性更换周期决策；可靠性工程。

\textbf{核心原理}：对威布尔、指数、正态、对数正态四种分布分别做极大似然估计，用AIC选最优。威布尔形状参数$\beta>1$=磨损失效，$\beta\approx1$=随机失效，$\beta<1$=早期失效。B10寿命=90\%产品仍未失效的时间。

\textbf{操作要点}：候选分布auto自动选择，删失事件值设1，Bootstrap重复999次，评估分位数0.50。

\textbf{结果关键}：最优分布威布尔，$\beta=2.59$（磨损期），B10=4721小时，MTTF=10005小时。

\textbf{注意}：数据含右删失（还没坏的设备）时不能直接算均值，必须用生存分析方法。
\end{algobrief}

%% ========== U控制图 ==========
\advancedalgo{U控制图（单位缺陷数图）}{u_chart}

\begin{algobrief}
\textbf{一句话}：监控\textbf{不等样本量}下的单位缺陷率——每批抽检数量不同，但要看每单位缺陷数是否稳定。

\textbf{美赛场景区}：汽车漆面缺陷监控（每天抽检台数不同）。

\textbf{核心原理}：缺陷数服从泊松分布，$u_i=c_i/n_i$为第i批单位缺陷数。控制限随$n_i$变化（$UCL=\bar{u}+3\sqrt{\bar{u}/n_i}$），所以是阶梯状的。

\textbf{操作要点}：缺陷数列、检验单位数列、分组列指定；3$\sigma$控制限；过离散时启用Laney修正。

\textbf{结果关键}：中心线$\bar{u}$=0.66（剔除失控后），7个子组失控，过程未受控。

\textbf{注意}：U图用于等单位数不固定的计数数据；样本量固定时用C图。
\end{algobrief}

%% ========== 属性一致性分析 ==========
\advancedalgo{属性一致性分析}{attribute_agreement}

\begin{algobrief}
\textbf{一句话}：评估多名检验员对同一样本的判定是否一致（如零件合格/返修/报废的目视检验）。

\textbf{美赛场景区}：质量检验系统的可靠性评估，判断检验员之间、检验员与标准之间是否一致。

\textbf{核心原理}：用一致率和Kappa系数衡量。Kappa扣除了"碰巧一致"的部分：$Kappa=(P_o-P_e)/(1-P_e)$。$Kappa>0.75$良好，$0.4$--$0.75$中等，$<0.4$差。

\textbf{操作要点}：指定样本列、评价人列、判定列、标准列；Bootstrap 1000次算置信区间。

\textbf{结果关键}：总体一致率90\%，Kappa=0.74（良好）；评价人丙一致率仅76\%需培训。

\textbf{注意}：Kappa为0说明一致率和随机猜差不多；Kappa=1说明完全一致。
\end{algobrief}

%% ========== C控制图 ==========
\advancedalgo{C控制图（缺陷数图）}{c_chart}

\begin{algobrief}
\textbf{一句话}：监控\textbf{固定检查单元}上的缺陷总数——每块电路板的焊接缺陷数是否稳定。

\textbf{美赛场景区}：SMT贴片车间AOI检测的缺陷数监控。

\textbf{核心原理}：缺陷数服从泊松分布，中心线$\bar{c}$，控制限$\bar{c}\pm3\sqrt{\bar{c}}$。样本量固定时用C图，不固定用U图。

\textbf{操作要点}：指定缺陷数列，3$\sigma$控制限，游程检验判随机性。

\textbf{结果关键}：中心线$\bar{c}=4.35$，UCL=10.6，第118块板缺陷15个超限→失控。

\textbf{注意}：C图和U图的区别——C图要求每个检查单元大小相同。
\end{algobrief}

%% ========== 加速寿命试验ALT ==========
\advancedalgo{加速寿命试验（ALT）}{alt}

\begin{algobrief}
\textbf{一句话}：在高温/高压等加速条件下做寿命试验，外推正常使用条件下的产品寿命。

\textbf{美赛场景区}：电容器、LED等长寿命产品的可靠性评估——常温下要几十年才坏，等不起。

\textbf{核心原理}：用威布尔分布描述寿命，用阿伦尼斯模型描述温度对寿命的加速效应（温度越高寿命越短）。在几个高温水平做试验，拟合模型后外推到常温。

\textbf{操作要点}：分布weibull，加速模型arrhenius，应力单位℃，目标B10寿命。

\textbf{结果关键}：40℃下B10寿命约28560小时；125℃比40℃加速349倍。

\textbf{注意}：外推距离越远越不可靠——正常温度离试验温度越远，置信区间越宽。
\end{algobrief}

%% ========== 量具R&R分析 ==========
\advancedalgo{量具R\&R分析（重复性与再现性）}{gauge_rr}

\begin{algobrief}
\textbf{一句话}：评估测量系统本身准不准——同一人反复测同一零件的波动（重复性）和不同人测同一件的差异（再现性）。

\textbf{美赛场景区}：新购量具投入使用前的验证；六西格玛DMAIC的Measure阶段。

\textbf{核心原理}：总波动=零件间真实差异+测量系统波动（GRR）。用ANOVA分解方差，\%GRR<10\%优秀，10--30\%有条件接受，>30\%不可用。ndc（可区分类别数）$\geq$5才算合格。

\textbf{操作要点}：3个检验员、10个零件、每人每件测5次；完整交叉ANOVA；研究变差倍数6。

\textbf{结果关键}：\%GRR=14.32\%（有条件接受），ndc=9（可区分），零件间占97.95\%。

\textbf{注意}：如果零件间差异太小（抽样范围窄），\%GRR会被放大失真，应扩大抽样范围。
\end{algobrief}

%% ========== 量具线性与偏倚研究 ==========
\advancedalgo{量具线性与偏倚研究}{gauge_linearity}

\begin{algobrief}
\textbf{一句话}：检查量具读数是否整体偏高/偏低（偏倚），以及偏差是否随量程变化（线性）。

\textbf{美赛场景区}：推拉力计、卡尺等量具的校准验证。

\textbf{核心原理}：在量程内选多个标准件，重复测量。每个水平做t检验看平均偏倚是否为0；再对偏倚vs参考值做线性回归，看斜率是否显著。\%线性$<$10\%良好。

\textbf{操作要点}：参考值列、测量值列指定；每水平重复多次；$\alpha=0.05$。

\textbf{结果关键}：偏倚随参考值线性上升（\%线性≈11.2\%），有条件接受，建议校准。

\textbf{注意}：偏倚是常数偏移（截距），线性是偏差随量程变化（斜率）。
\end{algobrief}

%% ========== Xbar-S控制图 ==========
\advancedalgo{Xbar-S控制图（均值-标准差图）}{xbar_s_chart}

\begin{algobrief}
\textbf{一句话}：Xbar-R的升级版——用子组标准差S代替极差R估计波动，子组较大（$n>10$）时更准确。

\textbf{美赛场景区}：饮料灌装净含量监控，子组大小3--7不等。

\textbf{核心原理}：与Xbar-R类似，但用$c_4$常数从子组标准差估计$\sigma$，控制限随子组大小逐组计算。

\textbf{操作要点}：子组容量3--7，合并估计$\sigma$，判异规则=超限+连续7点同侧。

\textbf{结果关键}：Xbar图失控（2个超限+9个游程），S图受控→均值偏移问题。

\textbf{注意}：子组小（$n<10$）用R图即可，子组大用S图更精确。
\end{algobrief}

%% ========== 公差区间 ==========
\advancedalgo{公差区间}{tolerance_interval}

\begin{algobrief}
\textbf{一句话}：估计"未来95\%的产品测量值落在什么范围内"——比均值$\pm$标准差更严格的区间。

\textbf{美赛场景区}：注塑件孔径和翘曲量的正常波动范围估计。

\textbf{核心原理}：置信区间估计的是"均值在哪"，公差区间估计的是"\textbf{未来个体值}覆盖范围"。正态时用$\bar{x}\pm k\cdot s$（k比z大），非正态时用顺序统计量的非参数法。

\textbf{操作要点}：覆盖比例95\%，置信水平95\%，双侧；正态性auto判断。

\textbf{结果关键}：孔径正态法区间[9.94, 10.09]mm；翘曲量非正态区间[0.018, 0.364]mm。

\textbf{注意}：公差区间比预测区间更宽——它要覆盖"未来大量个体"而非"单个预测"。
\end{algobrief}

% 第14章 空间计量
\chapter{空间计量}

\begin{chapteroverview}
\textbf{这个分类是干什么的？}

空间计量经济学研究\textbf{有地理位置信息的数据}。普通回归假设样本之间互相独立，但现实中"近朱者赤"——相邻地区的房价、污染、经济增长往往相互影响。空间计量就是把这种\textbf{空间依赖关系}量化并纳入模型，避免普通回归得出错误结论。

\textbf{什么数据需要空间分析？}

只要你的数据有\textbf{地理位置信息}（经纬度、区县、城市、邮编），就应该考虑空间分析！典型如：各区县的房价/污染/GDP/创新产出——这些指标天然存在"相邻地区互相影响"。如果忽略空间效应，回归系数和显著性都会失真。

\textbf{美赛什么题型常用？}

美赛C题（大数据分析）涉及区域经济、环境监测、城市规划时常用。Moran指数判断"污染是不是连片聚集"，空间回归量化"邻居效应有多大"。注意：\textbf{SPSS原生对空间计量支持有限}，主要通过Dabbit AI SPSS工具箱或Python的pysal库实现。

\textbf{学习路径建议}

先学\textbf{空间权重矩阵构造}（定义谁和谁是邻居）和\textbf{Moran指数}（判断有没有空间聚集），再学\textbf{空间OLS回归}（基准+空间诊断），然后是\textbf{空间滞后模型SAR}（邻居因变量影响本地）。进阶了解SEM、SDM等更复杂的空间模型。
\end{chapteroverview}

\section{基础级算法}

%% ========== 空间权重矩阵构造 ==========
\basicalgo{空间权重矩阵构造}{spatial_weights}

\begin{algointro}
空间权重矩阵W是空间分析的\textbf{地基}——它定义了"谁和谁是邻居"。没有W，后面所有空间模型都无从谈起。它是一个 $n\times n$ 矩阵，$W_{ij}$ 表示区域 $i$ 和区域 $j$ 之间的空间关联强度。
\end{algointro}

\begin{algoscene}
\begin{itemize}
    \item 任何空间分析的\textbf{第一步}——先定义邻居关系
    \item 门店网络、城市网络、区域经济分析的前置步骤
    \item 需要判断"距离越近影响越大"还是"共享边界才算邻居"
\end{itemize}
\end{algoscene}

\begin{algoprinciple}
想象你在画一张城市地图，要定义"相邻"：

\begin{itemize}
    \item \textbf{K近邻法（KNN）}：每个点找最近的K个点当邻居（如K=4）——不管实际距离多远，每个点都有4个邻居
    \item \textbf{反距离法（IDW）}：距离越近权重越大，权重 $=1/d^\alpha$，远处的点也有微弱影响
    \item \textbf{邻接法}：共享边界才算邻居（如两个区县接壤）
\end{itemize}

构造后通常做\textbf{行标准化}：每行除以行和，使权重变成"邻居的加权平均"。这样 $Wy$ 就是"邻居的平均取值"，好解释。

就像定义社交网络里"好友"的标准——是通讯录里前4个？还是距离1公里内？还是互相关注？不同定义影响后续分析。
\end{algoprinciple}

\begin{algosposs}
SPSS原生不支持空间权重矩阵，在\textbf{Dabbit AI SPSS工具箱}中操作：

\begin{enumerate}
    \item 上传含坐标的数据（x坐标列、y坐标列）
    \item 权重类型选"反距离"或"k近邻"
    \item 距离度量选"欧氏距离"（平面坐标）或"haversine"（经纬度）
    \item 近邻数K设为4，距离衰减幂设为1
    \item 标准化方式选"行标准化"
    \item 点击运行，输出连接密度、邻居数、孤立单元数等结构统计
\end{enumerate}
\end{algosposs}

\begin{algoresult}
\begin{itemize}
    \item \textbf{连接密度}：非零连接占比，1.0表示所有点都有邻居
    \item \textbf{孤立单元数}：=0最好——有孤立单元说明该点没有邻居，空间滞后无法计算
    \item \textbf{邻居数分布}：每个点平均几个邻居，是否均匀
    \item \textbf{W·y相关}：y与邻居均值的相关系数，正相关说明高值点倾向聚集
\end{itemize}
\end{algoresult}

\begin{algonotes}
\textbf{什么时候用？}任何空间分析的前置步骤。

\textbf{常见坑}
\begin{itemize}
    \item \textbf{坑1}：不做行标准化——权重直接是距离倒数，不同区域权重和差异巨大
    \item \textbf{坑2}：K设太大或太小——K=4是常用起点，可敏感性分析
\end{itemize}

\textbf{注意}：空间权重矩阵的选择会影响结论，论文中应说明选了哪种W并做敏感性检验。
\end{algonotes}

%% ========== 空间OLS回归 ==========
\basicalgo{空间OLS回归}{spatial_ols}

\begin{algointro}
先做普通最小二乘回归（OLS），然后\textbf{检验残差有没有空间自相关}。如果有，说明普通回归不可靠，需要改用空间模型。它是空间分析的"基准诊断"步骤。
\end{algointro}

\begin{algoscene}
\begin{itemize}
    \item 任何空间回归分析的\textbf{第一步}：先做OLS基准，再诊断
    \item 判断"普通回归的结论可不可信"
    \item 为选择SAR还是SEM提供依据（LM检验）
\end{itemize}
\end{algoscene}

\begin{algoprinciple}
普通回归假设 $y = X\beta + \varepsilon$，且 $\varepsilon$ 之间互相独立。但现实中相邻地区的残差往往相关（比如一个政策同时影响相邻几个区县）。

步骤：
\begin{enumerate}
    \item 先跑普通OLS，得到系数和残差
    \item 对残差做\textbf{Moran's I检验}：如果残差Moran's I显著，说明残差有空间聚集→OLS假设被破坏
    \item 做\textbf{LM检验}：LM-error显著→选SEM；LM-lag显著→选SAR；都显著→看稳健形式
\end{enumerate}

就像体检：先做基础检查（OLS），再做进阶筛查（空间诊断），决定要不要进一步治疗（空间模型）。
\end{algoprinciple}

\begin{algosposs}
在\textbf{Dabbit AI SPSS工具箱}中操作：

\begin{enumerate}
    \item 上传区域数据（因变量、自变量、坐标）
    \item 指定因变量（如地均GDP）和自变量（投资强度、研发强度等）
    \item 权重方法选KNN（k=4），距离度量haversine，行标准化
    \item 稳健标准误选"否"
    \item 运行后先看OLS系数表，再看空间依赖诊断表
\end{enumerate}
\end{algosposs}

\begin{algoresult}
\begin{itemize}
    \item \textbf{OLS系数}：如投资强度系数2.09（$p<0.001$），研发强度1.51（$p<0.001$）
    \item \textbf{$R^2=0.63$}：模型解释力
    \item \textbf{残差Moran's I=0.31（$p<0.001$）}：残差显著空间自相关→OLS不可靠
    \item \textbf{LM检验}：LM-Error显著、稳健LM-Error显著→建议用空间误差模型SEM
\end{itemize}
\end{algoresult}

\begin{algonotes}
\textbf{什么时候用？}空间分析的起点，必须先做。

\textbf{什么时候不用？}如果残差Moran's I不显著，说明没有空间效应，OLS即可，不需要后续空间模型。

\textbf{常见坑}：跳过空间诊断直接做空间模型——应先证明"需要"空间模型。
\end{algonotes}

%% ========== Moran指数（空间自相关） ==========
\basicalgo{{\starmark}Moran指数（空间自相关）}{moran_i}

\begin{algointro}
Moran's I是衡量\textbf{空间聚集程度}的核心指标。它回答一个问题：高值区域是不是和高值区域相邻？低值和低值相邻？还是完全随机散布？比如PM2.5是不是连片污染？房价是不是核心区高、外围低？
\end{algointro}

\begin{algoscene}
\begin{itemize}
    \item \textbf{环境科学}：PM2.5污染是否连片聚集？
    \item \textbf{城市经济}：高GDP区县是否空间聚集？
    \item \textbf{房地产}：高价小区是否扎堆？
    \item 任何"地理上相近的观测值是否相似"的问题
\end{itemize}
\end{algoscene}

\begin{algoprinciple}
Moran's I的公式本质是\textbf{变量值与其空间滞后的相关系数}：

\[ I = \frac{n}{\sum_i\sum_j w_{ij}} \cdot \frac{\sum_i\sum_j w_{ij}(x_i-\bar{x})(x_j-\bar{x})}{\sum_i(x_i-\bar{x})^2} \]

\begin{itemize}
    \item $I>0$且显著：\textbf{正空间自相关}——高值邻高值、低值邻低值（连片聚集）
    \item $I<0$且显著：\textbf{负空间自相关}——高值邻低值（高低交错）
    \item $I\approx 0$：空间随机分布，没有聚集
\end{itemize}

\textbf{Moran散点图}：横轴是变量值（标准化），纵轴是邻居均值。一、三象限=高-高/低-低聚集；二、四象限=高-低/低-高离群。

就像看地图：如果红色（高值）都连成片、蓝色（低值）都连成片，Moran's I就大；如果红蓝交错，I就接近0。
\end{algoprinciple}

\begin{algosposs}
SPSS原生不支持Moran指数，在\textbf{Dabbit AI SPSS工具箱}中操作：

\begin{enumerate}
    \item 上传含坐标的数据
    \item 变量列选要检验的变量（如PM2\_5年均浓度）
    \item 权重构建方式选"k近邻"，k=4
    \item 行标准化选"是"
    \item 置换检验次数设为999，显著性水平0.05
    \item 固定随机种子42（可复现）
    \item 点击运行
\end{enumerate}
\end{algosposs}

\begin{algoresult}
\begin{itemize}
    \item \textbf{Moran's I=0.84}（期望$E[I]=-0.007$）：远大于期望
    \item \textbf{z值=15.2，置换p=0.002 $<0.05$}：显著正空间自相关
    \item \textbf{象限占比}：高-高38\%、低-低54\%——大部分是同类聚集
    \item \textbf{结论}：PM2.5浓度呈连片聚集分布，不是随机散布
\end{itemize}

\textbf{注意}：全局Moran's I只告诉你"整体有没有聚集"，不能定位\textbf{哪里}是热点（那需要局部Moran's I/LISA）。
\end{algoresult}

\begin{algonotes}
\textbf{什么时候用？}数据有地理位置，想判断"是不是空间聚集"时。

\textbf{什么时候不用？}
\begin{itemize}
    \item 数据没有空间坐标——无法算Moran's I
    \item 已经在回归残差上做过Moran检验——那是诊断残差，这里是检验因变量本身
\end{itemize}

\textbf{常见坑}
\begin{itemize}
    \item \textbf{坑1}：把Moran's I当相关系数解读——I的取值范围不是[-1,1]，而是近似$[-1/(n-1), 1]$
    \item \textbf{坑2}：全局I显著就说"所有区域都聚集"——全局I只说整体趋势，个别区域可能是离群点
\end{itemize}

\textbf{和空间回归的关系}：Moran's I是"要不要做空间模型"的诊断；空间回归是"做了之后怎么解释效应"。
\end{algonotes}

%% ========== 空间滞后模型（SLM/SAR） ==========
\basicalgo{{\starmark}空间滞后模型（SLM/SAR）}{sar_model}

\begin{algointro}
空间滞后模型（SAR/SLM）在普通回归中加入\textbf{邻居因变量} $Wy$——不仅看本地自变量，还看邻居的因变量对本地的影响。比如：一个小区涨价会带动周边小区跟涨。
\end{algointro}

\begin{algoscene}
\begin{itemize}
    \item \textbf{房地产估价}：小区挂牌价受周边小区价格影响
    \item \textbf{区域经济}：一个城市的GDP受相邻城市增长带动
    \item \textbf{创新扩散}：一个地区的创新产出受邻近地区知识溢出影响
\end{itemize}
\end{algoscene}

\begin{algoprinciple}
普通回归：$y = X\beta + \varepsilon$

空间滞后模型：$y = \rho W y + X\beta + \varepsilon$

\begin{itemize}
    \item $Wy$是"邻居因变量的加权平均"
    \item $\rho$（空间滞后系数）：邻居每涨1单位，本地涨$\rho$单位
    \item $\rho>0$显著：存在\textbf{正向空间溢出}（邻居好我也好）
\end{itemize}

关键：因为$Wy$和$y$互相影响（你影响我、我影响你），形成了\textbf{空间反馈循环}，所以不能直接用OLS估计，要用极大似然（ML）。

效应分解：
\begin{itemize}
    \item \textbf{直接效应}：本地自变量对本地因变量的影响
    \item \textbf{间接效应（溢出）}：本地自变量变化通过空间网络传导到邻居
    \item \textbf{总效应}=直接+间接
\end{itemize}
\end{algoprinciple}

\begin{algosposs}
在\textbf{Dabbit AI SPSS工具箱}中操作：

\begin{enumerate}
    \item 上传含坐标和变量的数据
    \item 因变量选"avg\_price"，自变量选"metro\_km, school\_score, building\_age"
    \item 坐标选lon/lat，权重类型knn近邻，k=4，行标准化
    \item 估计方法选极大似然ML
    \item 点击运行，输出系数表和效应分解表
\end{enumerate}
\end{algosposs}

\begin{algoresult}
\begin{itemize}
    \item \textbf{$\rho=0.563$（$p<0.001$）}：显著正向空间溢出——邻居均价每涨1000元，本地涨约563元
    \item \textbf{LR检验}：$\rho=0$被拒绝→SAR比OLS显著更优
    \item \textbf{效应分解}：如距地铁距离的直接效应-0.53、间接效应-0.58、总效应-1.11——间接占比52\%！
    \item \textbf{间接占比}：超过50\%说明空间溢出效应非常强，不能忽略
\end{itemize}
\end{algoresult}

\begin{algonotes}
\textbf{什么时候用？}Moran's I显著，且LM-lag检验显著时。

\textbf{什么时候不用？}
\begin{itemize}
    \item LM-error显著而LM-lag不显著→用SEM
    \item Moran's I不显著→用普通OLS即可
\end{itemize}

\textbf{常见坑}
\begin{itemize}
    \item \textbf{坑1}：把SAR系数直接当OLS系数解读——SAR中自变量系数只是直接效应，总效应要算间接效应
    \item \textbf{坑2}：忽略$Wy$的内生性——$Wy$和$y$互为因果，OLS估计有偏，必须用ML或IV
\end{itemize}

\textbf{和SEM的区别}：SAR是\textbf{因变量}之间的空间溢出（邻居y影响本地y）；SEM是\textbf{误差项}之间的空间相关（未观测因素在相邻地区间传导）。
\end{algonotes}

\section{进阶级算法速览}

%% ========== 空间误差模型（SEM） ==========
\advancedalgo{空间误差模型（SEM）}{spatial_error_model}

\begin{algobrief}
\textbf{一句话}：空间依赖体现在\textbf{误差项}中——未观测的扰动在相邻地区间传导，而非因变量本身有溢出。

\textbf{美赛场景区}：房价分析中，相邻片区共享未观测的地段口碑、学区预期等因素，导致残差空间相关。

\textbf{核心原理}：$y=X\beta+u$，$u=\lambda Wu+\varepsilon$。$\lambda$是空间误差系数。LM-error显著而LM-lag不显著时选SEM。

\textbf{操作要点}：KNN权重k=4，ML估计，行标准化。

\textbf{结果关键}：$\lambda=0.537$（$p<0.001$）显著，说明误差存在空间扩散；系数与OLS接近但标准误更准。

\textbf{注意}：SEM中自变量没有空间溢出项，其效应就是直接效应，不需要分解。
\end{algobrief}

%% ========== 空间杜宾模型（SDM） ==========
\advancedalgo{空间杜宾模型（SDM）}{sdm_model}

\begin{algobrief}
\textbf{一句话}：最全面的空间模型——同时包含因变量空间滞后$Wy$和自变量空间滞后$WX$，既看邻居y也看邻居x。

\textbf{美赛场景区}：区域创新研究——本地研发投入影响本地创新，邻居研发投入也通过知识溢出影响本地。

\textbf{核心原理}：$y=\rho Wy+X\beta+WX\theta+\varepsilon$。可以检验能否退化为SAR（$\theta=0$）或SEM。SDM是最一般化的设定。

\textbf{操作要点}：ML估计，wx列选"全部解释变量"，行标准化。

\textbf{结果关键}：$\rho=0.45$显著；研发投入直接效应1.12、间接效应0.81——溢出效应很大。LR检验SDM显著优于SAR和SEM。

\textbf{注意}：SDM最全面但参数最多，小样本时估计不稳定；大样本首选SDM。
\end{algobrief}

%% ========== 空间滞后误差模型（SAC） ==========
\advancedalgo{空间滞后误差模型（SAC）}{sac_model}

\begin{algobrief}
\textbf{一句话}：同时包含空间滞后$\rho Wy$和空间误差$\lambda Wu$，SAR和SEM的合体。

\textbf{美赛场景区}：住宅价格中既有邻居房价的实质溢出（$\rho$），又有未观测区位因素的误差相关（$\lambda$）。

\textbf{核心原理}：$y=\rho Wy+X\beta+u$，$u=\lambda Wu+\varepsilon$。两个空间参数都显著说明两种机制都存在。

\textbf{操作要点}：ML估计，KNN权重，行标准化。

\textbf{结果关键}：$\rho=0.479$和$\lambda=0.401$都显著→两种空间机制同时存在。

\textbf{注意}：SAC参数多，对样本量要求高；小样本时$\rho$和$\lambda$可能难以区分。
\end{algobrief}

%% ========== 空间面板模型 ==========
\advancedalgo{空间面板模型}{spatial_panel}

\begin{algobrief}
\textbf{一句话}：空间模型+面板数据（多地区×多年份），同时控制时间不变的地区固定效应和空间溢出。

\textbf{美赛场景区}：15个城市10年的经济增长面板，研究数字经济对GDP的拉动及跨城溢出。

\textbf{核心原理}：在SAR/SEM/SDM基础上加入个体固定效应（吸收各城市不随时间变的禀赋差异），用组内去均值后做ML估计。

\textbf{操作要点}：effects选fixed，model type选sar，KNN权重k=4。

\textbf{结果关键}：$\rho=0.424$显著（城市间经济联动）；数字人才直接效应15.5、间接溢出10.6——近40\%效应通过邻国外溢。

\textbf{注意}：面板数据要有足够时间维度（$T>5$），固定效应吸收了不随时间变的因素。
\end{algobrief}

%% ========== 空间杜宾误差模型（SDEM） ==========
\advancedalgo{空间杜宾误差模型（SDEM）}{sdem_model}

\begin{algobrief}
\textbf{一句话}：包含自变量空间滞后$WX$和误差空间相关$\lambda$，但\textbf{不包含}因变量空间滞后$Wy$。

\textbf{美赛场景区}：创新产出研究——邻居研发投入有溢出（WX），但误差项也有空间关联（$\lambda$）。

\textbf{核心原理}：$y=X\beta+WX\theta+u$，$u=\lambda Wu+\varepsilon$。与SDM的区别：SDEM没有$Wy$，所以效应分解简单——$\beta$是直接效应，$\theta$就是间接效应，不需要矩阵求逆。

\textbf{操作要点}：KNN权重k=4，ML估计，行标准化。

\textbf{结果关键}：$\lambda=0.502$显著（误差空间相关）；研发投入直接效应1.10、间接效应0.78。

\textbf{注意}：SDEM比SDM简单（不需要$(I-\rho W)^{-1}$），如果LM-lag不显著但LM-ERR显著，SDEM是好选择。
\end{algobrief}

%% ========== 自变量空间滞后模型（SLX） ==========
\advancedalgo{自变量空间滞后模型（SLX）}{slx_model}

\begin{algobrief}
\textbf{一句话}：最简单的空间模型——只加邻居自变量$WX$，不加$Wy$也不加误差项，用普通OLS就能估计。

\textbf{美赛场景区}：便利店研究——邻居门店的营销投入是否影响本店营业额（客流共享/分流）。

\textbf{核心原理}：$y=X\beta+WX\theta+\varepsilon$。$\beta$=本地自变量效应（直接），$\theta$=邻居自变量效应（间接溢出）。因为没有$Wy$内生性，直接OLS即可。

\textbf{操作要点}：knn k=4，lag mode=all，稳健标准误，行标准化。

\textbf{结果关键}：营销投入直接效应1.56、间接效应1.08——邻居营销也能带动本店（品牌集聚效应）。

\textbf{注意}：SLX是入门级空间模型，如果残差Moran's I仍显著，说明还需要加误差项（SDEM）。
\end{algobrief}

% 第15章 试验设计
\chapter{试验设计}

\begin{chapteroverview}
\textbf{这个分类是干什么的？}

试验设计（DOE）解决一个问题：\textbf{怎么用最少的试验次数，找到最优工艺参数组合}。比如化工生产中，温度、时间、催化剂用量都影响收率——怎么安排试验才能既省钱又高效地找到最佳配方？试验设计用统计方法科学安排试验点，系统地分析各因素的主效应、交互作用，最后预测最优条件。

\textbf{美赛什么题型常用？}

美赛实验/化工/材料类题目中非常实用——只要题目涉及"调参数找最优"（反应温度、时间、配比），就用试验设计。正交设计用少量试验筛选因素，响应面分析精确找最优点。SPSS通过"分析 → 试验设计"菜单可完成部分操作。

\textbf{学习路径建议}

先学\textbf{全析因设计}（理解因素与交互作用的基础），再学\textbf{正交实验设计}（用少量试验做多因素筛选，美赛高频），然后学\textbf{响应面分析}（精确寻优）。CCD和BBD是响应面的两种常用设计方案。部分析因和Plackett-Burman用于早期因子筛选。
\end{chapteroverview}

\section{基础级算法}

%% ========== 全析因设计 ==========
\basicalgo{全析因设计}{full_factorial}

\begin{algointro}
全析因设计就是\textbf{所有因素的所有水平组合都做一遍试验}。比如2个温度水平$\times$3个时间水平$=$6个组合，每个组合重复几次。它能最完整地估计主效应和交互作用，但因素多了试验次数爆炸。
\end{algointro}

\begin{algoscene}
\begin{itemize}
    \item 因素少（$\leq$4个）、水平少（$\leq$3个）的精细试验
    \item 需要研究因素间交互作用时
    \item 美赛中因素少但要搞清楚每个因素影响时
\end{itemize}
\end{algoscene}

\begin{algoprinciple}
想象做蛋糕：烤温有高/中/低3档，时间有60/90/120/150/180分钟5档。全析因就是$3\times5=15$种组合\textbf{全部做一遍}。

然后用方差分析（ANOVA）分解响应的变异：
\begin{itemize}
    \item \textbf{主效应}：温度单独的影响有多大？时间单独的影响有多大？
    \item \textbf{交互作用}：温度高时延长时间更好，还是温度低时延长时间更好？如果折线图交叉明显，说明有交互
\end{itemize}

关键洞察：当交互显著时，不能单独看"温度越高越好"——因为温度的效果取决于时间水平。
\end{algoprinciple}

\begin{algosposs}
\begin{enumerate}
    \item SPSS菜单：\textbf{分析 → 试验设计 → 析因}
    \item 响应变量选"抗折强度"，因素选"烧成温度"和"保温时间"
    \item 模型选全模型（主效应+所有交互）
    \item 多重比较选Tukey HSD
    \item ANOVA类型选III型
    \item 点击确定
\end{enumerate}
\end{algosposs}

\begin{algoresult}
\begin{itemize}
    \item \textbf{方差分析表}：烧成温度$F=111$（$p<0.001$）、保温时间$F=26$（$p<0.001$）、交互$F=6.2$（$p<0.001$）
    \item \textbf{偏$\eta^2$}：温度0.62（解释62\%变异）、时间0.44、交互0.27
    \item \textbf{交互作用图}：折线明显交叉→交互显著，不能单独下结论
    \item \textbf{最优组合}：1300℃$\times$180min，均值38.48
\end{itemize}
\end{algoresult}

\begin{algonotes}
\textbf{什么时候用？}因素少（$\leq$4个）、水平少，且需要完整研究交互作用。

\textbf{什么时候不用？}因素多（$\geq$5个）——全析因试验次数呈指数增长（$2^5=32$次还能接受，$2^7=128$次就太多了），改用部分析因或正交设计。

\textbf{常见坑}
\begin{itemize}
    \item \textbf{坑1}：忽略交互作用只看主效应——交互显著时主效应会误导
    \item \textbf{坑2}：不做重复试验——没有重复就无法估计误差，无法做ANOVA
\end{itemize}
\end{algonotes}

%% ========== 正交实验设计（田口设计） ==========
\basicalgo{{\starmark}正交实验设计（田口设计）}{orthogonal_design}

\begin{algointro}
正交实验设计用\textbf{正交表}安排试验——因素很多（如6个参数各5个水平），但不用做$5^6=15625$次，只做25次就能系统估计各因素的主效应。它是美赛中\textbf{最常用的实验设计方法}。
\end{algointro}

\begin{algoscene}
\begin{itemize}
    \item 注塑工艺优化：6个参数（料温、模温、保压、注射速度、冷却时间、背压）各5水平
    \item 化工配方筛选：多个因素各多个水平，试验成本高
    \item 美赛中"多因素参数优化"类题目
\end{itemize}
\end{algoscene}

\begin{algoprinciple}
正交表的巧妙之处在于\textbf{均衡分散、整齐可比}：

\begin{itemize}
    \item 任意两个因素的水平组合\textbf{均匀出现}——每个因素的每个水平都和其他因素的每个水平搭配过相同次数
    \item 这样每个因素的效应可以\textbf{独立估计}，不受其他因素干扰
    \item 虽然只做了25次（而非$5^6=15625$次），但信息效率很高
\end{itemize}

\textbf{分析步骤}：
\begin{enumerate}
    \item \textbf{极差分析}：算每个因素各水平的响应均值，极差R越大=该因素影响越大
    \item \textbf{方差分析}：检验因素是否显著，算贡献率
    \item \textbf{选优}：显著因素取最优水平组合，非显著因素按成本选
\end{enumerate}

就像田忌赛马：不需要每匹马都和每匹马赛一遍，用巧妙的对阵安排就能比较出马的优劣。
\end{algoprinciple}

\begin{algosposs}
在\textbf{Dabbit AI SPSS工具箱}中操作（SPSS也可通过"分析 → 试验设计 → 正交设计"实现）：

\begin{enumerate}
    \item 上传试验数据，因素列自动识别，响应列自动识别
    \item 信噪比类型选"望大特性"（拉伸强度越大越好）
    \item 优水平选择基准选"信噪比优先"（不仅均值高，还要波动小）
    \item 显著性水平0.05
    \item 点击运行，输出极差分析表和方差分析表
\end{enumerate}
\end{algosposs}

\begin{algoresult}
\begin{itemize}
    \item \textbf{极差排序}：料筒温度R=6.46$>$保压压力3.58$>$模具温度3.75$>$冷却时间3.48$>$注射速度1.61$>$背压0.34
    \item \textbf{方差分析}：料筒温度贡献率48.6\%（最关键），背压不显著（$p=0.31$）
    \item \textbf{最优组合}：料温260℃、模温70℃、保压40MPa、注射速度70\%、冷却35s
    \item \textbf{预测响应}：约63.27MPa（需验证试验确认）
\end{itemize}
\end{algoresult}

\begin{algonotes}
\textbf{什么时候用？}因素多（3--7个）、水平数相同、试验成本高、只需估计主效应。

\textbf{什么时候不用？}
\begin{itemize}
    \item 需要研究交互作用——正交表主效应可能和交互混叠
    \item 因素水平不等——需用混合水平正交表
\end{itemize}

\textbf{常见坑}
\begin{itemize}
    \item \textbf{坑1}：把预测最优组合直接当最终结论——必须追加\textbf{验证试验}确认
    \item \textbf{坑2}：忽略信噪比，只看均值——田口方法的核心是"稳健性"，不仅要均值高，还要波动小
\end{itemize}

\textbf{和全析因的区别}：正交设计用少量试验换高效筛选，但不能估计高阶交互；全析因信息完整但试验多。
\end{algonotes}

%% ========== 响应面分析 ==========
\basicalgo{{\starmark}响应面分析}{rsm}

\begin{algointro}
响应面分析（RSM）在正交/筛选找到关键因素后，用\textbf{二次多项式}拟合"因素→响应"的曲面，然后在曲面上找最优点。它不仅知道"哪个因素重要"，还能精确算出"最佳参数是多少"。
\end{algointro}

\begin{algoscene}
\begin{itemize}
    \item 化工浸提工艺优化：温度、时间、液料比三因素找最大得率
    \item 食品配方优化：蛋糕比容最大化
    \item 美赛中已经筛选出关键因素，需要精确寻优时
\end{itemize}
\end{algoscene}

\begin{algoprinciple}
先用正交/PB设计筛选出3--5个关键因素，然后用CCD或BBD设计安排试验点，拟合二次模型：

\[ y = \beta_0 + \sum\beta_i x_i + \sum\beta_{ii} x_i^2 + \sum\beta_{ij} x_i x_j + \varepsilon \]

\begin{itemize}
    \item \textbf{线性项}：因素增大，响应线性增大/减小
    \item \textbf{平方项}：存在最优点（开口向下=极大值，开口向上=极小值）
    \item \textbf{交互项}：两个因素联合作用
\end{itemize}

然后对二次模型求导找驻点，用二阶导数矩阵（Hessian）的特征值判断是极大点、极小点还是鞍点。等高线图和3D响应面图直观展示最优点位置。
\end{algoprinciple}

\begin{algosposs}
\begin{enumerate}
    \item SPSS菜单：\textbf{分析 → 试验设计 → 响应面}
    \item 响应列选"茶多酚得率"，因素选"浸提温度、浸提时间、液料比"
    \item 设计类型选CCD（中心复合设计）
    \item 优化目标选maximize（最大化得率）
    \item 模型简化开启（自动剔除不显著项）
    \item 点击确定
\end{enumerate}
\end{algosposs}

\begin{algoresult}
\begin{itemize}
    \item \textbf{模型$R^2=0.98$}，失拟检验$p=0.51$（不显著=模型充分）
    \item \textbf{最优条件}：温度83.8℃、时间51.7min、液料比25.8
    \item \textbf{预测响应}：得率83.54\%（95\%置信区间83.2--83.9）
    \item \textbf{曲面类型}：所有特征值为负→极大值点
\end{itemize}
\end{algoresult}

\begin{algonotes}
\textbf{什么时候用？}已经筛选出3--5个关键因素，需要精确找最优点。

\textbf{什么时候不用？}
\begin{itemize}
    \item 还不知道哪些因素重要——先做PB/正交筛选
    \item 因素太多（$>$5个）——二次模型参数太多
\end{itemize}

\textbf{常见坑}
\begin{itemize}
    \item \textbf{坑1}：不做失拟检验直接用模型预测——失拟显著说明二次模型不够
    \item \textbf{坑2}：驻点在试验域外——此时不能报驻点值，应在边界上找最优
\end{itemize}

\textbf{和正交设计的区别}：正交设计\textbf{筛选因素}（哪个重要），响应面\textbf{精确寻优}（最佳值是多少）。两者是接力关系。
\end{algonotes}

%% ========== 中心复合设计（CCD） ==========
\basicalgo{中心复合设计（CCD）}{ccd}

\begin{algointro}
CCD是响应面最常用的试验设计方案——在2水平全因子点（立方点）基础上，加轴向点（星号点）和中心点，共三类点。它能拟合完整的二次模型。
\end{algointro}

\begin{algoscene}
\begin{itemize}
    \item 植物提取工艺优化（5个因素）
    \item 任何需要响应面建模且因素数$2\leq k\leq5$时
\end{itemize}
\end{algoscene}

\begin{algoprinciple}
CCD由三部分组成：
\begin{itemize}
    \item \textbf{立方点}：$2^k$个全因子点（各因素高低水平组合）
    \item \textbf{轴向点}：$2k$个点（一个因素取$\pm\alpha$，其他取0）——用来估计平方项
    \item \textbf{中心点}：重复多次（估计纯误差和失拟）
\end{itemize}

$\alpha$选"可旋转"模式时，各方向预测方差相等，是最常用选择。
\end{algoprinciple}

\begin{algosposs}
在Dabbit AI工具箱中选择"中心复合设计"，指定各因素的高低水平，$\alpha$模式选rotatable，响应列选yield\_pct。
\end{algosposs}

\begin{algoresult}
\begin{itemize}
    \item 模型$R^2=0.99$，失拟检验不显著
    \item 最优编码条件换算为实际工艺参数（如温度66.3℃、时间36.2min等）
    \item 预测响应87.59（95\%预测区间86.3--88.9）
\end{itemize}
\end{algoresult}

\begin{algonotes}
\textbf{CCD vs BBD}：CCD有轴向点（试验范围比因子范围大），适合想探索边界外的情况；BBD没有角点，试验次数更少但不探索极端组合。
\end{algonotes}

\section{进阶级算法速览}

%% ========== 均匀设计 ==========
\advancedalgo{均匀设计}{uniform_design}

\begin{algobrief}
\textbf{一句话}：中国数学家方开泰、王元提出的试验设计方法——让试验点在因素空间中\textbf{均匀散布}，用极少试验覆盖整个多维空间。

\textbf{美赛场景区}：化工工艺4因素150水平，试验代价极高，无法做全因子。

\textbf{核心原理}：与正交表"整齐可比"不同，均匀设计追求"均匀散布"——每个因素的每个水平只出现一次，点在空间中均匀分布。配合二次回归模型预测最优组合。

\textbf{操作要点}：因素数和试验次数匹配均匀表，二次模型=主效应+平方项，优化方向最大化响应。

\textbf{结果关键}：CD2（中心化L2偏差）越小越均匀；二次模型$R^2=0.996$预测最优参数。

\textbf{注意}：均匀设计点散布均匀但没有重复，无法估计纯误差，需配合建模使用。
\end{algobrief}

%% ========== 拉丁方设计 ==========
\advancedalgo{拉丁方设计}{latin_square}

\begin{algobrief}
\textbf{一句话}：用$k\times k$方阵同时控制两个无关变量（如工位、时段），只关心处理因素（如配方）的差异。

\textbf{美赛场景区}：涂装产线5种配方对比，但5个工位和5个时段都有系统差异，需要同时扣除两类干扰。

\textbf{核心原理}：每个处理在每行每列恰好出现一次——就像数独，行=工位、列=时段、数字=配方。这样工位和时段的系统偏差被区组吸收，误差更纯净。

\textbf{操作要点}：指定行区组、列区组、处理因素、因变量；Tukey HSD多重比较。

\textbf{结果关键}：处理$F=113$（$p<0.001$），F3配方最好（均值5.98），F4最差（4.61）。

\textbf{注意}：拉丁方要求行、列、处理水平数相等；不能估计行×列交互。
\end{algobrief}

%% ========== Box-Behnken设计（BBD） ==========
\advancedalgo{Box-Behnken设计（BBD）}{bbd}

\begin{algobrief}
\textbf{一句话}：响应面的另一种设计方案——没有角点，只取边中点和中心点，试验次数比CCD少。

\textbf{美赛场景区}：蛋糕配方优化（加水量、打发时间、焙烤温度3因素）。

\textbf{核心原理}：BBD的点在边的中点（两因素取端点、第三因素取中心），不做极端角点组合——适合"极端组合不安全/不现实"的场景（如温度太高会烤糊）。

\textbf{操作要点}：3因素BBD只需15次试验（加中心点重复），拟合二次响应面。

\textbf{结果关键}：比容最优条件（加水43.7\%、打发9.8min、温度172℃），预测比容3.81mL/g。

\textbf{注意}：BBD不探索角点——如果最优点在角上附近，BBD会漏掉。CCD和BBD根据实际选择。
\end{algobrief}

%% ========== 部分析因设计 ==========
\advancedalgo{部分析因设计}{fractional_factorial}

\begin{algobrief}
\textbf{一句话}：全析因的"1/2"或"1/4"——$2^{k-p}$设计，用部分试验筛选主效应，牺牲高阶交互换试验次数。

\textbf{美赛场景区}：7个两水平因子全因子需128次太贵，用$2^{7-3}=16$次筛选关键因子。

\textbf{核心原理}：分辨率IV设计——主效应不与二阶交互混叠，但二阶交互之间互相混叠。用Lenth伪标准误判断哪些因子显著。

\textbf{操作要点}：模型项选main，显著性方法选Lenth，响应列选收率。

\textbf{结果关键}：7个因子中3个显著（温度、催化剂量、养晶时间），R$^2$=0.84。

\textbf{注意}：部分析因是筛选设计——目的是"找到关键少数"，不是精确寻优。显著因子后续做响应面。
\end{algobrief}

%% ========== Plackett-Burman设计 ==========
\advancedalgo{Plackett-Burman设计}{pb_design}

\begin{algobrief}
\textbf{一句话}：最经济的因子筛选设计——$N$次试验最多筛$N-1$个两水平因子，专门用来从很多候选因子中快速找出关键的。

\textbf{美赛场景区}：8个工艺因子（温度、时间、催化剂量等）想快速找出4个关键的。

\textbf{核心原理}：PB设计是分辨率III——主效应可能与二阶交互混叠，但用极少试验就能估计所有主效应。效应$=$高水平均值$-$低水平均值。

\textbf{操作要点}：因子列选8个两水平因子，响应列选收率，显著性规则auto。

\textbf{结果关键}：8个因子中4个显著（温度+3.9、时间-3.6、催化剂量+3.0、保温时间+1.9）。

\textbf{注意}：PB是纯筛选设计——找到显著因子后，下一步要做全因子或响应面，不能直接当最终结论。
\end{algobrief}

% 第16章 Meta分析
\chapter{Meta分析}

\begin{chapteroverview}
\textbf{这个分类是干什么的？}

Meta分析（元分析）是一种"站在巨人肩膀上"的统计方法。当同一个问题已经有很多人各自做过研究、结论却大小不一时，Meta分析不重新做实验，而是把这些\textbf{已发表的研究结果}当成数据，用一套加权的数学方法把它们"合并"成一个更可靠的总体结论。它回答的核心问题是：\textbf{把所有同类研究放在一起看，效应到底有多大？各研究结论一致吗？结论会不会被个别研究带偏？}

\textbf{美赛什么题型常用？}

Meta分析在美赛中不是高频题，但在两类场景下非常加分：
\begin{itemize}
    \item \textbf{文献综述与证据合成}：E题（环境科学）、F题（政策）中，题目往往要求你"综合已有研究/报告的估计值"，比如把各地不同研究报告的某政策减排效果、某疾病发生率合并成一个区域总体估计。
    \item \textbf{多来源数据合并}：当你手里有多份独立统计（如各城市的相关系数、各试点的发生率），需要汇总出一个带置信区间的总体数字时，Meta分析就是最规范的做法。
\end{itemize}

\textbf{学习路径建议}

先吃透\textbf{平均值Meta分析}和\textbf{二分类数据Meta分析}两个明星算法——它们覆盖了"连续结局"和"事件结局"两种最常见的数据形态，异质性判断、固定/随机效应模型选择、森林图这些核心概念都在这两个算法里讲透。再了解\textbf{相关系数Meta分析}。进阶部分的\textbf{一般倒方差法}是"万能底座"（只要别人给的是效应值加标准误就能合并），\textbf{单个率Meta分析}、\textbf{连续性数据Meta分析（SMD）}、\textbf{OR/HR值Meta分析}是它的特化版本，\textbf{P值合并}则用于"只有p值、没有效应量"的尴尬情形。
\end{chapteroverview}

\section{基础级算法}

%% ========== 平均值Meta分析 ==========
\basicalgo{{\starmark}平均值Meta分析}{meta_mean}

\begin{algointro}
平均值Meta分析，就是把多项研究里"两组均值之差"加权合并成一个总的平均差值。比如有150项研究都测了"运动康复组 vs 常规组的6分钟步行距离"，每个研究算出一个差值，Meta分析把它们按精度加权平均，告诉你综合来看干预到底提升了多少米。
\end{algointro}

\begin{algoscene}
\begin{itemize}
    \item 多项研究都报告了\textbf{两组的均值、标准差、样本量}，且结局\textbf{量纲一致}（都是米、mmHg、分），要合并组间差异
    \item 美赛E/F题中，汇总各地区/各试点报告的政策前后均值差异，估计总体效应
    \item 需要给出"合并效应量 + 置信区间 + 是否显著"的循证结论时
\end{itemize}
\end{algoscene}

\begin{algoprinciple}
想象你要给一场有很多评委打分的比赛算总分。有的评委很专业（样本量大、打分很稳），有的评委只是随手一看（样本量小、波动大）。合理的做法不是简单平均，而是\textbf{让专业评委权重更高}。这就是Meta分析的核心思想。

\begin{itemize}
    \item \textbf{效应量MD（均数差）}：每个研究贡献一个"干预组均值$-$对照组均值"。量纲一致时直接用MD，单位和原始结局一样（米、分）。
    \item \textbf{逆方差加权}：权重 $w_i = 1/v_i$，其中$v_i$是该研究效应量的方差。样本越大、标准差越小，方差越小，权重越大。
    \item \textbf{异质性}：各研究结果如果只是围绕同一个真值随机波动，叫"同质"；如果研究之间真的不一样，叫"异质"。用三个指标判断：
    \begin{itemize}
        \item Cochran $Q$检验：$Q$的$p<0.10$提示有异质性；
        \item $I^2$：总变异中有多少百分比来自研究间差异，$<25\%$低、$50\%$中等、$>75\%$高；
        \item $\tau^2$：研究间方差的绝对大小。
    \end{itemize}
    \item \textbf{固定效应 vs 随机效应}：异质小（$p\geq0.10$且$I^2<50\%$）用固定效应（假设大家都在测同一个真值）；异质大改用随机效应（承认研究间有真差异，额外加上$\tau^2$，区间更宽更稳）。
\end{itemize}
\end{algoprinciple}

\begin{algosposs}
\begin{enumerate}
    \item 在Dabbit AI SPSS工具箱中选择"平均值Meta分析"（SPSS原生无此模块，需借助工具箱或R扩展）
    \item 数据按"宽表两臂"布局：每个研究一行，录入两组各自的\textbf{样本量$n$、均值、标准差}
    \item 设置参数：显著性水准$\alpha=0.05$；效应量选\textbf{MD（均数差）}，仅结局量纲一致时用
    \item 合并模型选\textbf{auto}：由$Q$检验$p$值与$I^2$自动决定固定/随机效应
    \item $\tau^2$估计选\textbf{DL（DerSimonian-Laird矩法）}；结局方向设为\textbf{high\_is\_better}（数值越大越优）
    \item 勾选输出森林图、漏斗图、留一法敏感性分析、Egger发表偏倚检验
    \item 点击运行，查看合并结果与异质性诊断
\end{enumerate}

\textbf{默认参数参考}：$\alpha=0.05$，模型auto自动选择，$\tau^2$用DL法，效应量MD，结局越大越优。
\end{algosposs}

\begin{algoresult}
以某运动康复案例为例（150项研究，结局为6分钟步行距离，单位米）：

\begin{itemize}
    \item \textbf{异质性表}：$Q=241.48$（$df=149$，$p<0.001$），$I^2=38.3\%$，$\tau^2=62.83$。因$Q$检验$p<0.10$，自动选用\textbf{随机效应模型}。
    \item \textbf{合并效应（核心）}：随机效应合并$MD=31.55$，$95\%CI=(29.39,\ 33.70)$，$Z=28.65$，$p<0.001$。\textbf{置信区间不含0}，说明干预显著提升了步行距离。
    \item \textbf{预测区间}：$15.73\sim47.36$，表示再做一项同类研究，其真实效应大概率落在此范围，比置信区间宽，体现了研究间差异。
    \item \textbf{森林图}：每个研究是一个方块（大小=权重）加横线（=置信区间），底部菱形是合并效应。菱形不跨过0线即显著。
    \item \textbf{敏感性分析}：逐篇剔除后合并值围绕31.55波动，方向不反转，说明结果稳健。
    \item \textbf{发表偏倚}：Egger回归截距$p<0.10$提示可能有小样本偏倚，结论需谨慎。
\end{itemize}

\textbf{怎么判断好坏}：合并$p<0.05$且置信区间不跨0为"显著"；$I^2$越高说明研究越不一致，合并值只能解释为"平均趋势"，不能当成唯一真值。
\end{algoresult}

\begin{algonotes}
\textbf{什么时候用？}多项研究都给了两组均值与标准差、且结局单位一致，要合并出一个总体均值差时。

\textbf{什么时候不用？}
\begin{itemize}
    \item 各研究结局\textbf{量纲/量表不统一}时（如一个用米、一个用评分），不能用MD，要用标准化均数差SMD（见进阶"连续性数据Meta分析"）。
    \item 手里只有别人报告的效应值和标准误、没有原始两组均值时，用"一般倒方差Meta分析"。
\end{itemize}

\textbf{常见坑}
\begin{itemize}
    \item \textbf{坑1}：异质性很大还硬用固定效应模型，导致置信区间过窄、过度自信。异质时务必上随机效应。
    \item \textbf{坑2}：忽略发表偏倚。只搜到阳性结果的研究会让合并效应被夸大，一定要看漏斗图/Egger检验。
    \item \textbf{坑3}：把Meta分析当成简单平均。不加权的平均会让小样本研究和大样本研究"一人一票"，严重失真。
\end{itemize}

\textbf{和相似算法的区别}
\begin{itemize}
    \item vs \textbf{连续性数据Meta分析（SMD）}：本算法用原始单位的MD，要求各研究结局量纲一致；SMD把结果标准化成无量纲数，适合量表不同的研究。
    \item vs \textbf{独立样本t检验}：t检验是\textbf{一次实验内}两组比均值；Meta分析是\textbf{跨多个研究}把多个均值差合并。
\end{itemize}
\end{algonotes}

%% ========== 相关系数Meta分析 ==========
\basicalgo{相关系数Meta分析}{meta_corr}

\begin{algointro}
相关系数Meta分析，就是把多项研究各自报告的"两个变量之间的相关系数$r$"合并成一个更可靠的总体相关估计。比如150项研究都测了"组织支持感 vs 工作投入"的相关，有的报0.2、有的报0.5，Meta分析告诉你综合起来到底有多相关。
\end{algointro}

\begin{algoscene}
\begin{itemize}
    \item 多项研究分别报告了\textbf{同一对变量的相关系数$r$和样本量$n$}，结论强弱不一
    \item 美赛中需要综合多项实证研究，判断两个变量是否稳定相关、相关多强
    \item 用于文献综述中"某关系是否稳健"的定量总结
\end{itemize}
\end{algoscene}

\begin{algoprinciple}
相关系数$r$有个毛病：它的分布在接近$\pm1$时很歪，不能直接拿来加权平均。所以要先做一个"扶正"变换。

\begin{itemize}
    \item \textbf{Fisher z变换}：把$r$转成近似正态的 $z = \frac{1}{2}\ln\frac{1+r}{1-r}$。在$z$尺度上做加权合并、算置信区间，最后再反变换回$r$方便解释。
    \item \textbf{权重}：每个研究的方差约为 $1/(n-3)$，样本量$n$越大权重越高。
    \item \textbf{异质性与模型}：同样用$Q$、$I^2$判断。案例中$Q=298.13$（$p<0.001$），$I^2=50.0\%$，研究间差异明显，故用随机效应。
    \item \textbf{效应强弱解读}：$|r|$约0.1弱、0.3中等、0.5以上强。
\end{itemize}
\end{algoprinciple}

\begin{algosposs}
\begin{enumerate}
    \item 在工具箱中选择"相关系数Meta分析"
    \item 数据每研究一行，录入\textbf{样本量$n$}和\textbf{相关系数$r$}两列
    \item 输入效应量选\textbf{r}；合并模型选\textbf{随机效应}；置信水平95\%
    \item $\tau^2$估计选DerSimonian-Laird；最小纳入研究数设为2
    \item 勾选\textbf{Egger发表偏倚检验}与\textbf{留一法敏感性分析}
    \item Knapp-Hartung小样本校正默认关闭（研究数很多时可不开）
    \item 运行并查看合并$r$、置信区间与异质性
\end{enumerate}

\textbf{默认参数参考}：随机效应模型，95\%置信，DL法估计$\tau^2$，Egger与留一法开启。
\end{algosposs}

\begin{algoresult}
以"组织支持感 vs 工作投入"案例为例（150项研究，合计样本11828）：

\begin{itemize}
    \item \textbf{合并结果（核心）}：随机效应合并 $r=0.322$，$95\%CI=(0.297,\ 0.346)$，$Z=23.62$，$p<0.001$。属\textbf{中等强度正相关}，显著但不等于因果。
    \item \textbf{异质性}：$Q=298.13$（$p<0.001$），$I^2=50.0\%$，$\tau^2=0.0132$，研究间差异明显超出抽样误差。
    \item \textbf{发表偏倚}：Egger截距检验不显著，未发现明显小研究效应。
    \item \textbf{敏感性}：剔除任意一篇后合并$r$在$0.319\sim0.325$间波动，方向与显著性不变，结论稳健。
\end{itemize}

\textbf{怎么看}：置信区间不含0即相关显著；区间越窄估计越准。务必记住——\textbf{相关不等于因果}。
\end{algoresult}

\begin{algonotes}
\textbf{什么时候用？}多篇文献都报告了同一对变量的相关系数及其样本量时。

\textbf{什么时候不用？}研究报告的是回归系数、OR等其他效应量时，应改用对应的Meta方法，不能硬塞进来。

\textbf{常见坑}
\begin{itemize}
    \item \textbf{坑1}：直接对$r$做加权平均，跳过Fisher z变换，导致两端的研究权重失真。
    \item \textbf{坑2}：把显著的合并$r$解读成"一个变量导致另一个"，相关Meta分析只能说明关联稳定。
    \item \textbf{坑3}：纳入了样本量$n\leq3$或$r$越界的脏数据，导致方差算错。
\end{itemize}

\textbf{和相似算法的区别}
\begin{itemize}
    \item vs \textbf{第2章相关分析}：那是在\textbf{一份数据集}里算两个变量的相关；本算法是\textbf{跨多项研究}把多个相关系数合并。
    \item vs \textbf{一般倒方差Meta分析}：本算法专门处理$r$（自动做Fisher z变换）；倒方差法更通用，只要给效应值和标准误即可。
\end{itemize}
\end{algonotes}

%% ========== 二分类数据Meta分析 ==========
\basicalgo{{\starmark}二分类数据Meta分析}{meta_binary}

\begin{algointro}
二分类数据Meta分析，合并的是"事件/没事件"这种结局。比如每个研究都报了"试验组多少人发生了静脉血栓、对照组多少人发生"，Meta分析把它们合并成一个总的风险比（OR或RR），回答干预到底把风险降了多少。
\end{algointro}

\begin{algoscene}
\begin{itemize}
    \item 多项研究都是\textbf{两臂对照}，结局是\textbf{二分类事件}（发生/未发生、达标/未达标）
    \item 美赛中合并各试点的政策实施率、发生率、合格率等事件型结果
    \item 需要估计"干预 vs 对照"的\textbf{风险倍数}并判断是否显著
\end{itemize}
\end{algoscene}

\begin{algoprinciple}
二分类结局的效应量不能直接减人数，要用"比值"。

\begin{itemize}
    \item \textbf{OR（比值比）}：试验组事件 odds 除以对照组事件 odds。$OR<1$表示试验组事件更少（好事）；$OR=1$表示两组无差别。
    \item \textbf{对数尺度合并}：OR是偏态的，要先取$\ln(OR)$再合并，算完置信区间再取指数还原。
    \item \textbf{异质性与模型}：$Q$检验$p<0.10$或$I^2>50\%$时用随机效应（$\tau^2$用DL估计），否则用固定效应，两个模型结果都会报告。
    \item \textbf{亚组分析}：可按地区/人群分组，用$Q_b$检验组间差异。
    \item \textbf{零事件校正}：某组0事件时OR算不出来，用0.5连续性校正。
\end{itemize}
\end{algoprinciple}

\begin{algosposs}
\begin{enumerate}
    \item 在工具箱中选择"二分类数据Meta分析"
    \item 数据每研究一行，录入两组各自的\textbf{事件数}与\textbf{总例数}；如有分组列（地区）一并录入
    \item 效应量选\textbf{OR}；置信水平0.95；Q检验显著性水准0.10
    \item 合并模型选\textbf{按异质性自动选择}；随机效应$\tau^2$用DL法
    \item 零事件校正选\textbf{0.5连续性校正}；发表偏倚用Egger回归
    \item 亚组分析选"自动识别分组列并分析"；勾选留一法敏感性分析
    \item 运行并查看固定/随机两模型结果、亚组表与森林图、漏斗图
\end{enumerate}

\textbf{默认参数参考}：效应量OR，异质性阈值$I^2=50\%$，零事件0.5校正，Egger检验。
\end{algosposs}

\begin{algoresult}
以"新型口服抗凝药预防术后VTE"案例为例（150项研究，多地区医疗中心）：

\begin{itemize}
    \item \textbf{合并效应（核心）}：随机效应 $OR=0.614$，$95\%CI=(0.575,\ 0.656)$，$p<0.001$。意思是\textbf{试验组事件风险约为对照组的0.61倍}，即风险下降约39\%，显著有效。固定效应$OR=0.631$作对照。
    \item \textbf{异质性}：$Q=272.96$（$p<0.001$），$I^2=45.4\%$，$\tau^2=0.0749$。按自动规则（$p<0.10$）选随机效应。
    \item \textbf{预测区间}：$0.358\sim1.055$，比置信区间宽，体现研究间差异。
    \item \textbf{亚组分析}：$Q_b=66.10$（$p<0.001$），说明华北/华东/华南/西南各地区效应确有差别，合并成一个总效应会掩盖地区差异。
    \item \textbf{森林图}：无效应线在$OR=1$处，合并菱形不跨过1线即显著；在1线左侧表示试验组更优。
    \item \textbf{敏感性}：剔除任意一篇后OR仍在$0.61\sim0.62$，结论稳健。
\end{itemize}

\textbf{怎么看}：OR的95\%CI\textbf{不含1}即显著；$OR<1$试验组事件少（如干预好），$OR>1$试验组事件多。
\end{algoresult}

\begin{algonotes}
\textbf{什么时候用？}两臂对照、结局为"事件发生与否"时，这是医学/政策评价最常见的Meta场景。

\textbf{什么时候不用？}结局是连续数值（如血压、评分）时，用平均值Meta分析；只给了OR和置信区间的现成数据，用"OR/HR值Meta分析"。

\textbf{常见坑}
\begin{itemize}
    \item \textbf{坑1}：把OR当RR读。OR是"比值比"，当事件率高时OR会比RR夸大，解读时要说清是比值比。
    \item \textbf{坑2}：忽略零事件研究。某组0事件会让OR无穷大，必须做0.5连续性校正或剔除。
    \item \textbf{坑3}：异质性明显还不做亚组分析，一个总效应掩盖了地区/人群差异。
\end{itemize}

\textbf{和相似算法的区别}
\begin{itemize}
    \item vs \textbf{单个率Meta分析}：本算法是\textbf{两组对照}（算OR/RR）；单个率Meta分析只合并\textbf{一组自身的发生率}，没有对照。
    \item vs \textbf{OR/HR值Meta分析}：本算法从\textbf{事件数/总例数}现场算OR；后者直接吃别人已算好的OR和置信区间。
\end{itemize}
\end{algonotes}

\section{进阶级算法速览}

%% ========== 一般倒方差Meta分析 ==========
\advancedalgo{一般倒方差Meta分析}{meta_invvar}

\begin{algobrief}
\textbf{一句话}：最通用的Meta合并底座——只要每个研究给了"效应值"和"标准误"两列，不管它是均差、回归系数还是对数OR，都能按 $1/SE^2$ 加权合并。

\textbf{美赛场景区}：各研究报告的效应形式五花八门（有的给均差、有的给回归系数），但都附带标准误，需要统一合并出总体净效应。

\textbf{核心原理}：以标准误平方的倒数 $w_i=1/SE_i^2$ 为权重，先算固定效应下的Cochran $Q$，结合$I^2$、$\tau^2$决定固定/随机效应（DL法估计$\tau^2$），再加权合并并做$Z$检验。若输入是对数尺度，取指数还原成OR/HR/RR。

\textbf{操作要点}：工具箱选"一般倒方差Meta分析"；指定效应值列effect、标准误列se、研究标签列study；默认随机效应、双侧检验、95\%置信、异质性水准0.10。

\textbf{结果关键}：合并效应值及其95\%CI、$p$值（CI不含0即显著）；案例中合并值4.20（CI 3.88$\sim$4.52，$p<0.001$），$I^2=67.1\%$异质性较高。配合去一法与Egger检验查稳健性与小样本偏倚。

\textbf{注意}：这是其他Meta算法的"母版"——平均值、二分类、SMD都是它在特定数据布局下的特化。当你拿不到原始两组数据、只有别人报告的效应估计时，就用它。
\end{algobrief}

%% ========== 单个率Meta分析 ==========
\advancedalgo{单个率Meta分析}{meta_single_rate}

\begin{algobrief}
\textbf{一句话}：只有一组数据（没有对照组）时，把多家单位各自报告的"发生率/阳性率"合并成一个总体率。

\textbf{美赛场景区}：评估区域医院感染率、某病总体患病率、某政策整体覆盖率——各单位只报"事件数/总数"，想得到区域平均水平。

\textbf{核心原理}：率本身在接近0或100\%时不正态，先做Freeman-Tukey双反正弦（或logit）变换，再加权合并，最后回变换到百分比刻度。随机效应先估$\tau^2$再加权，并用$Q$、$I^2$、Egger、留一法评估异质性与偏倚。

\textbf{操作要点}：工具箱选"单个率Meta分析"；每研究录事件数与样本量；变换选freeman\_tukey，模型random，$\tau^2$用DL；留一法与Egger自动。

\textbf{结果关键}：案例（150家医院）合并感染率3.7\%（95\%CI 3.4\%$\sim$4.0\%），预测区间1.1\%$\sim$7.9\%；$I^2=90\%$异质性极高，说明各医院差异很大，合并值只是平均水平。

\textbf{注意}：异质性极高时，"一个总率"意义有限，更该做的是找出率为何在医院间差异这么大（按规模/地区分层）。
\end{algobrief}

%% ========== 连续性数据Meta分析（SMD） ==========
\advancedalgo{连续性数据Meta分析}{meta_smd}

\begin{algobrief}
\textbf{一句话}：当各研究的连续结局\textbf{量表或单位不统一}时，不用原始均差，而用标准化均数差（SMD，Hedges g）把结果放到同一把尺子上再合并。

\textbf{美赛场景区}：各研究用不同工具测量同一构念（一个用mmHg、一个用评分），要合并它们的干预效应大小。

\textbf{核心原理}：$SMD = $两组均差$/$合并标准差，并做小样本校正得Hedges g，这样量纲消失、可跨研究比较。$|g|$约0.2小、0.5中、0.8大。固定与随机效应（DL）各算一遍，用$Q$、$I^2$、$\tau^2$诊异质，Egger查小样本偏倚。

\textbf{操作要点}：工具箱选"连续性数据Meta分析"；效应量类型选SMD-Hedges g；随机效应-DL为主并报告固定效应；输出森林图+漏斗图；异质性阈值50\%。

\textbf{结果关键}：案例（低钠盐对收缩压）随机效应 $SMD=-0.458$（CI $-0.498\sim-0.418$，$p<0.001$），属中等效应，干预组血压更低；$I^2=67.4\%$异质性偏高。

\textbf{注意}：本算法与基础"平均值Meta分析"的唯一区别就是效应量——结局单位一致用MD（可解释为多少米/mmHg），单位/量表不同用SMD（只说效应大小，不说实际单位）。
\end{algobrief}

%% ========== P值合并 ==========
\advancedalgo{P值合并}{meta_pvalue}

\begin{algobrief}
\textbf{一句话}：当你手头只有一堆研究的\textbf{p值}（连效应量、原始数据都没有），把这些p值合并成一个"整体证据是否显著"的结论。

\textbf{美赛场景区}：文献只报告了各检验的p值和方向，想据此判断"综合起来到底有没有效应"。这是信息最匮乏时的无奈之选。

\textbf{核心原理}：常用三套方法——Fisher卡方合并（把$-2\ln p$求和，服从卡方）、Stouffer z合并（可加权，还能利用方向）、Tippett最小p合并。配合KS诊断p值分布、留一法查是否被单一研究主导。

\textbf{操作要点}：工具箱选"P值合并"；指定p值列、方向列、权重列（可空）；方法选all（同时跑Fisher/Stouffer加权/Tippett）；显著性水准0.05。

\textbf{结果关键}：案例（150项研究）Fisher合并$p=3.3\times10^{-52}<0.05$，整体显著偏离无效假设；Stouffer两尾$p=8.8\times10^{-89}$支持正向为主。剔除任一研究结论方向不变。

\textbf{注意}：P值合并只能告诉你"整体是否有效应"，\textbf{既不估计效应有多大，也不量化异质性}，更不能当结论性的效应量来写。它是最后手段，能用效应量合并就别用它。
\end{algobrief}

%% ========== OR/HR值Meta分析 ==========
\advancedalgo{OR/HR值Meta分析}{meta_or_hr}

\begin{algobrief}
\textbf{一句话}：当别人已经把每项研究的OR（比值比）或HR（风险比）连同95\%置信区间算好给你时，直接拿这些"成品效应"做加权合并。

\textbf{美赛场景区}：观察性流行病学证据合成，如"睡眠不足 vs 2型糖尿病风险"——各研究报告的是OR及CI，要合并平均关联强度。

\textbf{核心原理}：把OR/HR取对数，由CI反推标准误，按逆方差$1/SE^2$加权；用$Q$、$I^2$、$\tau^2$选固定/随机效应（DL）；输出合并OR、预测区间，Egger查发表偏倚，留一法查稳健性。注意观察性研究只能谈关联、不能谈因果。

\textbf{操作要点}：工具箱选"OR/HR值Meta分析"；录效应值or、区间上下限ci\_low/ci\_high、样本量n；模型auto、95\%CI、$\alpha=0.05$、异质性水准0.10；开Egger、留一法、预测区间。

\textbf{结果关键}：案例（睡眠不足与糖尿病）合并$OR=1.24$（CI 1.19$\sim$1.28，$p<0.001$），睡眠不足者糖尿病风险高约24\%；$I^2=90.6\%$异质性很高，用随机效应；预测区间0.81$\sim$1.89；Egger $p=0.115$无明显偏倚。

\textbf{注意}：合并OR/HR的解读务必标注"观察性证据、关联不等于因果"。$I^2$很高时，结论只是各研究的平均关联，要解释差异来源需做亚组或Meta回归。
\end{algobrief}

% 第17章 文本分析
\chapter{文本分析}

\begin{chapteroverview}
\textbf{这个分类是干什么的？}

现实世界里超过80\%的数据是"文字"——顾客评论、政策文件、新闻报道、问卷开放题、社交媒体留言。文本分析就是教你\textbf{把这些非结构化的文字变成能统计、能画图、能建模的数字}。它不要求你读完全部几千条评论，而是自动告诉你：大家最常说什么词？情绪是正面还是负面？这些文字能归纳成哪几个话题？

\textbf{美赛什么题型常用？}

文本分析在美赛中是\textbf{C题（大数据/数据挖掘）}和\textbf{E/F题（环境/政策）}的常客：
\begin{itemize}
    \item C题经常给你一份CSV，里面一列是评论文本、政策文本或新闻标题，要求你做情感倾向、主题挖掘、热点识别。
    \item E/F题要分析政策文件、公众意见时，词频、LDA主题模型能帮你快速归纳各方关注点。
    \item 拿到一堆文字却"不知道里面在讲什么"时，从本章入手。
\end{itemize}

\textbf{学习路径建议}

按"从简单到复杂、从描述到建模"的顺序学：先\textbf{词频统计}（数一数词）和\textbf{词云图}（把词频画出来）——这是一切文本分析的地基；再学\textbf{文本情感分析}（判断正负情绪，是有监督学习）和\textbf{LDA主题分析}（自动归纳话题，是无监督学习）这两个明星算法。进阶部分的\textbf{TF-IDF}是"给词加权"的核心武器，\textbf{文本聚类}把相似文本自动分组，\textbf{共词网络}看词与词的关系，\textbf{新词发现}则用来挖词典里没有的新说法。
\end{chapteroverview}

\section{基础级算法}

%% ========== 词频统计 ==========
\basicalgo{词频统计}{word_freq}

\begin{algointro}
词频统计就是把一段段文字"切碎成词"，然后数一数每个词出现了多少次。它是所有文本分析的第一步——你连大家在提什么词都不知道，后面什么情感、主题都是空中楼阁。
\end{algointro}

\begin{algoscene}
\begin{itemize}
    \item 美赛C题给了评论文本、政策文本后，\textbf{第一步}：先看高频词是什么，摸清语料在讲什么
    \item 按类别分组对比：产品使用类客户在说什么词？物流售后类客户在说什么词？
    \item 论文中"文本数据概况"部分的基础统计
\end{itemize}
\end{algoscene}

\begin{algoprinciple}
想象你有一篮子碎纸，每张写着一句话。词频统计做两件事：

\begin{itemize}
    \item \textbf{分词}：把"这个手机拍照很清晰"切成["这个","手机","拍照","很","清晰"]。中文没有空格，要用词典做最大匹配切分；英文按空格切、统一小写。
    \item \textbf{过滤停用词}："的、了、很、the、a"这种到处都是、却没有信息量的词叫停用词，先删掉，否则榜单会被它们霸占。
    \item \textbf{计数与相对频率}：统计每个词出现次数，并算\textbf{相对频率}（该词频次$/$总词元数），以及\textbf{累计覆盖率}（前$N$个高频词占了总词次的百分之多少）。
\end{itemize}

生活类比：累计覆盖率就像盘点仓库——如果前10个词就覆盖了17\%的内容，说明话题比较集中；如果要几百个词才覆盖一半，说明内容很分散。

\textbf{分组比较注意}：两组文本条数不同，不能直接比绝对次数，要比\textbf{相对频率}（组内占比），否则人多的那组什么词都"更多"。
\end{algoprinciple}

\begin{algosposs}
\begin{enumerate}
    \item 在Dabbit AI SPSS工具箱中选择"词频统计"（SPSS原生无中文分词，需工具箱或Python扩展）
    \item 上传含文本列的数据，指定\textbf{文本字段}；如有分组列（如咨询类别），选入\textbf{分组字段}
    \item 停用词范围选\textbf{chinese\_and\_english}（中英停用词都过滤）
    \item 分词模式选auto；设置最短词长（通常$\geq2$），是否去数字
    \item 设定top N（看前多少个高频词）、最小文档频次min\_doc\_freq
    \item 运行，查看全局词频表与分组词频表
\end{enumerate}

\textbf{默认参数参考}：停用词中英双语，自动分词，最短词长2，分组间只比较相对频率。
\end{algosposs}

\begin{algoresult}
以"客服意见反馈"案例为例（150条文本，2933个有效词元）：

\begin{itemize}
    \item \textbf{全局高频词}：出现最多的是"清晰"（79次，占2.69\%）；前10个高频词累计覆盖17.56\%——说明话题不算特别集中。
    \item \textbf{分组结果}：
    \begin{itemize}
        \item "产品使用反馈"组（75篇）前三高频：清晰、稳定、优秀；
        \item "物流售后反馈"组（75篇）前三高频：服务、到位、处理。
    \end{itemize}
    \item \textbf{怎么读}：产品组在夸产品属性（清晰、稳定），售后组在谈服务环节（服务、处理），组间用词侧重一目了然。
\end{itemize}

\textbf{注意}：组间差异只反映\textbf{用词侧重}，不直接代表态度好坏，要下结论得回到原文语境。
\end{algoresult}

\begin{algonotes}
\textbf{什么时候用？}拿到任何文本语料后的第一步，先摸高频词。

\textbf{什么时候不用？}词频只告诉你"谁说得多"，不告诉你"说得对不对、情绪好不好"，也不区分一词多义——要判断倾向用情感分析，要归纳话题用LDA。

\textbf{常见坑}
\begin{itemize}
    \item \textbf{坑1}：不过滤停用词，结果榜单第一名是"的、了、是"，毫无信息量。
    \item \textbf{坑2}：直接比两组绝对频次。条数不同时必须比相对频率。
    \item \textbf{坑3}：只看词频不看原文。"问题"出现100次到底是在夸还是在骂，必须抽查原文。
\end{itemize}

\textbf{和相似算法的区别}
\begin{itemize}
    \item vs \textbf{词云图}：词频统计给精确数字表；词云图只是它的可视化版本。
    \item vs \textbf{TF-IDF}：词频只数出现次数；TF-IDF会惩罚到处都出现的"通用词"，更能找出某组独有的代表词。
\end{itemize}
\end{algonotes}

%% ========== 词云图 ==========
\basicalgo{词云图（词频可视化）}{wordcloud}

\begin{algointro}
词云图就是把词频统计的结果画成一张图：词越大，出现得越多。它的作用不是给精确数字，而是\textbf{一眼看出大家最关心什么}，特别适合放在论文里做直观展示。
\end{algointro}

\begin{algoscene}
\begin{itemize}
    \item 论文或汇报中需要\textbf{直观展示}文本热点时，词云图视觉冲击力强
    \item 美赛C题要求快速呈现"用户最常讨论的主题"
    \item 配合词频统计表使用：图负责好看，表负责准确
\end{itemize}
\end{algoscene}

\begin{algoprinciple}
词云图背后没有新算法，就是\textbf{词频统计的图形化}：

\begin{itemize}
    \item 先做分词、去停用词、计数（和词频统计完全一样的预处理）；
    \item 按词频排序，词频越高字号越大、位置越居中；
    \item 可自定义停用词（把"无意义但高频"的词手动加进黑名单）、设置最大词数与画布尺寸。
\end{itemize}

生活类比：词云就像海报——字号大小就是声音大小，谁嗓门大谁就被放在中间。它帮你\textbf{快速抓重点}，但别指望从图上读出精确数字。
\end{algoprinciple}

\begin{algosposs}
\begin{enumerate}
    \item 工具箱选"词云图（词频可视化）"
    \item 指定\textbf{文本字段}；若已有词频表，可直接选count字段与text字段
    \item 打开\textbf{remove stopwords}过滤停用词；自定义停用词可手动追加
    \item 设置最短词长min\_word\_len（默认2）、最大词数max\_words、画布宽高
    \item 运行后导出图片；建议同时导出高频词条形图以备核对
\end{enumerate}

\textbf{默认参数参考}：过滤停用词，最短词长2，自定义停用词按需添加。
\end{algosposs}

\begin{algoresult}
以"电商购物体验开放题"案例为例（150条文本，657个有效词次）：

\begin{itemize}
    \item 词云显示讨论主要围绕\textbf{包装、问题、客服}，三者合计约占词次7.6\%；
    \item 其中"包装"最突出，出现17次、占2.6\%；
    \item \textbf{条形图核对}：前若干高频词画成条形并标注数值，与词云结论一致、互相印证。
\end{itemize}

\textbf{怎么看}：图上最大的几个词就是核心话题；讨论集中在少数几个大词上说明用户关注点明确。
\end{algoresult}

\begin{algonotes}
\textbf{什么时候用？}需要可视化展示文本热点、放在论文配图里时。

\textbf{什么时候不用？}需要精确数字排名时，别从词云上估，回去看词频统计表。

\textbf{常见坑}
\begin{itemize}
    \item \textbf{坑1}：词云里全是"手机、东西"这种泛词，因为没加自定义停用词。
    \item \textbf{坑2}：用词云当正式结论，忽视了它只是示意图、没有置信度概念。
    \item \textbf{坑3}：把不同情感的词混在一张图里，"好"和"差"都很大，却没法区分情绪。
\end{itemize}

\textbf{和相似算法的区别}：词云是词频统计的\textbf{可视化产物}，二者预处理完全相同，区别只在一个给表、一个给图。
\end{algonotes}

%% ========== 文本情感分析 ==========
\basicalgo{{\starmark}文本情感分析}{sentiment}

\begin{algointro}
文本情感分析就是让机器给每段文字判"情绪"——这条评论是夸（正面）还是骂（负面）？它把定性的文字变成了可统计的情绪比例，是美赛处理评论/舆情文本的\textbf{核心武器}。
\end{algointro}

\begin{algoscene}
\begin{itemize}
    \item 美赛C题给了成百上千条顾客评论、社媒留言，要\textbf{量化口碑}：正面占多少、负面占多少
    \item 分析不同门店、不同时期的情绪倾向变化
    \item 找出"哪些词最能区分好评和差评"，定位满意/不满的具体原因
\end{itemize}
\end{algoscene}

\begin{algoprinciple}
情感分析本质是一个\textbf{有监督二分类问题}：先有人工标注好的一批文本（正面/负面），机器学着自己判断。

\begin{itemize}
    \item \textbf{文本向量化（TF-IDF）}：把每段文字变成一串数字。中文不用分词，用\textbf{字符级n-gram}（char\_wb）——相邻两三个字切成特征，比如"太差"、"好吃"，避免了中文分词的麻烦。
    \item \textbf{分类器}：用\textbf{逻辑回归}（logistic），简单可解释。类别不均衡时开\textbf{class weight=balanced}自动加权。
    \item \textbf{为什么可解释}：逻辑回归的系数告诉你哪些词是"正面信号"（系数大正）、哪些是"负面信号"（系数大负）。
    \item \textbf{怎么评估}：把数据切成训练集/测试集（test size=0.2），训练集上做5折交叉验证估稳定性，测试集上看准确率、宏平均F1、AUC。
\end{itemize}

生活类比：就像老师批改作文——先拿标准答案（标注数据）教学生（训练模型），再让学生做一套没见过的卷子（测试集），看他和标准答案对得上多少。
\end{algoprinciple}

\begin{algosposs}
\begin{enumerate}
    \item 工具箱选"文本情感分析"
    \item 指定\textbf{文本列}与\textbf{情感标签列}（必须有预先标注的正负标签，这是有监督学习）
    \item 模式选auto；分析器选\textbf{char\_wb}（字符级，中文免分词）
    \item 分类器选\textbf{logistic}；类别权重balanced；max\_features=5000（保留信息量最大的5000个词特征）
    \item 测试集比例test size=0.2，交叉验证5折，随机种子42（保证可复现）
    \item 运行，查看效果指标与高区分度关键词
\end{enumerate}

\textbf{默认参数参考}：char\_wb字符n-gram，逻辑回归，类别加权，5折CV，测试集20\%，最大特征5000。
\end{algosposs}

\begin{algoresult}
以"连锁快餐顾客评价"案例为例（150条文本，测试集30条）：

\begin{itemize}
    \item \textbf{效果指标（核心）}：宏平均$F1=0.861$，准确率$0.867$，$AUC=0.940$。
    \item \textbf{和基线比}：多数类基线$0.600$（即全猜多数类的准确率）。模型0.867远高于0.600，说明\textbf{正负情感能被可靠识别}。
    \item \textbf{混淆矩阵}：看有多少正面被误判成负面、反之亦然，找出易错类型。
    \item \textbf{关键词证据}：模型系数会列出支持正面/负面的高区分度词（如"好吃、干净"推正面，"慢、凉、错"推负面）。
\end{itemize}

\textbf{怎么判断好坏}：$F1$和AUC都明显高于"多数类基线"才算有用；AUC>0.8算好，>0.9算很好。
\end{algoresult}

\begin{algonotes}
\textbf{什么时候用？}有一批\textbf{已标注正负情感}的文本，要批量自动判断海量新文本的情绪倾向。

\textbf{什么时候不用？}
\begin{itemize}
    \item \textbf{没有标注数据时}：逻辑回归是有监督方法，必须有标签。无标签时只能用词典法或先人工标注一小批。
    \item 情感是多分类（喜/怒/哀/乐）或需要细粒度时，二分类模型不够。
\end{itemize}

\textbf{常见坑}
\begin{itemize}
    \item \textbf{坑1}：没有标注数据就硬上有监督模型。标签是这个方法的命根子。
    \item \textbf{坑2}：类别严重不平衡（好评140条、差评10条）却不调类别权重，模型会"全猜好评"拿高准确率却毫无用处——所以要看$F1$和AUC，不能只看准确率。
    \item \textbf{坑3}：把测试集混进训练。必须独立留出测试集，否则效果指标虚高。
\end{itemize}

\textbf{和相似算法的区别}
\begin{itemize}
    \item vs \textbf{词频统计}：词频只数词；情感分析判断\textbf{情绪极性}，是分类任务。
    \item vs \textbf{LDA主题分析}：情感分析回答"情绪是正还是负"（有监督）；LDA回答"在讲什么话题"（无监督，不关心褒贬）。
\end{itemize}
\end{algonotes}

%% ========== LDA主题分析 ==========
\basicalgo{{\starmark}LDA主题分析}{lda_topic}

\begin{algointro}
LDA（隐含狄利克雷分配）是文本分析里最经典的"话题归纳器"。面对几千条没人读过的评论，它自动告诉你：这些文字大致分成$K$个话题，每个话题由哪些词构成、每条评论主要在聊哪个话题。它是\textbf{无监督}的——不需要人工标注。
\end{algointro}

\begin{algoscene}
\begin{itemize}
    \item 美赛C题给了大段无标注评论/政策文本，要\textbf{自动归纳话题结构}
    \item 顾客评论自动分成"质量、物流、价格、客服、尺码"等话题，并看各话题占比
    \item E/F题政策文本分析：归纳各方关注点、议题演变
\end{itemize}
\end{algoscene}

\begin{algoprinciple}
LDA用两个"骰子"的比喻来理解：

\begin{itemize}
    \item \textbf{文档--主题骰子}：每篇评论不是只聊一件事，而是按比例混合多个主题。比如某评论60\%在讲价格、40\%在讲物流。
    \item \textbf{主题--词骰子}：每个主题又是一堆词的概率分布。"价格主题"里最可能出现price、total、final；"物流主题"里最可能出现shipping、delivery、pack。
    \item \textbf{核心思想}：反复\textbf{共同出现}的词会被归到同一个主题。算法就是来回调整这两个骰子，让"每个主题的词"和"每篇文档的主题混合"都最自洽。
\end{itemize}

\textbf{怎么定主题数$K$}：不是拍脑袋。用\textbf{困惑度（perplexity）}扫描多个$K$，选困惑度下降变平缓的位置，再结合\textbf{主题可解释性}人工判断——主题词读起来像不像一个有意义的话题。

生活类比：LDA就像给一堆博客自动打标签。它不会预先告诉你标签是什么，而是自己发现"总被一起提到的词构成一类"，再让你给这些类起名字。
\end{algoprinciple}

\begin{algosposs}
\begin{enumerate}
    \item 工具箱选"LDA主题分析"
    \item 指定\textbf{文本列}feedback\_text、\textbf{文档ID列}review\_id
    \item 预处理：分词、去停用词（可加extra停用词）；设min\_df=2（至少出现在2篇文档才保留）、max\_df（剔除过于通用的词）
    \item 主题数：若不确定，开启\textbf{tune topics}让它扫描推荐$K$；确定后填num\_topics
    \item 学习方法选batch；doc\_topic prior与topic\_word prior设auto；随机状态固定可复现
    \item 运行，查看主题--词分布、文档--主题占比与代表文档
\end{enumerate}

\textbf{默认参数参考}：batch学习，自动先验，min\_df=2/max\_df过滤，主题数自动调优。
\end{algosposs}

\begin{algoresult}
以"跨境电商服装英文评论"案例为例（150篇有效文档）：

\begin{itemize}
    \item \textbf{主题数}：自动归纳为\textbf{5个主题}。
    \item \textbf{最强主题}：占比最高的是「price total final」（主题强度22.6\%），核心词price、total、final——显然在讲\textbf{价格/费用}。
    \item \textbf{怎么解读每个主题}：看该主题下概率最高的若干词，给它起名字。比如出现shipping/delivery/pack的就叫"物流"，出现size/fit/color的就叫"尺码版型"。
    \item \textbf{文档--主题分布}：每条评论有一个"主导主题"，据此统计各话题在全体评论中的占比。
    \item \textbf{代表文档}：每个主题抽几条最典型的原文回读，验证主题名取得对不对。
\end{itemize}

\textbf{怎么判断好坏}：主题词\textbf{语义清晰、彼此区分}就是好结果；如果两个主题词几乎一样，说明$K$太多或数据太稀，需要合并或调整。
\end{algoresult}

\begin{algonotes}
\textbf{什么时候用？}有大量\textbf{无标注}文本，想自动归纳话题结构、看讨论焦点分布。

\textbf{什么时候不用？}
\begin{itemize}
    \item 你已经知道文本该分哪几类、且有标签——直接做分类更准。
    \item 文本太短（每条就几个字）时，LDA效果差，短文本建议先用文本聚类。
\end{itemize}

\textbf{常见坑}
\begin{itemize}
    \item \textbf{坑1}：主题数$K$乱设。一定要结合困惑度扫描和人工可解释性，别只挑数字最小的。
    \item \textbf{坑2}：不看代表文档就给主题起名。主题词可能有歧义，必须回读原文确认。
    \item \textbf{坑3}：把LDA当因果。它只发现"词常一起出现"，不代表这些主题之间有因果关系。
\end{itemize}

\textbf{和相似算法的区别}
\begin{itemize}
    \item vs \textbf{文本聚类}：两者都是无监督分组。LDA给出\textbf{软分配}（一篇评论可同时属于多个主题，带比例）且有可读的主题词；K-Means聚类是\textbf{硬分配}（一篇只进一簇），适合短文本。
    \item vs \textbf{情感分析}：情感分析看"褒贬"，LDA看"话题"，回答的是完全不同的问题，常配合使用（先LDA分话题，再对每个话题做情感）。
\end{itemize}
\end{algonotes}

\section{进阶级算法速览}

%% ========== 社会网络关系图（共词矩阵） ==========
\advancedalgo{社会网络关系图（共词矩阵）}{co_word_network}

\begin{algobrief}
\textbf{一句话}：把"经常出现在同一篇文献里的关键词"连成一张网络图，看哪些词是研究热点核心、哪些词对总被搭配使用、整个领域分成几个主题簇。

\textbf{美赛场景区}：分析某领域文献/政策文件的关键词共现，识别研究热点结构与主题聚类，适合做领域综述、技术情报分析。

\textbf{核心原理}：以关键词为节点、两词在同一篇中共现的次数为连边权重，构建无向加权网络；计算度中心性（连边最多的核心词）、中介中心性（连接不同主题的"桥"），再做社区发现自动分成若干主题簇。

\textbf{操作要点}：工具箱选"社会网络关系图（共词矩阵）"；指定关键词列；布局spring；权重用共现次数；中心性度量度中心性；开启社区检测。用分隔符[ ,;，；、/|\\s]+拆分多关键词。

\textbf{结果关键}：案例（150篇智能制造文献）30个高频词节点、230条连边、5个社区、网络密度0.529；"智能制造"是核心节点，"智能制造--数字化转型"共现最强（21篇同现）。节点越大越居中越核心，连边越粗共现越强。

\textbf{注意}：共现只表示"同篇出现"，\textbf{不等于相关或因果}。中介中心性高的词即使频率不高，也可能是连接不同主题的关键桥接。
\end{algobrief}

%% ========== TF-IDF分析 ==========
\advancedalgo{TF-IDF分析}{tfidf}

\begin{algobrief}
\textbf{一句话}：在词频基础上加权——一个词在本文档出现越多、但在整个语料中出现越少，权重越高，从而揪出真正"代表这一类"的特色词。

\textbf{美赛场景区}：美赛C题分好组（如配送/质量/售后）后，找出每组\textbf{独有}的代表关键词，而不是各组都提到的通用词。

\textbf{核心原理}：TF（词频）$=$词在本文档出现次数；IDF（逆文档频率）$=\log\frac{\text{文档总数}}{\text{含该词的文档数}}$。TF-IDF$=TF\times IDF$。像"的、是"处处都有，IDF接近0，权重被压下去；只在某组反复出现的词，IDF高、权重突出。常用l2归一化。

\textbf{操作要点}：工具箱选"TF-IDF分析"；norm选l2；ngram range(1,1)；设min\_df/max\_df过滤；按分组质心提取各组代表词。

\textbf{结果关键}：案例中商品质量组突出"手感、紧密"，售后服务组突出"售后、客服、问题"，配送服务组突出"快递、配送、包装"——三组关注点被关键词清晰区分。

\textbf{注意}：TF-IDF是文本向量化的标准武器，也是情感分析、文本聚类、LDA之前的常用特征步骤。它找的是"区分性"，不是"重要性"。
\end{algobrief}

%% ========== 文本聚类分析 ==========
\advancedalgo{文本聚类分析}{text_cluster}

\begin{algobrief}
\textbf{一句话}：不预设任何类别标签，把"诉求相似"的文本自动归成几堆，适合客服留言这种口语化、长短不一、没人预先分类的数据。

\textbf{美赛场景区}：C题中对售后留言、投诉文本做自动分诊，归纳出几类典型问题及占比，供人工复核。

\textbf{核心原理}：先用TF-IDF字符n-gram把文本转成数值矩阵，再跑K-Means；从$k=2$到上限逐个试，以\textbf{轮廓系数}最高（同分取小$k$）选簇数；最后用每簇质心的高权重词命名主题，用轮廓系数、CH指数、簇内余弦相似度评估质量。

\textbf{操作要点}：工具箱选"文本聚类分析"；指定文本列；vectorizer选tfidf，分析器char\_wb，ngram 2+1；k设auto让其自动选；固定随机种子。

\textbf{结果关键}：案例（生鲜售后留言）聚成4簇，轮廓系数0.107（文本聚类普遍偏低，0.1左右即基本可分）。用每簇高频词人工概括主题（如"配送压坏""售后回复"等）。

\textbf{注意}：文本聚类轮廓系数天然偏低，别拿数值变量聚类的高标准来要求。聚类后\textbf{必须人工读每簇代表词来命名}，机器只负责分堆，主题含义要人来解释。
\end{algobrief}

%% ========== 新词发现 ==========
\advancedalgo{新词发现}{new_word}

\begin{algobrief}
\textbf{一句话}：不依赖任何词典，直接从海量文本里自动找出"大家新造出来、词典里还没有"的词，比如"半固态电池""零重力座椅"。

\textbf{美赛场景区}：分析新兴领域（新能源、AI、网络热词）时，业务词库更新跟不上，用它自动挖掘候选新词补充词表。

\textbf{核心原理}：三重证据筛选——(1)\textbf{凝固度（PMI互信息）}：字之间结合得紧不紧，切一刀算最小PMI；(2)\textbf{自由度（左右熵）}：这个词能不能出现在不同语境（左右邻居越多样越能独立成词）；(3)\textbf{词频}。三者都过阈值才入选，按综合得分排序。

\textbf{操作要点}：工具箱选"新词发现"；候选字串长度2$\sim$6；设最小词频、最小PMI、最小左右熵；去重、去纯数字、去非中文。

\textbf{结果关键}：案例（车主评论）挖出36个候选新词，头部如"半固态电池、零重力座椅、能量回收、城市领航、哨兵模式"，复用强度与成词证据最充分，可交人工确认后入词库。

\textbf{注意}：这是辅助工具，输出只是\textbf{候选清单}，必须人工复核再入库，机器可能把高频搭配误判成词。它解决的是"分词词典太旧"的问题。
\end{algobrief}

% 第18章 功效分析
\chapter{功效分析}

\begin{chapteroverview}
\textbf{这个分类是干什么的？}

功效分析回答的是一个\textbf{研究设计阶段}的问题：\textbf{样本量到底够不够大？}它在你正式收集数据之前就算清楚：要想以足够把握检出某个真实差异，需要招多少样本；反过来，现有样本量下，你有多大把握能发现效应。它不做假设检验，而是告诉你"实验要花多少钱、招多少人才不至于白做"。

\textbf{美赛什么题型常用？}

功效分析在美赛中\textbf{很少直接作为解题模型}，但它在两个地方很加分：
\begin{itemize}
    \item 题目要求你\textbf{评估数据收集方案、说明样本量是否充足}时，功效分析是最规范的论据。
    \item 论文中论证"我们采集的数据量足以支撑后续检验"时，引用功效分析结果能体现严谨性。
\end{itemize}
如果你拿到的是现成大数据集，基本用不到本章；只有涉及"规划抽样规模"时才用。

\textbf{学习路径建议}

先学\textbf{功效分析之均值差}（明星算法，最经典，讲透"样本量、功效、效应量、显著性水平"四者关系），再了解\textbf{功效分析之线性回归}。进阶部分按检验对象区分：\textbf{相关性}、\textbf{率差}、\textbf{方差（多组比较）}、\textbf{广义模型}（处理二项/计数结局），原理相通，只是背后的分布不同。
\end{chapteroverview}

\section{基础级算法}

%% ========== 功效分析之均值差 ==========
\basicalgo{{\starmark}功效分析之均值差}{power_mean}

\begin{algointro}
功效分析之均值差，就是在做两（或多）组均值比较之前，先算清楚：要检出预设的均值差异，每组需要多少样本，或者现有样本能达到多大的"检出把握"。
\end{algointro}

\begin{algoscene}
\begin{itemize}
    \item 规划对照试验、问卷调查时，确定\textbf{需要招募多少样本}
    \item 美赛中论证"现有数据量是否足以检出组间差异"
    \item 评估已有预试验数据的功效够不够
\end{itemize}
\end{algoscene}

\begin{algoprinciple}
功效分析绕着四个互相牵制的量转：

\begin{itemize}
    \item \textbf{显著性水平$\alpha$}：容忍"误报"的概率，通常0.05。
    \item \textbf{功效$1-\beta$}：真实差异存在时，能成功检出它的概率，通常目标0.8。
    \item \textbf{效应量$d$（Cohen's d）}：均值差除以合并标准差，$d\approx0.2$小、$0.5$中、$0.8$大。效应越小，越难检出，需要样本越多。
    \item \textbf{样本量$n$}：其他三个量定了，就能反解$n$。
\end{itemize}

生活类比：功效就像捕鱼。网眼大小（$\alpha$）、鱼到底多大（效应量）决定了要撒多大面积的网（样本量）才能有80\%把握捕到这条鱼。鱼越小（效应小），网要撒得越大。

\textbf{注意}：功效是"长期检出概率"，不是"这次一定显著"——达到80\%功效也意味着仍有20\%的可能检不出。
\end{algoprinciple}

\begin{algosposs}
\begin{enumerate}
    \item 工具箱选"功效分析之均值差"
    \item 检验类型选\textbf{两独立样本}（也支持配对、单样本、单因素方差分析）；检验方向\textbf{双侧}
    \item 求解目标选\textbf{样本量}；$\alpha=0.05$，目标功效$0.80$，两组分配比ratio$=1$
    \item 效应量选\textbf{数据估计优先}（用预试验数据自动算Cohen's d，没有时手动给值）
    \item 运行后把结果\textbf{向上取整}，再用取整规模反算实际功效
    \item 看灵敏度分析：效应量上下浮动25\%时样本量怎么变
\end{enumerate}

\textbf{默认参数参考}：两独立样本、双侧、$\alpha=0.05$、目标功效0.8、1:1分配，效应量由数据估计。
\end{algosposs}

\begin{algoresult}
以"新旧客服话术对比"案例为例（预试验每组75人）：

\begin{itemize}
    \item \textbf{效应量}：$d=0.33$（小到中等效应）。
    \item \textbf{核心结论}：在双侧$\alpha=0.05$、目标功效0.8下，每组约需\textbf{147例}（两组共294例）才能较有把握检出该均值差。
    \item \textbf{现有数据功效}：按预试验每组75例，实际功效仅约0.516，明显低于目标——说明试点规模只够粗估效应量，正式研究要按294例规划。
    \item \textbf{灵敏度}：效应量下调25\%（$d=0.247$）需260例/组；上调25\%（$d=0.411$）只需94例/组。效应量估得越小、要的样本越多。
\end{itemize}

\textbf{怎么看}：报告"达到80\%功效每组需$n$例"；若担心试点高估了效应，建议按保守值（下调25\%）招募。
\end{algoresult}

\begin{algonotes}
\textbf{什么时候用？}设计两/多组均值比较的研究前，定样本量。

\textbf{什么时候不用？}已经收集完数据、做假设检验时不需要功效分析（那是第3章的事）；现成大数据集通常样本量已远超需求。

\textbf{常见坑}
\begin{itemize}
    \item \textbf{坑1}：把"达到目标功效"理解成"必然显著"。功效是概率，仍有漏检可能。
    \item \textbf{坑2}：效应量拍脑袋给大了，导致样本量算得偏小、实际检不出。务必用预试验或文献估计效应量，并做灵敏度分析。
    \item \textbf{坑3}：忽略前提。功效公式近似要求数据近似正态、方差齐，案例中已用Shapiro与Levene检验确认。
\end{itemize}

\textbf{和相似算法的区别}
\begin{itemize}
    \item vs \textbf{第3章t检验/ANOVA}：t检验是拿到数据后\textbf{检验差异是否存在}；功效分析是\textbf{之前}算要多少样本。
    \item vs \textbf{功效分析之方差}：均值差针对两组；方差（ANOVA）功效针对三组及以上。
\end{itemize}
\end{algonotes}

%% ========== 功效分析之线性回归 ==========
\basicalgo{功效分析之线性回归}{power_regression}

\begin{algointro}
功效分析之线性回归，回答的是：要想让回归模型（或某个自变量的贡献）在统计上站得住，正式研究需要收集多少条样本。
\end{algointro}

\begin{algoscene}
\begin{itemize}
    \item 规划问卷调研：已知大概的$R^2$，算需要多少份有效问卷
    \item 论证回归分析的样本量是否充足
\end{itemize}
\end{algoscene}

\begin{algoprinciple}
回归功效基于\textbf{非中心$F$分布}：

\begin{itemize}
    \item 先用预试验数据拟合OLS，得到\textbf{调整$R^2$}，换算成效应量 $f^2 = R^2/(1-R^2)$。
    \item 非中心参数 $\lambda = n\cdot f^2$，自由度 $v = n-k-1$（$k$为预测变量数）。
    \item 功效$=1-F_{nc}(F_{\text{临界}}; u, v, \lambda)$，再二分反推满足目标功效的最小$n$。
    \item Cohen基准：$f^2\approx0.02$小、$0.15$中、$0.35$大。
\end{itemize}
\end{algoprinciple}

\begin{algosposs}
\begin{enumerate}
    \item 工具箱选"功效分析之线性回归"
    \item 指定目标列与预测变量列；效应来源选\textbf{data\_observed}（用数据估计$R^2$）
    \item 分析模式选\textbf{sample\_size}（求样本量）；$\alpha=0.05$，目标功效0.80
    \item 运行，查看所需最小样本量与功效曲线
\end{enumerate}

\textbf{默认参数参考}：$\alpha=0.05$、目标功效0.8，效应量由观测数据估计（用调整$R^2$，更保守）。
\end{algosposs}

\begin{algoresult}
以"顾客复购意愿"案例为例（4个预测变量，预试验150条）：

\begin{itemize}
    \item 观测$R^2=0.183$，调整$R^2=0.161$，换算 $f^2=0.191$（中等效应）。
    \item \textbf{核心结论}：$\alpha=0.05$、功效0.8下，正式研究至少需\textbf{68条}有效样本，此时功效80.5\%。
    \item 现有150条样本对应功效约99.5\%，绰绰有余。
\end{itemize}

\textbf{怎么看}：报告"至少需$n$条有效样本"；$n$对$\alpha$和目标功效敏感，收紧要求会明显增加样本量。
\end{algoresult}

\begin{algonotes}
\textbf{什么时候用？}规划回归类调研、论证问卷样本量。

\textbf{常见坑}
\begin{itemize}
    \item \textbf{坑1}：用未调整的$R^2$算效应量，会高估、样本量算少。应使用调整$R^2$。
    \item \textbf{坑2}：忽略预测变量个数$k$。变量越多，自由度损失越大，需要样本越多。
\end{itemize}

\textbf{和相似算法的区别}：本算法针对回归模型的整体或某个自变量的检验；均值差功效针对两组均值比较。
\end{algonotes}

\section{进阶级算法速览}

%% ========== 功效分析之相关性 ==========
\advancedalgo{功效分析之相关性}{power_corr}

\begin{algobrief}
\textbf{一句话}：要检验两个连续变量的相关是否显著（$H_0:\rho=0$），需要多少对成对观测，或现有样本能检出多强的相关。

\textbf{美赛场景区}：规划需要配对观测的相关分析（如培训学时与满意度），论证样本是否足够检出目标相关。

\textbf{核心原理}：检验统计量 $t=r\sqrt{(n-2)/(1-r^2)}$，原假设下服从$t(n-2)$，备择下按非中心参数$\delta=\rho\sqrt{n/(1-\rho^2)}$的非中心$t$分布精确算功效，Fisher z近似交叉验证。

\textbf{操作要点}：工具箱选"功效分析之相关性"；选x、y两列；相关类型pearson；求解目标sample\_size；双侧检验。

\textbf{结果关键}：案例（培训学时vs满意度）观测$r=0.357$，达80\%功效至少需约\textbf{56对}；现有150对功效约99.6\%。三问：现样本功效多大？需多少对？现规模能检出多小的$|r|$？

\textbf{注意}：效应量越小、目标功效越高，所需对数越多。
\end{algobrief}

%% ========== 功效分析之率差 ==========
\advancedalgo{功效分析之率差}{power_proportion}

\begin{algobrief}
\textbf{一句话}：两组合格率/发生率比较（二分类结局），在给定两组率的前提下，需要多少样本才能检出这个率差。

\textbf{美赛场景区}：规划临床试验、政策达标率对比，如降压治疗后达标率从30\%提到50\%需要多少患者。

\textbf{核心原理}：对两组率$p_1,p_2$做反正弦变换得到效应量Cohen's $h$，按两率$z$检验的解析功效公式求每组样本量；向上取整后回代实际功效，并用蒙特卡洛模拟复核正态近似。

\textbf{操作要点}：工具箱选"功效分析之率差"；填对照组率$p_1$、试验组率$p_2$；求解目标样本量；$\alpha=0.05$、双侧、1:1分配。

\textbf{结果关键}：案例$p_1=0.30$、$p_2=0.50$，$\alpha=0.05$、功效0.8下共需\textbf{186例}（每组93）；蒙特卡洛经验功效0.798与解析0.801吻合。

\textbf{注意}：样本量对$p_1,p_2$假设很敏感——率越接近50\%、两组差越小，需要的样本越多。务必做$p_1$扰动敏感性分析。
\end{algobrief}

%% ========== 功效分析之方差 ==========
\advancedalgo{功效分析之方差}{power_anova}

\begin{algobrief}
\textbf{一句话}：单因素方差分析（三组及以上均值比较）前，规划每组需要多少样本才能检出组间差异。

\textbf{美赛场景区}：如比较三套策略、三套方案的效果，预试验后规划正式对比实验的样本量。

\textbf{核心原理}：四要素互相牵制——效应量$f$（由预试验$\eta^2$估计，$f\approx0.1$小、$0.25$中、$0.4$大）、总样本量$N$、$\alpha$、功效$1-\beta$。基于非中心$F$分布$\lambda=N f^2$，按均衡分配反推每组样本量，取整后重算实际功效。

\textbf{操作要点}：工具箱选"功效分析之方差"；填组数$k$；效应来源数据估计；求解目标所需样本量；$\alpha$与目标功效默认0.05/0.8。

\textbf{结果关键}：案例（三套拣货策略，$k=3$）预试验得$\eta^2=0.083$、$f\approx0.30$（中效应），需每组\textbf{37例、共111例}，取整后功效约0.80。

\textbf{注意}：需先做正态性与方差齐性（Levene）检验；若异方差严重，当前结论不适用，应改用Welch方差分析的模拟功效。
\end{algobrief}

%% ========== 功效分析之广义模型 ==========
\advancedalgo{功效分析之广义模型}{power_glm}

\begin{algobrief}
\textbf{一句话}：当结局不是正态连续变量，而是二分类、计数等更复杂的结局时，用广义线性模型的功效分析来定样本量。是均值差功效的"升级版"。

\textbf{美赛场景区}：规划安全性/不良反应对照研究——如联合用药使肝功能异常率从20\%升到40\%（$OR=2$），需要多少患者。

\textbf{核心原理}：先按结局选分布族与链接函数（二项/泊松/负二项）；给定效应量、基线率、$\alpha$、分配比，用Wald检验的渐近功效二分求解总样本量；再固定种子生成数据、拟合对应GLM做蒙特卡洛模拟，在解点和1.25倍样本点核对解析近似是否靠谱。

\textbf{操作要点}：工具箱选"功效分析之广义模型"；family选binomial；solve选n；填效应（OR=2）、基线率、$\alpha=0.05$、双侧、功效0.8；指定group与outcome列。

\textbf{结果关键}：案例（二项结局、$OR=2$、基线20\%）需总样本\textbf{352例}（对照176、暴露176）；蒙特卡洛实测功效81.1\%，与解析80.1\%吻合。

\textbf{注意}：这是功效分析里最灵活但也最复杂的一类。结果对$\alpha$、基线率假设敏感；资源有限时可评估放宽$\alpha$或用单侧检验以降样本量。美赛中用到时，重点是把"样本量怎么来的"讲清楚。
\end{algobrief}


% ---------- 附录 ----------
\appendix
% 附录A 美赛真题题型与算法速查表
\chapter{美赛真题题型与算法速查表}

本章是给你的"考场急救手册"。比赛三天，拿到题目往往一时不知道从何下手——这张表就是帮你\textbf{先判断题型、再反查该用哪些算法、对应本书哪一章}。建议赛前打印出来贴在桌边。

\section{六大题型与常用算法}

美赛（MCM/ICM）按字母分为六类，特点和常用算法如下。

\subsection{A题：连续型问题}
A题常涉及连续变化的物理量、微分方程、连续优化（如流体、扩散、生态、物理过程）。
\textbf{常用算法}：描述性统计（第1章）摸清数据；回归分析（第4章）拟合连续关系；时间序列与计量经济（第6章）处理动态过程；预测模型（第7章）外推；规划求解（第12章）做连续优化。

\subsection{B题：离散型问题}
B题常涉及离散决策、调度、路径、分配、图论优化（如排班、选址、网络流）。
\textbf{常用算法}：规划求解（第12章，整数规划/0-1规划是主力）；聚类与降维（第5章）做分组；相关分析（第2章）找影响因素；差异比较（第3章）验证方案优劣。

\subsection{C题：大数据 / 数据挖掘}
C题给你一个较大的数据集（常含文本列、多变量、多表），要求清洗、探索、建模、预测。这是本书最对口的题型。
\textbf{常用算法}：描述性统计（第1章）起步；相关分析（第2章）与差异比较（第3章）探索关系；回归分析（第4章）与机器学习（第11章）建模预测；\textbf{文本分析（第17章）}处理评论/政策文本（词频、情感分析、LDA主题分析）；聚类与降维（第5章）做客户/样本分群；综合评价（第8章）打分排序。

\subsection{D题：网络科学}
D题关注节点与连接构成的网络（如社交网络、交通网、传播网络）。
\textbf{常用算法}：\textbf{社会网络关系图（共词矩阵，第17章）}的中心性与社区发现思路可迁移；聚类与降维（第5章）做网络社团识别；相关分析（第2章）找节点属性关系；规划求解（第12章）做网络流/最短路。

\subsection{E题：环境科学}
E题涉及生态、污染、可持续发展、资源环境，常需综合多方证据、做政策评估。
\textbf{常用算法}：综合评价（第8章）构建环境指标体系；空间计量（第14章）分析区域分布；时间序列（第6章）看趋势；\textbf{Meta分析（第16章）}汇总多地区/多研究的环境效应估计；文本分析（第17章）分析政策文件与公众意见。

\subsection{F题：政策建模}
F题关注社会经济政策、公平性、伦理，数据常偏宏观、定性与定量混合。
\textbf{常用算法}：综合评价（第8章）做政策效果打分；回归分析（第4章）与机器学习（第11章）找政策影响因素；\textbf{Meta分析（第16章）}综合多项政策研究证据；\textbf{功效分析（第18章）}论证样本量是否充足；问卷研究（第9章）设计调研。

\section{问题类型 $\rightarrow$ 推荐算法速查表}

下表按"你拿到题目后会遇到的典型任务"反查算法。左列是任务信号，右列是首选算法与章节。

\begin{longtable}{p{4.2cm} p{6.2cm} p{2.6cm}}
\rowcolor{algoBlue!15}
\textbf{题目信号 / 任务} & \textbf{首选算法} & \textbf{章节} \\
\endfirsthead
\rowcolor{algoBlue!15}
\textbf{题目信号 / 任务} & \textbf{首选算法} & \textbf{章节} \\
\endhead

拿到任何数据，第一步 & 描述性统计、数据概览 & 第1章 \\
数据偏态严重、要不要换方法 & 正态性检验 & 第1章 \\
两类样本谁高谁低、差多少 & 两独立样本t检验 / Mann-Whitney & 第3章 \\
三组及以上比较、看组间差异 & 单因素方差分析 / Kruskal-Wallis & 第3章 \\
两个变量是否一起变动、多强 & 相关分析（Pearson/Spearman） & 第2章 \\
用一个或多个变量预测另一变量 & 线性回归 / 多元回归 & 第4章 \\
分类预测（是否违约、好评差评） & 逻辑回归 / 机器学习分类 & 第11章 \\
给对象打分排序、选最优方案 & 综合评价（AHP/熵权/TOPSIS等） & 第8章 \\
把样本自动分成几类 & 聚类分析（K-Means等） & 第5章 \\
变量太多、想降维 & 主成分分析 / 因子分析 & 第5章 \\
有时间顺序、要预测未来 & 时间序列（ARIMA等） & 第6章 \\
要做更准的数值预测 & 机器学习预测模型 & 第7、11章 \\
文本里大家最常说什么词 & 词频统计、词云图 & 第17章 \\
评论文本是好评还是差评 & 文本情感分析 & 第17章 \\
一大段文字里讲了哪几个话题 & LDA主题分析 & 第17章 \\
想找某类文本独有的关键词 & TF-IDF分析 & 第17章 \\
无标签留言自动分诊归类 & 文本聚类分析 & 第17章 \\
各研究结果不一，要合并成总结论 & 平均值/二分类/相关系数Meta分析 & 第16章 \\
只有一堆p值，想判断整体是否有效应 & P值合并 & 第16章 \\
规划抽样：要招多少样本才够 & 功效分析之均值差/回归 & 第18章 \\
多城市/多区域数据有空间聚集 & 空间计量 & 第14章 \\
流程/工序稳定性、合格率监控 & 质量控制 & 第13章 \\
实验分组怎么设计最省样本 & 试验设计 & 第15章 \\
问卷数据信效度、因子结构 & 问卷研究 & 第9章 \\
\end{longtable}

\section{推荐建模流程}

拿到题目后，绝大多数数据驱动的题目（尤其C/E/F题）都可以按下面这条主线走，\textbf{不要一上来就上复杂模型}：

\begin{enumerate}
    \item \textbf{第一步：描述性统计（第1章）}
    先把数据看透——均值、分布、缺失、异常值。这一步决定你后面能不能用参数方法，论文的"数据描述"部分也靠它。\textbf{没有这一步，后面都是空中楼阁。}

    \item \textbf{第二步：差异比较 / 相关探索（第2、3章）}
    看各组之间有没有差别、变量之间有没有关系。这一阶段帮你\textbf{发现规律、选对变量}，也为论文的"探索性分析"提供图表。

    \item \textbf{第三步：回归 / 评价建模（第4、8章）}
    在前两步基础上，建立解释或评价模型：想找影响因素用回归，想打分排序用综合评价。这是论文的\textbf{核心模型部分}。

    \item \textbf{第四步：预测 / 深化（第7、11、17章）}
    题目要求预测未来时，用时间序列或机器学习；有文本数据时，用第17章的词频、情感分析、LDA补充定性洞察。

    \item \textbf{贯穿始终：稳健性与敏感性检验}
    换个样本量、换个参数结论还成立吗？功效分析（第18章）说明样本够不够，敏感性分析说明结论扎不扎实。\textbf{评委最看重这一步。}
\end{enumerate}

\textbf{一句话口诀}：先描述、再比较、后建模、终预测，每步都做稳健性。

\section{使用提醒}

\begin{itemize}
    \item 本表给的是\textbf{首选算法}，实际题目常需要2$\sim$3个算法组合使用（如C题：描述统计 $\rightarrow$ 回归建模 $\rightarrow$ 情感分析 $\rightarrow$ 综合评价）。
    \item 选算法前先看\textbf{数据类型}：连续数值用回归/方差分析，分类结局用逻辑回归，文本用第17章，多研究证据合并用第16章。
    \item 任何模型跑完都要\textbf{回原文/回数据核对}，别迷信输出数字——这是全书反复强调的原则。
\end{itemize}

\chapter{优秀论文模型拆解}

本章从 2020--2025 年 MCM/ICM 的 C 题（离散建模题）与 E 题（可持续性环境题）共 12 道真题中，每道题挑选 2--3 篇代表性获奖论文，重点拆解它们的\textbf{模型组合流水线}：先用什么做数据清洗，再用什么建模，最后用什么评价与预测。读者可以把这一章当作一张"模型搭配菜谱"，在自己参赛时按图索骥。文中括号里标注的章节号对应本书正文 18 个分类，方便回查算法细节。

% ============================================================
\section{2020 年 C 题：亚马逊商品评论与评分分析}

\textbf{题目背景}：Sunshine 公司希望根据吹风机、安抚奶嘴、微波炉三类商品在亚马逊上的历史评分（1--5 星）与文字评论，判断商品口碑、预测未来销量，并给新品上市提建议。核心难点是\textbf{同时利用结构化的星级数据和非结构化的评论文本}。

\subsection{论文 2020C-2002116：CE-VADER 情感分析 + AR 时序预测}
\textbf{使用模型}：情感分析（第 17 章 文本分析）、LDA 主题模型（第 17 章）、自回归模型 AR（第 6 章 时间序列与计量经济）、微分方程口碑模型（第 12 章 规划求解中的连续动力学）、描述性统计（第 1 章）。

\textbf{模型组合逻辑}：论文先对评论文本做清洗、分词、词形还原，再用改进的 VADER 词典法把每条评论打到五档情感分（强正/弱正/中/弱负/强负）；接着把五档情感分与五颗星做交叉表，验证文本情感与星级是否对齐；然后用微分方程定义"声誉率"随时间的衰减与累积；最后用 AR 自回归模型对未来声誉率做外推，并用 RMSE 在验证集上检验。

\textbf{为什么这么选}：文本情感负责把"消费者说了什么"量化成连续分数，微分方程负责刻画口碑随时间的传播规律，AR 模型则利用声誉率的自相关性做短期预测。三者分别解决"量化""机理""外推"三件事，互不替代。

\textbf{可借鉴亮点}：把情感分析当作\textbf{数据预处理}而不是终点——情感分最终被喂给时序模型做预测，这种"文本→数值→预测"的流水线是处理评论类数据的标准套路。

\subsection{论文 2020C-2009116：词典法情感打分 + 主成分分析 + ARIMA}
\textbf{使用模型}：描述性统计（第 1 章）、缺失值/异常值处理（第 1 章）、词典法情感分析（第 17 章）、主成分分析 PCA（第 5 章 聚类与降维）、ARIMA 时间序列（第 6 章）、Pearson/Spearman 相关分析（第 2 章）、Mann-Whitney U 检验（第 3 章 差异比较）。

\textbf{模型组合逻辑}：先做缺失值、异常值、重复值清洗，并用描述性统计和箱线图初步摸清三类商品的评分分布；再用基于词典的方法给每条评论打情感分；为了把多维的评论指标（长度、有用票数、星级、情感分等）压缩，用 PCA 提取主成分；最后对聚合后的日声誉序列建 ARIMA 模型做未来趋势预测。

\textbf{为什么这么选}：当自变量维度多、相关性强时，PCA 能先去冗余再喂给时序模型，避免 ARIMA 因输入维度爆炸而失稳；非参数 Mann-Whitney 检验则用来比较不同商品间评分分布是否有显著差异，弥补均值比较的粗糙。

\textbf{可借鉴亮点}：把 PCA 放在"特征工程"环节，而不是当成最终评价手段——这种"先用降维降噪、再用时序外推"的两段式结构，在数据维度高时非常实用。

\subsection{论文 2020C-2010638：Logistic 回归 + LDA 主题模型 + AHP 评价}
\textbf{使用模型}：NLTK 文本分词与情感打分（第 17 章）、多分类 Logistic 回归（第 4 章 回归分析）、LDA 主题模型（第 17 章）、层次分析法 AHP（第 8 章 综合评价）、聚类分析（第 5 章）、相关分析（第 2 章）。

\textbf{模型组合逻辑}：先用 NLTK 把评论转成 [-1,1] 的情感分；再用多分类 Logistic 回归分析"有用票数"与评论长度、星级、情感分之间的非线性关系（发现倒 U 型）；接着用 LDA 主题模型自动抽出评论者最关心的几个话题；最后用 AHP 对多个商品维度赋权，给出销售策略建议。

\textbf{为什么这么选}：Logistic 回归擅长处理"有用票数"这种分类型目标，LDA 负责回答"消费者在聊什么"，AHP 则把前面的量化结果汇总成一个可解释的打分体系。三个模型对应\textbf{预测、解释、决策}三层任务。

\textbf{可借鉴亮点}：一篇论文里同时出现"统计回归+无监督主题+主观赋权评价"，正好演示了定量与定性方法如何在同一道题里各司其职。

% ============================================================
\section{2020 年 E 题：一次性塑料废弃物减排}

\textbf{题目背景}：一次性塑料废弃物对海洋、陆地和人类健康造成持续损害。题目要求建立模型估算"在不造成进一步环境损害的前提下，每年最多可以安全处理多少一次性塑料废弃物"，并设计减排政策。

\subsection{论文 2020E-2010035：PWH 系统动力学 + 动态规划 + K-means/AHP}
\textbf{使用模型}：系统动力学/物质流模型（第 12 章 规划求解）、动态规划（第 12 章）、K-means 聚类（第 5 章）、层次分析法 AHP（第 8 章）、一元线性回归（第 4 章）、ARMA 时序（第 6 章）。

\textbf{模型组合逻辑}：先构建 PWH（Production-Waste-Harm）物质流模型，把塑料按年产量、寿命、回收率、焚烧率分成若干状态，逐年滚动计算废弃物量和环境损害；然后在 PWH 外层嵌套动态规划，在"替代材料、政策干预、回收技术"三个决策变量下求最小损害路径；为了给不同国家/地区分类施策，再用 K-means 把样本国聚成几类，用 AHP 给政策方案赋权。

\textbf{为什么这么选}：塑料问题是典型的\textbf{存量-流量}问题，必须用微分/差分方程刻画逐年累积；动态规划负责在多年时间轴上找最优决策；聚类和 AHP 则把全球样本分而治之，避免一个模型套所有国家。

\textbf{可借鉴亮点}：把机理模型（PWH）当作仿真器，把动态规划当作优化器，把聚类+AHP 当作政策解释层——这种"仿真+优化+评价"三段式是环境政策类题目的经典组合。

\subsection{论文 2020E-2016287：环境承载力阈值模型 + PCA + 指数平滑}
\textbf{使用模型}：环境承载力概念模型（第 12 章）、主成分分析 PCA（第 5 章）、指数平滑/Holt-Winters（第 6 章）、熵权法（第 8 章）、线性回归（第 4 章）。

\textbf{模型组合逻辑}：先定义环境承载力量 EBQ 与环境承载力容量 EBC，二者相等即为"安全阈值"；为了从众多环境指标中提取综合承载力指标，用 PCA 降维；然后用指数平滑法外推未来废弃物产生量；最后用熵权法给各影响因子赋权，得到不同减排政策的综合得分。

\textbf{为什么这么选}：阈值类问题需要一个清晰的"等式"作为决策依据（EBQ=EBC），PCA 解决指标多而杂的问题，指数平滑则在数据季节性强时比 ARIMA 更稳健。

\textbf{可借鉴亮点}：用"承载力容量 vs. 实际负荷"的差值作为预警信号，这种\textbf{阈值思想}在所有环境/安全/金融风控题里都可以套用。

\subsection{论文 2020E-2021103：以海鸟为指示物种的种群微分方程}
\textbf{使用模型}：种群动力学微分方程（第 12 章 连续模型）、生物指示物法（第 8 章 综合评价）、线性回归（第 4 章）。

\textbf{模型组合逻辑}：论文没有直接给"塑料污染指数"打分，而是选择海鸟种群作为海洋健康的代理指标，建立含出生率、自然死亡率、塑料致死亡率的常微分方程，把累积塑料量作为驱动变量代入，求解未来海鸟种群规模，反推当前塑料产量的"安全上限"。

\textbf{为什么这么选}：直接量化"环境损害"很难，但海鸟数量是可观测的；用一个机理清晰的 ODE 把"塑料产量→鸟类死亡"这条因果链写出来，比拍脑袋加权更有说服力。

\textbf{可借鉴亮点}：当题目要求"评估环境影响"时，找一个\textbf{可观测的生物/物理代理指标}再建机理模型，比堆砌评价指标更有说服力。

% ============================================================
\section{2021 年 C 题：亚洲大黄蜂入侵监测}

\textbf{题目背景}：亚洲大黄蜂（Vespa mandarinia）在美国华盛顿州被发现，威胁本地蜜蜂和养蜂业。公众上报了大量目击记录，真假混杂。题目要求预测其扩散、给目击记录分优先级、并判断是否已被剿灭。

\subsection{论文 2021C-2107870：K-means 半监督 + 神经网络分类}
\textbf{使用模型}：K-means 聚类（第 5 章）、半监督机器学习（第 11 章 机器学习）、神经网络分类（第 11 章）、ROC/AUC 评价（第 11 章）、描述性统计 EDA（第 1 章）。

\textbf{模型组合逻辑}：面对大量"未确认"的目击图片，先用 K-means 把未标注图片聚成若干簇，再把簇中心与已知正例/负例比对，自动给未标注样本打伪标签；然后用这些带伪标签的数据训练神经网络分类器，给每条目击记录输出一个"误报概率"，按概率从高到低安排调查优先级。

\textbf{为什么这么选}：标注数据极少（只有少数 PositiveID/NegativeID），纯监督学习会过拟合；半监督聚类+伪标签是解决"少标签、多数据"的标准武器。

\textbf{可借鉴亮点}：把"聚类"用在\textbf{自动打标签}上，而不是用来做客户分群——这种用法在数据科学竞赛里非常常见，值得借鉴。

\subsection{论文 2021C-2109298：元胞自动机 + SVM + 逻辑增长}
\textbf{使用模型}：元胞自动机 Cellular Automata（第 12 章 仿真模型）、支持向量机 SVM（第 11 章）、逻辑增长模型 Logistic（第 4/12 章）、GIS 空间分析（第 14 章 空间计量）、蒙特卡洛模拟（第 12 章）。

\textbf{模型组合逻辑}：把华盛顿州切成 2925 个 12 公里见方的栅格，每个栅格作为元胞，根据温度、季节、邻近栅格的大黄蜂数量按自定义规则更新状态，模拟未来 10 年的扩散；用 Logistic 模型拟合种群总量的 S 型增长；用 SVM 对目击记录做真假二分类；最后用蒙特卡洛扰动参数，给出扩散范围的置信区间。

\textbf{为什么这么选}：生物入侵本质是\textbf{空间扩散过程}，元胞自动机天然适合刻画"从一个格子跳到邻近格子"的局部规则；SVM 则在高维小样本的图片/文本分类上表现稳定。

\textbf{可借鉴亮点}：元胞自动机+GIS 栅格化是处理"地理扩散"类题目的利器，比纯回归模型更能画出未来分布图。

\subsection{论文 2021C-2107815：Logistic 增长 + 高斯扩散 + 空间适宜性}
\textbf{使用模型}：Logistic 增长模型（第 4 章）、高斯概率分布（第 1 章）、因子分析（第 5 章）、GIS 空间可视化（第 14 章）、MAE 误差评价（第 11 章）。

\textbf{模型组合逻辑}：先用因子分析从温度、海拔、植被等环境变量中提取"环境适宜性"公因子；然后用 Logistic 方程描述蜂群数量随时间的 S 型增长；再用高斯分布把每个目击事件的时间和位置展开成概率云，叠加到栅格上得到巢穴位置概率图。

\textbf{为什么这么选}：Logistic 解决"数量怎么长"，高斯概率解决"在哪里出现"，因子分析解决"环境因子怎么浓缩"——时间、空间、环境三个维度各用一个模型。

\textbf{可借鉴亮点}：用高斯概率云把"点目击"转化成"面概率"，再叠加到 GIS 地图上，这种可视化方式在所有空间分布题里都可以复用。

% ============================================================
\section{2021 年 E 题：全球粮食系统优化}

\textbf{题目背景}：现有粮食系统追求效率和利润，却造成环境破坏和饥饿不平等。题目要求在"环境可持续、分配公平、经济效益"三者间重新优化全球粮食系统。

\subsection{论文 2021E-2122175：熵权法 + 变异系数法 + 粒子群优化}
\textbf{使用模型}：熵权法 EWM（第 8 章 综合评价）、变异系数法 CVM（第 8 章）、粒子群优化 PSO（第 12 章 规划求解）、灰色系统预测（第 6 章）、MAPE 误差（第 11 章）、敏感性分析（第 1 章）。

\textbf{模型组合逻辑}：先把粮食系统拆成"活动—结果—目标"三层指标体系；用熵权法和变异系数法\textbf{客观地}给每个指标赋权（避免 AHP 的主观随意性），合成食品系统结果指数 FOI 和目标指数 FGI；然后用粒子群算法调整各指标取值，在满足约束下最小化环境足迹、最大化公平性；最后用 MAPE 和敏感性分析检验权重扰动下结论是否稳健。

\textbf{为什么这么选}：指标多且权重难定时，两种客观赋权法互相校验比单押一种更稳；PSO 则适合处理这种变量多、约束模糊的连续优化问题。

\textbf{可借鉴亮点}：\textbf{双客观赋权法交叉验证}（熵权+变异系数）是应对"主观赋权被质疑"的标准做法，比单用 AHP 更严谨。

\subsection{论文 2021E-2119031：运输优化 + 线性规划 + 遗传算法}
\textbf{使用模型}：线性规划 LP（第 12 章）、非线性规划（第 12 章）、遗传算法 GA（第 12 章）、微分方程（第 12 章）、熵权法（第 8 章）。

\textbf{模型组合逻辑}：把粮食系统拆成生产、运输、销售三段；在运输段建立"利润驱动的运输选择算法"，本质是一个带产地/销地约束的线性规划问题；当生产函数引入规模报酬非线性时，改用遗传算法全局求解；最后用熵权法对各国方案打分。

\textbf{为什么这么选}：粮食调拨本质是\textbf{运输问题}，LP 是标准答案；当目标函数非线性或整数约束多时，遗传算法比单纯形法更灵活。

\textbf{可借鉴亮点}：先从最简单的 LP 入手，发现不够再升级到 GA——这种"先易后难、按需升级"的建模思路比一上来就上深度学习更稳妥。

\subsection{论文 2021E-2102057：EEE 三系统耦合模型 + Logistic 回归}
\textbf{使用模型}：三系统耦合综合评价（第 8 章）、Logistic 回归（第 4 章）、GIS 空间分析（第 14 章）、结构方程 SEM（第 5 章）、因子分析 EFA（第 5 章）。

\textbf{模型组合逻辑}：把粮食系统抽象为环境（E）、公平（E）、经济（E）三个互锁子模型；每个子模型用若干指标加权合成；用 Logistic 回归分析"改变膳食结构"对温室气体排放的影响；用 GIS 把各国结果画在地图上；用 SEM 检验三个子系统之间的因果路径。

\textbf{为什么这么选}：题目要求同时优化三个互相冲突的目标，把系统拆成子模块再耦合，比建一个大而全的黑箱模型更可控、可解释。

\textbf{可借鉴亮点}：\textbf{模块化建模}——每个子问题单独建一个小模型，再通过耦合变量连起来，这是处理复杂系统题目的通用范式。

% ============================================================
\section{2022 年 C 题：黄金与比特币每日交易策略}

\textbf{题目背景}：给交易员 1000 美元初始资金，根据黄金和比特币的每日历史价格，决定每天买/卖/持有，最大化收益。每次交易要付手续费。

\subsection{论文 2022C-2200688：Prophet 预测 + XGBoost 决策 + 模型对比}
\textbf{使用模型}：Facebook Prophet 时序预测（第 6 章）、XGBoost 机器学习（第 11 章）、LSTM 深度学习（第 7 章 预测模型）、ARIMA（第 6 章）、EDA 探索性分析（第 1 章）。

\textbf{模型组合逻辑}：先对价格序列做 EDA，观察趋势、节假日效应、波动聚类；然后并行用 ARIMA、LSTM、Prophet 三种方法预测下一日价格，用 MAE/RMSE 对比后选择 Prophet（因为它对节假日和突变点鲁棒、可解释）；再把预测值、过去 N 日收益率、波动率作为特征喂给 XGBoost，输出"买/卖/持"三分类决策；最后回测五年数据。

\textbf{为什么这么选}：金融价格噪声大、突变多，ARIMA 抓非线性弱，LSTM 容易过拟合，Prophet 自带趋势+季节+节假日分解，更适合这种有明显日历效应的金融序列；XGBoost 则负责把预测信号翻译成交易动作。

\textbf{可借鉴亮点}：\textbf{多模型并行对比再选优}，而不是凭直觉选一个——这种做法在论文里显得严谨，也能避免被评委质疑"为什么用这个模型"。

\subsection{论文 2022C-2208834：协整配对交易 + LSTM/XGBoost}
\textbf{使用模型}：协整检验（第 6 章 计量经济）、误差修正模型 ECM（第 6 章）、LSTM（第 7 章）、XGBoost（第 11 章）、蒙特卡洛模拟（第 12 章）。

\textbf{模型组合逻辑}：先检验黄金和比特币价格序列是否协整（即长期存在稳定价差）；发现协整后，构建配对交易策略——当价差偏离均值时做多低估、做空高估；用改进的 Dollar Neutral 策略控制仓位；再用 LSTM/XGBoost 预测价差短期走势来辅助开平仓信号；最后用蒙特卡洛模拟扰动参数做稳健性检验。

\textbf{为什么这么选}：两个风险资产如果协整，就可以\textbf{对冲掉市场风险}，只赚均值回归的钱——这是量化金融里最经典的套利思路之一。

\textbf{可借鉴亮点}：把计量经济学的协整/ECM 用在交易策略里，比纯时序外推更有金融直觉，论文立刻显得"懂行"。

\subsection{论文 2022C-2218743：ARIMA 预测 + CVaR 风险度量 + NSGA-II}
\textbf{使用模型}：ARIMA（第 6 章）、CVaR 条件风险价值（第 8/12 章）、NSGA-II 多目标遗传算法（第 12 章）、移动平均趋势判断（第 6 章）、敏感性分析（第 1 章）。

\textbf{模型组合逻辑}：先用 ARIMA 预测下一日价格；用移动平均长排列法判断当前是牛市还是熊市，决定是否交易；然后构建"收益-CVaR"双目标优化模型，用改进的 NSGA-II 求一组 Pareto 最优仓位；再用下行半方差约束过滤极端方案。

\textbf{为什么这么选}：只看收益不看风险是散户思维，CVaR 比 VaR 更能衡量尾部损失；NSGA-II 专门处理"收益 vs. 风险"这种本质冲突的双目标问题。

\textbf{可借鉴亮点}：\textbf{多目标优化画 Pareto 前沿}，把"最佳策略"变成"一系列可供选择的方案"，比强行给一个答案更专业。

% ============================================================
\section{2022 年 E 题：森林管理与碳汇}

\textbf{题目背景}：森林既能固碳，又能提供木材。题目要求设计一份森林管理方案，在"适度砍伐收获木材"和"保留森林继续固碳"之间找到平衡。

\subsection{论文 2022E-2202838：换算因子法 + Logistic 生长 + 动态目标规划}
\textbf{使用模型}：森林生物量换算因子法（第 8 章 综合评价）、Logistic 生长方程（第 4/12 章）、动态目标规划（第 12 章）、AHP 层次分析（第 8 章）、GIS 栅格（第 14 章）、敏感性分析（第 1 章）。

\textbf{模型组合逻辑}：先按"换算因子法"用森林蓄积量乘以碳转换系数，分别计算活立木、林下植物、土壤三部分的碳储量；用 Logistic 方程描述树木生物量随年龄的 S 型增长；然后把森林分成多个区域，以"是否砍伐、树龄、死亡率"为状态变量，建立动态目标规划模型，在"碳汇最大、经济收益最大"两个目标下求最优轮伐期；用 AHP 给不同目标赋权。

\textbf{为什么这么选}：森林生长有明确的生物学 S 曲线，Logistic 比线性外推更贴合；动态规划则天然处理"多年决策路径"问题。

\textbf{可借鉴亮点}：把林业领域成熟的\textbf{换算因子法}拿来即用，再嵌套优化模型——跨学科借用领域公式是获奖论文的常见打法。

\subsection{论文 2022E-2223495：微分方程 + 多目标规划 + 演化博弈}
\textbf{使用模型}：常微分方程（第 12 章）、多目标规划（第 12 章）、演化博弈（第 12 章）、模拟退火（第 12 章）、社会网络分析 SNA（第 17 章）。

\textbf{模型组合逻辑}：先建碳汇微分方程，把森林面积、树龄、气候、木材用途、碳分解率都写成参数，模拟不同砍伐率下的碳汇轨迹；然后用多目标规划同时优化碳汇量和经济收益；再引入演化博弈，刻画政府、林农、企业三方在不同政策下的策略演化；用模拟退火求解。

\textbf{为什么这么选}：森林管理不只是技术问题，还是\textbf{利益相关方博弈问题}——引入演化博弈能解释"为什么政策落地难"，让方案更现实。

\textbf{可借鉴亮点}：在环境类题目里引入博弈论分析利益相关方，能显著提升论文的政策深度。

\subsection{论文 2022E-2209336：AHP+EWM 组合赋权 + DEA 效率评价}
\textbf{使用模型}：AHP（第 8 章）、熵权法 EWM（第 8 章）、DEA 数据包络分析（第 8 章）、BP 神经网络（第 11 章）、GIS（第 14 章）。

\textbf{模型组合逻辑}：先分别用 AHP（主观）和 EWM（客观）给森林价值的五个指标赋权，再把两组权重加权融合，得到综合得分；用 DEA 评价不同管理方案的投入产出效率；最后用 BP 神经网络对新方案做预测打分。

\textbf{为什么这么选}：AHP 反映专家经验，EWM 反映数据本身的离散程度，二者融合比单押一种更稳健；DEA 则擅长评价"多投入多产出"的方案效率。

\textbf{可借鉴亮点}：\textbf{主客观组合赋权}（AHP+熵权）是综合评价题里出现频率极高的套路，几乎可以套用到所有"打分排序"类问题。

% ============================================================
\section{2023 年 C 题：Wordle 单词游戏数据分析}

\textbf{题目背景}：Wordle 是 NYT 推出的每日猜词游戏，玩家在 Twitter 上分享自己的猜测次数分布。题目要求预测未来分享人数、分析单词难度、并给游戏设计提建议。

\subsection{论文 2023C-2307946：ARIMA + LightGBM + GBDT/MMoE}
\textbf{使用模型}：ARIMA 时序（第 6 章）、LightGBM（第 11 章）、GBDT（第 11 章）、MMoE 多任务学习（第 11 章）、K-means（第 5 章）、EDA（第 1 章）。

\textbf{模型组合逻辑}：先对每日分享人数做时序分解，确定 ARIMA 的阶数后用 ARIMA(0,1,1) 给出 80\% 置信预测区间；然后定义单词词频、重复字母、元音数等属性，用 LightGBM 预测"困难模式比例"；最后用 GBDT/MMoE 多任务模型同时预测 (1,2,3,4,5,6,X) 七个猜测次数的百分比分布。

\textbf{为什么这么选}：分享人数是单变量时序，ARIMA 足够；猜测次数分布是多标签回归，MMoE 多任务学习能共享信息、提升小样本表现。

\textbf{可借鉴亮点}：\textbf{按子问题选模型}——时序问题用 ARIMA，多标签回归用 GBDT/MMoE，不强行用一个模型包打天下。

\subsection{论文 2023C-2310767：ARIMA-LSTM 混合 + Spearman + Lasso + GMM}
\textbf{使用模型}：ARIMA（第 6 章）、LSTM（第 7 章）、Spearman 相关（第 2 章）、Lasso 回归（第 4 章）、高斯混合模型 GMM（第 5 章）、XGBoost（第 11 章）。

\textbf{模型组合逻辑}：发现 ARIMA 只能抓线性趋势、LSTM 只能抓非线性，于是把二者残差相加构建 ARIMA-LSTM 混合模型，RMSE 降到 0.0432；然后用 Spearman 相关分析单词属性与困难模式比例的单调关系；用 Lasso 回归做变量选择；用 GMM 把单词按难度聚成几类。

\textbf{为什么这么选}：金融/时序数据往往同时含线性和非线性成分，\textbf{线性模型抓趋势+残差用深度学习拟合}是经典混合建模套路。

\textbf{可借鉴亮点}：ARIMA-LSTM 残差组合法简单有效，比纯 LSTM 更稳、比纯 ARIMA 更准，是时序预测的首选 baseline。

\subsection{论文 2023C-2307166：SIR 传染病模型 + 遗传算法 + Bootstrap}
\textbf{使用模型}：SIR 传染病微分方程（第 12 章）、遗传算法 GA（第 12 章）、Bootstrap 重抽样（第 1 章）、KNN 近邻（第 11 章）、交叉验证（第 11 章）。

\textbf{模型组合逻辑}：把 Wordle 的传播类比成传染病——玩家=感染者，长期不玩=康复者，Twitter 分享=传播途径；用 SIR 微分方程拟合每日分享人数曲线，用遗传算法最小化 MSE 来标定参数；用 Bootstrap 重抽样 1000 次给出预测置信区间；用 KNN 对单词难度做分类。

\textbf{为什么这么选}：Wordle 的流行曲线和传染病爆发曲线形状高度相似（先爆发后衰减），SIR 模型有明确的机理可解释，比纯黑箱时序模型更有故事性。

\textbf{可借鉴亮点}：\textbf{跨领域类比建模}——把"社交传播"看成"传染病"，是传播类题目（热搜、病毒视频、新游戏）的通用思路。

% ============================================================
\section{2023 年 E 题：光污染风险评估与干预}

\textbf{题目背景}：人工灯光造成的光污染影响野生动物、植物和人类健康。题目要求建立光污染风险评估模型，对不同类型地区（城市、郊区、乡村、保护区）分级并设计干预策略。

\subsection{论文 2023E-2307336：EWM-TOPSIS 风险评价}
\textbf{使用模型}：熵权法 EWM（第 8 章）、TOPSIS 逼近理想解排序（第 8 章）、GIS 夜间遥感（第 14 章）、微分方程（第 12 章）、敏感性分析（第 1 章）。

\textbf{模型组合逻辑}：从人类、野生生物、植物、能源浪费四个维度选指标，用熵权法客观赋权，再用 TOPSIS 计算每个地区与"最优/最劣方案"的距离，得到风险得分；按得分把地区分成四级（脆弱/较差/一般/良好）；用夜间遥感影像把结果画在深圳四类区域上；最后用微分方程模拟干预措施实施后的光污染衰减过程。

\textbf{为什么这么选}：多指标综合排序问题，TOPSIS 是最经典的解法之一，配合熵权法避免主观赋权；GIS 遥感数据则让评价有空间可视化。

\textbf{可借鉴亮点}：\textbf{EWM-TOPSIS 是环境/评价类题目的"黄金组合"}，几乎可以套用到所有"多指标打分排序"问题。

\subsection{论文 2023E-2305598：AHP+熵权+变异系数 + 模糊聚类}
\textbf{使用模型}：AHP（第 8 章）、熵权法（第 8 章）、变异系数法（第 8 章）、模糊聚类（第 5 章）、BP 神经网络（第 11 章）、Spearman 相关（第 2 章）。

\textbf{模型组合逻辑}：把 10 个指标分成光线、社会、自然三类；分别用 AHP（主观）、熵权法（客观）、变异系数法（客观）三套权重，再融合成综合权重，构建光污染指数 LPI；用模糊聚类把 55 个地区分成四类；用 BP 神经网络预测干预后的风险等级变化。

\textbf{为什么这么选}：三种赋权法交叉验证，结论稳健；模糊聚类比硬聚类更适合"分级边界模糊"的环境问题。

\textbf{可借鉴亮点}：\textbf{三套权重融合+模糊聚类分级}，把"打分"和"分类"两步都做扎实，是环境评价题的标准答案范式。

\subsection{论文 2023E-2312411：EWM-TOPSIS + ArcGIS + 克里金插值}
\textbf{使用模型}：EWM-TOPSIS（第 8 章）、ArcGIS 空间分析（第 14 章）、克里金插值 Kriging（第 14 章）、ARIMA 时序（第 6 章）、逻辑回归（第 4 章）。

\textbf{模型组合逻辑}：先用 EWM-TOPSIS 给陕西 7 个指标打分；用 ArcGIS 提取地块几何中心；用克里金插值把离散点的光污染指数插值成连续面；再用 ARIMA 预测未来光污染趋势；用逻辑回归判断干预措施是否达标。

\textbf{为什么这么选}：环境监测站点稀疏，克里金插值能把点数据扩展成连续地图；空间插值+综合评价的组合让结果既能打分又能画图。

\textbf{可借鉴亮点}：\textbf{GIS+克里金插值}是处理"区域环境评价"类题目的利器，比单纯列表打分更直观。

% ============================================================
\section{2024 年 C 题：网球比赛中的"动量"}

\textbf{题目背景}：2023 温网决赛 Alcaraz 逆转 Djokovic，球迷热议"动量"（momentum）。题目要求量化动量、判断动量是否真实存在、并预测动量何时切换。

\subsection{论文 2024C-2406324：EWMA 动量 + LSTM+XGBoost + K-means 迁移学习 + SHAP}
\textbf{使用模型}：指数加权移动平均 EWMA（第 6 章）、LSTM（第 7 章）、XGBoost（第 11 章）、K-means 聚类（第 5 章）、SHAP 可解释性（第 11 章）、敏感性分析（第 1 章）、游程检验（第 3 章）。

\textbf{模型组合逻辑}：先用 EWMA 给每个时刻的胜率波动做平滑，定义"动量分"；用游程检验判断连胜连败是否随机；然后用 LSTM 预测动量时序、用 XGBoost 做动量贡献因子分解；引入 K-means 迁移学习把女子比赛数据辅助男子比赛预测；用 SHAP 值解释哪些因素最影响动量切换。

\textbf{为什么这么选}：动量本质是\textbf{时间序列上的局部趋势}，EWMA 是最自然的定义方式；LSTM 抓时序依赖，XGBoost+SHAP 兼顾精度和可解释性。

\textbf{可借鉴亮点}：\textbf{SHAP 可解释性}是机器学习论文的加分项——不仅预测，还能告诉评委"为什么是这个因素"。

\subsection{论文 2024C-2401919：PCA-TOPSIS 实时评价 + KMO/Bartlett 检验}
\textbf{使用模型}：PCA 主成分（第 5 章）、TOPSIS（第 8 章）、KMO 检验与 Bartlett 球形检验（第 5 章）、白噪声检验（第 6 章）、描述性统计（第 1 章）。

\textbf{模型组合逻辑}：从 5 个维度 18 个指标构建球员实时表现评分体系；用 KMO 和 Bartlett 检验确认数据适合做因子分析后，用 PCA 降维；再用 TOPSIS 计算每个时刻两名球员与理想解的距离，得到动量曲线；用白噪声检验判断动量是否随机波动。

\textbf{为什么这么选}：指标多达 18 个，直接加权会重复，PCA 先浓缩；TOPSIS 把浓缩后的指标合成单一得分；统计检验则用来回答"动量是不是真的"。

\textbf{可借鉴亮点}：\textbf{PCA+TOPSIS 组合}非常适合"指标多、要打分"的场景，KMO/Bartlett 检验则让降维步骤显得规范。

\subsection{论文 2024C-2403774：XGBoost 动量量化 + 时间卷积 + IQR 异常值}
\textbf{使用模型}：XGBoost（第 11 章）、时间卷积网络 TCN（第 7 章）、IQR 异常值检测（第 1 章）、Pearson/Spearman 相关（第 2 章）、交叉验证（第 11 章）、RMSE/MAE（第 11 章）。

\textbf{模型组合逻辑}：先用 IQR 法识别并剔除异常值；按不同时间窗口构造"近期表现""历史交锋""发球优势"等特征；用 XGBoost 回归实时胜率；用时间卷积网络捕捉更长时间尺度的动量模式；用交叉验证调参。

\textbf{为什么这么选}：体育数据噪声大，IQR 比均值±3σ 更稳健；XGBoost 在表格数据上训练快、可解释，比深度网络更适合这种中小数据集。

\textbf{可借鉴亮点}：\textbf{先做特征工程再上模型}——按时间窗口手工构造特征，往往比端到端深度学习更有效。

% ============================================================
\section{2024 年 E 题：极端天气下的财产保险可持续性}

\textbf{题目背景}：极端天气频发导致保险公司赔付飙升。题目要求评估保险公司在高风险地区是否可持续经营、如何定价保费、以及如何保护历史地标建筑。

\subsection{论文 2024E-2413552：多尺度地理加权回归 MGWR + PCA}
\textbf{使用模型}：多尺度地理加权回归 MGWR（第 14 章 空间计量）、PCA（第 5 章）、GIS（第 14 章）、熵权法（第 8 章）、优化模型（第 12 章）。

\textbf{模型组合逻辑}：把气候数据、地理位置、保险赔付数据整合到一起；用 MGWR 回归揭示"极端天气概率→赔付额"的关系如何随空间变化（不同地区斜率不同）；用 PCA 把多维气候指标压缩；用熵权法给地区风险打分；最后在风险定价模型中优化保费。

\textbf{为什么这么选}：普通 OLS 假设"全国一套系数"，但飓风在佛罗里达和在阿拉斯加的影响完全不同；MGWR 允许每个地区有自己的回归斜率，是空间异质性问题的标准答案。

\textbf{可借鉴亮点}：\textbf{MGWR/GWR 是空间类题目的"降维打击"武器}——一旦数据带地理坐标，就应该考虑它。

\subsection{论文 2024E-2407414：POT 极值理论 + 泊松分布 + 破产理论}
\textbf{使用模型}：超阈值峰值 POT 极值统计（第 6/12 章）、泊松分布（第 1 章）、破产理论 Ruin Theory（第 12 章）、AHP（第 8 章）、GIS（第 14 章）。

\textbf{模型组合逻辑}：用 POT 方法拟合灾害损失超过阈值的尾部分布（广义帕累托分布）；用泊松分布拟合灾害发生频率；据此计算保险公司的破产概率（POR）；用破产理论确定最优自留额和投资策略；用 AHP 给历史地标的保护优先级打分。

\textbf{为什么这么选}：保险赔付是典型的\textbf{厚尾分布}——平时小赔、偶尔巨赔，普通正态假设会严重低估风险；极值理论专门处理这种尾部问题。

\textbf{可借鉴亮点}：\textbf{POT 极值理论+破产概率}是精算/风控题目的专业组合，比简单线性回归显得更懂保险。

\subsection{论文 2024E-2401102：ARIMA 风险预测 + Lane AFC 保费模型}
\textbf{使用模型}：ARIMA 时序（第 6 章）、Lane AFC 保费定价模型（第 8/12 章）、保险需求曲线（第 12 章）、AHP（第 8 章）、优化模型（第 12 章）。

\textbf{模型组合逻辑}：按"地区-灾害类型"配对，用 ARIMA 预测未来极端天气发生频率；据此计算预期损失；用 Lane AFC 模型把预期损失转换成保费；再引入保险需求曲线，刻画保费高低如何影响客户数量；最后在"保费收入-赔付支出-客户留存"之间优化。

\textbf{为什么这么选}：ARIMA 负责"未来风险多大"，Lane AFC 负责"该收多少钱"，需求曲线负责"客户会不会买"——三步闭环正好对应保险经营的三个环节。

\textbf{可借鉴亮点}：把行业经典定价模型（Lane AFC）直接拿来用，比自己拍脑袋定保费更有依据。

% ============================================================
\section{2025 年 C 题：2028 洛杉矶奥运会奖牌预测}

\textbf{题目背景}：根据历届奥运数据，预测 2028 洛杉矶奥运会各国奖牌数、识别首次夺牌国家、评估"伟大教练"对成绩的影响。

\subsection{论文 2025C-2501869：Hurdle-Tobit 融合模型 + RCA 显示性比较优势}
\textbf{使用模型}：Hurdle 两阶段模型（第 4 章 回归分析）、Tobit 删失回归（第 4 章）、显示性比较优势指数 RCA（第 8 章）、随机森林/ XGBoost 对比（第 11 章）、EDA（第 1 章）。

\textbf{模型组合逻辑}：奖牌数有大量零值（很多国家没拿过牌）且右删失，普通 OLS 会有偏；于是用 Hurdle 模型分两步——先用 Logistic 回归判断"是否夺牌"，再用截断计数模型预测"夺多少块"；对右删失部分用 Tobit 模型；引入 RCA 指数衡量各国在某个项目上的传统优势；最后与随机森林、XGBoost 对比，发现 Hurdle-Tobit 在小样本和预测覆盖率上更优。

\textbf{为什么这么选}：奖牌数据是典型的\textbf{零膨胀+删失计数数据}，Hurdle/Tobit 这类计量经济模型比通用机器学习更贴合数据生成过程。

\textbf{可借鉴亮点}：\textbf{针对数据特点选模型}——零膨胀用 Hurdle、删失用 Tobit，而不是无脑上深度学习，这是获奖论文的重要特征。

\subsection{论文 2025C-2500759：PCA 降维 + LSTM + XGBoost-Bootstrap 双通道}
\textbf{使用模型}：PCA（第 5 章）、LSTM（第 7 章）、XGBoost（第 11 章）、Bootstrap 重抽样（第 1 章）、Spearman 相关（第 2 章）、ARMA（第 6 章）。

\textbf{模型组合逻辑}：先用 PCA 对赛事多维数据降维降噪；用 LSTM 挖掘各国奖牌数的时序演化趋势并纳入东道主效应；再用 XGBoost-Bootstrap 双通道输出带 95\% 置信区间的预测；用 Spearman 相关分析项目优势与奖牌数的关系。

\textbf{为什么这么选}：PCA 解决奥运指标维度爆炸问题；LSTM 抓时间依赖；Bootstrap 给出置信区间——三者组合覆盖了"降噪-预测-不确定性量化"全流程。

\textbf{可借鉴亮点}：\textbf{双通道模型+Bootstrap 区间估计}，比单点预测更符合奥运奖牌这种不确定事件的本质。

\subsection{论文 2025C-2514362：混合效应负二项回归 + 历史阶段修正}
\textbf{使用模型}：混合效应负二项回归（第 4 章）、Bootstrap（第 1 章）、交叉验证（第 11 章）、SVR（第 11 章）、敏感性分析（第 1 章）。

\textbf{模型组合逻辑}：先做时空数据清洗，发现 1952--1991 年间奖牌分布因冷战存在结构性断点，引入"冷战修正因子"；然后构建混合效应负二项回归，把国家作为随机效应，把东道主效应、参赛人数、历史成绩作为固定效应；用 Bootstrap 和交叉验证检验稳健性。

\textbf{为什么这么选}：奖牌数是离散计数且过度离散，负二项回归比 Poisson 更合适；混合效应处理"每个国家是一个面板"的结构；历史断点修正则让模型不会把冷战时期的异常趋势外推到未来。

\textbf{可借鉴亮点}：\textbf{对历史数据做结构性诊断}——发现断点后专门引入修正变量，这种细节是区分优秀论文和平庸论文的关键。

% ============================================================
\section{2025 年 E 题：从森林到农田的生态系统稳定}

\textbf{题目背景}：森林被砍伐后转为有机农田，如何在这一转型过程中维持生态系统稳定？题目要求建模预测农药、化肥、季节变化对食物链/网的影响。

\subsection{论文 2025E-2508861：Lotka-Volterra + 广义 Logistic + 空间扩散}
\textbf{使用模型}：Lotka-Volterra 捕食者-猎物方程（第 12 章）、广义 Logistic 增长（第 12 章）、空间扩散项（第 14 章 空间计量）、季节变化函数（第 12 章）、线性规划（第 12 章）。

\textbf{模型组合逻辑}：把农田生态系统抽象为作物、害虫、天敌三级食物链；用 Lotka-Volterra 方程描述捕食关系；加入广义 Logistic 项刻画环境承载力；加入空间扩散项描述物种迁移；加入周期函数模拟季节性农药/化肥施用；对比"不用杀虫剂""不用化肥""三者都用"三种情景下的生态稳定性。

\textbf{为什么这么选}：生态系统是典型的\textbf{多物种耦合动力系统}，Lotka-Volterra 是最经典的机理模型；空间扩散和季节项则让模型从"实验室"走向"真实农田"。

\textbf{可借鉴亮点}：\textbf{机理模型+情景对比}——不直接预测一个数字，而是比较多种管理方案下的系统轨迹，这是环境生态类题目的标准做法。

\subsection{论文 2025E-2515136：Petri 网食物网模型 + Runge-Kutta 求解}
\textbf{使用模型}：Petri 网（第 12 章）、Lotka-Volterra（第 12 章）、能量流理论（第 8 章）、Runge-Kutta 数值求解（第 12 章）、模糊综合评价（第 8 章）。

\textbf{模型组合逻辑}：用 Petri 网把食物网抽象为"库所（物种）+变迁（捕食/死亡）"的网络；在 Petri 网框架下嵌入 Lotka-Volterra 方程和能量流分配；用四阶 Runge-Kutta 数值求解逐年生物量轨迹；最后用模糊综合评价对不同管理方案打分。

\textbf{为什么这么选}：当物种数多、食物关系复杂时，纯 ODE 方程组会变得很难写；Petri 网用图形化方式组织状态和转移，更适合复杂食物网。

\textbf{可借鉴亮点}：\textbf{Petri 网+动力学方程}的组合，既保留了机理模型的严谨，又用网络结构处理了复杂关系。

\subsection{论文 2025E-2502355：递进式微分方程建模}
\textbf{使用模型}：常微分方程（第 12 章）、系统动力学（第 12 章）、AHP（第 8 章）、熵权法（第 8 章）、敏感性分析（第 1 章）。

\textbf{模型组合逻辑}：论文按"从简到繁"的思路逐步升级：先建最简单的食物链 ODE，再加入农药和除草剂的影响，再加入农业周期和季节波动，最后加入温度、施肥等参数；每加一层都与上一层对比，观察生物量曲线如何变化。

\textbf{为什么这么选}：生态系统太复杂，一上来就建大模型容易失控；\textbf{递进式建模}让每一步的贡献都清晰可辨，也方便评委理解。

\textbf{可借鉴亮点}：\textbf{"积木式"建模}——每加一个变量就重新跑一遍并对比结果，这种做法在论文里会显得思路极其清晰。

% ============================================================
\section{本章小结}

通览 12 道题、36 篇获奖论文，可以总结出几条\textbf{模型组合的通用规律}：

\begin{enumerate}
  \item \textbf{数据预处理层}：几乎所有论文都先用描述性统计（第 1 章）和 EDA 摸清数据，再做缺失值/异常值处理；文本类题目（2020C、2021C 图片）还会加一层文本/图像特征工程。
  \item \textbf{降维/赋权层}：指标多时用 PCA/因子分析（第 5 章）降维，需要打分时用 AHP/熵权法/TOPSIS（第 8 章）做综合评价。
  \item \textbf{核心建模层}：时序问题用 ARIMA/Prophet/LSTM（第 6、7 章），空间问题用 GIS/GWR/Moran（第 14 章），机理问题用微分方程/系统动力学/元胞自动机（第 12 章），分类预测问题用 XGBoost/RF/神经网络（第 11 章）。
  \item \textbf{优化决策层}：涉及"找最优策略"的题目，再嵌套线性规划、动态规划、遗传算法或 NSGA-II（第 12 章）。
  \item \textbf{评价检验层}：最后用敏感性分析、交叉验证、Bootstrap、ROC/AUC 证明结论稳健。
\end{enumerate}

记住一个核心思想：\textbf{没有万能模型，只有合适的模型组合}。获奖论文的共同点不是用了最复杂的算法，而是每个模型都在流水线里承担了一件清晰、不可替代的任务。


\end{document}
